% === START OF CHUNK preamble+ch1-2 ===
% ============================================================
%  Differential Topology — A Graduate Course
%  Preamble + Chapter 1 + Chapter 2
% ============================================================

\documentclass[12pt,a4paper,openany]{book}

% ── Encoding and fonts ──────────────────────────────────────
\usepackage[utf8]{inputenc}
\usepackage[T1]{fontenc}
\usepackage{lmodern}
\usepackage[margin=2.5cm]{geometry}

% ── Mathematics ─────────────────────────────────────────────
\usepackage{amsmath,amssymb,amsthm,mathrsfs,mathtools}
\usepackage{esint}

% ── Graphics ────────────────────────────────────────────────
\usepackage{tikz}
\usetikzlibrary{arrows.meta,calc,positioning,patterns,decorations.markings,decorations.pathmorphing,cd}
\usepackage{pgfplots}
\pgfplotsset{compat=1.18}

% ── Colored boxes ───────────────────────────────────────────
\usepackage[most]{tcolorbox}
\tcbuselibrary{theorems,skins,breakable}

% ── Hyperlinks ──────────────────────────────────────────────
\usepackage[colorlinks=true,linkcolor=blue!60!black,citecolor=green!50!black]{hyperref}
\usepackage{cleveref}

% ── Index ───────────────────────────────────────────────────
\usepackage{makeidx}
\makeindex

% ── Lists ───────────────────────────────────────────────────
\usepackage{enumitem}

% ── Headers ─────────────────────────────────────────────────
\usepackage{fancyhdr}

% ── Tables and captions ────────────────────────────────────
\usepackage{caption}
\usepackage{booktabs}
\usepackage{array}
\usepackage{longtable}

% ── Mini table of contents ─────────────────────────────────
\usepackage{minitoc}

% ── Micro-typographic refinements ──────────────────────────
\usepackage{microtype}

% ── Colors ──────────────────────────────────────────────────
\usepackage{xcolor}

% ── Commutative diagrams ───────────────────────────────────
\usepackage{tikz-cd}

% ============================================================
%  Theorem-like environments
% ============================================================

\newtheorem{theorem}{Theorem}[section]
\newtheorem{proposition}[theorem]{Proposition}
\newtheorem{lemma}[theorem]{Lemma}
\newtheorem{corollary}[theorem]{Corollary}

\theoremstyle{definition}
\newtheorem{definition}[theorem]{Definition}
\newtheorem{example}[theorem]{Example}
\newtheorem{exercise}{Exercise}[chapter]
\newtheorem{notation}[theorem]{Notation}

\theoremstyle{remark}
\newtheorem{remark}[theorem]{Remark}

% ── tcolorbox styling for theorem environments ─────────────

\tcolorboxenvironment{theorem}{
  enhanced, breakable,
  colback=blue!3, colframe=blue!50!black,
  boxrule=0.8pt, arc=2pt,
  left=6pt, right=6pt, top=4pt, bottom=4pt,
  fonttitle=\bfseries
}

\tcolorboxenvironment{proposition}{
  enhanced, breakable,
  colback=cyan!3, colframe=cyan!50!black,
  boxrule=0.8pt, arc=2pt,
  left=6pt, right=6pt, top=4pt, bottom=4pt
}

\tcolorboxenvironment{lemma}{
  enhanced, breakable,
  colback=green!3, colframe=green!50!black,
  boxrule=0.8pt, arc=2pt,
  left=6pt, right=6pt, top=4pt, bottom=4pt
}

\tcolorboxenvironment{corollary}{
  enhanced, breakable,
  colback=violet!3, colframe=violet!50!black,
  boxrule=0.8pt, arc=2pt,
  left=6pt, right=6pt, top=4pt, bottom=4pt
}

\tcolorboxenvironment{definition}{
  enhanced, breakable,
  colback=orange!4, colframe=orange!70!black,
  boxrule=0.8pt, arc=2pt,
  left=6pt, right=6pt, top=4pt, bottom=4pt
}

\tcolorboxenvironment{example}{
  enhanced, breakable,
  colback=yellow!5, colframe=yellow!60!black,
  boxrule=0.6pt, arc=2pt,
  left=6pt, right=6pt, top=4pt, bottom=4pt
}

\tcolorboxenvironment{exercise}{
  enhanced, breakable,
  colback=red!3, colframe=red!50!black,
  boxrule=0.6pt, arc=2pt,
  left=6pt, right=6pt, top=4pt, bottom=4pt
}

\tcolorboxenvironment{remark}{
  enhanced, breakable,
  colback=gray!5, colframe=gray!50!black,
  boxrule=0.5pt, arc=2pt,
  left=6pt, right=6pt, top=4pt, bottom=4pt
}

\tcolorboxenvironment{notation}{
  enhanced, breakable,
  colback=teal!4, colframe=teal!50!black,
  boxrule=0.5pt, arc=2pt,
  left=6pt, right=6pt, top=4pt, bottom=4pt
}

% ============================================================
%  Custom commands
% ============================================================

\newcommand{\N}{\mathbb{N}}
\newcommand{\Z}{\mathbb{Z}}
\newcommand{\Q}{\mathbb{Q}}
\newcommand{\R}{\mathbb{R}}
\newcommand{\C}{\mathbb{C}}
\newcommand{\RP}{\mathbb{R}\mathrm{P}}
\newcommand{\CP}{\mathbb{C}\mathrm{P}}
\newcommand{\Id}{\mathrm{Id}}
\DeclareMathOperator{\im}{im}
\newcommand{\GL}{\mathrm{GL}}
\newcommand{\SL}{\mathrm{SL}}
\newcommand{\SO}{\mathrm{SO}}
\newcommand{\SU}{\mathrm{SU}}
\newcommand{\Orth}{\mathrm{O}}
\newcommand{\Uni}{\mathrm{U}}
\newcommand{\Diff}{\mathrm{Diff}}
\newcommand{\dd}{\,\mathrm{d}}
\newcommand{\abs}[1]{\left\lvert #1 \right\rvert}
\newcommand{\norm}[1]{\left\lVert #1 \right\rVert}
\newcommand{\set}[1]{\left\{ #1 \right\}}
\newcommand{\rel}{\sim}
\newcommand{\Gr}{\mathrm{Gr}}
\DeclareMathOperator{\rank}{rank}
\DeclareMathOperator{\codim}{codim}
\DeclareMathOperator{\sgn}{sgn}
\DeclareMathOperator{\Hom}{Hom}

% ── Page style ──────────────────────────────────────────────
\pagestyle{fancy}
\fancyhf{}
\fancyhead[LE]{\leftmark}
\fancyhead[RO]{\rightmark}
\fancyfoot[C]{\thepage}
\renewcommand{\headrulewidth}{0.4pt}

% ============================================================
\begin{document}
% ============================================================

\dominitoc

\begin{titlepage}
\begin{tikzpicture}[remember picture, overlay]
  \fill[left color=blue!15!white, right color=blue!5!white]
    (current page.south west) rectangle (current page.north east);
  \fill[color=orange!70!yellow]
    ([yshift=-2cm]current page.north west) rectangle ([yshift=-2.3cm]current page.north east);
  \fill[color=blue!80!black, rounded corners=8pt]
    ([xshift=3cm, yshift=-4cm]current page.north west) rectangle
    ([xshift=-3cm, yshift=-10cm]current page.north east);
  \node[anchor=center, text=white, font=\Huge\bfseries]
    at ([yshift=-6cm]current page.north) {Differential Topology};
  \node[anchor=center, text=orange!80!yellow, font=\Large]
    at ([yshift=-7.5cm]current page.north) {Notes de Cours / Lecture Notes};
  \node[anchor=center, text=white!80, font=\large]
    at ([yshift=-8.8cm]current page.north) {Graduate Level --- 2025--2026};
  \node[anchor=center, text=blue!80!black, font=\Large\itshape]
    at ([yshift=-12cm]current page.north) {Ya\'e Ulrich Gaba};
  \draw[orange!70!yellow, line width=2pt]
    ([xshift=5cm, yshift=-13.5cm]current page.north west) --
    ([xshift=-5cm, yshift=-13.5cm]current page.north east);
  \node[anchor=center, text width=12cm, align=center, text=blue!60!black, font=\itshape]
    at ([yshift=-16cm]current page.north) {Manifolds, embeddings, and Morse theory: smooth structures at the heart of geometry.};
  \node[anchor=south, text=gray, font=\small]
    at ([yshift=1.5cm]current page.south) {\today};
\end{tikzpicture}
\end{titlepage}
\tableofcontents

% ============================================================
%  Notation Table
% ============================================================

\chapter*{Table of Notation}
\addcontentsline{toc}{chapter}{Table of Notation}
\markboth{Table of Notation}{Table of Notation}

\begin{longtable}{>{\raggedright}p{3.8cm} p{9.5cm}}
\toprule
\textbf{Symbol} & \textbf{Meaning} \\
\midrule
\endfirsthead
\toprule
\textbf{Symbol} & \textbf{Meaning} \\
\midrule
\endhead
\bottomrule
\endfoot
$\N, \Z, \Q, \R, \C$ & Natural numbers, integers, rationals, reals, complex numbers \\
$\R^n$ & Euclidean $n$-space \\
$S^n$ & Unit $n$-sphere in $\R^{n+1}$ \\
$\RP^n$ & Real projective $n$-space \\
$\CP^n$ & Complex projective $n$-space \\
$T^n$ & $n$-torus $S^1 \times \cdots \times S^1$ \\
$\GL(n,\R)$, $\GL(n,\C)$ & General linear groups \\
$\SL(n,\R)$, $\SL(n,\C)$ & Special linear groups \\
$\Orth(n)$, $\SO(n)$ & Orthogonal and special orthogonal groups \\
$\Uni(n)$, $\SU(n)$ & Unitary and special unitary groups \\
$\Gr(k,n)$ & Grassmannian of $k$-planes in $\R^n$ \\
$\Diff(M)$ & Diffeomorphism group of $M$ \\
$C^\infty(M)$ & Smooth real-valued functions on $M$ \\
$C^\infty(M,N)$ & Smooth maps from $M$ to $N$ \\
$T_pM$ & Tangent space at $p$ \\
$TM$ & Tangent bundle of $M$ \\
$\dd f_p$, $Df(p)$ & Differential of $f$ at $p$ \\
$\Id_M$ & Identity map on $M$ \\
$\partial M$ & Boundary of the manifold $M$ \\
$\operatorname{Int}(M)$ & Interior of $M$ \\
$\abs{\cdot}$, $\norm{\cdot}$ & Absolute value, norm \\
$\rel$ & Equivalence relation \\
\bottomrule
\end{longtable}

% ============================================================
%  Preface
% ============================================================

\chapter*{Preface}
\addcontentsline{toc}{chapter}{Preface}
\markboth{Preface}{Preface}

Differential topology studies the properties of smooth manifolds and smooth
maps that are invariant under diffeomorphisms.  Unlike differential geometry,
which equips manifolds with additional structure such as Riemannian metrics
or connections, differential topology works with the bare smooth structure
alone.  The subject is therefore closer in spirit to algebraic topology, yet
its methods are overwhelmingly analytic and geometric.

This text is addressed to graduate students who have completed a solid course
in point-set topology and possess working knowledge of linear algebra,
multivariable calculus, and basic group theory.  Some familiarity with
algebraic topology (fundamental group, covering spaces) is helpful but not
strictly required for the first half of the book.

The first two chapters establish the language of smooth manifolds and smooth
maps.  We take care to present full proofs of the foundational results ---
the existence of partitions of unity, the inverse and implicit function
theorems on manifolds, and the constant rank theorem --- as they underpin
everything that follows.  The exposition is supplemented by numerous examples
drawn from geometry and Lie group theory, as well as exercises that range
from routine verifications to more demanding problems.

\medskip
\noindent\textit{Prerequisites.}\enspace
Second-year graduate standing in mathematics; point-set topology
(Hausdorff spaces, compactness, paracompactness, quotient topology);
multivariable analysis (inverse function theorem in $\R^n$,
implicit function theorem); linear algebra.

% ============================================================
% ============================================================
%
%         CHAPTER 1  —  SMOOTH MANIFOLDS
%
% ============================================================
% ============================================================

\chapter{Smooth Manifolds --- Definitions and Examples}
\label{ch:smooth-manifolds}
\minitoc
\bigskip

\noindent
The notion of a \emph{smooth manifold}\index{manifold!smooth} lies at the
very foundation of differential topology.  Intuitively, a smooth manifold
is a topological space that ``locally looks like'' Euclidean space and on
which the tools of calculus can be applied in a consistent, coordinate-free
manner.  In this chapter we make this idea precise and develop the basic
vocabulary that will accompany us throughout the course.

% ============================================================
\section{Topological Preliminaries}
\label{sec:top-prelim}
% ============================================================

Before introducing manifolds, we recall the topological notions that appear
in their definition.

\begin{definition}[Second-countable space]\label{def:second-countable}
\index{second-countable}
A topological space $X$ is \emph{second-countable} if its topology admits a
countable base, i.e.\ there exists a countable collection
$\mathcal{B} = \{B_i\}_{i \in \N}$ of open sets such that every open set
$U \subseteq X$ can be written as a union of members of $\mathcal{B}$.
\end{definition}

\begin{definition}[Hausdorff space]\label{def:hausdorff}
\index{Hausdorff}
A topological space $X$ is \emph{Hausdorff} if for every pair of distinct
points $p,q \in X$ there exist disjoint open sets $U \ni p$ and $V \ni q$.
\end{definition}

\begin{definition}[Paracompact space]\label{def:paracompact}
\index{paracompact}
A topological space $X$ is \emph{paracompact} if every open cover of $X$
admits a locally finite open refinement.
\end{definition}

\begin{proposition}[Second-countable implies paracompact]
\label{prop:second-countable-paracompact}
\index{paracompact!from second-countability}
Every second-countable Hausdorff space is paracompact.
\end{proposition}

\begin{proof}
Let $\{B_i\}_{i\in\N}$ be a countable base for the topology.  Given an open
cover $\{U_\alpha\}$, for each $x \in X$ choose some $U_\alpha$ containing
$x$ and then a basis element $B_{i(x)} \subseteq U_\alpha$ with
$x \in B_{i(x)}$.  The countable subcover $\{B_{i(x)}\}$ can be refined to
a locally finite cover by a standard exhaustion argument using compact sets
(which exist because a second-countable Hausdorff space is Lindel\"of and
every open cover has a countable subcover).
\end{proof}

\begin{remark}\label{rem:why-hausdorff-second-countable}
We require manifolds to be Hausdorff and second-countable to exclude
pathologies: the line with two origins is locally Euclidean but not
Hausdorff, and the long line is Hausdorff and locally Euclidean but not
second-countable.  Both conditions together guarantee that manifolds are
metrizable and paracompact, which is essential for the existence of
partitions of unity.
\end{remark}

% ============================================================
\section{Charts, Atlases, and Smooth Structures}
\label{sec:charts-atlases}
% ============================================================

\begin{definition}[Topological manifold]\label{def:top-manifold}
\index{manifold!topological}
A \emph{topological manifold of dimension $n$} is a topological space $M$
that is:
\begin{enumerate}[label=(\roman*)]
  \item Hausdorff,
  \item second-countable, and
  \item locally Euclidean of dimension $n$: every point $p \in M$ possesses
        an open neighbourhood $U$ homeomorphic to an open subset of $\R^n$.
\end{enumerate}
\end{definition}

\begin{definition}[Chart]\label{def:chart}
\index{chart}
A \emph{chart} on a topological manifold $M^n$ is a pair $(U,\varphi)$
where $U \subseteq M$ is open and $\varphi \colon U \to \varphi(U)
\subseteq \R^n$ is a homeomorphism onto an open subset of $\R^n$.  We call
$U$ the \emph{chart domain}\index{chart!domain} and $\varphi$ the
\emph{coordinate map}\index{coordinate map}.  If $\varphi(p) = 0$ we say
the chart is \emph{centred}\index{chart!centred} at $p$.
\end{definition}

\begin{definition}[Transition map]\label{def:transition}
\index{transition map}
Given two charts $(U_\alpha, \varphi_\alpha)$ and
$(U_\beta, \varphi_\beta)$ on $M$ with $U_\alpha \cap U_\beta \neq
\varnothing$, the \emph{transition map} is
\[
  \varphi_{\beta\alpha}
  \;=\;
  \varphi_\beta \circ \varphi_\alpha^{-1}
  \colon
  \varphi_\alpha(U_\alpha \cap U_\beta)
  \longrightarrow
  \varphi_\beta(U_\alpha \cap U_\beta).
\]
This is a homeomorphism between open subsets of $\R^n$.
\end{definition}

% ── TikZ: chart and transition map ────────────────────────

\begin{center}
\begin{tikzpicture}[scale=1.0,
  every node/.style={font=\small},
  >={Stealth[length=5pt]}]
  % Manifold
  \draw[thick, fill=blue!5, rounded corners=20pt]
    (-1.2,1.8) -- (2.8,2.2) -- (4.2,1.0) -- (3.0,-0.8)
    -- (-0.2,-0.6) -- cycle;
  \node at (1.5,2.5) {$M$};
  % U_alpha
  \draw[thick, blue!60, fill=blue!10, opacity=0.6]
    (0.2,0.3) ellipse (1.2 and 0.9);
  \node[blue!70!black] at (-0.3,1.3) {$U_\alpha$};
  % U_beta
  \draw[thick, red!60, fill=red!10, opacity=0.6]
    (2.2,0.5) ellipse (1.1 and 0.85);
  \node[red!70!black] at (3.0,1.4) {$U_\beta$};
  % Point p
  \fill (1.2,0.4) circle (1.5pt) node[above] {$p$};
  % phi_alpha
  \draw[->, thick, blue!60!black]
    (-0.5,-0.5) -- (-2.0,-2.2) node[midway, left] {$\varphi_\alpha$};
  % phi_beta
  \draw[->, thick, red!60!black]
    (3.0,-0.3) -- (4.5,-2.2) node[midway, right] {$\varphi_\beta$};
  % R^n left
  \draw[thick, blue!30, fill=blue!5]
    (-3.8,-2.5) rectangle (-0.5,-4.5);
  \node[blue!70!black] at (-2.1,-2.3) {$\R^n$};
  \draw[blue!50, fill=blue!15, opacity=0.5]
    (-2.8,-3.5) ellipse (0.8 and 0.6);
  % R^n right
  \draw[thick, red!30, fill=red!5]
    (3.0,-2.5) rectangle (6.3,-4.5);
  \node[red!70!black] at (4.6,-2.3) {$\R^n$};
  \draw[red!50, fill=red!15, opacity=0.5]
    (5.0,-3.5) ellipse (0.8 and 0.6);
  % Transition map
  \draw[->, thick, purple!70!black, decorate,
    decoration={snake, amplitude=1pt, segment length=6pt}]
    (-1.8,-3.5) -- (4.0,-3.5)
    node[midway, below=4pt]
    {$\varphi_{\beta\alpha} = \varphi_\beta \circ \varphi_\alpha^{-1}$};
\end{tikzpicture}
\captionof{figure}{Two overlapping charts and their transition map.}
\label{fig:transition-map}
\end{center}

\begin{definition}[Smooth atlas]\label{def:smooth-atlas}
\index{atlas!smooth}
A \emph{smooth atlas} (or $C^\infty$ atlas) on $M$ is a collection
$\mathcal{A} = \{(U_\alpha, \varphi_\alpha)\}_{\alpha \in A}$ of charts
such that:
\begin{enumerate}[label=(\roman*)]
  \item $\{U_\alpha\}_{\alpha \in A}$ covers $M$;
  \item for every $\alpha, \beta \in A$ with $U_\alpha \cap U_\beta \neq
        \varnothing$, the transition map $\varphi_{\beta\alpha}$ is a
        smooth ($C^\infty$) diffeomorphism between open subsets of $\R^n$.
\end{enumerate}
\end{definition}

\begin{definition}[Compatible charts and maximal atlas]\label{def:maximal-atlas}
\index{atlas!maximal}\index{smooth structure}
Two smooth atlases $\mathcal{A}$ and $\mathcal{A}'$ on $M$ are
\emph{compatible} if $\mathcal{A} \cup \mathcal{A}'$ is again a smooth
atlas.  Compatibility is an equivalence relation.  A \emph{smooth
structure}\index{smooth structure} on $M$ is a maximal smooth atlas
$\overline{\mathcal{A}}$, i.e.\ the union of all atlases compatible with a
given atlas $\mathcal{A}$.  A \emph{smooth manifold}\index{manifold!smooth}
is a pair $(M, \overline{\mathcal{A}})$.
\end{definition}

\begin{remark}\label{rem:determined-by-atlas}
In practice, we specify a smooth manifold by giving \emph{any} smooth atlas;
the maximal atlas is then determined.  Two atlases define the same smooth
structure if and only if they are compatible.
\end{remark}

\begin{proposition}[Dimension is well-defined]\label{prop:dim-well-defined}
\index{dimension}
If $M$ is a connected topological manifold that is locally Euclidean of
dimension $n$ and also of dimension $m$, then $n = m$.
\end{proposition}

\begin{proof}
If $U \subseteq M$ lies in the domain of charts to both $\R^n$ and $\R^m$,
then there is a homeomorphism between open subsets of $\R^n$ and $\R^m$.
By Brouwer's invariance of domain\index{invariance of domain}, this forces
$n = m$.
\end{proof}

% ============================================================
\section{Examples of Smooth Manifolds}
\label{sec:examples-manifolds}
% ============================================================

\begin{example}[Euclidean space $\R^n$]\label{ex:Rn}
\index{Euclidean space}
The space $\R^n$ with the single chart $(\R^n, \Id_{\R^n})$ is a smooth
$n$-manifold.  More generally, any open subset $U \subseteq \R^n$ inherits
a smooth structure from the inclusion chart.
\end{example}

\begin{example}[The $n$-sphere $S^n$]\label{ex:sphere}
\index{sphere!$S^n$}
Define $S^n = \{x \in \R^{n+1} : \norm{x} = 1\}$.  We equip $S^n$ with a
smooth structure via \emph{stereographic projection}\index{stereographic projection}.
Let $N = (0,\ldots,0,1)$ and $S = (0,\ldots,0,-1)$ be the north and south
poles.  Define
\begin{align*}
  \varphi_N \colon S^n \setminus \{N\} &\longrightarrow \R^n, &
  \varphi_N(x_1,\ldots,x_{n+1}) &= \frac{1}{1 - x_{n+1}}
    (x_1, \ldots, x_n), \\
  \varphi_S \colon S^n \setminus \{S\} &\longrightarrow \R^n, &
  \varphi_S(x_1,\ldots,x_{n+1}) &= \frac{1}{1 + x_{n+1}}
    (x_1, \ldots, x_n).
\end{align*}
The transition map on $S^n \setminus \{N,S\}$ is
\[
  \varphi_S \circ \varphi_N^{-1}(y) = \frac{y}{\norm{y}^2},
\]
which is smooth (indeed real-analytic) on $\R^n \setminus \{0\}$.  Hence
$\{(S^n \setminus \{N\}, \varphi_N),\; (S^n \setminus \{S\}, \varphi_S)\}$
is a smooth atlas on $S^n$.
\end{example}

% ── TikZ: stereographic projection ────────────────────────

\begin{center}
\begin{tikzpicture}[scale=1.4, >={Stealth[length=5pt]}]
  % Circle (S^1 for illustration)
  \draw[thick] (0,0) circle (1.5);
  \fill (0,1.5) circle (1.5pt) node[above right] {$N$};
  \fill (0,-1.5) circle (1.5pt) node[below right] {$S$};
  % Point on circle
  \coordinate (P) at ({1.5*cos(50)},{1.5*sin(50)});
  \fill (P) circle (1.5pt) node[above right] {$p$};
  % Line from N through P to x-axis
  \draw[thick, blue!70, dashed] (0,1.5) -- (P) -- ({1.5*cos(50)/(1-sin(50))},0);
  \fill[blue!70] ({1.5*cos(50)/(1-sin(50))},0) circle (1.5pt);
  \node[blue!70!black, below] at ({1.5*cos(50)/(1-sin(50))},0) {$\varphi_N(p)$};
  % Equatorial line
  \draw[thick, gray!50] (-3,0) -- (3,0);
  \node[gray!70!black] at (3.0,-0.3) {$\R^n$};
\end{tikzpicture}
\captionof{figure}{Stereographic projection from the north pole.}
\label{fig:stereographic}
\end{center}

\begin{example}[Real projective space $\RP^n$]\label{ex:RPn}
\index{projective space!real}
Define $\RP^n = ({\R^{n+1} \setminus \{0\}})/{\rel}$ where
$x \rel \lambda x$ for all $\lambda \in \R \setminus \{0\}$.
Denote the equivalence class of $(x_0,\ldots,x_n)$ by
$[x_0 : \cdots : x_n]$.  For $i = 0,\ldots,n$ set
\[
  U_i = \{[x_0:\cdots:x_n] : x_i \neq 0\}, \qquad
  \varphi_i([x_0:\cdots:x_n])
  = \Bigl(\frac{x_0}{x_i},\ldots,\widehat{\frac{x_i}{x_i}},\ldots,
    \frac{x_n}{x_i}\Bigr) \in \R^n,
\]
where the hat denotes omission.  The transition maps are rational functions
with non-vanishing denominators, hence smooth.  The resulting atlas gives
$\RP^n$ the structure of a compact smooth $n$-manifold.
\end{example}

\begin{example}[Complex projective space $\CP^n$]\label{ex:CPn}
\index{projective space!complex}
Analogously, $\CP^n = (\C^{n+1} \setminus \{0\})/{\rel}$ with
$z \rel \lambda z$ for $\lambda \in \C \setminus \{0\}$.  Charts
$U_i = \{[z_0:\cdots:z_n] : z_i \neq 0\}$ with
$\varphi_i([z]) = (z_0/z_i, \ldots, \widehat{z_i/z_i}, \ldots, z_n/z_i)
\in \C^n \cong \R^{2n}$ yield a smooth atlas.  Thus $\CP^n$ is a compact
smooth manifold of (real) dimension $2n$.
\end{example}

\begin{example}[The torus $T^n$]\label{ex:torus}
\index{torus}
The \emph{$n$-torus} is $T^n = \underbrace{S^1 \times \cdots \times S^1}_{n}$.
As a product of smooth manifolds it inherits a product smooth structure.
Equivalently, $T^n \cong \R^n / \Z^n$ via the quotient by the lattice,
and the local sections of the projection $\R^n \to \R^n/\Z^n$ serve as
charts.
\end{example}

% ── TikZ: the 2-torus ────────────────────────────────────

\begin{center}
\begin{tikzpicture}[scale=1.2]
  % Outer ellipse (visible part of torus)
  \draw[thick] (0,0) ellipse (2.5 and 1.0);
  % Inner hole - front part
  \draw[thick] (0,0.15) ellipse (0.9 and 0.35);
  % Top curve
  \draw[thick, blue!60, dashed]
    plot[smooth, tension=0.7] coordinates
    {(-2.5,0) (-1.8,0.65) (-0.9,0.85) (0,0.88) (0.9,0.85) (1.8,0.65) (2.5,0)};
  % Label
  \node at (3.0,0) {$T^2$};
  % Arrows indicating two generating circles
  \draw[-{Stealth[length=4pt]}, red!70, thick]
    (2.5,0) arc[start angle=0, end angle=-40, x radius=2.5, y radius=1.0];
  \node[red!70!black] at (2.8,-0.7) {$\theta$};
  \draw[-{Stealth[length=4pt]}, green!60!black, thick]
    (1.7,0.58) arc[start angle=90, end angle=30, x radius=0.3, y radius=0.6];
  \node[green!60!black] at (2.2,0.8) {$\phi$};
\end{tikzpicture}
\captionof{figure}{The 2-torus $T^2$ with its two angular coordinates
$(\theta, \phi)$.}
\label{fig:torus}
\end{center}

\begin{example}[General linear group $\GL(n,\R)$]\label{ex:GLn}
\index{Lie group!$\GL(n,\R)$}\index{general linear group}
The set $\GL(n,\R)$ of invertible $n \times n$ real matrices is an open
subset of $\R^{n^2}$ (as $\det \colon \R^{n^2} \to \R$ is continuous and
$\GL(n,\R) = \det^{-1}(\R \setminus \{0\})$).  It is therefore a smooth
manifold of dimension $n^2$.  Matrix multiplication
$\GL(n,\R) \times \GL(n,\R) \to \GL(n,\R)$ and inversion
$A \mapsto A^{-1}$ are smooth maps (the entries of $A^{-1}$ are rational
functions of the entries of $A$), so $\GL(n,\R)$ is a
\emph{Lie group}\index{Lie group}.
\end{example}

\begin{example}[Special linear group $\SL(n,\R)$]\label{ex:SLn}
\index{Lie group!$\SL(n,\R)$}\index{special linear group}
The subgroup $\SL(n,\R) = \{A \in \GL(n,\R) : \det A = 1\}$ is the
preimage of the regular value $1$ under $\det \colon \GL(n,\R) \to \R^*$.
By the preimage theorem (which we prove in Chapter~\ref{ch:smooth-maps}),
$\SL(n,\R)$ is a smooth submanifold of dimension $n^2 - 1$, and is a Lie
group.
\end{example}

\begin{example}[Orthogonal group $\Orth(n)$ and $\SO(n)$]\label{ex:On}
\index{Lie group!$\Orth(n)$}\index{Lie group!$\SO(n)$}%
\index{orthogonal group}\index{special orthogonal group}
The orthogonal group $\Orth(n) = \{A \in \GL(n,\R) : A^T A = I_n\}$ is the
preimage of $I_n$ under the smooth map $F(A) = A^T A$ from $\GL(n,\R)$ to
the space $\mathrm{Sym}(n)$ of symmetric matrices.  Since $I_n$ is a
regular value of $F$, $\Orth(n)$ is a compact smooth manifold of dimension
$n^2 - \frac{n(n+1)}{2} = \frac{n(n-1)}{2}$.  Its identity component is
$\SO(n) = \{A \in \Orth(n) : \det A = 1\}$.
\end{example}

\begin{example}[Unitary group $\Uni(n)$ and $\SU(n)$]\label{ex:Un}
\index{Lie group!$\Uni(n)$}\index{Lie group!$\SU(n)$}%
\index{unitary group}\index{special unitary group}
Similarly, $\Uni(n) = \{A \in \GL(n,\C) : A^*A = I_n\}$ is a compact Lie
group of (real) dimension $n^2$, and
$\SU(n) = \{A \in \Uni(n) : \det A = 1\}$ has dimension $n^2 - 1$.
In particular $\SU(2)$ is diffeomorphic to $S^3$.
\end{example}

\begin{example}[Grassmannian $\Gr(k,n)$]\label{ex:grassmannian}
\index{Grassmannian}
The \emph{Grassmannian} $\Gr(k,n)$ is the set of all $k$-dimensional linear
subspaces of $\R^n$.  For each subset $I \subseteq \{1,\ldots,n\}$ of size
$k$, let $U_I$ consist of those $k$-planes $V$ such that the projection
$\R^n \to \R^I$ restricts to an isomorphism on $V$.  Identifying each such
$V$ with the unique $k \times (n-k)$ matrix expressing the complementary
coordinates in terms of the $I$-coordinates gives a chart
$\varphi_I \colon U_I \xrightarrow{\;\sim\;} \R^{k(n-k)}$.
The transition maps are rational, hence smooth, and $\Gr(k,n)$ is a compact
smooth manifold of dimension $k(n-k)$.  Note that
$\Gr(1, n+1) \cong \RP^n$.
\end{example}

% ============================================================
\section{Smooth Functions on Manifolds}
\label{sec:smooth-functions}
% ============================================================

Before turning to manifolds with boundary, we single out the important
special case of smooth \emph{real-valued} functions.

\begin{definition}[Smooth function]\label{def:smooth-function}
\index{smooth function}\index{$C^\infty(M)$}
Let $M$ be a smooth manifold.  A function $f \colon M \to \R$ is
\emph{smooth} if for every chart $(U,\varphi)$ in the smooth structure
of $M$, the composition $f \circ \varphi^{-1} \colon \varphi(U) \to \R$
is smooth in the ordinary sense of calculus.  The set of all smooth
functions on $M$ is denoted $C^\infty(M)$.
\end{definition}

\begin{proposition}\label{prop:Cinfty-algebra}
\index{algebra of smooth functions}
The set $C^\infty(M)$ is a commutative $\R$-algebra under pointwise
operations: for $f, g \in C^\infty(M)$ and $\lambda \in \R$,
\[
  (f+g)(p) = f(p) + g(p), \qquad
  (\lambda f)(p) = \lambda \cdot f(p), \qquad
  (fg)(p) = f(p)\,g(p).
\]
\end{proposition}

\begin{proof}
The sum, scalar multiple, and product of smooth functions on open subsets
of $\R^n$ are smooth.  The claim follows by checking in local coordinates.
\end{proof}

\begin{remark}\label{rem:smooth-determines-structure}
A fundamental fact (which we shall not prove here) is that the smooth
structure on a manifold $M$ is completely determined by the algebra
$C^\infty(M)$.  More precisely, if $M$ and $N$ are smooth manifolds and
$\Phi \colon C^\infty(N) \to C^\infty(M)$ is an $\R$-algebra isomorphism,
then there is a unique diffeomorphism $f \colon M \to N$ with
$\Phi(g) = g \circ f$ for all $g \in C^\infty(N)$.  This is the content
of Milnor's exercise, and it foreshadows the sheaf-theoretic viewpoint on
manifolds.
\end{remark}

\begin{example}[Height function on the sphere]\label{ex:height-function}
\index{height function}
The function $h \colon S^n \to \R$ defined by
$h(x_0, x_1, \ldots, x_n) = x_n$ is smooth.  Indeed, in the stereographic
chart $\varphi_S$ from Example~\ref{ex:sphere}, one computes
\[
  h \circ \varphi_S^{-1}(y) = \frac{\norm{y}^2 - 1}{\norm{y}^2 + 1},
\]
which is a smooth function on $\R^n$.  The critical points of $h$ are the
north and south poles.
\end{example}

\begin{example}[Smooth functions on $\RP^n$]\label{ex:smooth-on-RPn}
\index{projective space!smooth functions}
A function $f \colon \RP^n \to \R$ is smooth if and only if
$f \circ \pi \colon S^n \to \R$ is smooth, where $\pi \colon S^n \to \RP^n$
is the quotient map.  Moreover, $f \circ \pi$ must satisfy
$(f \circ \pi)(-x) = (f \circ \pi)(x)$ for all $x \in S^n$.  Hence
$C^\infty(\RP^n)$ identifies with the even smooth functions on $S^n$.
\end{example}

% ============================================================
\section{Manifolds with Boundary}
\label{sec:manifolds-boundary}
% ============================================================

\begin{notation}\label{not:half-space}
\index{half-space}
We denote the closed upper half-space by
$\mathbb{H}^n = \{(x_1,\ldots,x_n) \in \R^n : x_n \geq 0\}$.
Its boundary is $\partial \mathbb{H}^n = \{x_n = 0\} \cong \R^{n-1}$.
\end{notation}

\begin{definition}[Smooth manifold with boundary]\label{def:manifold-boundary}
\index{manifold!with boundary}
A \emph{smooth $n$-manifold with boundary} is a second-countable Hausdorff
space $M$ equipped with a maximal smooth atlas whose charts map open subsets
of $M$ homeomorphically to open subsets of $\mathbb{H}^n$, with smooth
transition maps.

The \emph{boundary}\index{boundary} $\partial M$ consists of all points
mapped to $\partial\mathbb{H}^n$ by some (hence every) chart.  The
\emph{interior} is $\operatorname{Int}(M) = M \setminus \partial M$.
\end{definition}

\begin{proposition}\label{prop:boundary-invariance}
The boundary $\partial M$ is well-defined: if $p$ maps to
$\partial \mathbb{H}^n$ under one chart, it does so under every chart
containing $p$.
\end{proposition}

\begin{proof}
Suppose charts $(U,\varphi)$ and $(V,\psi)$ both contain $p$, with
$\varphi(p) \in \partial \mathbb{H}^n$ but
$\psi(p) \in \operatorname{Int}(\mathbb{H}^n)$.  Then the transition map
$\psi \circ \varphi^{-1}$ sends a boundary point of $\mathbb{H}^n$ to an
interior point, contradicting the smooth invariance of domain
\index{invariance of domain} (a diffeomorphism between open subsets of
$\mathbb{H}^n$ maps boundary to boundary).
\end{proof}

\begin{proposition}\label{prop:boundary-manifold}
If $M$ is a smooth $n$-manifold with boundary, then $\partial M$ is a
smooth $(n-1)$-manifold without boundary.
\end{proposition}

\begin{proof}
For each boundary chart $(U, \varphi)$ the restriction
$\varphi|_{U \cap \partial M}$ maps $U \cap \partial M$ homeomorphically
onto an open subset of $\R^{n-1} \times \{0\} \cong \R^{n-1}$.  The
transition maps restrict to smooth diffeomorphisms between open subsets of
$\R^{n-1}$.
\end{proof}

\begin{example}[Closed disk]\label{ex:closed-disk}
\index{disk, closed}
The closed unit disk $\overline{D}^n = \{x \in \R^n : \norm{x} \leq 1\}$
is a smooth $n$-manifold with boundary $\partial \overline{D}^n = S^{n-1}$.
\end{example}

\begin{example}[Cylinder and M\"obius band]\label{ex:cylinder-moebius}
\index{cylinder}\index{M\"obius band}
The cylinder $S^1 \times [0,1]$ is a $2$-manifold with boundary
$S^1 \sqcup S^1$.  The M\"obius band is a non-orientable $2$-manifold
with boundary $S^1$.
\end{example}

\begin{example}[Half-spaces and quadrants]\label{ex:half-space-examples}
\index{half-space!examples}
The upper half-plane $\mathbb{H}^2 = \{(x,y) \in \R^2 : y \geq 0\}$ is a
$2$-manifold with boundary $\partial \mathbb{H}^2 = \R \times \{0\}$.
Note that the quadrant $Q = \{(x,y) : x \geq 0,\, y \geq 0\}$ is
\emph{not} a manifold with boundary (the origin has no neighbourhood
homeomorphic to an open subset of $\mathbb{H}^2$); rather, $Q$ is a
\emph{manifold with corners}\index{manifold!with corners}, a notion we do
not develop here.
\end{example}

% ── TikZ: manifold with boundary ──────────────────────────

\begin{center}
\begin{tikzpicture}[scale=1.1, >={Stealth[length=5pt]}]
  % M region
  \draw[thick, fill=green!8]
    plot[smooth cycle, tension=0.7] coordinates
    {(-2,0.5) (-1,2) (1.5,2.2) (3,1.5) (3.5,0) (2.5,-0.5) (0,-0.5)};
  \node at (0.8,1.8) {$M$};
  % Boundary highlighted
  \draw[very thick, red!70!black]
    plot[smooth, tension=0.8] coordinates
    {(-2,0.5) (-1.5,-0.1) (0,-0.5) (2.5,-0.5) (3.5,0)};
  \node[red!70!black] at (1.0,-1.0) {$\partial M$};
  % Chart to half-space
  \draw[->, thick] (4.0,0.5) -- (5.5,0.5) node[midway, above] {$\varphi$};
  % Half-space
  \draw[thick, fill=green!5] (6,-0.8) rectangle (9,2);
  \draw[very thick, red!60!black] (6,-0.8) -- (9,-0.8);
  \node at (7.5,1.5) {$\mathbb{H}^n$};
  \node[red!60!black] at (7.5,-1.2) {$\partial \mathbb{H}^n$};
\end{tikzpicture}
\captionof{figure}{A boundary chart maps a neighbourhood of a boundary
point to an open subset of the half-space $\mathbb{H}^n$.}
\label{fig:boundary-chart}
\end{center}

% ============================================================
\section{Partitions of Unity}
\label{sec:partitions-unity}
% ============================================================

Partitions of unity\index{partition of unity} are the single most important
technical tool in the theory of smooth manifolds.  They allow one to patch
together local constructions into global ones.

\begin{definition}[Support]\label{def:support}
\index{support}
The \emph{support} of a continuous function $f \colon M \to \R$ is
$\operatorname{supp}(f) = \overline{\{p \in M : f(p) \neq 0\}}$.
\end{definition}

\begin{definition}[Partition of unity]\label{def:partition-unity}
\index{partition of unity}
Let $\{U_\alpha\}_{\alpha \in A}$ be an open cover of a smooth manifold
$M$.  A \emph{smooth partition of unity subordinate to $\{U_\alpha\}$} is a
family $\{\rho_\alpha\}_{\alpha \in A}$ of smooth functions
$\rho_\alpha \colon M \to [0,1]$ such that:
\begin{enumerate}[label=(\roman*)]
  \item $\operatorname{supp}(\rho_\alpha) \subseteq U_\alpha$ for every
        $\alpha$;
  \item the collection $\{\operatorname{supp}(\rho_\alpha)\}$ is locally
        finite;
  \item $\sum_{\alpha} \rho_\alpha(p) = 1$ for every $p \in M$ (the sum
        is finite at each point by~(ii)).
\end{enumerate}
\end{definition}

We need a bump function\index{bump function} construction.

\begin{lemma}[Smooth bump function on $\R^n$]\label{lem:bump}
For any $0 < r < R$ there exists a smooth function
$\beta \colon \R^n \to [0,1]$ with $\beta \equiv 1$ on
$\overline{B}_r(0)$ and $\operatorname{supp}(\beta) \subseteq B_R(0)$.
\end{lemma}

\begin{proof}
Define $h \colon \R \to \R$ by
\[
  h(t) =
  \begin{cases}
    e^{-1/t} & t > 0,\\
    0 & t \leq 0.
  \end{cases}
\]
Then $h$ is smooth.  Set $g(t) = h(t)/(h(t) + h(1-t))$; this is a smooth
non-decreasing function with $g(t) = 0$ for $t \leq 0$ and $g(t) = 1$ for
$t \geq 1$.  Now define
\[
  \beta(x) = g\!\left(\frac{R^2 - \norm{x}^2}{R^2 - r^2}\right).
\]
This satisfies the stated properties.
\end{proof}

% ── TikZ: bump function ──────────────────────────────────

\begin{center}
\begin{tikzpicture}
  \begin{axis}[
    width=10cm, height=5.5cm,
    axis lines=middle,
    xlabel={$t$},
    ylabel={$\beta(t)$},
    xmin=-2.5, xmax=2.5,
    ymin=-0.15, ymax=1.25,
    xtick={-2,-1,1,2},
    xticklabels={$-R$,$-r$,$r$,$R$},
    ytick={0,1},
    yticklabels={$0$,$1$},
    samples=200,
    domain=-2.5:2.5,
    every axis x label/.style={at={(ticklabel* cs:1)}, anchor=west},
    every axis y label/.style={at={(ticklabel* cs:1)}, anchor=south},
  ]
    % Bump function (1D analogue)
    \addplot[thick, blue!70!black] {
      (abs(x) <= 1) * 1 +
      (abs(x) > 1) * (abs(x) < 2) * exp(-1/(1 - ((abs(x)-1))^2 / 1)) +
      (abs(x) >= 2) * 0
    };
    % Shading for support
    \addplot[fill=blue!10, draw=none, domain=-2:2] {1.2} \closedcycle;
    \node at (axis cs:0,1.12) {\footnotesize $\operatorname{supp}(\beta) \subseteq B_R(0)$};
  \end{axis}
\end{tikzpicture}
\captionof{figure}{A smooth bump function: $\beta \equiv 1$ on
$\overline{B}_r(0)$ and $\operatorname{supp}(\beta) \subseteq B_R(0)$.}
\label{fig:bump-function}
\end{center}

\begin{theorem}[Existence of partitions of unity]\label{thm:partition-unity}
\index{partition of unity!existence}
Let $M$ be a smooth manifold (possibly with boundary) and let
$\{U_\alpha\}_{\alpha \in A}$ be an open cover of $M$.  Then there exists a
smooth partition of unity subordinate to $\{U_\alpha\}$.
\end{theorem}

\begin{proof}
The proof proceeds in several steps.

\medskip
\noindent\textbf{Step 1: Exhaustion by compact sets.}
Since $M$ is second-countable and locally compact Hausdorff, there exists an
exhaustion $K_1 \subseteq K_2 \subseteq \cdots$ by compact sets with
$K_j \subseteq \operatorname{Int}(K_{j+1})$ and $M = \bigcup_{j} K_j$.
Set $K_0 = \varnothing$.

\medskip
\noindent\textbf{Step 2: Locally finite refinement.}
For each $p \in K_{j+1} \setminus \operatorname{Int}(K_j)$, choose
$\alpha(p)$ such that $p \in U_{\alpha(p)}$, and a chart
$(V_p, \varphi_p)$ centred at $p$ with
\[
  V_p \subseteq U_{\alpha(p)} \cap
  \bigl(\operatorname{Int}(K_{j+2}) \setminus K_{j-1}\bigr).
\]
Choose $r_p > 0$ so that $\varphi_p^{-1}(\overline{B}_{r_p}(0)) \subseteq V_p$.
By compactness of $K_{j+1} \setminus \operatorname{Int}(K_j)$, finitely
many of these balls cover this annulus.  Combining over all $j$ gives a
countable, locally finite refinement $\{W_i\}_{i \in \N}$ of the original
cover, together with smooth bump functions $\psi_i \colon M \to [0,1]$ with
$\operatorname{supp}(\psi_i) \subseteq W_i$ and $\psi_i > 0$ on a
neighbourhood of the corresponding compact piece.

\medskip
\noindent\textbf{Step 3: Construction of the partition.}
For each $i$, associate $W_i$ to some $U_{\alpha(i)}$ with
$W_i \subseteq U_{\alpha(i)}$.  For each $\alpha \in A$ define
\[
  \tilde{\rho}_\alpha = \sum_{\{i : \alpha(i)=\alpha\}} \psi_i.
\]
Since $\{\operatorname{supp}(\psi_i)\}$ is locally finite, each
$\tilde{\rho}_\alpha$ is smooth with
$\operatorname{supp}(\tilde{\rho}_\alpha) \subseteq U_\alpha$.
The sum $\Psi = \sum_\alpha \tilde{\rho}_\alpha = \sum_i \psi_i$ is
everywhere positive (the bumps $\psi_i$ cover $M$).  Setting
\[
  \rho_\alpha = \frac{\tilde{\rho}_\alpha}{\Psi}
\]
yields the desired partition of unity.
\end{proof}

\begin{corollary}\label{cor:smooth-urysohn}
\index{Urysohn-type lemma}
Let $A \subseteq M$ be closed and $U \supseteq A$ open in a smooth manifold
$M$.  Then there exists $f \in C^\infty(M)$ with $0 \leq f \leq 1$,
$f|_A = 1$, and $\operatorname{supp}(f) \subseteq U$.
\end{corollary}

\begin{proof}
Apply Theorem~\ref{thm:partition-unity} to the cover $\{U, M \setminus A\}$
and take $f = \rho_U$.
\end{proof}

\begin{corollary}\label{cor:extension-lemma}
\index{extension lemma}
Let $A \subseteq M$ be closed and $f \colon A \to \R$ a function that
extends smoothly to a neighbourhood of each point in $A$.  Then $f$ extends
to a smooth function $\tilde{f} \colon M \to \R$.
\end{corollary}

\begin{proof}
Cover $A$ by open sets $\{V_i\}$ on which smooth extensions $f_i$ exist.
Let $\{U_i\}$ be the cover $\{V_i\} \cup \{M \setminus A\}$ and
$\{\rho_i\}$ a subordinate partition of unity.  Define
$\tilde{f} = \sum_{i} \rho_i f_i$, where we set $f_i = 0$ on
$M \setminus A$.
\end{proof}

% ============================================================
\section{Exercises}
\label{sec:exercises-ch1}
% ============================================================

\begin{exercise}\label{exo:circle-atlas}
Construct a smooth atlas on $S^1$ using four charts (open arcs) and verify
that the transition maps are smooth.
\end{exercise}

\begin{exercise}\label{exo:RP1-circle}
\index{projective space!$\RP^1$}
Show that $\RP^1$ is diffeomorphic to $S^1$.
\textit{Hint:} exhibit an explicit smooth bijection with smooth inverse,
or use the map $[x_0:x_1] \mapsto e^{2i\arctan(x_1/x_0)}$.
\end{exercise}

\begin{exercise}\label{exo:CP1-S2}
\index{projective space!$\CP^1$}
Prove that $\CP^1$ is diffeomorphic to $S^2$.
\end{exercise}

\begin{exercise}\label{exo:product-manifold}
\index{product manifold}
Let $M^m$ and $N^n$ be smooth manifolds.  Show that $M \times N$ has a
natural smooth structure of dimension $m + n$, with charts
$(U_\alpha \times V_\beta,\; \varphi_\alpha \times \psi_\beta)$.
\end{exercise}

\begin{exercise}\label{exo:torus-quotient}
Verify explicitly that the atlas on $T^2 = \R^2/\Z^2$ obtained from local
sections of the projection $\pi \colon \R^2 \to \R^2/\Z^2$ is a smooth
atlas, and that $T^2$ is diffeomorphic to $S^1 \times S^1$.
\end{exercise}

\begin{exercise}\label{exo:GL-two-components}
\index{general linear group!components}
Show that $\GL(n,\R)$ has exactly two connected components:
$\GL^+(n,\R) = \{A : \det A > 0\}$ and $\GL^-(n,\R) = \{A : \det A < 0\}$.
\end{exercise}

\begin{exercise}\label{exo:SU2-S3}
\index{Lie group!$\SU(2)$}
Prove that $\SU(2)$ is diffeomorphic to $S^3$ by showing that every element
of $\SU(2)$ can be written as
$\bigl(\begin{smallmatrix} \alpha & -\bar\beta \\ \beta & \bar\alpha
\end{smallmatrix}\bigr)$ with $\abs{\alpha}^2 + \abs{\beta}^2 = 1$.
\end{exercise}

\begin{exercise}\label{exo:grassmann-compact}
\index{Grassmannian!compactness}
Prove that $\Gr(k,n)$ is compact.
\textit{Hint:} exhibit it as a quotient of a compact group, e.g.\
$\Gr(k,n) \cong \Orth(n)/(\Orth(k) \times \Orth(n-k))$.
\end{exercise}

\begin{exercise}[Partition of unity application]\label{exo:pou-metric}
\index{partition of unity!application}
Use a partition of unity to show that every smooth manifold admits a
Riemannian metric.  (You need only show the existence of a smooth,
positive-definite symmetric $(0,2)$-tensor field.)
\end{exercise}

\begin{exercise}\label{exo:boundary-product}
Let $M$ be a smooth manifold with boundary and $N$ a smooth manifold without
boundary.  Show that $M \times N$ is a smooth manifold with boundary
$\partial(M \times N) = (\partial M) \times N$.
\end{exercise}


% ============================================================
% ============================================================
%
%         CHAPTER 2  —  SMOOTH MAPS AND DIFFEOMORPHISMS
%
% ============================================================
% ============================================================

\chapter{Smooth Maps and Diffeomorphisms}
\label{ch:smooth-maps}
\minitoc
\bigskip

\noindent
Having defined smooth manifolds, we now study the maps between them.  The
central idea is that smoothness of a map $f \colon M \to N$ can be tested
in local coordinates; the compatibility of the smooth structures ensures
that this notion is independent of the charts chosen.  We develop the
fundamental analytic tools --- the inverse and implicit function theorems
on manifolds, the constant rank theorem --- and introduce the key notions
of immersion, submersion, and diffeomorphism.

% ============================================================
\section{Smooth Maps Between Manifolds}
\label{sec:smooth-maps}
% ============================================================

\begin{definition}[Smooth map]\label{def:smooth-map}
\index{smooth map}\index{map!smooth}
Let $M^m$ and $N^n$ be smooth manifolds.  A continuous map
$f \colon M \to N$ is \emph{smooth} (or $C^\infty$) if for every
$p \in M$ there exist charts $(U,\varphi)$ on $M$ around $p$ and
$(V,\psi)$ on $N$ around $f(p)$ with $f(U) \subseteq V$ such that the
\emph{coordinate representation}
\[
  \hat{f} = \psi \circ f \circ \varphi^{-1}
  \colon \varphi(U) \longrightarrow \psi(V)
\]
is smooth as a map between open subsets of Euclidean spaces.

We write $C^\infty(M,N)$\index{$C^\infty(M,N)$} for the set of all smooth
maps from $M$ to $N$, and $C^\infty(M) = C^\infty(M,\R)$.
\end{definition}

% ── TikZ: coordinate representation ──────────────────────

\begin{center}
\begin{tikzpicture}[scale=0.95, >={Stealth[length=5pt]},
  every node/.style={font=\small}]
  % M
  \draw[thick, fill=blue!5, rounded corners=15pt]
    (-3,1.5) -- (-0.5,2) -- (0.8,0.8) -- (0,-0.5) -- (-2.5,-0.3) -- cycle;
  \node at (-1.2,2.1) {$M$};
  \draw[blue!50, fill=blue!12] (-1.3,0.7) ellipse (0.9 and 0.6);
  \node[blue!70!black] at (-1.3,0.1) {$U$};
  \fill (-1.3,0.7) circle (1.2pt) node[above right] {$p$};
  % N
  \draw[thick, fill=red!5, rounded corners=15pt]
    (4,1.5) -- (7,2) -- (8,0.8) -- (7,-0.5) -- (4.5,-0.3) -- cycle;
  \node at (5.5,2.1) {$N$};
  \draw[red!50, fill=red!12] (6,0.7) ellipse (0.9 and 0.6);
  \node[red!70!black] at (6,0.1) {$V$};
  \fill (6,0.7) circle (1.2pt) node[above right] {$f(p)$};
  % f arrow
  \draw[->, thick] (0.5,1.0) -- (4.2,1.0) node[midway, above] {$f$};
  % phi arrow
  \draw[->, thick, blue!60!black] (-2.0,-0.4) -- (-3.5,-2.2)
    node[midway, left] {$\varphi$};
  % psi arrow
  \draw[->, thick, red!60!black] (6.8,-0.2) -- (7.5,-2.2)
    node[midway, right] {$\psi$};
  % R^m
  \draw[thick, blue!30, fill=blue!5] (-5.0,-2.5) rectangle (-1.5,-4.5);
  \node[blue!60!black] at (-3.2,-2.3) {$\R^m$};
  % R^n
  \draw[thick, red!30, fill=red!5] (5.5,-2.5) rectangle (9.0,-4.5);
  \node[red!60!black] at (7.3,-2.3) {$\R^n$};
  % hat f arrow
  \draw[->, thick, purple!70!black]
    (-2.0,-3.5) -- (6.0,-3.5)
    node[midway, below=3pt]
    {$\hat{f} = \psi \circ f \circ \varphi^{-1}$};
\end{tikzpicture}
\captionof{figure}{The coordinate representation of a smooth map $f$.}
\label{fig:coord-representation}
\end{center}

\begin{remark}\label{rem:chart-independence}
The definition is independent of the choice of charts: if $(U',\varphi')$
and $(V',\psi')$ are other charts, then
$\psi' \circ f \circ (\varphi')^{-1}
= (\psi' \circ \psi^{-1}) \circ \hat{f} \circ (\varphi \circ (\varphi')^{-1})$
is a composition of smooth maps, hence smooth.
\end{remark}

\begin{proposition}[Composition of smooth maps]\label{prop:composition-smooth}
\index{smooth map!composition}
If $f \colon M \to N$ and $g \colon N \to P$ are smooth, then
$g \circ f \colon M \to P$ is smooth.
\end{proposition}

\begin{proof}
In local coordinates, $\widehat{g \circ f} = \hat{g} \circ \hat{f}$,
which is a composition of smooth maps between Euclidean open sets.
\end{proof}

\begin{definition}[The smooth category]\label{def:smooth-category}
\index{category of smooth manifolds}
Smooth manifolds and smooth maps form a category, denoted $\mathbf{Man}$ or
$\mathbf{Diff}$.  The identity morphism on $M$ is $\Id_M$, and composition
is associative by the previous proposition.
\end{definition}

The following commutative diagram illustrates the ``functorial'' nature of
the coordinate representation:

\begin{center}
\begin{tikzcd}[column sep=large, row sep=large]
  U \arrow[r, "f"] \arrow[d, "\varphi"'] & V \arrow[r, "g"] \arrow[d, "\psi"] & W \arrow[d, "\chi"] \\
  \varphi(U) \arrow[r, "\hat{f}"'] & \psi(V) \arrow[r, "\hat{g}"'] & \chi(W)
\end{tikzcd}
\captionof{figure}{Composition in local coordinates.}
\label{fig:composition-diagram}
\end{center}

% ============================================================
\section{Diffeomorphisms}
\label{sec:diffeomorphisms}
% ============================================================

\begin{definition}[Diffeomorphism]\label{def:diffeomorphism}
\index{diffeomorphism}
A smooth map $f \colon M \to N$ is a \emph{diffeomorphism} if it is
bijective and its inverse $f^{-1} \colon N \to M$ is also smooth.  We say
$M$ and $N$ are \emph{diffeomorphic}\index{diffeomorphic} and write
$M \cong N$ (or $M \approx_{\mathrm{diff}} N$ when we wish to distinguish
from homeomorphism).
\end{definition}

\begin{remark}\label{rem:diffeo-vs-homeo}
A smooth bijection need not be a diffeomorphism: $f \colon \R \to \R$,
$f(t) = t^3$ is a smooth bijection but $f^{-1}(s) = s^{1/3}$ is not smooth
at $s = 0$.
\end{remark}

\begin{proposition}\label{prop:diffeo-equivalence}
Diffeomorphism is an equivalence relation on the class of smooth manifolds.
\end{proposition}

\begin{proof}
Reflexivity: $\Id_M$ is a diffeomorphism.  Symmetry: if $f$ is a
diffeomorphism then so is $f^{-1}$.  Transitivity: if $f \colon M \to N$
and $g \colon N \to P$ are diffeomorphisms, then $g \circ f$ is smooth
with smooth inverse $f^{-1} \circ g^{-1}$.
\end{proof}

\begin{definition}[Diffeomorphism group]\label{def:Diff-group}
\index{diffeomorphism group}
The group of all diffeomorphisms $f \colon M \to M$ under composition is
denoted $\Diff(M)$.
\end{definition}

% ============================================================
\section{The Inverse Function Theorem on Manifolds}
\label{sec:inverse-function-thm}
% ============================================================

We first recall the differential of a smooth map, which we shall define
properly through the tangent space in a later chapter.  For the present
purposes, the following coordinate formulation suffices.

\begin{definition}[Differential in coordinates]\label{def:differential-coords}
\index{differential}
Let $f \colon M^m \to N^n$ be smooth and $p \in M$.  Choose charts
$(U,\varphi)$ around $p$ and $(V,\psi)$ around $f(p)$.  The
\emph{differential of $f$ at $p$} (in these coordinates) is the linear map
\[
  \dd f_p \colon \R^m \longrightarrow \R^n, \qquad
  \dd f_p = D\hat{f}(\varphi(p)),
\]
where $D\hat{f}$ denotes the usual (Euclidean) derivative of the coordinate
representation $\hat{f} = \psi \circ f \circ \varphi^{-1}$.  The rank of
$\dd f_p$ is independent of the charts chosen.
\end{definition}

\begin{theorem}[Inverse function theorem --- manifold version]
\label{thm:inverse-function}
\index{inverse function theorem}
Let $f \colon M^n \to N^n$ be a smooth map between manifolds of the same
dimension.  If $\dd f_p \colon \R^n \to \R^n$ is an isomorphism at some
point $p \in M$, then there exist open neighbourhoods $U$ of $p$ and $V$ of
$f(p)$ such that $f|_U \colon U \to V$ is a diffeomorphism.
\end{theorem}

\begin{proof}
Choose charts $(U_0, \varphi)$ centred at $p$ and $(V_0, \psi)$ centred at
$f(p)$.  The coordinate representation
$\hat{f} = \psi \circ f \circ \varphi^{-1}$ satisfies
$D\hat{f}(0) = \dd f_p$, which is invertible by hypothesis.  By the
classical inverse function theorem\index{inverse function theorem!classical}
in $\R^n$, there exist open sets $\hat{U} \ni 0$ and $\hat{V} \ni 0$ such
that $\hat{f}|_{\hat{U}} \colon \hat{U} \to \hat{V}$ is a smooth
diffeomorphism.  Setting $U = \varphi^{-1}(\hat{U})$ and
$V = \psi^{-1}(\hat{V})$, we get
\[
  f|_U = \psi^{-1} \circ \hat{f}|_{\hat{U}} \circ \varphi,
\]
which is a diffeomorphism from $U$ to $V$ as a composition of
diffeomorphisms.
\end{proof}

\begin{corollary}\label{cor:local-diffeo-open}
\index{open map}
A local diffeomorphism (i.e.\ a smooth map whose differential is everywhere
an isomorphism) is an open map.
\end{corollary}

% ============================================================
\section{The Implicit Function Theorem on Manifolds}
\label{sec:implicit-function-thm}
% ============================================================

\begin{theorem}[Implicit function theorem --- manifold version]
\label{thm:implicit-function}
\index{implicit function theorem}
Let $f \colon M^m \to N^n$ be a smooth map with $m \geq n$.  Suppose
$q \in N$ is a \emph{regular value}\index{regular value} of $f$, meaning
that $\dd f_p$ is surjective for every $p \in f^{-1}(q)$.  Then
$f^{-1}(q)$ is a smooth submanifold of $M$ of dimension $m - n$.
\end{theorem}

\begin{proof}
Let $p \in f^{-1}(q)$.  Choose charts $(U, \varphi)$ centred at $p$ and
$(V, \psi)$ centred at $q$ with $f(U) \subseteq V$.  The coordinate
representation $\hat{f} = \psi \circ f \circ \varphi^{-1}$ satisfies
$D\hat{f}(0) \colon \R^m \to \R^n$ surjective, hence of rank $n$.

By the classical implicit function theorem, after permuting coordinates if
necessary, the set $\hat{f}^{-1}(0)$ is locally the graph of a smooth
function $g \colon W \subseteq \R^{m-n} \to \R^n$.  More precisely, there
is an open neighbourhood $\hat{U}$ of $0$ in $\R^m$ such that
\[
  \hat{f}^{-1}(0) \cap \hat{U}
  = \{(x', g(x')) : x' \in W\}.
\]
The map $x' \mapsto \varphi^{-1}(x', g(x'))$ gives a chart for
$f^{-1}(q)$ near $p$.  Since transition maps between such charts are smooth,
$f^{-1}(q)$ is a smooth submanifold of dimension $m - n$.
\end{proof}

\begin{corollary}[Preimage theorem]\label{cor:preimage}
\index{preimage theorem}
If $f \colon M \to N$ is a smooth map and $q$ is a regular value, then
$f^{-1}(q)$ is either empty or a closed smooth submanifold of $M$ of
codimension $n = \dim N$.
\end{corollary}

\begin{example}[The sphere revisited]\label{ex:sphere-preimage}
\index{sphere!as preimage}
Consider $f \colon \R^{n+1} \to \R$, $f(x) = \norm{x}^2$.  Then
$Df(x) = 2x$, which is surjective for $x \neq 0$.  Hence every $c > 0$ is a
regular value, and $f^{-1}(1) = S^n$ is a smooth $n$-dimensional
submanifold of $\R^{n+1}$.
\end{example}

\begin{example}[$\Orth(n)$ as a smooth manifold]\label{ex:On-preimage}
\index{orthogonal group!as preimage}
Define $F \colon M(n,\R) \to \mathrm{Sym}(n)$ by $F(A) = A^T A$.  Then
$\Orth(n) = F^{-1}(I_n)$.  The derivative is
$DF(A) \cdot H = H^T A + A^T H$.  At $A \in \Orth(n)$ this becomes
$DF(A) \cdot H = H^T A + A^T H$, and for any symmetric $S$ the matrix
$H = \tfrac{1}{2}AS$ satisfies $DF(A) \cdot H = S$.  Hence $I_n$ is a
regular value, confirming that $\Orth(n)$ is a smooth submanifold of
dimension $n^2 - \tfrac{n(n+1)}{2} = \tfrac{n(n-1)}{2}$.
\end{example}

% ============================================================
\section{Immersions and Submersions}
\label{sec:immersions-submersions}
% ============================================================

\begin{definition}[Immersion, submersion]\label{def:immersion-submersion}
\index{immersion}\index{submersion}
Let $f \colon M^m \to N^n$ be smooth.
\begin{enumerate}[label=(\roman*)]
  \item $f$ is an \emph{immersion} at $p$ if $\dd f_p$ is injective
        (i.e.\ $\rank \dd f_p = m$).
  \item $f$ is a \emph{submersion} at $p$ if $\dd f_p$ is surjective
        (i.e.\ $\rank \dd f_p = n$).
  \item $f$ is an \emph{immersion} (resp.\ \emph{submersion}) if it is one
        at every point.
\end{enumerate}
\end{definition}

\begin{definition}[Local diffeomorphism]\label{def:local-diffeomorphism}
\index{local diffeomorphism}
A smooth map $f \colon M^n \to N^n$ is a \emph{local diffeomorphism} if
$\dd f_p$ is an isomorphism for all $p \in M$ (equivalently, if $f$ is both
an immersion and a submersion).
\end{definition}

\begin{proposition}[Canonical form for immersions]\label{prop:immersion-canonical}
\index{immersion!canonical form}
If $f \colon M^m \to N^n$ is an immersion at $p$ (so $m \leq n$), then
there exist charts $(U,\varphi)$ centred at $p$ and $(V,\psi)$ centred at
$f(p)$ such that the coordinate representation takes the form
\[
  \hat{f}(x_1, \ldots, x_m) = (x_1, \ldots, x_m, 0, \ldots, 0).
\]
\end{proposition}

\begin{proof}
In coordinates, $D\hat{f}(0)$ has rank $m$.  After reordering the
coordinates of $\R^n$ we may assume the first $m$ rows of the $n \times m$
Jacobian form an invertible $m \times m$ matrix.  Define
$F \colon \varphi(U) \times \R^{n-m} \to \R^n$ by
\[
  F(x, y) = \hat{f}(x) + (0, y).
\]
Then $DF(0,0)$ is invertible.  By the inverse function theorem, $F$ is a
local diffeomorphism near $(0,0)$.  Setting $\psi' = F^{-1} \circ \psi$
(restricted appropriately), we get $\psi' \circ f \circ \varphi^{-1}(x) =
F^{-1}(\hat{f}(x)) = F^{-1}(F(x,0)) = (x, 0)$.
\end{proof}

\begin{proposition}[Canonical form for submersions]\label{prop:submersion-canonical}
\index{submersion!canonical form}
If $f \colon M^m \to N^n$ is a submersion at $p$ (so $m \geq n$), then
there exist charts $(U,\varphi)$ centred at $p$ and $(V,\psi)$ centred at
$f(p)$ such that
\[
  \hat{f}(x_1, \ldots, x_m) = (x_1, \ldots, x_n).
\]
\end{proposition}

\begin{proof}
The Jacobian $D\hat{f}(0)$ has rank $n$.  After reordering, assume the first
$n$ columns are linearly independent.  Define
$G \colon \varphi(U) \to \R^m$ by $G(x) = (\hat{f}(x), x_{n+1},\ldots,x_m)$.
Then $DG(0)$ is invertible, so $G$ is a local diffeomorphism.  Setting
$\varphi' = G \circ \varphi$, the coordinate representation becomes
$\psi \circ f \circ (\varphi')^{-1}(y_1,\ldots,y_m) = (y_1,\ldots,y_n)$.
\end{proof}

% ── TikZ: immersion and submersion ────────────────────────

\begin{center}
\begin{tikzpicture}[scale=0.85, >={Stealth[length=5pt]},
  every node/.style={font=\small}]
  % Immersion picture
  \begin{scope}[shift={(-5,0)}]
    \node[above] at (0,2.5) {\textbf{Immersion} ($m < n$)};
    % Domain (line)
    \draw[thick, blue!70] (-1.5,0) -- (1.5,0);
    \node[blue!70!black, below] at (0,-0.3) {$M^m$};
    % Arrow
    \draw[->, thick] (2.0,0) -- (3.0,0) node[midway, above] {$f$};
    % Codomain with embedded curve
    \draw[thick, fill=red!5, rounded corners=10pt]
      (3.5,-1.5) rectangle (7.5,1.5);
    \node[red!70!black] at (7.0,1.2) {$N^n$};
    \draw[thick, blue!70, decorate,
      decoration={snake, amplitude=3pt, segment length=12pt}]
      (4.0,0) -- (7.0,0);
  \end{scope}
  % Submersion picture
  \begin{scope}[shift={(5,0)}]
    \node[above] at (0,2.5) {\textbf{Submersion} ($m > n$)};
    % Domain (rectangle)
    \draw[thick, fill=blue!5, rounded corners=8pt]
      (-2,-1.5) rectangle (2,1.5);
    \node[blue!70!black] at (1.5,1.2) {$M^m$};
    % Vertical fibers
    \foreach \x in {-1.2,-0.4,0.4,1.2}
      \draw[gray!60, thick] (\x,-1.2) -- (\x,1.2);
    % Arrow
    \draw[->, thick] (2.5,0) -- (3.5,0) node[midway, above] {$f$};
    % Codomain (line)
    \draw[thick, red!70] (4.0,0) -- (7.0,0);
    \node[red!70!black, below] at (5.5,-0.3) {$N^n$};
  \end{scope}
\end{tikzpicture}
\captionof{figure}{Schematic depiction of an immersion (left) and a
submersion (right).  The fibres of the submersion are shown as vertical
lines.}
\label{fig:immersion-submersion}
\end{center}

% ============================================================
\section{The Constant Rank Theorem}
\label{sec:constant-rank}
% ============================================================

The canonical forms for immersions and submersions are special cases of the
following fundamental result.

\begin{theorem}[Constant rank theorem]\label{thm:constant-rank}
\index{constant rank theorem}
Let $f \colon M^m \to N^n$ be a smooth map of constant rank $r$ on an open
set containing $p$.  Then there exist charts $(U,\varphi)$ centred at $p$
and $(V,\psi)$ centred at $f(p)$ in which $f$ has the coordinate
representation
\[
  \hat{f}(x_1, \ldots, x_m) = (x_1, \ldots, x_r, 0, \ldots, 0).
\]
\end{theorem}

\begin{proof}
By the Euclidean constant rank theorem (see, e.g., Boothby), working in any
pair of charts we may find smooth coordinate changes
$\alpha$ on the domain side and $\beta$ on the target side such that
$\beta \circ \hat{f} \circ \alpha^{-1}$ has the stated form.  Replacing
$\varphi$ by $\alpha \circ \varphi$ and $\psi$ by $\beta \circ \psi$ gives
charts with the desired property.

In more detail: after a linear change of coordinates we may assume
$D\hat{f}(0)$ has its upper-left $r \times r$ block invertible.  Write
$\hat{f}(x) = (\hat{f}_1(x), \hat{f}_2(x))$ with
$\hat{f}_1 \colon \R^m \to \R^r$ and $\hat{f}_2 \colon \R^m \to \R^{n-r}$.
Define $\alpha(x) = (\hat{f}_1(x), x_{r+1}, \ldots, x_m)$; then
$D\alpha(0)$ is invertible.  In the new coordinates $y = \alpha(x)$,
the map becomes $\hat{f} \circ \alpha^{-1}(y) = (y_1, \ldots, y_r,
h(y))$ for some smooth $h$.  The constant rank condition forces
$\partial h / \partial y_j = 0$ for $j = r+1, \ldots, m$ (otherwise the
rank would exceed $r$).  Hence $h$ depends only on $(y_1, \ldots, y_r)$.
Define $\beta \colon \R^n \to \R^n$ by
$\beta(z_1,\ldots,z_n) = (z_1,\ldots,z_r, z_{r+1} - h(z_1,\ldots,z_r),
\ldots, z_n - h_{n-r}(z_1,\ldots,z_r))$.  Then
$\beta \circ \hat{f} \circ \alpha^{-1}(y) = (y_1,\ldots,y_r,0,\ldots,0)$.
\end{proof}

\begin{corollary}\label{cor:constant-rank-level-set}
If $f \colon M \to N$ has constant rank $r$ on all of $M$, then for every
$q \in f(M)$ the level set $f^{-1}(q)$ is a closed smooth submanifold of
$M$ of dimension $m - r$.
\end{corollary}

\begin{proof}
In the canonical coordinates of Theorem~\ref{thm:constant-rank},
$f^{-1}(q)$ locally coincides with $\{(0,\ldots,0,x_{r+1},\ldots,x_m)\}$,
a smooth submanifold of dimension $m - r$.
\end{proof}

% ============================================================
\section{Examples}
\label{sec:examples-ch2}
% ============================================================

\begin{example}[Projection maps]\label{ex:projections}
\index{projection}
Let $M \times N$ be a product manifold.  The projections
$\pi_M \colon M \times N \to M$ and $\pi_N \colon M \times N \to N$ are
smooth submersions.  In product charts
$(\varphi \times \psi)$, $\hat{\pi}_M(x,y) = x$.
\end{example}

\begin{example}[Inclusion of submanifolds]\label{ex:inclusion}
\index{inclusion map}
If $S \subseteq M$ is a smooth submanifold, the inclusion
$\iota \colon S \hookrightarrow M$ is a smooth immersion (indeed an
embedding).
\end{example}

\begin{example}[The Hopf map]\label{ex:hopf}
\index{Hopf map}
The \emph{Hopf map} $h \colon S^3 \to S^2$ is defined as follows.
View $S^3 \subseteq \C^2$ as $\{(z_1,z_2) : \abs{z_1}^2 + \abs{z_2}^2 = 1\}$
and $S^2 \cong \CP^1$.  Define
\[
  h(z_1, z_2) = [z_1 : z_2] \in \CP^1 \cong S^2.
\]
Explicitly, using the identification $S^2 \subseteq \R^3$,
\[
  h(z_1,z_2) = \bigl(2\operatorname{Re}(z_1 \bar{z}_2),\;
  2\operatorname{Im}(z_1 \bar{z}_2),\;
  \abs{z_1}^2 - \abs{z_2}^2\bigr).
\]
The Hopf map is a smooth submersion.  Each fibre $h^{-1}(p)$ is a great
circle $S^1 \subseteq S^3$, and $h$ exhibits $S^3$ as an $S^1$-bundle over
$S^2$.
\end{example}

% ── TikZ: Hopf fibration schematic ──────────────────────

\begin{center}
\begin{tikzpicture}[scale=1.0, >={Stealth[length=5pt]}]
  % S^3 (schematic)
  \draw[thick, fill=blue!5] (0,2) ellipse (2 and 1);
  \node at (0,3.3) {$S^3$};
  % fibres drawn as small circles
  \draw[blue!60, thick] (-1.0,2) ellipse (0.25 and 0.45);
  \draw[blue!60, thick] (0,2) ellipse (0.25 and 0.45);
  \draw[blue!60, thick] (1.0,2) ellipse (0.25 and 0.45);
  \node[blue!70!black, right] at (1.5,2) {$S^1$ fibres};
  % arrow
  \draw[->, thick] (0,0.7) -- (0,-0.3) node[midway, right] {$h$};
  % S^2
  \draw[thick, fill=red!5] (0,-1.5) circle (1.0);
  \node at (0,-2.8) {$S^2$};
  \fill[red!50] (-0.5,-1.5) circle (1.5pt);
  \fill[red!50] (0,-1.5) circle (1.5pt);
  \fill[red!50] (0.5,-1.5) circle (1.5pt);
\end{tikzpicture}
\captionof{figure}{Schematic of the Hopf fibration $h \colon S^3 \to S^2$
with $S^1$ fibres.}
\label{fig:hopf}
\end{center}

\begin{example}[The determinant map]\label{ex:det-submersion}
\index{determinant}
The map $\det \colon \GL(n,\R) \to \R^*$ is a smooth submersion.  Its
derivative at $A$ in the direction $H$ is
\[
  D(\det)(A) \cdot H = \det(A)\, \operatorname{tr}(A^{-1}H).
\]
For $A \in \GL(n,\R)$ and any $c \in \R^*$, choosing $H = \frac{c}{\det(A)
\cdot n} A$ shows surjectivity.  The level set $\det^{-1}(1) = \SL(n,\R)$.
\end{example}

\begin{example}[Smooth maps between spheres]\label{ex:sphere-map}
\index{sphere!maps between}
The antipodal map $a \colon S^n \to S^n$, $a(x) = -x$, is a
diffeomorphism.  For $n$ even, $a$ is orientation-reversing; for $n$ odd, it
is orientation-preserving (as we shall see when we discuss orientations).
\end{example}

\begin{example}[Covering maps as local diffeomorphisms]\label{ex:covering}
\index{covering map}\index{local diffeomorphism!example}
The exponential map $e \colon \R \to S^1$, $e(t) = e^{2\pi i t}$, is a
smooth surjective local diffeomorphism (and a covering map).  More generally,
the quotient map $\R^n \to T^n = \R^n/\Z^n$ is a smooth local
diffeomorphism.
\end{example}

\begin{example}[The Gauss map]\label{ex:gauss-map}
\index{Gauss map}
Let $S \subseteq \R^3$ be a smooth surface (a $2$-dimensional submanifold).
If $S$ is orientable, there exists a smooth unit normal field
$\nu \colon S \to S^2$, called the \emph{Gauss map}.  For instance, for
the sphere $S^2 = \{x \in \R^3 : \norm{x}=1\}$ the outward normal is
$\nu(x) = x$, so the Gauss map is the identity.  For the torus
$T^2 \subseteq \R^3$ parametrized by
\[
  \Phi(\theta, \phi) = \bigl((R + r\cos\phi)\cos\theta,\;
  (R + r\cos\phi)\sin\theta,\; r\sin\phi\bigr),
\]
the Gauss map sends each point to its outward unit normal, and is a smooth
map $T^2 \to S^2$ of degree zero.
\end{example}

\begin{example}[Power maps on $S^1$]\label{ex:power-maps}
\index{power map}
For each integer $k \geq 1$, the map $p_k \colon S^1 \to S^1$ defined by
$p_k(z) = z^k$ (viewing $S^1 \subseteq \C$) is smooth.  It is a
$k$-sheeted covering map for $k \geq 2$ and a diffeomorphism for $k = 1$.
In real coordinates, writing $z = e^{i\theta}$, the map is
$\theta \mapsto k\theta$, and the differential is multiplication by $k$,
which is everywhere an isomorphism.  Hence $p_k$ is a local diffeomorphism
for all $k \geq 1$.
\end{example}

% ============================================================
\section{Local Diffeomorphisms and Global Properties}
\label{sec:local-diffeo}
% ============================================================

\begin{proposition}\label{prop:local-diffeo-injective}
\index{local diffeomorphism!injective implies diffeomorphism}
A smooth map $f \colon M \to N$ that is both a local diffeomorphism and a
bijection is a diffeomorphism.
\end{proposition}

\begin{proof}
Since $f$ is a local diffeomorphism, for each $q \in N$ there is a
neighbourhood $V_q$ of $q$ on which $f^{-1}|_{V_q}$ is smooth.  As $f$ is
bijective, these local inverses patch together to give the global inverse
$f^{-1}$, which is therefore smooth.
\end{proof}

\begin{proposition}\label{prop:local-diffeo-covering}
\index{local diffeomorphism!as covering map}
Let $f \colon M \to N$ be a local diffeomorphism with $M$ compact and $N$
connected.  Then $f$ is surjective and a (finite-sheeted) covering map.
In particular, if $N$ is simply connected, $f$ is a diffeomorphism.
\end{proposition}

\begin{proof}
Since $f$ is a local diffeomorphism, it is an open map
(Corollary~\ref{cor:local-diffeo-open}).  Since $M$ is compact, $f(M)$ is
compact, hence closed in $N$ (as $N$ is Hausdorff).  As $f(M)$ is both open
and closed in the connected space $N$, we get $f(M) = N$.

For each $q \in N$ the fibre $f^{-1}(q)$ is discrete (by the local
diffeomorphism condition) and compact (as a closed subset of $M$), hence
finite.  A standard argument using the local inverse charts shows $f$ is a
covering map.  If $N$ is simply connected, covering theory gives that $f$ is
a homeomorphism, hence a diffeomorphism by
Proposition~\ref{prop:local-diffeo-injective}.
\end{proof}

\begin{remark}\label{rem:local-not-global}
The converse fails: a local diffeomorphism from a non-compact manifold need
not be a covering map.  For instance, the inclusion of an open interval
$(0,1) \hookrightarrow \R$ is a local diffeomorphism that is not surjective.
A more interesting example: the restriction of $e \colon \R \to S^1$ to
$(0, 2)$ is a local diffeomorphism that is neither injective nor a covering.
\end{remark}

% ── TikZ: local diffeomorphism vs diffeomorphism ──────────

\begin{center}
\begin{tikzpicture}[scale=0.9, >={Stealth[length=5pt]},
  every node/.style={font=\small}]
  % Left: covering map
  \begin{scope}[shift={(-4,0)}]
    \node[above] at (0,2.8) {\textbf{Covering map}};
    % Helix (schematic)
    \draw[thick, blue!70, decorate,
      decoration={coil, segment length=8pt, amplitude=4pt}]
      (0,2.2) -- (0,-1.5);
    \node[blue!70!black, left] at (-0.5,0.3) {$\R$};
    \draw[->, thick] (0,-1.8) -- (0,-2.5) node[midway, right] {$e$};
    % Circle
    \draw[thick, red!70] (0,-3.5) circle (0.7);
    \node[red!70!black, right] at (0.9,-3.5) {$S^1$};
  \end{scope}
  % Right: diffeomorphism
  \begin{scope}[shift={(4,0)}]
    \node[above] at (0,2.8) {\textbf{Diffeomorphism}};
    % Domain
    \draw[thick, fill=blue!8, rounded corners=12pt]
      (-1.5,-0.5) rectangle (1.5,2);
    \node[blue!70!black] at (0,2.3) {$M$};
    \draw[->, thick] (0,-0.8) -- (0,-1.8) node[midway, right] {$f \cong$};
    % Codomain
    \draw[thick, fill=red!8, rounded corners=12pt]
      (-1.5,-2.2) rectangle (1.5,-4.2);
    \node[red!70!black] at (0,-4.5) {$N$};
  \end{scope}
\end{tikzpicture}
\captionof{figure}{A covering map (local diffeomorphism, left) versus a
global diffeomorphism (right).}
\label{fig:local-vs-global-diffeo}
\end{center}

% ============================================================
\section{Embeddings}
\label{sec:embeddings}
% ============================================================

\begin{definition}[Smooth embedding]\label{def:embedding}
\index{embedding}
A smooth map $f \colon M \to N$ is a \emph{smooth embedding} if it is an
injective immersion that is a homeomorphism onto its image $f(M)$ (with
the subspace topology from $N$).
\end{definition}

\begin{remark}\label{rem:embedding-vs-immersion}
Every embedding is an injective immersion, but the converse fails.  The
standard example is an injective immersion of $\R$ into $T^2$ whose image
is a dense subset of the torus (an irrational-slope line).  This map is an
injective immersion but not an embedding because the image, with the
subspace topology, is not homeomorphic to $\R$.
\end{remark}

\begin{proposition}[Compact domain implies embedding]\label{prop:compact-embedding}
\index{embedding!from compact domain}
If $M$ is compact and $f \colon M \to N$ is an injective immersion, then
$f$ is a smooth embedding.
\end{proposition}

\begin{proof}
An injective continuous map from a compact space to a Hausdorff space is
a homeomorphism onto its image (a closed subset).  Since $f$ is also an
immersion, it is a smooth embedding.
\end{proof}

\begin{proposition}[Image of an embedding is a submanifold]
\label{prop:embedding-image}
\index{embedding!image is submanifold}
If $f \colon M^m \to N^n$ is a smooth embedding, then $f(M)$ is a smooth
submanifold of $N$ of dimension $m$, and $f \colon M \to f(M)$ is a
diffeomorphism.
\end{proposition}

\begin{proof}
Since $f$ is an immersion, the canonical form
(Proposition~\ref{prop:immersion-canonical}) provides, for each $p \in M$,
charts in which $f$ looks like the standard inclusion
$\R^m \hookrightarrow \R^n$.  This shows that $f(M)$ is locally a ``slice''
of $N$, hence a submanifold.  Since $f$ is a homeomorphism onto $f(M)$, and
the smooth structures agree via the canonical form, $f$ is a diffeomorphism
onto this submanifold.
\end{proof}

% ── TikZ: embedding vs non-embedding ─────────────────────

\begin{center}
\begin{tikzpicture}[scale=0.85, >={Stealth[length=5pt]},
  every node/.style={font=\small}]
  % Left: proper embedding
  \begin{scope}[shift={(-4.5,0)}]
    \node[above] at (0,2.2) {\textbf{Embedding}};
    \draw[thick, fill=yellow!8, rounded corners=10pt]
      (-2,-1.5) rectangle (2,1.5);
    \node at (1.5,1.2) {$N$};
    \draw[very thick, blue!70]
      plot[smooth, tension=0.6] coordinates
      {(-1.2,-0.8) (-0.3,0.5) (0.6,-0.2) (1.3,0.7)};
    \node[blue!70!black, below] at (0,-1.8) {$f(M)$: submanifold};
  \end{scope}
  % Right: non-embedding (figure-eight)
  \begin{scope}[shift={(4.5,0)}]
    \node[above] at (0,2.2) {\textbf{Injective immersion,}};
    \node[above] at (0,1.7) {\textbf{not an embedding}};
    \draw[thick, fill=yellow!8, rounded corners=10pt]
      (-2,-1.5) rectangle (2,1.5);
    \node at (1.5,1.2) {$N$};
    % Figure eight
    \draw[very thick, red!70]
      (0,0) .. controls (1.2,1.2) and (1.2,-1.2) .. (0,0)
             .. controls (-1.2,1.2) and (-1.2,-1.2) .. (0,0);
    \fill[red!70] (0,0) circle (2pt);
    \node[red!70!black, below] at (0,-1.8) {Self-intersection point};
  \end{scope}
\end{tikzpicture}
\captionof{figure}{An embedding (left) versus a non-injective immersion
producing a figure-eight (right).}
\label{fig:embedding-vs-immersion}
\end{center}

\begin{example}[Standard embeddings]\label{ex:standard-embeddings}
\index{embedding!standard examples}
The following are smooth embeddings:
\begin{enumerate}[label=(\alph*)]
  \item The inclusion $\iota \colon S^n \hookrightarrow \R^{n+1}$.
  \item The map $f \colon \R \to \R^2$, $f(t) = (t, t^2)$ (a parabola).
  \item The Segre embedding\index{Segre embedding}
    $\sigma \colon \CP^m \times \CP^n \hookrightarrow \CP^{mn+m+n}$
    defined by
    $\sigma([x_0:\cdots:x_m], [y_0:\cdots:y_n]) = [x_i y_j]_{0 \leq i \leq m,\, 0 \leq j \leq n}$.
  \item The Veronese embedding\index{Veronese embedding}
    $v_d \colon \CP^n \hookrightarrow \CP^N$
    ($N = \binom{n+d}{d}-1$) sending $[x_0:\cdots:x_n]$ to the list of
    all monomials of degree $d$.
\end{enumerate}
\end{example}

% ============================================================
\section{Smooth Maps on Manifolds with Boundary}
\label{sec:smooth-maps-boundary}
% ============================================================

\begin{definition}\label{def:smooth-map-boundary}
\index{smooth map!on manifolds with boundary}
Let $M$ and $N$ be smooth manifolds, possibly with boundary.  A map
$f \colon M \to N$ is \emph{smooth} if for every $p \in M$ there exist
boundary-type charts $(U,\varphi)$ and $(V,\psi)$ with
$f(U) \subseteq V$ such that $\psi \circ f \circ \varphi^{-1}$ is smooth
in the sense that it extends to a smooth map on an open subset of $\R^m$.
\end{definition}

\begin{proposition}\label{prop:smooth-boundary-to-boundary}
If $f \colon M \to N$ is a diffeomorphism between manifolds with boundary,
then $f(\partial M) = \partial N$ and
$f|_{\partial M} \colon \partial M \to \partial N$ is a diffeomorphism.
\end{proposition}

\begin{proof}
If $p \in \partial M$, then every chart around $p$ sends it to
$\partial \mathbb{H}^m$.  If $f(p) \notin \partial N$, an interior chart
around $f(p)$ composed with $f$ and the inverse of a boundary chart would
give a diffeomorphism mapping a boundary point to an interior point of
$\mathbb{H}^n$, contradicting smooth invariance of domain.  The same
argument applied to $f^{-1}$ shows $f^{-1}(\partial N) = \partial M$.
The restriction $f|_{\partial M}$ is smooth with smooth inverse
$f^{-1}|_{\partial N}$.
\end{proof}

% ============================================================
\section{Exercises}
\label{sec:exercises-ch2}
% ============================================================

\begin{exercise}\label{exo:smooth-iff-components}
\index{smooth map!component criterion}
Let $f \colon M \to \R^n$ be a map.  Show that $f$ is smooth if and only
if each component function $f_i = \pi_i \circ f \colon M \to \R$ is smooth,
where $\pi_i$ is the $i$-th coordinate projection.
\end{exercise}

\begin{exercise}\label{exo:graph-submanifold}
\index{graph of a smooth map}
Let $f \colon M \to N$ be smooth.  Prove that its graph
$\Gamma_f = \{(p, f(p)) : p \in M\} \subseteq M \times N$ is a smooth
submanifold diffeomorphic to $M$.
\end{exercise}

\begin{exercise}\label{exo:smooth-bijection-not-diffeo}
Give a detailed proof that $f \colon \R \to \R$, $f(t) = t^3$, is a smooth
bijection that is \emph{not} a diffeomorphism.  Where does the inverse
function theorem fail to apply?
\end{exercise}

\begin{exercise}\label{exo:hopf-submersion}
\index{Hopf map!is a submersion}
Verify that the Hopf map $h \colon S^3 \to S^2$ is a smooth submersion by
computing its differential in suitable coordinates.
\end{exercise}

\begin{exercise}\label{exo:quotient-RPn}
\index{projective space!quotient map}
Show that the quotient map $\pi \colon S^n \to \RP^n$,
$\pi(x) = [x]$, is a smooth two-sheeted covering map and a local
diffeomorphism.
\end{exercise}

\begin{exercise}\label{exo:submersion-open}
\index{submersion!is an open map}
Prove that every submersion is an open map.
\textit{Hint:} use the canonical form for submersions
(Proposition~\ref{prop:submersion-canonical}).
\end{exercise}

\begin{exercise}\label{exo:immersion-injective}
\index{immersion!injective}
Give an example of an immersion that is not injective, and an example of
an injective immersion that is not an embedding (i.e.\ not a homeomorphism
onto its image).  \textit{Hint for the second:} consider the figure-eight
curve or a dense line on the torus.
\end{exercise}

\begin{exercise}\label{exo:constant-rank-composition}
Let $f \colon M \to N$ and $g \colon N \to P$ be smooth maps.
Show that $\rank \dd(g \circ f)_p \leq \min(\rank \dd f_p, \rank \dd g_{f(p)})$.
Give an example where strict inequality holds.
\end{exercise}

\begin{exercise}\label{exo:product-submersion}
\index{submersion!product}
Let $f \colon M \to N$ and $g \colon M' \to N'$ be submersions.  Prove that
$f \times g \colon M \times M' \to N \times N'$ is a submersion.  State
and prove the analogous result for immersions.
\end{exercise}

\begin{exercise}\label{exo:det-regular-values}
Determine all critical points and critical values of the determinant map
$\det \colon M(n,\R) \to \R$ (here $M(n,\R) \cong \R^{n^2}$ is the space
of all $n \times n$ real matrices, not just invertible ones).
\textit{Answer:} the only critical value is $0$; equivalently, $c \neq 0$
is always a regular value.
\end{exercise}

\begin{exercise}\label{exo:embedding-figure-eight}
\index{figure-eight curve}
Let $\gamma \colon (-\pi, \pi) \to \R^2$ be defined by
$\gamma(t) = (\sin 2t, \sin t)$.
\begin{enumerate}[label=(\alph*)]
  \item Show that $\gamma$ is an immersion.
  \item Show that $\gamma$ is injective.
  \item Prove that $\gamma$ is \emph{not} an embedding.
  \textit{Hint:} examine what happens as $t \to \pm \pi$.
\end{enumerate}
\end{exercise}

\begin{exercise}\label{exo:diffeomorphism-Rn}
\index{diffeomorphism!of $\R^n$}
\begin{enumerate}[label=(\alph*)]
  \item Show that any open ball $B_r(0) \subseteq \R^n$ is diffeomorphic
    to $\R^n$.
    \textit{Hint:} use the map $x \mapsto x / (r - \norm{x})$ or a similar
    radial rescaling.
  \item Conclude that every connected smooth $n$-manifold has an open
    subset diffeomorphic to $\R^n$.
\end{enumerate}
\end{exercise}

\begin{exercise}[Ehresmann's lemma, special case]\label{exo:ehresmann}
\index{Ehresmann's lemma}
Let $f \colon M \to \R$ be a smooth proper submersion (proper means
preimages of compact sets are compact).  Show that for any two regular
values $a < b$ of $f$, the preimages $f^{-1}(a)$ and $f^{-1}(b)$ are
diffeomorphic.
\textit{Hint:} this requires tools from Chapter~\ref{ch:smooth-manifolds}
(partitions of unity) and the flow of a vector field; give a proof sketch
assuming the existence of flows for compactly supported vector fields.
\end{exercise}

\begin{exercise}\label{exo:smooth-retraction}
\index{retraction}
Let $M$ be a smooth manifold and $A \subseteq M$ a smooth submanifold.  A
\emph{smooth retraction} onto $A$ is a smooth map $r \colon M \to A$ with
$r|_A = \Id_A$.
\begin{enumerate}[label=(\alph*)]
  \item Show that $r$ is a submersion at every point of $A$.
  \item Conclude that if $M$ is connected and $\dim A < \dim M$, then no
    smooth retraction $M \to A$ can exist when $M$ is compact without
    boundary.  \textit{Hint:} what is $r^{-1}(q)$ for a regular value $q$?
\end{enumerate}
\end{exercise}


% ============================================================
% === END OF CHUNK preamble+ch1-2 ===
% === START OF CHUNK ch3-4 ===
%%%%%%%%%%%%%%%%%%%%%%%%%%%%%%%%%%%%%%%%%%%%%%%%%%%%%%%%%%%%%
%% DIFFERENTIAL TOPOLOGY — A Graduate Course
%% Chapter 3: Tangent and Cotangent Bundles
%% Chapter 4: Submanifolds and the Regular Value Theorem
%%%%%%%%%%%%%%%%%%%%%%%%%%%%%%%%%%%%%%%%%%%%%%%%%%%%%%%%%%%%%

%%%%%%%%%%%%%%%%%%%%%%%%%%%%%%%%%%%%%%%%%%%%%%%%%%%%%%%%%%%%%
%% CHAPTER 3
%%%%%%%%%%%%%%%%%%%%%%%%%%%%%%%%%%%%%%%%%%%%%%%%%%%%%%%%%%%%%
\chapter{Tangent and Cotangent Bundles}\label{ch:tangent-cotangent}
\minitoc\bigskip

The notion of tangent vector is fundamental in differential topology.
On an open subset $U \subset \R^n$, tangent vectors are simply elements
of~$\R^n$ attached at a point.  On a smooth manifold, one must define
tangent vectors \emph{intrinsically}, without reference to an ambient
Euclidean space.  We present two classical approaches---via derivations
and via equivalence classes of curves---and prove their equivalence.
We then assemble the tangent spaces into the \emph{tangent bundle}, a
smooth manifold in its own right, and dualise to obtain the cotangent
bundle.  Vector fields, the Lie bracket, and flows round out the chapter.

%%%%%%%%%%%%%%%%%%%%%%%%%%%%%%%%%%%%%%%%%%%%%%%%%%%%%%%%%%%%%
\section{The tangent space}\label{sec:tangent-space}
%%%%%%%%%%%%%%%%%%%%%%%%%%%%%%%%%%%%%%%%%%%%%%%%%%%%%%%%%%%%%

Throughout this chapter, $M$ denotes a smooth manifold of dimension~$n$
and $p \in M$.

%% ── Derivations approach ──────────────────────────────────
\subsection{Derivations}\label{ssec:derivations}

\begin{definition}[Tangent vector as derivation]\label{def:derivation}%
\index{tangent vector!derivation}\index{derivation}
Let $C^\infty_p(M)$ denote the $\R$-algebra of germs\index{germ} of smooth
functions at~$p$.  A \textbf{derivation at~$p$} is an $\R$-linear map
$v \colon C^\infty_p(M) \to \R$ satisfying the \emph{Leibniz rule}:
\[
  v(fg) \;=\; f(p)\,v(g) \;+\; g(p)\,v(f)
  \qquad \forall\, f,g \in C^\infty_p(M).
\]
The set of all derivations at~$p$ is denoted $T_pM$\index{tangent space}
and called the \textbf{tangent space of $M$ at~$p$}.
\end{definition}

\begin{lemma}\label{lem:derivation-constant}
If $v \in T_pM$ and $c \in \R$ is a constant function, then $v(c) = 0$.
\end{lemma}

\begin{proof}
Write $c = c \cdot 1$.  Then $v(c) = v(c \cdot 1) = c\, v(1) + 1 \cdot v(c)
= c\, v(1) + v(c)$, so $c\, v(1) = 0$.  Applying this with $c = 1$ gives
$v(1) = 0$, hence $v(c) = c\, v(1) = 0$ for all $c$.
\end{proof}

\begin{proposition}\label{prop:tangent-space-basis}%
\index{tangent space!basis}
Let $(U, \varphi = (x^1, \dots, x^n))$ be a chart around~$p$.  Define, for
$i = 1, \dots, n$,
\[
  \left.\frac{\partial}{\partial x^i}\right|_p \colon C^\infty_p(M) \to \R,
  \qquad
  f \;\longmapsto\;
  \frac{\partial (f \circ \varphi^{-1})}{\partial r^i}\bigg|_{\varphi(p)},
\]
where $r^1, \dots, r^n$ are the standard coordinates on~$\R^n$.  Then
$\left\{\left.\dfrac{\partial}{\partial x^1}\right|_p, \dots,
\left.\dfrac{\partial}{\partial x^n}\right|_p\right\}$
is a basis of $T_pM$.  In particular, $\dim T_pM = n$.
\end{proposition}

\begin{proof}
\emph{Step~1: Spanning.}
By a standard lemma (Hadamard's lemma)\index{Hadamard's lemma}, every
$f \in C^\infty_p(M)$ can be written in local coordinates as
\[
  f \circ \varphi^{-1}(r)
  = f(p) + \sum_{i=1}^{n} (r^i - r^i_0)\, g_i(r),
\]
where $r_0 = \varphi(p)$ and $g_i$ are smooth with
$g_i(r_0) = \frac{\partial(f \circ \varphi^{-1})}{\partial r^i}(r_0)$.
For any derivation~$v$, Lemma~\ref{lem:derivation-constant} gives
$v(f(p)) = 0$, so
\[
  v(f) = \sum_{i=1}^{n} v(x^i - x^i(p))\, g_i(\varphi(p))
  = \sum_{i=1}^{n} v(x^i)\,
  \left.\frac{\partial}{\partial x^i}\right|_p(f).
\]
Hence $v = \sum_i v(x^i)\, \left.\frac{\partial}{\partial x^i}\right|_p$.

\emph{Step~2: Linear independence.}
If $\sum_i a^i \left.\frac{\partial}{\partial x^i}\right|_p = 0$, apply
this to the coordinate function~$x^j$ to get $a^j = 0$.
\end{proof}

%% ── Curves approach ───────────────────────────────────────
\subsection{Curves approach}\label{ssec:curves}

\begin{definition}[Tangent vector as equivalence class of curves]%
\label{def:tangent-curve}\index{tangent vector!curves approach}
Two smooth curves $\gamma_1, \gamma_2 \colon (-\varepsilon, \varepsilon) \to M$
with $\gamma_1(0) = \gamma_2(0) = p$ are \textbf{equivalent} if for some
(equivalently, every) chart $(U, \varphi)$ around~$p$,
\[
  (\varphi \circ \gamma_1)'(0) = (\varphi \circ \gamma_2)'(0).
\]
Denote the equivalence class of~$\gamma$ by $[\gamma]$ and the set of all
such classes by $\widetilde{T}_pM$.
\end{definition}

\begin{remark}\label{rem:chart-independence-curves}
That the definition is chart-independent follows from the chain rule
in~$\R^n$: if $(V, \psi)$ is another chart, then
$(\psi \circ \gamma)'(0) = D(\psi \circ \varphi^{-1})_{\varphi(p)}
\cdot (\varphi \circ \gamma)'(0)$, and the transition map
$\psi \circ \varphi^{-1}$ is a diffeomorphism.
\end{remark}

The set $\widetilde{T}_pM$ becomes a real vector space by defining, in a
chart $(U, \varphi)$,
\[
  [\gamma_1] + [\gamma_2]
  \;=\; \bigl[\varphi^{-1}\bigl(\varphi(p)
    + (\varphi \circ \gamma_1)'(0)\,t
    + (\varphi \circ \gamma_2)'(0)\,t\bigr)\bigr],
  \qquad
  \lambda\,[\gamma] \;=\;
  \bigl[\varphi^{-1}\bigl(\varphi(p)
    + \lambda\,(\varphi \circ \gamma)'(0)\,t\bigr)\bigr].
\]

%% ── Equivalence of the two definitions ────────────────────
\subsection{Equivalence of the two definitions}%
\label{ssec:equivalence-tangent}

\begin{theorem}\label{thm:tangent-equiv}%
\index{tangent space!equivalence of definitions}
The map
\[
  \Phi \colon \widetilde{T}_pM \longrightarrow T_pM,
  \qquad
  [\gamma] \;\longmapsto\;
  \Bigl( f \mapsto (f \circ \gamma)'(0) \Bigr)
\]
is a well-defined linear isomorphism.
\end{theorem}

\begin{proof}
\emph{Well-definedness.}
If $[\gamma_1] = [\gamma_2]$, then for any smooth germ~$f$ at~$p$,
\[
  (f \circ \gamma_1)'(0)
  = D(f \circ \varphi^{-1})_{\varphi(p)}\, (\varphi \circ \gamma_1)'(0)
  = D(f \circ \varphi^{-1})_{\varphi(p)}\, (\varphi \circ \gamma_2)'(0)
  = (f \circ \gamma_2)'(0),
\]
so $\Phi([\gamma_1]) = \Phi([\gamma_2])$.  A direct computation shows that
$\Phi([\gamma])$ satisfies the Leibniz rule, hence lies in~$T_pM$.

\emph{Linearity.}
We verify $\Phi([\gamma_1] + [\gamma_2]) = \Phi([\gamma_1]) +
\Phi([\gamma_2])$ using the chain rule in coordinates; the argument for
scalar multiplication is analogous.

\emph{Surjectivity.}
Given $v = \sum_i a^i \left.\frac{\partial}{\partial x^i}\right|_p$, let
$\gamma(t) = \varphi^{-1}(\varphi(p) + t\, a)$, where $a = (a^1, \dots,
a^n)$.  Then $\Phi([\gamma]) = v$.

\emph{Injectivity.}
If $\Phi([\gamma]) = 0$, then $(f \circ \gamma)'(0) = 0$ for all~$f$.
Applying this to the coordinate functions $x^i$ gives
$(\varphi \circ \gamma)'(0) = 0$, so $[\gamma] = 0$.
\end{proof}

Henceforth, we identify $\widetilde{T}_pM$ with $T_pM$ via~$\Phi$
without further comment.

\begin{example}\label{ex:tangent-Rn}%
\index{tangent space!of Euclidean space}
For $M = \R^n$ with the identity chart, $T_p\R^n \cong \R^n$ via
$\sum_i a^i \left.\frac{\partial}{\partial r^i}\right|_p \mapsto
(a^1, \dots, a^n)$.
\end{example}

\begin{example}\label{ex:tangent-Sn}%
\index{tangent space!of sphere}
Viewing $S^{n-1} \subset \R^n$ as a level set (see
Chapter~\ref{ch:submanifolds}), one obtains
$T_pS^{n-1} \cong \set{v \in \R^n : \langle v, p \rangle = 0}$,
the orthogonal complement of~$p$.
\end{example}

%%%%%%%%%%%%%%%%%%%%%%%%%%%%%%%%%%%%%%%%%%%%%%%%%%%%%%%%%%%%%
\section{The differential}\label{sec:differential}
%%%%%%%%%%%%%%%%%%%%%%%%%%%%%%%%%%%%%%%%%%%%%%%%%%%%%%%%%%%%%

\begin{definition}[Differential of a smooth map]\label{def:differential}%
\index{differential}\index{pushforward|see{differential}}
Let $f \colon M \to N$ be a smooth map between manifolds and $p \in M$.
The \textbf{differential}\footnote{Also called the \emph{pushforward},
\emph{tangent map}, or \emph{derivative} of~$f$ at~$p$.} of $f$ at~$p$
is the linear map
\[
  \dd f_p \colon T_pM \longrightarrow T_{f(p)}N,
  \qquad
  (\dd f_p\, v)(g) \;=\; v(g \circ f),
\]
for $v \in T_pM$ and $g \in C^\infty_{f(p)}(N)$.
\end{definition}

In the curves picture, $\dd f_p([\gamma]) = [f \circ \gamma]$.

\begin{proposition}[Coordinate expression]\label{prop:diff-coords}%
\index{differential!coordinate expression}
Let $(U, \varphi = (x^i))$ and $(V, \psi = (y^j))$ be charts around~$p$
and $f(p)$ respectively.  Write $\hat{f} = \psi \circ f \circ \varphi^{-1}$.
Then
\[
  \dd f_p\!\left(\left.\frac{\partial}{\partial x^i}\right|_p\right)
  \;=\; \sum_{j=1}^m
  \frac{\partial \hat{f}^j}{\partial r^i}(\varphi(p))\;
  \left.\frac{\partial}{\partial y^j}\right|_{f(p)}.
\]
Thus the matrix of $\dd f_p$ in these coordinate bases is the
\textbf{Jacobian matrix}\index{Jacobian matrix}
$\left[\frac{\partial \hat{f}^j}{\partial r^i}(\varphi(p))\right]_{j,i}$.
\end{proposition}

\begin{proof}
For any smooth germ $g$ at $f(p)$,
\[
  \dd f_p\!\left(\left.\frac{\partial}{\partial x^i}\right|_p\right)(g)
  = \left.\frac{\partial}{\partial x^i}\right|_p(g \circ f)
  = \frac{\partial (g \circ f \circ \varphi^{-1})}{\partial r^i}(\varphi(p))
  = \sum_j \frac{\partial (g \circ \psi^{-1})}{\partial r^j}(\psi(f(p)))\,
    \frac{\partial \hat{f}^j}{\partial r^i}(\varphi(p)),
\]
by the chain rule in~$\R^n$.  The right-hand side equals
$\sum_j \frac{\partial \hat{f}^j}{\partial r^i}(\varphi(p))\,
\left.\frac{\partial}{\partial y^j}\right|_{f(p)}(g)$.
\end{proof}

%% ── Chain rule and functoriality ──────────────────────────
\subsection{Chain rule and functoriality}\label{ssec:chain-rule}

\begin{theorem}[Chain rule]\label{thm:chain-rule}%
\index{chain rule}
Let $f \colon M \to N$ and $g \colon N \to P$ be smooth maps.  Then for
every $p \in M$,
\[
  \dd(g \circ f)_p \;=\; \dd g_{f(p)} \circ \dd f_p.
\]
\end{theorem}

\begin{proof}
For $v \in T_pM$ and $h \in C^\infty_{g(f(p))}(P)$,
\[
  \dd(g \circ f)_p(v)(h)
  = v(h \circ g \circ f)
  = \dd f_p(v)(h \circ g)
  = \dd g_{f(p)}(\dd f_p(v))(h).
  \qedhere
\]
\end{proof}

\begin{corollary}\label{cor:diffeo-iso}%
\index{diffeomorphism!induces isomorphism}
If $f \colon M \to N$ is a diffeomorphism, then $\dd f_p$ is a linear
isomorphism for all $p \in M$.
\end{corollary}

\begin{proof}
Since $f^{-1} \circ f = \Id_M$, the chain rule gives
$\dd(f^{-1})_{f(p)} \circ \dd f_p = \dd(\Id_M)_p = \Id_{T_pM}$.
Similarly, $\dd f_p \circ \dd(f^{-1})_{f(p)} = \Id_{T_{f(p)}N}$.
\end{proof}

\begin{remark}[Functoriality]\label{rem:functoriality}%
\index{functoriality!tangent functor}
The assignments $M \mapsto TM$ and $f \mapsto \dd f$ define a (covariant)
functor $T$ from the category $\mathbf{Man}$ of smooth manifolds and
smooth maps to the category $\mathbf{Vect}_\R$ of real vector spaces
and linear maps.  The chain rule is precisely the statement that~$T$
preserves composition, and $\dd(\Id_M)_p = \Id_{T_pM}$ says that~$T$
preserves identities.
\end{remark}

The following commutative diagram summarises the chain rule:
\[
\begin{tikzcd}[column sep=large, row sep=large]
  T_pM \arrow[r, "\dd f_p"] \arrow[dr, "\dd(g \circ f)_p"'] &
  T_{f(p)}N \arrow[d, "\dd g_{f(p)}"] \\
  & T_{g(f(p))}P
\end{tikzcd}
\]

%%%%%%%%%%%%%%%%%%%%%%%%%%%%%%%%%%%%%%%%%%%%%%%%%%%%%%%%%%%%%
\section{The tangent bundle}\label{sec:tangent-bundle}
%%%%%%%%%%%%%%%%%%%%%%%%%%%%%%%%%%%%%%%%%%%%%%%%%%%%%%%%%%%%%

\begin{definition}[Tangent bundle]\label{def:tangent-bundle}%
\index{tangent bundle}
The \textbf{tangent bundle} of~$M$ is the disjoint union
\[
  TM \;=\; \bigsqcup_{p \in M} T_pM
  \;=\; \set{(p, v) : p \in M,\; v \in T_pM},
\]
together with the natural projection
$\pi \colon TM \to M$, $(p, v) \mapsto p$.
\end{definition}

\begin{theorem}\label{thm:TM-manifold}%
\index{tangent bundle!smooth structure}
If $M$ is a smooth $n$-manifold, then $TM$ carries a natural smooth
structure making it a smooth $2n$-manifold, with respect to which
$\pi \colon TM \to M$ is a smooth surjective submersion.
\end{theorem}

\begin{proof}
\emph{Step~1: Topology and charts.}
Let $\mathcal{A} = \{(U_\alpha, \varphi_\alpha)\}$ be a smooth atlas
for~$M$.  For each~$\alpha$, define
\[
  \widetilde{\varphi}_\alpha \colon \pi^{-1}(U_\alpha)
  \longrightarrow \varphi_\alpha(U_\alpha) \times \R^n
  \;\subset\; \R^{2n},
\]
by
\[
  \widetilde{\varphi}_\alpha(p, v)
  \;=\;
  \Bigl(\varphi_\alpha(p),\;
  v(x^1_\alpha),\, \dots,\, v(x^n_\alpha)\Bigr),
\]
where $x^i_\alpha = r^i \circ \varphi_\alpha$ are the coordinate functions.
Concretely, if $v = \sum_i a^i \left.\frac{\partial}{\partial
x^i_\alpha}\right|_p$, then
$\widetilde{\varphi}_\alpha(p,v) = (\varphi_\alpha(p), a^1, \dots, a^n)$.
Since $\varphi_\alpha(U_\alpha)$ is open in~$\R^n$,
$\varphi_\alpha(U_\alpha) \times \R^n$ is open in~$\R^{2n}$, and
$\widetilde{\varphi}_\alpha$ is a bijection onto this open set.
We give $TM$ the topology generated by declaring a set
$W \subset TM$ open if and only if
$\widetilde{\varphi}_\alpha(W \cap \pi^{-1}(U_\alpha))$ is open in
$\R^{2n}$ for every~$\alpha$.

\emph{Step~2: Transition maps.}
On $\pi^{-1}(U_\alpha \cap U_\beta)$, the transition map is
\begin{align*}
  \widetilde{\varphi}_\beta \circ \widetilde{\varphi}_\alpha^{-1}
  \colon
  \varphi_\alpha(U_\alpha \cap U_\beta) \times \R^n
  &\longrightarrow
  \varphi_\beta(U_\alpha \cap U_\beta) \times \R^n, \\
  (x, a) &\longmapsto
  \Bigl(\tau_{\beta\alpha}(x),\;
  D\tau_{\beta\alpha}(x)\, a\Bigr),
\end{align*}
where $\tau_{\beta\alpha} = \varphi_\beta \circ \varphi_\alpha^{-1}$ is
the transition map of~$M$.  Since $\tau_{\beta\alpha}$ is a smooth
diffeomorphism, the map $(x, a) \mapsto (\tau_{\beta\alpha}(x),\,
D\tau_{\beta\alpha}(x)\, a)$ is smooth (in fact, the second component
is linear in~$a$ and smooth in~$x$).

\emph{Step~3: Hausdorff and second countability.}
Because $M$ is Hausdorff and second-countable, and the fibres~$\R^n$ have
these properties, the topology on $TM$ is Hausdorff and second-countable.
(Alternatively, $TM$ is covered by countably many of the sets
$\pi^{-1}(U_\alpha)$, each homeomorphic to an open subset of~$\R^{2n}$.)

\emph{Step~4: The projection $\pi$.}
In the chart $\widetilde{\varphi}_\alpha$, the map $\pi$ is represented by
the projection $(x, a) \mapsto x$, which is a smooth submersion.
\end{proof}

\begin{example}\label{ex:tangent-bundle-Rn}
$T\R^n \cong \R^n \times \R^n = \R^{2n}$ as a smooth manifold.
\end{example}

\begin{example}\label{ex:tangent-bundle-S1}%
\index{tangent bundle!of the circle}
$TS^1 \cong S^1 \times \R$, since $S^1$ is a Lie group and every Lie
group has trivial tangent bundle.  More generally, $T(S^1 \times \cdots
\times S^1) \cong T^n \times \R^n$.
\end{example}

\begin{example}\label{ex:tangent-bundle-S2}%
\index{tangent bundle!of $S^2$}
The tangent bundle $TS^2$ is \emph{non-trivial}: $TS^2 \not\cong
S^2 \times \R^2$.  This is a consequence of the \textbf{hairy ball
theorem}\index{hairy ball theorem}, which states that there is no
non-vanishing continuous vector field on $S^2$.  If $TS^2$ were trivial,
then the constant section $p \mapsto (p, e_1)$ (for a fixed
$e_1 \in \R^2$) would give a non-vanishing vector field, a contradiction.
More generally, $TS^n$ is trivial if and only if $n \in \{1, 3, 7\}$
(Adams's theorem, 1962).
\end{example}

\begin{remark}[The tangent functor on morphisms]\label{rem:tangent-functor-maps}%
\index{tangent functor!on maps}
If $f \colon M \to N$ is smooth, the \textbf{global differential}
$\dd f \colon TM \to TN$ defined by $\dd f(p, v) = (f(p), \dd f_p\, v)$
is a smooth map.  In fact, if $\widetilde{\varphi}_\alpha$ and
$\widetilde{\psi}_\beta$ are the induced charts on $TM$ and $TN$, then
\[
  \widetilde{\psi}_\beta \circ \dd f \circ \widetilde{\varphi}_\alpha^{-1}
  (x, a) = \bigl(\hat{f}(x),\; D\hat{f}(x)\, a\bigr),
\]
which is smooth in $(x, a)$.  The diagram
\[
\begin{tikzcd}[column sep=large, row sep=large]
  TM \arrow[r, "\dd f"] \arrow[d, "\pi_M"'] &
  TN \arrow[d, "\pi_N"] \\
  M \arrow[r, "f"'] & N
\end{tikzcd}
\]
commutes, expressing $\dd f$ as a \textbf{bundle map}\index{bundle map}
over $f$.
\end{remark}

The following diagram illustrates the local trivialisation of $TM$:
\[
\begin{tikzcd}[column sep=huge, row sep=large]
  \pi^{-1}(U_\alpha)
    \arrow[r, "\widetilde{\varphi}_\alpha"]
    \arrow[d, "\pi"'] &
  \varphi_\alpha(U_\alpha) \times \R^n
    \arrow[d, "\mathrm{pr}_1"] \\
  U_\alpha \arrow[r, "\varphi_\alpha"'] &
  \varphi_\alpha(U_\alpha)
\end{tikzcd}
\]

%%%%%%%%%%%%%%%%%%%%%%%%%%%%%%%%%%%%%%%%%%%%%%%%%%%%%%%%%%%%%
\section{The cotangent bundle}\label{sec:cotangent-bundle}
%%%%%%%%%%%%%%%%%%%%%%%%%%%%%%%%%%%%%%%%%%%%%%%%%%%%%%%%%%%%%

\begin{definition}[Cotangent space and bundle]\label{def:cotangent}%
\index{cotangent space}\index{cotangent bundle}
The \textbf{cotangent space} at~$p$ is the dual vector space
\[
  T_p^*M \;=\; (T_pM)^* \;=\; \Hom_\R(T_pM, \R).
\]
Elements of $T_p^*M$ are called \textbf{covectors}\index{covector} at~$p$.
The \textbf{cotangent bundle}\index{cotangent bundle} is
\[
  T^*M \;=\; \bigsqcup_{p \in M} T_p^*M,
  \qquad \pi^* \colon T^*M \to M,\quad (p, \xi) \mapsto p.
\]
It is a smooth $2n$-manifold by the same construction as for~$TM$, with
charts
\[
  \widetilde{\varphi}_\alpha^*(p, \xi)
  = \bigl(\varphi_\alpha(p),\; \xi_1, \dots, \xi_n\bigr),
  \qquad\text{where } \xi = \sum_i \xi_i\, \dd x^i_\alpha\big|_p.
\]
\end{definition}

\begin{definition}[Differential of a function]\label{def:df}%
\index{differential!of a function}
For $f \in C^\infty(M)$, the \textbf{differential of $f$} is the smooth
section $\dd f \colon M \to T^*M$ defined by
\[
  \dd f_p(v) = v(f), \qquad v \in T_pM.
\]
In local coordinates $(x^1, \dots, x^n)$,
\[
  \dd f = \sum_{i=1}^n \frac{\partial f}{\partial x^i}\, \dd x^i.
\]
\end{definition}

\begin{proposition}\label{prop:cotangent-basis}%
\index{cotangent space!basis}
In a chart $(U, \varphi = (x^1, \dots, x^n))$, the differentials
$\dd x^1|_p, \dots, \dd x^n|_p$ form the basis of $T_p^*M$ dual to
$\left.\frac{\partial}{\partial x^1}\right|_p, \dots,
\left.\frac{\partial}{\partial x^n}\right|_p$:
\[
  \dd x^i\!\left(\frac{\partial}{\partial x^j}\right) = \delta^i_j.
\]
\end{proposition}

\begin{proof}
By definition,
$\dd x^i|_p\!\left(\left.\frac{\partial}{\partial x^j}\right|_p\right)
= \left.\frac{\partial}{\partial x^j}\right|_p(x^i)
= \frac{\partial (x^i \circ \varphi^{-1})}{\partial r^j}(\varphi(p))
= \frac{\partial r^i}{\partial r^j} = \delta^i_j$.
\end{proof}

\begin{definition}[Pullback of a covector]\label{def:pullback-covector}%
\index{pullback!covector}
If $f \colon M \to N$ is smooth, the \textbf{pullback} is
\[
  f^* \colon T^*_{f(p)}N \to T^*_pM,
  \qquad
  (f^*\xi)(v) = \xi(\dd f_p\, v).
\]
\end{definition}

Note that $f^*$ goes in the opposite direction to $\dd f$:
\[
\begin{tikzcd}[column sep=large]
  T_pM \arrow[r, "\dd f_p"] &
  T_{f(p)}N \\
  T_p^*M &
  T_{f(p)}^*N \arrow[l, "f^*"']
\end{tikzcd}
\]

%%%%%%%%%%%%%%%%%%%%%%%%%%%%%%%%%%%%%%%%%%%%%%%%%%%%%%%%%%%%%
\section{Vector fields}\label{sec:vector-fields}
%%%%%%%%%%%%%%%%%%%%%%%%%%%%%%%%%%%%%%%%%%%%%%%%%%%%%%%%%%%%%

\begin{definition}[Vector field]\label{def:vector-field}%
\index{vector field}
A \textbf{(smooth) vector field} on $M$ is a smooth section of the
tangent bundle, i.e., a smooth map $X \colon M \to TM$ such that
$\pi \circ X = \Id_M$.  Equivalently, $X$ assigns to each $p \in M$ a
tangent vector $X_p \in T_pM$, smoothly in~$p$.
The set of all smooth vector fields on $M$ is denoted
$\mathfrak{X}(M)$\index{$\mathfrak{X}(M)$}.
\end{definition}

In local coordinates $(x^1, \dots, x^n)$, a vector field is written
\[
  X = \sum_{i=1}^n X^i\, \frac{\partial}{\partial x^i},
\]
where $X^i \colon U \to \R$ are smooth functions (the \textbf{component
functions}\index{component functions} of~$X$).

\begin{remark}\label{rem:vf-module}%
\index{vector field!module structure}
The set $\mathfrak{X}(M)$ is a module over $C^\infty(M)$: for
$f \in C^\infty(M)$ and $X \in \mathfrak{X}(M)$, define
$(fX)_p = f(p)\, X_p$.  It is also a real vector space (taking~$f$
to be constant), but it is \emph{not} finite-dimensional
if $\dim M \geq 1$.
\end{remark}

\begin{remark}\label{rem:vf-derivation}
A vector field $X$ acts as a \textbf{derivation of $C^\infty(M)$},
i.e., an $\R$-linear map $X \colon C^\infty(M) \to C^\infty(M)$
satisfying $X(fg) = fX(g) + gX(f)$.  Conversely, every such derivation
arises from a unique smooth vector field.
\end{remark}

%%%%%%%%%%%%%%%%%%%%%%%%%%%%%%%%%%%%%%%%%%%%%%%%%%%%%%%%%%%%%
\section{The Lie bracket}\label{sec:lie-bracket}
%%%%%%%%%%%%%%%%%%%%%%%%%%%%%%%%%%%%%%%%%%%%%%%%%%%%%%%%%%%%%

\begin{definition}[Lie bracket]\label{def:lie-bracket}%
\index{Lie bracket}\index{bracket|see{Lie bracket}}
For $X, Y \in \mathfrak{X}(M)$, the \textbf{Lie bracket} $[X, Y]$ is
the vector field defined by
\[
  [X, Y](f) \;=\; X(Y(f)) - Y(X(f))
\]
for all $f \in C^\infty(M)$.
\end{definition}

\begin{proposition}[Coordinate formula]\label{prop:lie-bracket-coords}%
\index{Lie bracket!coordinate expression}
In local coordinates $(x^1, \dots, x^n)$, if
$X = \sum_i X^i \frac{\partial}{\partial x^i}$ and
$Y = \sum_j Y^j \frac{\partial}{\partial x^j}$, then
\[
  [X, Y] = \sum_{k=1}^n \left(
    \sum_{i=1}^n X^i \frac{\partial Y^k}{\partial x^i}
    - Y^i \frac{\partial X^k}{\partial x^i}
  \right) \frac{\partial}{\partial x^k}.
\]
\end{proposition}

\begin{proof}
Apply $[X, Y]$ to a coordinate function $x^k$:
\begin{align*}
  [X, Y](x^k)
  &= X(Y(x^k)) - Y(X(x^k))
  = X(Y^k) - Y(X^k) \\
  &= \sum_i X^i \frac{\partial Y^k}{\partial x^i}
   - \sum_i Y^i \frac{\partial X^k}{\partial x^i}.
\end{align*}
Since $[X,Y] = \sum_k [X,Y](x^k)\, \frac{\partial}{\partial x^k}$
locally, the formula follows.
\end{proof}

\begin{theorem}[Properties of the Lie bracket]\label{thm:lie-bracket-props}%
\index{Lie bracket!properties}
The Lie bracket satisfies, for all $X, Y, Z \in \mathfrak{X}(M)$ and
$f, g \in C^\infty(M)$:
\begin{enumerate}[label=(\roman*)]
  \item \textbf{$\R$-bilinearity:}
    $[\alpha X + \beta Y, Z] = \alpha[X, Z] + \beta[Y, Z]$ for
    $\alpha, \beta \in \R$.
  \item \textbf{Antisymmetry:} $[X, Y] = -[Y, X]$.
  \item \textbf{Jacobi identity\index{Jacobi identity}:}
    $[X, [Y, Z]] + [Y, [Z, X]] + [Z, [X, Y]] = 0$.
  \item \textbf{$C^\infty(M)$-Leibniz rule:}
    $[fX, gY] = fg\,[X, Y] + f(Xg)\,Y - g(Yf)\,X$.
\end{enumerate}
In particular, $(\mathfrak{X}(M), [\cdot,\cdot])$ is a
\textbf{Lie algebra}\index{Lie algebra} over~$\R$ (but not over
$C^\infty(M)$, by~(iv)).
\end{theorem}

\begin{proof}
Properties (i) and (ii) are immediate from the definition.

\emph{(iii) Jacobi identity.}  For $f \in C^\infty(M)$,
\begin{align*}
  [X,[Y,Z]](f)
  &= X(Y(Z(f)) - Z(Y(f))) - (Y(Z(X(f))) - Z(Y(X(f)))) \\
  &= X(Y(Z(f))) - X(Z(Y(f))) - Y(Z(X(f))) + Z(Y(X(f))).
\end{align*}
Summing over the three cyclic permutations of $(X, Y, Z)$, all twelve
terms cancel in pairs.

\emph{(iv) Leibniz rule.}  For $h \in C^\infty(M)$,
\begin{align*}
  [fX, gY](h)
  &= fX(gY(h)) - gY(fX(h)) \\
  &= fX(g)\, Y(h) + fg\, X(Y(h)) - gY(f)\, X(h) - gf\, Y(X(h)) \\
  &= fg\,[X,Y](h) + f(Xg)\, Y(h) - g(Yf)\, X(h).
  \qedhere
\end{align*}
\end{proof}

\begin{example}\label{ex:lie-bracket-R2}
On $\R^2$ with coordinates $(x, y)$, let $X = \frac{\partial}{\partial x}$
and $Y = x\,\frac{\partial}{\partial y}$.  Then
\[
  [X, Y] = \frac{\partial}{\partial x}\!\left(x\right)\,
  \frac{\partial}{\partial y}
  - x\, \frac{\partial}{\partial x}\!\left(1\right)\,
  \frac{\partial}{\partial x}
  = \frac{\partial}{\partial y}.
\]
\end{example}

\begin{proposition}[$f$-relatedness]\label{prop:f-related}%
\index{$f$-related vector fields}
Let $f \colon M \to N$ be smooth.  Vector fields $X \in \mathfrak{X}(M)$
and $Y \in \mathfrak{X}(N)$ are \textbf{$f$-related}\index{$f$-related}
(written $X \sim_f Y$) if $\dd f_p(X_p) = Y_{f(p)}$ for all $p \in M$.
If $X_1 \sim_f Y_1$ and $X_2 \sim_f Y_2$, then
$[X_1, X_2] \sim_f [Y_1, Y_2]$.
\end{proposition}

\begin{proof}
For $g \in C^\infty(N)$ and $p \in M$,
\begin{align*}
  \dd f_p([X_1, X_2]_p)(g)
  &= [X_1, X_2]_p(g \circ f)
  = X_{1,p}(X_2(g \circ f)) - X_{2,p}(X_1(g \circ f)) \\
  &= X_{1,p}((Y_2 g) \circ f) - X_{2,p}((Y_1 g) \circ f) \\
  &= Y_{1,f(p)}(Y_2 g) - Y_{2,f(p)}(Y_1 g)
  = [Y_1, Y_2]_{f(p)}(g).
  \qedhere
\end{align*}
\end{proof}

%%%%%%%%%%%%%%%%%%%%%%%%%%%%%%%%%%%%%%%%%%%%%%%%%%%%%%%%%%%%%
\section{Flows of vector fields}\label{sec:flows}
%%%%%%%%%%%%%%%%%%%%%%%%%%%%%%%%%%%%%%%%%%%%%%%%%%%%%%%%%%%%%

\begin{definition}[Integral curve]\label{def:integral-curve}%
\index{integral curve}
Let $X \in \mathfrak{X}(M)$.  An \textbf{integral curve} of $X$ through
$p \in M$ is a smooth curve $\gamma \colon J \to M$, defined on an open
interval $J \ni 0$, such that $\gamma(0) = p$ and
$\gamma'(t) = X_{\gamma(t)}$ for all $t \in J$.
\end{definition}

The existence and uniqueness of integral curves is a direct consequence
of the classical Cauchy--Lipschitz theorem\index{Cauchy--Lipschitz theorem}
(Picard--Lindel\"of theorem) applied in local coordinates.

\begin{theorem}[Existence and uniqueness of integral curves]%
\label{thm:integral-curve}\index{integral curve!existence and uniqueness}
Let $X \in \mathfrak{X}(M)$ and $p \in M$.  There exists a unique
maximal integral curve $\gamma_p \colon J_p \to M$ of $X$ with
$\gamma_p(0) = p$, where $J_p \subset \R$ is the largest open interval
on which the solution exists.
\end{theorem}

\begin{proof}[Proof sketch]
In a chart $(U, \varphi)$, the equation $\gamma'(t) = X_{\gamma(t)}$
becomes the ODE system $\dot{y}(t) = F(y(t))$ in~$\R^n$, where
$F = (\varphi_* X) \circ \varphi^{-1}$ is smooth and hence locally
Lipschitz.  The Cauchy--Lipschitz theorem guarantees a unique local
solution.  By covering $M$ with charts and using uniqueness on overlaps,
one extends to a unique maximal solution.
\end{proof}

\begin{definition}[Flow]\label{def:flow}%
\index{flow}
The \textbf{flow} of $X \in \mathfrak{X}(M)$ is the map
\[
  \theta \colon \mathcal{D} \to M,
  \qquad (t, p) \mapsto \theta_t(p) = \gamma_p(t),
\]
where $\mathcal{D} = \set{(t, p) \in \R \times M : t \in J_p}$ is the
\textbf{flow domain}\index{flow domain}.
\end{definition}

\begin{theorem}[Properties of the flow]\label{thm:flow-props}%
\index{flow!properties}
Let $\theta$ be the flow of $X \in \mathfrak{X}(M)$.
\begin{enumerate}[label=(\roman*)]
  \item $\mathcal{D}$ is an open subset of $\R \times M$ containing
    $\{0\} \times M$.
  \item $\theta \colon \mathcal{D} \to M$ is smooth.
  \item $\theta_0 = \Id_M$.
  \item \textbf{Group law:} $\theta_t \circ \theta_s(p) = \theta_{t+s}(p)$
    whenever both sides are defined.
\end{enumerate}
\end{theorem}

\begin{definition}[Complete vector field]\label{def:complete-vf}%
\index{vector field!complete}\index{complete vector field}
A vector field $X$ is \textbf{complete} if $J_p = \R$ for all $p \in M$,
i.e., $\mathcal{D} = \R \times M$.  In this case, $\theta_t \colon M \to M$
is a diffeomorphism for each~$t$, and $t \mapsto \theta_t$ is a smooth
group homomorphism from $(\R, +)$ to $\Diff(M)$\index{one-parameter group
of diffeomorphisms}, called the
\textbf{one-parameter group of diffeomorphisms} generated by~$X$.
\end{definition}

\begin{proposition}\label{prop:compact-complete}%
\index{vector field!complete!on compact manifold}
Every smooth vector field on a compact manifold is complete.
\end{proposition}

\begin{proof}
If $M$ is compact, the maximal interval~$J_p$ must be all of~$\R$ for
every~$p$: if $J_p = (a, b)$ with $b < \infty$, then
$\gamma_p(t)$ leaves every compact set as $t \to b^-$ (by a standard
escape lemma), contradicting the compactness of~$M$.
\end{proof}

\begin{example}\label{ex:flow-rotation}%
\index{flow!rotation on $S^1$}
On $S^1 \subset \R^2$, the vector field
$X_{(x,y)} = (-y, x)$ (the unit tangent to the circle) has flow
$\theta_t(x, y) = (x\cos t - y\sin t,\; x\sin t + y\cos t)$, which is
rotation by angle $t$.  This vector field is complete, and the
one-parameter group $t \mapsto \theta_t$ is the group of rotations
$\SO(2)$.
\end{example}

\begin{example}\label{ex:flow-incomplete}%
\index{flow!incomplete}
On $M = \R$, the vector field $X = x^2 \frac{\dd}{\dd x}$ has integral
curves $\gamma_p(t) = \frac{p}{1 - pt}$ (for $p \neq 0$), defined on
$(-\infty, 1/p)$ if $p > 0$ and $(1/p, +\infty)$ if $p < 0$.  The flow
is not complete: solutions blow up in finite time.  This cannot happen on a
compact manifold, confirming Proposition~\ref{prop:compact-complete}.
\end{example}

\begin{remark}[Lie derivative of a function]\label{rem:lie-derivative-fn}%
\index{Lie derivative!of a function}
If $X \in \mathfrak{X}(M)$ with flow $\theta$ and $f \in C^\infty(M)$,
the \textbf{Lie derivative} of $f$ along $X$ is
\[
  (\mathcal{L}_X f)(p)
  = \left.\frac{\dd}{\dd t}\right|_{t=0} f(\theta_t(p))
  = X_p(f) = (Xf)(p).
\]
Thus $\mathcal{L}_X f = Xf$: the Lie derivative of a function is simply
the action of $X$ as a derivation.
\end{remark}

The flow is depicted schematically below.

\begin{center}
\begin{tikzpicture}[scale=1.0,
  >=Stealth,
  flow/.style={thick, blue!70!black, decoration={markings,
    mark=at position 0.55 with {\arrow{>}}}, postaction={decorate}}]
  %% manifold blob
  \draw[thick, fill=gray!8] (0,0) ellipse (3.5 and 2);
  \node at (2.8, 1.6) {$M$};
  %% flow lines
  \draw[flow] (-2.5, -0.8) .. controls (-1.2, -0.4) and (0.3, 0.3) .. (2.0, 0.7);
  \draw[flow] (-2.5, -1.4) .. controls (-1.0, -1.0) and (0.5, -0.3) .. (2.2, 0.0);
  \draw[flow] (-2.3, 0.2) .. controls (-0.8, 0.5) and (0.7, 1.0) .. (2.4, 1.2);
  %% points
  \fill (-1.5, -0.55) circle (1.5pt) node[below] {$p$};
  \fill (0.8, 0.1) circle (1.5pt) node[above] {$\theta_t(p)$};
  %% vector at p
  \draw[->, thick, red!70!black] (-1.5, -0.55) -- (-0.7, -0.25)
    node[midway, above left] {$X_p$};
\end{tikzpicture}
\end{center}

%%%%%%%%%%%%%%%%%%%%%%%%%%%%%%%%%%%%%%%%%%%%%%%%%%%%%%%%%%%%%
\section{Exercises}
%%%%%%%%%%%%%%%%%%%%%%%%%%%%%%%%%%%%%%%%%%%%%%%%%%%%%%%%%%%%%

\begin{exercise}\label{exo:tangent-product}%
\index{tangent space!of product manifold}
Let $M$ and $N$ be smooth manifolds and $(p, q) \in M \times N$.
Prove that
\[
  T_{(p,q)}(M \times N) \;\cong\; T_pM \oplus T_qN,
\]
where the isomorphism is given by
$v \mapsto (\dd\mathrm{pr}_{1,(p,q)}\, v,\; \dd\mathrm{pr}_{2,(p,q)}\, v)$.
\end{exercise}

\begin{exercise}\label{exo:lie-bracket-compute}
On $\R^3$ with coordinates $(x, y, z)$, compute the Lie brackets
$[X, Y]$, $[Y, Z]$, and $[X, Z]$, where
\[
  X = y\,\frac{\partial}{\partial z} - z\,\frac{\partial}{\partial y},
  \qquad
  Y = z\,\frac{\partial}{\partial x} - x\,\frac{\partial}{\partial z},
  \qquad
  Z = x\,\frac{\partial}{\partial y} - y\,\frac{\partial}{\partial x}.
\]
These are the infinitesimal generators of rotations.  Verify the Jacobi
identity.
\end{exercise}

\begin{exercise}\label{exo:df-computation}
Let $f \colon \R^2 \to \R^3$ be given by
$f(u, v) = (u\cos v,\; u\sin v,\; v)$.
Compute $\dd f_{(u,v)}$ in the standard bases and find its rank.
\end{exercise}

\begin{exercise}\label{exo:tangent-bundle-trivial}%
\index{tangent bundle!trivialisable}
Prove that the tangent bundle $TS^1$ is trivial by exhibiting an explicit
global trivialisation $TS^1 \xrightarrow{\;\sim\;} S^1 \times \R$.
\end{exercise}

\begin{exercise}\label{exo:lie-bracket-leibniz}
Verify directly from the definition that $[X, Y]$ satisfies the Leibniz
rule, i.e., that $[X, Y]$ is a derivation of $C^\infty(M)$, not just
an endomorphism.
\end{exercise}

\begin{exercise}\label{exo:flow-linear}%
\index{flow!of linear vector field}
Let $A \in M_n(\R)$ and define $X_A \in \mathfrak{X}(\R^n)$ by
$(X_A)_x = Ax$ (identifying $T_x\R^n \cong \R^n$).  Show that the
flow of $X_A$ is $\theta_t(x) = e^{tA}x$.
\end{exercise}

\begin{exercise}\label{exo:bracket-flow}%
\index{Lie bracket!and flow}
Let $X, Y \in \mathfrak{X}(M)$ and let $\theta$ be the flow of $X$.
Show that the Lie bracket can be expressed as
\[
  [X, Y]_p
  = \lim_{t \to 0} \frac{(\dd\theta_{-t})_{\theta_t(p)}\,
    Y_{\theta_t(p)} - Y_p}{t}
  = \left.\frac{\dd}{\dd t}\right|_{t=0}
    (\theta_t^* Y)_p.
\]
This is the \textbf{Lie derivative}\index{Lie derivative!of a vector field}
$\mathcal{L}_X Y = [X, Y]$.
\end{exercise}

\begin{exercise}\label{exo:covector-pullback}
Let $f \colon \R^2 \to \R^2$ be given by $f(x, y) = (x^2 - y^2, 2xy)$.
Compute $f^*(\dd u)$ and $f^*(\dd v)$, where $(u, v)$ are coordinates on
the codomain.
\end{exercise}

\begin{exercise}\label{exo:tangent-RP}%
\index{tangent bundle!of projective space}
Let $\pi \colon S^n \to \RP^n$ be the canonical double covering.  Show
that for each $p \in S^n$, the differential
$\dd\pi_p \colon T_pS^n \to T_{\pi(p)}\RP^n$ is an isomorphism.
Use this to prove $\dim T_{[x]}\RP^n = n$.
\end{exercise}


%%%%%%%%%%%%%%%%%%%%%%%%%%%%%%%%%%%%%%%%%%%%%%%%%%%%%%%%%%%%%
%% CHAPTER 4
%%%%%%%%%%%%%%%%%%%%%%%%%%%%%%%%%%%%%%%%%%%%%%%%%%%%%%%%%%%%%
\chapter{Submanifolds and the Regular Value Theorem}%
\label{ch:submanifolds}
\minitoc\bigskip

With the tangent space machinery in place, we can now study submanifolds
and the key analytic result that produces them: the \emph{regular value
theorem}.  This single theorem, combined with its generalisations,
is responsible for establishing the smooth manifold structure of an
enormous number of classical examples---spheres, orthogonal groups,
special linear groups, and many more.

%%%%%%%%%%%%%%%%%%%%%%%%%%%%%%%%%%%%%%%%%%%%%%%%%%%%%%%%%%%%%
\section{Immersions, embeddings, and submanifolds}%
\label{sec:immersion-embedding}
%%%%%%%%%%%%%%%%%%%%%%%%%%%%%%%%%%%%%%%%%%%%%%%%%%%%%%%%%%%%%

\begin{definition}[Immersion and submersion]\label{def:immersion-submersion}%
\index{immersion}\index{submersion}
A smooth map $f \colon M \to N$ is:
\begin{itemize}
  \item an \textbf{immersion} if $\dd f_p$ is injective for every
    $p \in M$;
  \item a \textbf{submersion} if $\dd f_p$ is surjective for every
    $p \in M$.
\end{itemize}
\end{definition}

\begin{definition}[Immersed submanifold]\label{def:immersed-submanifold}%
\index{submanifold!immersed}
An \textbf{immersed submanifold} of $N$ is a subset $S \subset N$
together with a topology and smooth structure on $S$ making the
inclusion $\iota \colon S \hookrightarrow N$ a smooth immersion.
The \textbf{codimension}\index{codimension} is $\dim N - \dim S$.
\end{definition}

\begin{example}\label{ex:figure-eight}%
\index{figure-eight}
The figure-eight curve $\gamma \colon (-\pi, \pi) \to \R^2$,
$\gamma(t) = (\sin 2t, \sin t)$, is an injective immersion whose image
is not an embedded submanifold (the self-intersection at the origin has
no manifold neighbourhood).
\end{example}

\begin{center}
\begin{tikzpicture}[scale=1.4, >=Stealth]
  %% Figure-eight (lemniscate-like)
  \draw[thick, blue!70!black, domain=-0.97*pi:0.97*pi, samples=200,
    variable=\t]
    plot ({sin(2*\t r)}, {sin(\t r)});
  %% Mark the self-intersection
  \fill[red!70!black] (0, 0) circle (1.5pt)
    node[below right, font=\small] {self-intersection};
  %% Arrows showing direction
  \draw[->, thick, blue!70!black] ({sin(2*1.0 r)}, {sin(1.0 r)})
    -- ++(0.05, 0.03);
  \node[blue!70!black, anchor=south] at (0, 1.15) {$\gamma((-\pi,\pi))$};
\end{tikzpicture}
\end{center}

\begin{definition}[Embedding]\label{def:embedding}%
\index{embedding}
A smooth map $f \colon M \to N$ is a \textbf{(smooth) embedding} if it
is an injective immersion that is also a homeomorphism onto its image
$f(M) \subset N$ (with the subspace topology).
\end{definition}

\begin{definition}[Embedded submanifold]\label{def:embedded-submanifold}%
\index{submanifold!embedded}
An \textbf{embedded submanifold} (or \textbf{regular submanifold}) of $N$
is a subset $S \subset N$ such that for every $p \in S$ there exists a
chart $(U, \varphi)$ of $N$ with $p \in U$ and
\[
  \varphi(U \cap S) = \varphi(U) \cap (\R^k \times \{0\}),
\]
where $k = \dim S$ and we identify $\R^k \times \{0\} \subset \R^n$.
Such a chart is called a \textbf{slice chart}\index{slice chart} or
\textbf{adapted chart}\index{adapted chart}.
\end{definition}

\begin{proposition}\label{prop:embedded-iff-embedding}%
\index{submanifold!embedded!characterisation}
A subset $S \subset N$ is an embedded submanifold of dimension $k$ if and
only if the inclusion $\iota \colon S \hookrightarrow N$ is a smooth
embedding, when $S$ carries the subspace topology and the induced smooth
structure.
\end{proposition}

\begin{remark}\label{rem:immersed-vs-embedded}
Every embedded submanifold is an immersed submanifold, but not conversely.
The distinction lies in the topology: the topology of an immersed
submanifold need not agree with the subspace topology inherited from the
ambient manifold.  A classical example is a line of irrational slope in
the torus $T^2 = \R^2/\Z^2$: it is an injective immersion of $\R$, dense
in $T^2$, hence not embedded.
\end{remark}

\begin{proposition}\label{prop:compact-immersion-embedding}%
\index{immersion!compact domain}
If $M$ is compact and $f \colon M \to N$ is an injective immersion,
then $f$ is an embedding.
\end{proposition}

\begin{proof}
An injective continuous map from a compact space to a Hausdorff space is
a homeomorphism onto its image.  Since $f$ is also an immersion, it is
a smooth embedding.
\end{proof}

%%%%%%%%%%%%%%%%%%%%%%%%%%%%%%%%%%%%%%%%%%%%%%%%%%%%%%%%%%%%%
\section{Regular values and the preimage theorem}%
\label{sec:regular-value}
%%%%%%%%%%%%%%%%%%%%%%%%%%%%%%%%%%%%%%%%%%%%%%%%%%%%%%%%%%%%%

\begin{definition}[Regular and critical values]\label{def:regular-critical}%
\index{regular value}\index{critical value}\index{critical point}
Let $f \colon M \to N$ be a smooth map.  A point $p \in M$ is a
\textbf{regular point}\index{regular point} if $\dd f_p$ is surjective;
otherwise it is a \textbf{critical point}.  A point $q \in N$ is a
\textbf{regular value} if every $p \in f^{-1}(q)$ is a regular point
(in particular, $q$ is a regular value if $f^{-1}(q) = \varnothing$).
Otherwise, $q$ is a \textbf{critical value}.
\end{definition}

The main tool is the following local result, which is a smooth version
of the implicit function theorem.

\begin{theorem}[Submersion theorem / local normal form]%
\label{thm:submersion-local}\index{submersion!local normal form}
Let $f \colon M^m \to N^n$ be smooth with $\dd f_p$ surjective ($m \geq n$).
Then there exist charts $(U, \varphi)$ around $p$ and $(V, \psi)$ around
$f(p)$ such that
\[
  \psi \circ f \circ \varphi^{-1}(x^1, \dots, x^m)
  = (x^1, \dots, x^n).
\]
That is, in suitable coordinates, $f$ is just the projection onto the
first $n$ factors.
\end{theorem}

\begin{proof}
By replacing $M$ and $N$ with coordinate neighbourhoods, we may assume
$M \subset \R^m$ is open, $N \subset \R^n$ is open, and $p = 0$,
$f(0) = 0$.  Since $\dd f_0 \colon \R^m \to \R^n$ is surjective,
after reordering coordinates we may assume the last~$n$ columns of the
Jacobian $Df(0)$ form an invertible $n \times n$ matrix.  Write
$x = (x', x'') \in \R^{m-n} \times \R^n$ and define
$\Phi \colon M \to \R^m$ by $\Phi(x', x'') = (x', f(x', x''))$.
Then
\[
  D\Phi(0) = \begin{pmatrix}
    I_{m-n} & 0 \\ * & Df_0\big|_{\R^n}
  \end{pmatrix},
\]
which is invertible.  By the inverse function theorem,
$\Phi$ is a local diffeomorphism near~$0$.  Setting
$\varphi = \Phi$ and $\psi = \Id$, one verifies
\[
  f \circ \varphi^{-1}(y', y'')
  = f \circ \Phi^{-1}(y', y'')
  = y'',
\]
since $\Phi^{-1}(y', y'')$ maps to a point whose $f$-image is~$y''$.
Reindexing to place $y''$ as the first~$n$ coordinates gives the stated
normal form.
\end{proof}

\begin{theorem}[Regular value theorem / preimage theorem]%
\label{thm:regular-value}\index{regular value theorem}%
\index{preimage theorem|see{regular value theorem}}
Let $f \colon M^m \to N^n$ be a smooth map and $q \in N$ a regular value
of~$f$.  Then $f^{-1}(q)$ is either empty or an embedded submanifold
of~$M$ of dimension $m - n$.
\end{theorem}

\begin{proof}
Assume $S = f^{-1}(q) \neq \varnothing$ and let $p \in S$.
Since~$q$ is a regular value, $\dd f_p$ is surjective.  By the
submersion theorem (Theorem~\ref{thm:submersion-local}), there exist
charts $(U, \varphi)$ around $p$ with $\varphi(p) = 0$ and $(V, \psi)$
around $q$ with $\psi(q) = 0$ such that
\[
  \psi \circ f \circ \varphi^{-1}(x^1, \dots, x^m)
  = (x^1, \dots, x^n).
\]
Therefore
\begin{align*}
  \varphi(U \cap S)
  &= \varphi\bigl(U \cap f^{-1}(q)\bigr)
  = \set{x \in \varphi(U) : \psi \circ f \circ \varphi^{-1}(x) = 0} \\
  &= \set{x \in \varphi(U) : x^1 = \cdots = x^n = 0}
  = \varphi(U) \cap (\{0\} \times \R^{m-n}).
\end{align*}
After reindexing, this shows $(U, \varphi)$ is a slice chart for~$S$
of codimension~$n$.  Since $p \in S$ was arbitrary, $S$ is an embedded
submanifold of dimension $m - n$.
\end{proof}

\begin{corollary}[Tangent space of a level set]\label{cor:tangent-level-set}%
\index{tangent space!of a level set}
Under the hypotheses of Theorem~\ref{thm:regular-value}, if
$S = f^{-1}(q)$ and $p \in S$, then
\[
  T_pS \;=\; \ker \dd f_p
  \;=\; \set{v \in T_pM : \dd f_p(v) = 0}.
\]
\end{corollary}

\begin{proof}
Since $f|_S$ is the constant map~$q$, we have
$\dd(f \circ \iota)_p = \dd f_p \circ \dd\iota_p = 0$, where
$\iota \colon S \hookrightarrow M$ is the inclusion.  Hence
$\dd\iota_p(T_pS) \subset \ker \dd f_p$.  Since $\dd\iota_p$ is
injective (the inclusion is an embedding), we obtain
$T_pS \subset \ker \dd f_p$ (identifying $T_pS$ with its image in
$T_pM$).  By dimension count,
$\dim T_pS = m - n = m - \mathrm{rank}(\dd f_p)
= \dim \ker \dd f_p$, giving equality.
\end{proof}

The situation is depicted in the following figure:

\begin{center}
\begin{tikzpicture}[scale=1.1, >=Stealth]
  %% Ambient manifold M (a region of R^m)
  \draw[thick, fill=blue!5] (-3,-2) rectangle (3,2.5);
  \node[anchor=north east] at (3, 2.5) {$M$};

  %% Level set S = f^{-1}(q)
  \draw[very thick, red!70!black]
    (-2.2, -1.0) .. controls (-1.0, 0.5) and (0.5, 1.5) .. (2.5, 1.8);
  \node[red!70!black, anchor=south] at (0, 1.3) {$S = f^{-1}(q)$};

  %% Point p on S
  \fill (-0.3, 0.6) circle (1.8pt) node[below left] {$p$};

  %% Tangent space T_pS (a line tangent to S at p)
  \draw[thick, green!50!black, dashed]
    (-1.8, -0.2) -- (1.2, 1.4);
  \node[green!50!black, anchor=west] at (1.2, 1.4) {$T_pS = \ker\,\dd f_p$};

  %% Arrow indicating df_p direction (normal to S)
  \draw[->, thick, blue!70!black] (-0.3, 0.6) -- (-0.8, 1.8)
    node[midway, right] {$(\dd f_p)^\perp$};

  %% Right side: target N
  \draw[thick, fill=orange!5] (5,-0.5) rectangle (7,1.5);
  \node[anchor=north east] at (7, 1.5) {$N$};
  \fill (6, 0.5) circle (1.8pt) node[right] {$q$};

  %% Arrow f
  \draw[->, thick, shorten >=4pt, shorten <=4pt]
    (3.2, 1.0) -- (4.8, 0.7) node[midway, above] {$f$};
\end{tikzpicture}
\end{center}

%%%%%%%%%%%%%%%%%%%%%%%%%%%%%%%%%%%%%%%%%%%%%%%%%%%%%%%%%%%%%
\section{Applications of the regular value theorem}%
\label{sec:applications-rvt}
%%%%%%%%%%%%%%%%%%%%%%%%%%%%%%%%%%%%%%%%%%%%%%%%%%%%%%%%%%%%%

The following figure shows the level sets of $f(x, y) = x^2 + y^2$,
illustrating that regular values ($c > 0$) give smooth manifolds (circles)
while the critical value ($c = 0$) gives a single point.

\begin{center}
\begin{tikzpicture}[scale=0.9, >=Stealth]
  %% Axes
  \draw[->] (-3.2, 0) -- (3.2, 0) node[right] {$x$};
  \draw[->] (0, -3.2) -- (0, 3.2) node[above] {$y$};
  %% Level sets: circles of increasing radius
  \foreach \r/\lab in {0.8/$c=0.64$, 1.4/$c=1.96$, 2.0/$c=4$, 2.6/$c=6.76$} {
    \draw[thick, blue!60!black] (0,0) circle (\r);
  }
  \node[blue!60!black, font=\small] at (2.0, -2.5) {$f^{-1}(c)$};
  %% Critical point at origin
  \fill[red!70!black] (0, 0) circle (2pt)
    node[below right, font=\small] {$0$ (crit.\ pt.)};
  %% Labels for some level sets
  \node[font=\small, blue!60!black] at (1.8, 1.8) {$c=4$};
  \node[font=\small, blue!60!black] at (0.95, 0.95) {$c\!=\!1.96$};
\end{tikzpicture}
\end{center}

We now apply Theorem~\ref{thm:regular-value} to establish the manifold
structure of several classical spaces.

%% ── Spheres ───────────────────────────────────────────────
\subsection{Spheres}\label{ssec:sphere}

\begin{proposition}\label{prop:sphere-submanifold}%
\index{sphere!as submanifold}
The sphere $S^{n-1} = \set{x \in \R^n : \abs{x}^2 = 1}$ is a smooth
embedded submanifold of\/ $\R^n$ of dimension $n - 1$.
\end{proposition}

\begin{proof}
Define $f \colon \R^n \to \R$ by $f(x) = \abs{x}^2 = \sum_{i=1}^n (x^i)^2$.
Then $\dd f_x = 2\sum_i x^i\, \dd x^i$, so $\dd f_x = 0$ only if $x = 0$.
Since $0 \notin f^{-1}(1)$, every point of $S^{n-1} = f^{-1}(1)$ is a
regular point.  Hence $1$ is a regular value, and
$S^{n-1} = f^{-1}(1)$ is an embedded submanifold of dimension
$n - 1$.
\end{proof}

\begin{corollary}\label{cor:tangent-sphere}
For $p \in S^{n-1}$, $T_pS^{n-1} = \ker \dd f_p = \set{v \in \R^n :
\sum_i p^i v^i = 0} = p^\perp$.
\end{corollary}

%% ── Orthogonal group ──────────────────────────────────────
\subsection{The orthogonal group}\label{ssec:orthogonal}

\begin{proposition}\label{prop:On-submanifold}%
\index{orthogonal group}\index{$O(n)$}
The orthogonal group $O(n) = \set{A \in M_n(\R) : A^T A = I_n}$ is a
smooth embedded submanifold of $M_n(\R) \cong \R^{n^2}$ of dimension
$\frac{n(n-1)}{2}$.  In particular, $O(n)$ is a Lie group.
\end{proposition}

\begin{proof}
Let $\mathrm{Sym}_n(\R)$ denote the space of $n \times n$ real symmetric
matrices, which has dimension $\frac{n(n+1)}{2}$.  Define
$f \colon M_n(\R) \to \mathrm{Sym}_n(\R)$ by $f(A) = A^T A$.
Then $O(n) = f^{-1}(I_n)$.

We compute the differential.  For $A \in M_n(\R)$ and $B \in T_A M_n(\R)
\cong M_n(\R)$,
\[
  \dd f_A(B) = \lim_{t \to 0} \frac{(A+tB)^T(A+tB) - A^T A}{t}
  = B^T A + A^T B.
\]
We must show $\dd f_A$ is surjective at every $A \in O(n)$.
Let $S \in \mathrm{Sym}_n(\R)$.  Set $B = \frac{1}{2}AS$. Then
\[
  \dd f_A(B) = \frac{1}{2}(AS)^T A + \frac{1}{2}A^T(AS)
  = \frac{1}{2}S^T \underbrace{A^T A}_{I_n}
    + \frac{1}{2}\underbrace{A^T A}_{I_n} S
  = \frac{1}{2}S + \frac{1}{2}S = S.
\]
Hence $\dd f_A$ is surjective for all $A \in O(n)$, so $I_n$ is a
regular value.  The regular value theorem gives $O(n)$ as an
embedded submanifold of dimension
$n^2 - \frac{n(n+1)}{2} = \frac{n(n-1)}{2}$.
\end{proof}

\begin{corollary}\label{cor:tangent-On}%
\index{orthogonal group!tangent space}
The tangent space of $O(n)$ at $A$ is
\[
  T_A O(n) = \ker \dd f_A
  = \set{B \in M_n(\R) : B^T A + A^T B = 0}.
\]
At the identity, $T_I O(n) = \set{B \in M_n(\R) : B^T + B = 0}
= \mathfrak{o}(n)$\index{$\mathfrak{o}(n)$}, the space of skew-symmetric
matrices.
\end{corollary}

%% ── Special linear group ──────────────────────────────────
\subsection{The special linear group}\label{ssec:SLn}

\begin{proposition}\label{prop:SLn-submanifold}%
\index{special linear group}\index{$\SL(n, \R)$}
The special linear group $\SL(n, \R) = \set{A \in M_n(\R) : \det A = 1}$
is a smooth embedded submanifold of $M_n(\R)$ of dimension $n^2 - 1$.
\end{proposition}

\begin{proof}
Let $f = \det \colon M_n(\R) \to \R$.  We use the well-known formula
\[
  \dd(\det)_A(B) = \det(A)\, \mathrm{tr}(A^{-1}B)
\]
for invertible~$A$.  If $A \in \SL(n, \R)$, then $\det A = 1 \neq 0$,
so $\dd f_A(B) = \mathrm{tr}(A^{-1}B)$.  This is surjective
(onto~$\R$): taking $B = A$ gives
$\dd f_A(A) = \mathrm{tr}(I_n) = n \neq 0$.
Thus $1$ is a regular value of $\det$, and
$\SL(n, \R) = \det^{-1}(1)$ is an embedded submanifold of dimension
$n^2 - 1$.
\end{proof}

\begin{corollary}\label{cor:tangent-SLn}%
\index{special linear group!tangent space}
$T_I\,\SL(n, \R) = \ker \dd(\det)_I
= \set{B \in M_n(\R) : \mathrm{tr}(B) = 0}
= \mathfrak{sl}(n, \R)$\index{$\mathfrak{sl}(n, \R)$}.
\end{corollary}

%% ── Special orthogonal group ──────────────────────────────
\subsection{Other classical groups}\label{ssec:classical-groups}

\begin{example}[$\SO(n)$]\label{ex:SOn}%
\index{special orthogonal group}\index{$\SO(n)$}
The special orthogonal group
$\SO(n) = O(n) \cap \SL(n, \R) = \set{A \in O(n) : \det A = 1}$ is an
open (hence embedded) submanifold of $O(n)$, since it is a connected
component ($O(n)$ has exactly two connected components, distinguished
by the sign of the determinant).  In particular,
$\dim \SO(n) = \frac{n(n-1)}{2}$, and $T_I\,\SO(n) = \mathfrak{o}(n)$.
\end{example}

\begin{example}[$\SU(n)$]\label{ex:SUn}%
\index{special unitary group}\index{$\SU(n)$}
A similar argument applies to the unitary group
$U(n) = \set{A \in M_n(\C) : A^*A = I_n}$, using the map
$f(A) = A^*A$ valued in Hermitian matrices, to show $U(n)$ is a smooth
submanifold of $M_n(\C) \cong \R^{2n^2}$ of dimension $n^2$.  The
special unitary group $\SU(n) = \set{A \in U(n) : \det A = 1}$ is
then a submanifold of dimension $n^2 - 1$.
\end{example}

%%%%%%%%%%%%%%%%%%%%%%%%%%%%%%%%%%%%%%%%%%%%%%%%%%%%%%%%%%%%%
\section{The constant rank theorem}\label{sec:constant-rank}
%%%%%%%%%%%%%%%%%%%%%%%%%%%%%%%%%%%%%%%%%%%%%%%%%%%%%%%%%%%%%

The submersion theorem generalises to maps of constant rank.

\begin{theorem}[Constant rank theorem]\label{thm:constant-rank}%
\index{constant rank theorem}
Let $f \colon M^m \to N^n$ be a smooth map of constant rank~$r$ on an
open set containing~$p$.  Then there exist charts $(U, \varphi)$ around
$p$ and $(V, \psi)$ around $f(p)$ such that
\[
  \psi \circ f \circ \varphi^{-1}(x^1, \dots, x^m)
  = (x^1, \dots, x^r, 0, \dots, 0).
\]
\end{theorem}

\begin{proof}[Proof sketch]
The proof proceeds in three steps.  First, by the rank assumption,
after permuting coordinates we may assume the upper-left $r \times r$
block of $Df(p)$ is invertible.  Second, define $\varphi$ by including
$f^1, \dots, f^r$ among the coordinate functions and applying the
inverse function theorem.  Third, in the new coordinates the map has
the form $(x', x'') \mapsto (x', g(x'))$ for some smooth~$g$ (since
the rank is exactly~$r$, the remaining components cannot depend on~$x''$).
A final coordinate change $\psi$ on the target that subtracts $g(x')$
puts the map into the stated normal form.  We omit the details and refer
to \cite[\S4.12]{Lee2013}.
\end{proof}

\begin{corollary}[Immersion theorem / canonical form]%
\label{cor:immersion-local}%
\index{immersion!local normal form}
If $f \colon M^m \to N^n$ is an immersion ($r = m \leq n$), there exist
charts $(U, \varphi)$ around $p$ and $(V, \psi)$ around $f(p)$ in which
$f$ has the form
\[
  (x^1, \dots, x^m) \;\longmapsto\; (x^1, \dots, x^m, 0, \dots, 0).
\]
In particular, the image $f(U)$ is locally a ``flat'' $m$-dimensional
slice in $N$.
\end{corollary}

\begin{remark}\label{rem:local-forms-summary}%
\index{local normal forms!summary}
Summarising the local normal forms for smooth maps of constant rank~$r$:
\[
\begin{array}{lcll}
  \textbf{Type} & \textbf{Rank} & \textbf{Normal form} \\
  \midrule
  \text{Submersion} & r = n & (x^1, \dots, x^m) \mapsto (x^1, \dots, x^n)
    & (m \geq n) \\
  \text{Immersion} & r = m & (x^1, \dots, x^m) \mapsto
    (x^1, \dots, x^m, 0, \dots, 0) & (m \leq n) \\
  \text{General, rank $r$} & r \leq \min(m,n) &
    (x^1, \dots, x^m) \mapsto (x^1, \dots, x^r, 0, \dots, 0) & \\
\end{array}
\]
All three are instances of the constant rank theorem.
\end{remark}

\begin{corollary}\label{cor:constant-rank-level-set}%
\index{constant rank theorem!level set}
If $f \colon M \to N$ has constant rank~$r$ on a neighbourhood of
$f^{-1}(q)$, then $f^{-1}(q)$ (if non-empty) is an embedded submanifold
of dimension $m - r$.
\end{corollary}

%%%%%%%%%%%%%%%%%%%%%%%%%%%%%%%%%%%%%%%%%%%%%%%%%%%%%%%%%%%%%
\section{Submanifolds with boundary}\label{sec:submanifold-boundary}
%%%%%%%%%%%%%%%%%%%%%%%%%%%%%%%%%%%%%%%%%%%%%%%%%%%%%%%%%%%%%

\begin{definition}[Smooth manifold with boundary]\label{def:mfld-boundary}%
\index{manifold with boundary}\index{boundary!of manifold}
A \textbf{smooth $n$-manifold with boundary} is a second-countable
Hausdorff space $M$ locally modelled on the half-space
$\mathbb{H}^n = \set{x \in \R^n : x^n \geq 0}$, with smooth transition
maps.  The \textbf{boundary}
\[
  \partial M = \set{p \in M : \varphi(p) \in \partial\mathbb{H}^n
  = \R^{n-1} \times \{0\}}
\]
for any chart $(U, \varphi)$, is a smooth $(n-1)$-manifold (without
boundary).  The \textbf{interior} is
$\mathrm{Int}(M) = M \setminus \partial M$, an $n$-manifold without
boundary.
\end{definition}

\begin{proposition}\label{prop:boundary-submanifold}
If $M$ is a smooth $n$-manifold with boundary, then $\partial M$ is a
smooth embedded submanifold of $M$ of dimension $n-1$.
\end{proposition}

\begin{definition}[Neat submanifold]\label{def:neat-submanifold}%
\index{submanifold!neat}
A submanifold $S$ of a manifold with boundary $M$ is \textbf{neat} if
$\partial S = S \cap \partial M$ and $S$ is transverse to $\partial M$.
\end{definition}

\begin{example}\label{ex:half-disk}
Consider the closed half-disk
$D^n_+ = \set{x \in \R^n : \abs{x} \leq 1,\; x^n \geq 0}$.
Its boundary consists of two pieces:
$\partial D^n_+ = S^{n-1}_+ \cup D^{n-1}$, where
$S^{n-1}_+ = \set{x \in S^{n-1} : x^n \geq 0}$ and
$D^{n-1} = \set{x \in D^n : x^n = 0}$.
The equatorial disk $D^{n-1}$ sits neatly in $D^n_+$.
\end{example}

%%%%%%%%%%%%%%%%%%%%%%%%%%%%%%%%%%%%%%%%%%%%%%%%%%%%%%%%%%%%%
\section{Transversality}\label{sec:transversality}
%%%%%%%%%%%%%%%%%%%%%%%%%%%%%%%%%%%%%%%%%%%%%%%%%%%%%%%%%%%%%

\begin{definition}[Transversality]\label{def:transversality}%
\index{transversality}
Let $f \colon M \to N$ be smooth and $S \subset N$ an embedded
submanifold.  We say $f$ is \textbf{transverse to $S$}, written
$f \pitchfork S$, if for every $p \in f^{-1}(S)$,
\[
  \im(\dd f_p) + T_{f(p)}S \;=\; T_{f(p)}N.
\]
Two submanifolds $S_1, S_2 \subset N$ are \textbf{transverse}
($S_1 \pitchfork S_2$) if the inclusion $\iota_1 \colon S_1
\hookrightarrow N$ is transverse to $S_2$.
\end{definition}

\begin{theorem}[Transverse preimage theorem]%
\label{thm:transverse-preimage}%
\index{transversality!preimage theorem}
If $f \colon M \to N$ is smooth and $f \pitchfork S$ for an embedded
submanifold $S \subset N$ of codimension~$k$, then $f^{-1}(S)$ is
either empty or an embedded submanifold of $M$ of codimension~$k$.
Moreover, $T_p(f^{-1}(S)) = (\dd f_p)^{-1}(T_{f(p)}S)$ for
$p \in f^{-1}(S)$.
\end{theorem}

\begin{proof}
We reduce to the regular value theorem.  Locally near $q = f(p) \in S$,
choose a slice chart $(V, \psi)$ for $S$ in $N$ so that
$\psi(S \cap V) = \psi(V) \cap (\R^{n-k} \times \{0\})$.  Let
$\pi \colon \R^n \to \R^k$ be the projection onto the last $k$
coordinates and set $g = \pi \circ \psi \circ f \colon f^{-1}(V) \to \R^k$.
Then $f^{-1}(S) \cap f^{-1}(V) = g^{-1}(0)$.  The transversality
condition $\im(\dd f_p) + T_qS = T_qN$ translates to $\dd g_p$ being
surjective.  Hence $0$ is a regular value of $g$ and $g^{-1}(0)$ is an
embedded submanifold of codimension~$k$.
\end{proof}

\begin{corollary}\label{cor:transverse-submanifolds}
If $S_1, S_2 \subset N$ are embedded submanifolds with
$S_1 \pitchfork S_2$, then $S_1 \cap S_2$ is an embedded submanifold
of $N$ of dimension $\dim S_1 + \dim S_2 - \dim N$ (when non-empty).
\end{corollary}

\begin{example}\label{ex:transverse-planes}
In $\R^3$, the plane $S_1 = \set{z = 0}$ and the plane
$S_2 = \set{x = 0}$ are transverse: at any point $p \in S_1 \cap S_2$,
$T_pS_1 + T_pS_2 = \R^3$.  Their intersection is the $y$-axis, a
$1$-dimensional submanifold, consistent with $\dim S_1 + \dim S_2 -
\dim N = 2 + 2 - 3 = 1$.
\end{example}

\begin{example}\label{ex:non-transverse}
In $\R^2$, the $x$-axis $S_1 = \set{y = 0}$ and the curve
$S_2 = \set{y = x^2}$ are \emph{not} transverse at the origin, since
$T_0S_1 = T_0S_2 = \R \times \{0\}$.  The intersection $\{0\}$ is a
point but this is ``accidental''; a small perturbation of $S_2$ can
create two intersection points, one, or none.
\end{example}

The following illustration contrasts transverse and non-transverse
intersections:

\begin{center}
\begin{tikzpicture}[scale=1.0, >=Stealth]
  %% Transverse case
  \begin{scope}[shift={(-3.5, 0)}]
    \draw[thick, red!70!black] (-1.5, 0) -- (1.5, 0)
      node[right, font=\small] {$S_1$};
    \draw[thick, blue!70!black] (-0.5, -1.2) -- (0.5, 1.2)
      node[above, font=\small] {$S_2$};
    \fill (0, 0) circle (1.5pt);
    \node[below, font=\small] at (0, -1.5) {transverse};
  \end{scope}
  %% Non-transverse case
  \begin{scope}[shift={(2.0, 0)}]
    \draw[thick, red!70!black] (-1.5, 0) -- (1.5, 0)
      node[right, font=\small] {$S_1$};
    \draw[thick, blue!70!black, domain=-1.2:1.2, samples=60]
      plot (\x, {\x*\x});
    \node[blue!70!black, above, font=\small] at (0.8, 1.0) {$S_2$};
    \fill (0, 0) circle (1.5pt);
    \node[below, font=\small] at (0, -1.5) {not transverse};
  \end{scope}
\end{tikzpicture}
\end{center}

%%%%%%%%%%%%%%%%%%%%%%%%%%%%%%%%%%%%%%%%%%%%%%%%%%%%%%%%%%%%%
\section{Tubular neighbourhoods}\label{sec:tubular}
%%%%%%%%%%%%%%%%%%%%%%%%%%%%%%%%%%%%%%%%%%%%%%%%%%%%%%%%%%%%%

\begin{definition}[Normal bundle]\label{def:normal-bundle}%
\index{normal bundle}
Let $S \subset M$ be an embedded submanifold.  Choose a Riemannian metric
$g$ on $M$ (which always exists by a partition of unity argument).  The
\textbf{normal bundle} of $S$ in $M$ is
\[
  NS \;=\; \bigsqcup_{p \in S} N_pS,
  \qquad
  N_pS = (T_pS)^\perp \subset T_pM,
\]
where the orthogonal complement is taken with respect to~$g_p$.
It is a smooth vector bundle over $S$ of rank $\dim M - \dim S$.
\end{definition}

\begin{theorem}[Tubular neighbourhood theorem]\label{thm:tubular}%
\index{tubular neighbourhood theorem}
Let $S \subset M$ be a compact embedded submanifold.  Then there exists an
open neighbourhood $U \supset S$ in $M$ and a diffeomorphism
$\Phi \colon NS \supset V \xrightarrow{\;\sim\;} U$,
where $V$ is an open neighbourhood of the zero section in~$NS$, such that
$\Phi|_S = \Id_S$ (identifying $S$ with the zero section).

The image $U = \Phi(V)$ is called a
\textbf{tubular neighbourhood}\index{tubular neighbourhood} of $S$ in $M$.
\end{theorem}

The proof uses the exponential map of the Riemannian metric restricted
to normal vectors.  We omit it here and refer to
\cite[\S6.2]{Lee2013}.

\begin{center}
\begin{tikzpicture}[scale=1.0, >=Stealth]
  %% Submanifold S (a curve)
  \draw[very thick, red!70!black]
    (-3.0, 0) .. controls (-1.5, 0.6) and (1.5, -0.6) .. (3.0, 0);
  \node[red!70!black] at (3.3, 0) {$S$};

  %% Tubular neighbourhood (band around S)
  \draw[thick, blue!50!black, fill=blue!8, fill opacity=0.5]
    (-3.0, 0.9) .. controls (-1.5, 1.5) and (1.5, 0.3) .. (3.0, 0.9)
    -- (3.0, -0.9) .. controls (1.5, -1.5) and (-1.5, -0.3) .. (-3.0, -0.9)
    -- cycle;
  \node[blue!50!black] at (-3.5, 0.5) {$U$};

  %% Normal vectors at a few points
  \draw[->, thick, green!50!black] (-1.0, 0.25) -- (-1.0, 0.9)
    node[above, font=\small] {$N_pS$};
  \draw[->, thick, green!50!black] (-1.0, 0.25) -- (-1.0, -0.4);
  \fill (-1.0, 0.25) circle (1.5pt) node[below left, font=\small] {$p$};

  \draw[->, thick, green!50!black] (1.5, -0.25) -- (1.5, 0.4);
  \draw[->, thick, green!50!black] (1.5, -0.25) -- (1.5, -0.9);
  \fill (1.5, -0.25) circle (1.5pt);
\end{tikzpicture}
\end{center}

\begin{remark}\label{rem:tubular-non-compact}
The compactness assumption can be relaxed: the theorem holds for any
closed (i.e., properly embedded) submanifold, though the proof requires
more care with the width of the tubular neighbourhood.
\end{remark}

%%%%%%%%%%%%%%%%%%%%%%%%%%%%%%%%%%%%%%%%%%%%%%%%%%%%%%%%%%%%%
\section{Exercises}
%%%%%%%%%%%%%%%%%%%%%%%%%%%%%%%%%%%%%%%%%%%%%%%%%%%%%%%%%%%%%

\begin{exercise}\label{exo:Stiefel}%
\index{Stiefel manifold}
The \textbf{Stiefel manifold} $V_k(\R^n)$ is the set of orthonormal
$k$-frames in $\R^n$:
\[
  V_k(\R^n) = \set{A \in M_{n \times k}(\R) : A^T A = I_k}.
\]
Show that $V_k(\R^n)$ is a smooth embedded submanifold of
$M_{n \times k}(\R)$ of dimension $nk - \frac{k(k+1)}{2}$.
\textit{Hint:} mimic the proof for $O(n)$ using the map
$A \mapsto A^T A$.
\end{exercise}

\begin{exercise}\label{exo:graph-submanifold}%
\index{submanifold!graph}
Let $f \colon M \to N$ be smooth.  Prove that the graph
$\Gamma_f = \set{(p, f(p)) : p \in M} \subset M \times N$ is an
embedded submanifold diffeomorphic to $M$.
\end{exercise}

\begin{exercise}\label{exo:product-of-spheres}
Show that $S^m \times S^n \subset \R^{m+1} \times \R^{n+1}$ is an
embedded submanifold of $\R^{m+n+2}$, and determine its tangent space
at $(p, q)$.
\end{exercise}

\begin{exercise}\label{exo:hyperboloid}%
\index{hyperboloid}
Consider the hyperboloid
$H = \set{(x, y, z) \in \R^3 : x^2 + y^2 - z^2 = 1}$.
\begin{enumerate}[label=(\alph*)]
  \item Show that $H$ is a smooth $2$-dimensional submanifold of $\R^3$.
  \item Compute $T_{(1,0,0)}H$.
  \item Is $H$ compact?  Connected?
\end{enumerate}
\end{exercise}

\begin{exercise}\label{exo:transverse-curves}
Let $C_1, C_2 \subset \R^2$ be smooth curves (embedded $1$-dimensional
submanifolds).  Show that $C_1 \pitchfork C_2$ if and only if at every
point $p \in C_1 \cap C_2$, the tangent lines $T_pC_1$ and $T_pC_2$
are distinct.  In this case, $C_1 \cap C_2$ is a discrete set of points.
\end{exercise}

\begin{exercise}\label{exo:submersion-open}%
\index{submersion!is an open map}
Prove that every submersion $f \colon M \to N$ is an open map.
\textit{Hint:} use the submersion theorem.
\end{exercise}

\begin{exercise}\label{exo:regular-level-set-orientable}
Let $f \colon M \to \R$ be smooth with $c$ a regular value.  Show that
the level set $f^{-1}(c)$ has a trivial normal bundle and is therefore
orientable (assuming $M$ is orientable).
\end{exercise}

\begin{exercise}\label{exo:SLn-connected}%
\index{$\SL(n, \R)$!connected}
Show that $\SL(n, \R)$ is connected.
\textit{Hint:} use the fact that every matrix in $\SL(n, \R)$ is a
product of elementary matrices.
\end{exercise}

\begin{exercise}\label{exo:preimage-submanifold}
Let $f \colon \R^3 \to \R^2$ be defined by $f(x, y, z) = (x^2 + y^2, z)$.
\begin{enumerate}[label=(\alph*)]
  \item Find all critical points and critical values of $f$.
  \item For which $(a, b) \in \R^2$ is $f^{-1}(a, b)$ a smooth manifold?
    Describe the topology.
\end{enumerate}
\end{exercise}

\begin{exercise}\label{exo:transverse-intersection-codim}
Let $S_1$ and $S_2$ be compact embedded submanifolds of $\R^n$ with
$\dim S_1 + \dim S_2 < n$.  Prove that if $S_1 \pitchfork S_2$, then
$S_1 \cap S_2 = \varnothing$.
\end{exercise}

\bigskip
\begin{center}
\rule{0.5\textwidth}{0.4pt}
\end{center}

\noindent\textit{References for Chapters~\ref{ch:tangent-cotangent}
and~\ref{ch:submanifolds}:}
J.\,M.~Lee, \textit{Introduction to Smooth Manifolds}, 2nd ed.,
Springer, 2013 \cite{Lee2013};
V.~Guillemin and A.~Pollack, \textit{Differential Topology}, AMS Chelsea,
2010 \cite{GP2010};
L.\,W.~Tu, \textit{An Introduction to Manifolds}, 2nd ed., Springer,
2011 \cite{Tu2011};
M.\,W.~Hirsch, \textit{Differential Topology}, Springer GTM 33, 1976
\cite{Hirsch1976}.

% === END OF CHUNK ch3-4 ===
% === START OF CHUNK ch5-6 ===

%%%%%%%%%%%%%%%%%%%%%%%%%%%%%%%%%%%%%%%%%%%%%%%%%%%%%%%%%%%%%
%%%%%%%%%%%%%%%%%%%%%%%%%%%%%%%%%%%%%%%%%%%%%%%%%%%%%%%%%%%%%
%%
%%  CHAPTER 5 --- TRANSVERSALITY
%%
%%%%%%%%%%%%%%%%%%%%%%%%%%%%%%%%%%%%%%%%%%%%%%%%%%%%%%%%%%%%%
%%%%%%%%%%%%%%%%%%%%%%%%%%%%%%%%%%%%%%%%%%%%%%%%%%%%%%%%%%%%%

\chapter{Transversality}\label{ch:transversality}
\minitoc\vspace{1cm}

Transversality is one of the most powerful concepts in differential topology.
It provides a robust framework for understanding how submanifolds and maps
interact in \emph{generic} position. Where the implicit function theorem
tells us when a level set is a manifold, transversality vastly generalizes
this, showing that ``most'' geometric configurations are well-behaved.
The key insight, due to Thom, is that transversal intersections are not
only clean (yielding submanifolds) but also \emph{stable} and \emph{generic}.

% ============================================================
\section{Transverse maps and submanifolds}\label{sec:transverse-maps}
% ============================================================

\begin{definition}[Transversality of a map to a submanifold]\label{def:transverse-map}
\index{transversality!of a map}
Let $f \colon M \to N$ be a smooth map between smooth manifolds, and let
$W \subseteq N$ be a smooth submanifold. We say that $f$ is
\textbf{transverse}\index{transverse} to $W$, written
\[
  f \pitchfork W,
\]
if for every $x \in f^{-1}(W)$ we have
\[
  \dd f_x(T_x M) + T_{f(x)} W = T_{f(x)} N.
\]
In other words, the image of $\dd f_x$ together with the tangent space
of $W$ at $f(x)$ spans the entire tangent space of $N$ at $f(x)$.
\end{definition}

\begin{definition}[Transverse intersection of submanifolds]\label{def:transverse-submanifolds}
\index{transversality!of submanifolds}
Let $M$ be a smooth manifold, and let $A, B \subseteq M$ be smooth
submanifolds. We say that $A$ and $B$ \textbf{intersect transversally},
written
\[
  A \pitchfork B,
\]
if for every $x \in A \cap B$,
\[
  T_x A + T_x B = T_x M.
\]
\end{definition}

\begin{remark}\label{rem:transverse-inclusion}
The condition $A \pitchfork B$ is a special case of
Definition~\ref{def:transverse-map}: if $\iota \colon A \hookrightarrow M$
is the inclusion, then $A \pitchfork B$ if and only if $\iota \pitchfork B$.
\end{remark}

\begin{remark}\label{rem:transverse-empty}
If $f^{-1}(W) = \varnothing$, then $f \pitchfork W$ is satisfied
vacuously. In particular, disjoint submanifolds are always transverse.
\end{remark}

\begin{remark}\label{rem:transverse-submersion}
If $f \colon M \to N$ is a submersion, then $f \pitchfork W$ for every
submanifold $W \subseteq N$, since $\dd f_x(T_x M) = T_{f(x)} N$.
\end{remark}

% --- TikZ figure: transverse vs non-transverse intersection ---
\begin{figure}[ht]
\centering
\begin{tikzpicture}[scale=1.2]
  % Transverse case
  \begin{scope}[xshift=-3.5cm]
    \draw[thick, blue!70!black] (-1.5,0) .. controls (-0.5,0.3) and (0.5,-0.3) .. (1.5,0)
      node[right]{\small $A$};
    \draw[thick, red!70!black] (0,-1.3) .. controls (0.2,-0.3) and (-0.2,0.3) .. (0,1.3)
      node[above]{\small $B$};
    \filldraw[black] (0,0) circle (1.5pt);
    \node[below=8pt] at (0,-1.3) {\textbf{Transverse}: $A \pitchfork B$};
  \end{scope}
  % Non-transverse case
  \begin{scope}[xshift=3.5cm]
    \draw[thick, blue!70!black] (-1.5,0) .. controls (-0.5,0.4) and (0.5,0.4) .. (1.5,0)
      node[right]{\small $A$};
    \draw[thick, red!70!black] (-1.5,0.05) .. controls (-0.5,-0.35) and (0.5,-0.35) .. (1.5,0.05)
      node[right]{\small $B$};
    \filldraw[black] (-1.15,0.12) circle (1.5pt);
    \filldraw[black] (1.15,0.12) circle (1.5pt);
    \node[below=8pt] at (0,-0.6) {\textbf{Not transverse}};
  \end{scope}
\end{tikzpicture}
\caption{Transverse versus non-transverse intersections of curves in the plane.}
\label{fig:transverse-nontransverse}
\end{figure}

\begin{example}\label{ex:transverse-lines}
\index{transversality!examples}
In $\R^2$, two lines through the origin are transverse if and only if they
are distinct. If $A = \set{(x,0) : x \in \R}$ and
$B = \set{(x,mx) : x \in \R}$ with $m \neq 0$, then at the origin
$T_0 A + T_0 B = \R^2$, so $A \pitchfork B$. When $m = 0$, we have
$T_0 A = T_0 B$, so the intersection is not transverse.
\end{example}

\begin{example}\label{ex:transverse-spheres}
In $\R^3$, consider two spheres $S_1 = S^2(p_1, r_1)$ and
$S_2 = S^2(p_2, r_2)$ with $S_1 \cap S_2 \neq \varnothing$. They
are transverse if and only if they do not share a common tangent plane
at any intersection point, which happens precisely when they are not
internally tangent. When $S_1 \pitchfork S_2$, the intersection
$S_1 \cap S_2$ is a circle (a $1$-dimensional submanifold of $\R^3$).
\end{example}

\begin{example}\label{ex:transverse-graph}
Let $f \colon \R^n \to \R$ be smooth. The graph
$\Gamma_f = \set{(x, f(x)) : x \in \R^n} \subset \R^{n+1}$ is
transverse to the hyperplane $W = \R^n \times \set{c}$ if and only
if $c$ is a regular value of $f$. Indeed, at a point $(x_0, c)$
with $f(x_0) = c$, the tangent space $T_{(x_0,c)}\Gamma_f$ is spanned
by the vectors $(e_i, \partial_i f(x_0))$, and $T_{(x_0,c)}W$ is
$\R^n \times \set{0}$. These span $\R^{n+1}$ if and only if
$\dd f_{x_0} \neq 0$.
\end{example}

\begin{example}\label{ex:transverse-diagonal}
Let $\Delta = \set{(x,x) : x \in M} \subset M \times M$ be the diagonal.
If $f, g \colon N \to M$ are smooth, define
$h \colon N \to M \times M$ by $h(x) = (f(x), g(x))$.
Then $h \pitchfork \Delta$ if and only if for every $x \in N$ with
$f(x) = g(x)$,
\[
  \dd f_x(T_x N) + \dd g_x(T_x N) = T_{f(x)} M.
\]
This reformulation is central to intersection theory.
\end{example}

% ============================================================
\section{The Transversality Theorem}\label{sec:transversality-theorem}
% ============================================================

The following theorem is the fundamental structural result of transversality.
It generalizes the preimage theorem (regular value theorem) to the setting
of transverse maps.

\begin{theorem}[Transversality theorem]\label{thm:transversality}
\index{transversality theorem}
Let $f \colon M \to N$ be a smooth map and $W \subseteq N$ a smooth
submanifold of codimension $k$. If $f \pitchfork W$, then:
\begin{enumerate}[label=(\roman*)]
  \item $f^{-1}(W)$ is a smooth submanifold of $M$ of codimension $k$
    (hence $\dim f^{-1}(W) = \dim M - k$, provided it is nonempty).
  \item For every $x \in f^{-1}(W)$, the tangent space satisfies
    \[
      T_x(f^{-1}(W)) = (\dd f_x)^{-1}(T_{f(x)} W).
    \]
\end{enumerate}
\end{theorem}

\begin{proof}
Let $x_0 \in f^{-1}(W)$ and set $y_0 = f(x_0) \in W$.
Since $W$ is a submanifold of codimension $k$ in $N$, there exists a
neighborhood $V$ of $y_0$ in $N$ and a smooth submersion
$\varphi \colon V \to \R^k$ such that
\[
  W \cap V = \varphi^{-1}(0).
\]
Consider the composition $g = \varphi \circ f \colon f^{-1}(V) \to \R^k$.
We claim that $0$ is a regular value of $g$.

Let $x \in g^{-1}(0) = f^{-1}(W \cap V)$. We must show that
$\dd g_x \colon T_x M \to \R^k$ is surjective. By the chain rule,
$\dd g_x = \dd\varphi_{f(x)} \circ \dd f_x$. Let $w \in \R^k$. Since
$\varphi$ is a submersion, $\dd\varphi_{f(x)}$ is surjective; but we need
more. Since $\ker(\dd\varphi_{f(x)}) = T_{f(x)}W$ (as $W = \varphi^{-1}(0)$
locally), the map $\dd\varphi_{f(x)}$ induces an isomorphism
\[
  \overline{\dd\varphi}_{f(x)} \colon T_{f(x)}N / T_{f(x)}W
    \xrightarrow{\;\sim\;} \R^k.
\]
By the transversality condition $f \pitchfork W$, we have
$\dd f_x(T_x M) + T_{f(x)}W = T_{f(x)}N$. Passing to the quotient,
the map $\dd f_x$ composed with the projection
$\pi \colon T_{f(x)}N \to T_{f(x)}N / T_{f(x)}W$ is surjective:
\[
  \pi \circ \dd f_x \colon T_x M \to T_{f(x)}N / T_{f(x)}W
  \quad \text{is surjective.}
\]
Therefore
\[
  \dd g_x = \dd\varphi_{f(x)} \circ \dd f_x
  = \overline{\dd\varphi}_{f(x)} \circ \pi \circ \dd f_x
\]
is a composition of a surjection followed by an isomorphism, hence surjective.
Thus $0$ is a regular value of $g$.

By the preimage theorem (regular value theorem), $g^{-1}(0)$ is a smooth
submanifold of $f^{-1}(V)$ of codimension $k$. Since
$g^{-1}(0) = f^{-1}(W) \cap f^{-1}(V)$, this shows that $f^{-1}(W)$ is
locally (near $x_0$) a smooth submanifold of codimension $k$. As $x_0$ was
arbitrary, $f^{-1}(W)$ is a smooth submanifold of $M$ of codimension $k$.

For part (ii), the tangent space to $g^{-1}(0)$ at $x$ is
$\ker(\dd g_x) = \ker(\dd\varphi_{f(x)} \circ \dd f_x)$, which consists
of those $v \in T_x M$ such that $\dd f_x(v) \in \ker(\dd\varphi_{f(x)})
= T_{f(x)}W$. Hence $T_x(f^{-1}(W)) = (\dd f_x)^{-1}(T_{f(x)}W)$.
\end{proof}

\begin{corollary}\label{cor:transverse-submanifolds-intersection}
\index{transverse intersection}
If $A$ and $B$ are smooth submanifolds of $M$ with $A \pitchfork B$, then
$A \cap B$ is a smooth submanifold of $M$ with
\[
  \operatorname{codim}(A \cap B) = \operatorname{codim}(A)
    + \operatorname{codim}(B).
\]
Equivalently, $\dim(A \cap B) = \dim A + \dim B - \dim M$ (when $A \cap B
\neq \varnothing$).
\end{corollary}

\begin{proof}
Apply Theorem~\ref{thm:transversality} to the inclusion
$\iota \colon A \hookrightarrow M$ and the submanifold $B \subseteq M$.
Since $\iota \pitchfork B$, we have $\iota^{-1}(B) = A \cap B$ is a
submanifold of $A$ of codimension $\operatorname{codim}_M(B)$. Hence
$\operatorname{codim}_M(A \cap B) = \operatorname{codim}_A(A \cap B)
+ \operatorname{codim}_M(A) = \operatorname{codim}_M(B) + \operatorname{codim}_M(A)$.
\end{proof}

\begin{corollary}\label{cor:dimension-formula}
If $M$ has dimension $m$, $\dim A = a$, and $\dim B = b$ with $A \pitchfork B$,
then $\dim(A \cap B) = a + b - m$. In particular, if $a + b < m$, then
$A \cap B = \varnothing$ (since a manifold cannot have negative dimension).
\end{corollary}

\begin{example}\label{ex:transverse-curves-surface}
Two curves (dimension $1$) on a surface (dimension $2$) meeting
transversally intersect in a $0$-dimensional submanifold, i.e., a
discrete set of points. If the curves are compact, the intersection
is finite.
\end{example}

\begin{example}\label{ex:transverse-hyperplanes}
In $\R^n$, $k$ hyperplanes in general position (pairwise transverse
intersections propagating) meet in a submanifold of dimension $n - k$
(an affine subspace). Concretely, $k \leq n$ hyperplanes through the
origin given by linearly independent linear forms intersect in an
$(n-k)$-dimensional subspace.
\end{example}

% --- TikZ figure: preimage under transverse map ---
\begin{figure}[ht]
\centering
\begin{tikzpicture}[scale=1.1,>=Stealth]
  % Domain M
  \draw[thick, fill=blue!5, rounded corners=12pt]
    (-4.5,1.8) -- (-1.5,1.8) -- (-1.5,-1.8) -- (-4.5,-1.8) -- cycle;
  \node at (-3,2.1) {$M$};
  % Preimage curve
  \draw[very thick, green!50!black] (-4.0,-0.8) .. controls (-3.2,0.5)
    and (-2.8,-0.5) .. (-2.0,0.6);
  \node[green!50!black] at (-2.5,1.1) {\small $f^{-1}(W)$};

  % Arrow
  \draw[->,thick] (-1.2,0) -- (0.3,0) node[midway,above]{$f$};

  % Codomain N
  \draw[thick, fill=red!5, rounded corners=12pt]
    (0.6,1.8) -- (3.6,1.8) -- (3.6,-1.8) -- (0.6,-1.8) -- cycle;
  \node at (2.1,2.1) {$N$};
  % Submanifold W
  \draw[very thick, red!70!black] (1.0,-1.2) -- (3.2,1.4);
  \node[red!70!black] at (3.4,1.6) {\small $W$};
\end{tikzpicture}
\caption{When $f \pitchfork W$, the preimage $f^{-1}(W)$ is a submanifold
  of $M$ with the same codimension as $W$ in $N$.}
\label{fig:transverse-preimage}
\end{figure}

% ============================================================
\section{Stability of transversality}\label{sec:stability-transversality}
% ============================================================

Transversality is a robust condition: it persists under small perturbations.
This stability is crucial for applications and underpins the genericity results
that follow.

\begin{theorem}[Stability of transverse intersections]\label{thm:stability-transversality}
\index{stability!of transversality}
Let $f \colon M \to N$ be a smooth map with $M$ compact, and let
$W \subseteq N$ be a closed submanifold. If $f \pitchfork W$, then there
exists a neighborhood $\mathcal{U}$ of $f$ in $C^\infty(M, N)$
(in the $C^1$ topology) such that every $g \in \mathcal{U}$ satisfies
$g \pitchfork W$.
\end{theorem}

\begin{proof}[Proof sketch]
The transversality condition $\dd f_x(T_x M) + T_{f(x)}W = T_{f(x)}N$ is
an open condition on the $1$-jet of $f$ at each point. Since $M$ is compact,
$f^{-1}(W)$ is compact, and the finitely many local conditions can be
simultaneously satisfied in a $C^1$-neighborhood.
\end{proof}

\begin{remark}\label{rem:stability-noncompact}
When $M$ is not compact, stability still holds in the strong (Whitney)
$C^\infty$ topology, but not necessarily in the weak topology.
\end{remark}

% ============================================================
\section{The Thom Transversality Theorem}\label{sec:thom-transversality}
% ============================================================

The Thom transversality theorem asserts that transversality is the generic
situation. This is one of the most important results in differential
topology, providing a systematic way to achieve transversality by
arbitrarily small perturbations.

\begin{definition}[Residual and generic sets]\label{def:residual-generic}
\index{residual set}\index{generic property}
A subset of a topological space is \textbf{residual} if it contains a
countable intersection of open dense sets. A property of smooth maps
$M \to N$ is called \textbf{generic} if the set of maps satisfying it
is residual in $C^\infty(M, N)$ (with the strong topology). By the
Baire category theorem, residual subsets of complete metric spaces are
dense.
\end{definition}

\begin{theorem}[Thom Transversality Theorem]\label{thm:thom-transversality}
\index{Thom transversality theorem}
Let $M$ and $N$ be smooth manifolds and $W \subseteq N$ a smooth submanifold.
Then the set
\[
  \pitchfork(M, N; W) \;=\; \set{f \in C^\infty(M, N) : f \pitchfork W}
\]
is residual (and hence dense) in $C^\infty(M, N)$ equipped with the
strong Whitney $C^\infty$ topology. If $W$ is closed, then
$\pitchfork(M, N; W)$ is open and dense.
\end{theorem}

\begin{proof}[Proof sketch]
The proof uses Sard's theorem in a parametric fashion. We sketch the
main ideas.

\medskip\noindent
\textbf{Step 1: Local representation.}
Cover $M$ by countably many coordinate charts $\set{U_i}$. It suffices to
show that for each $i$, the set of maps $f$ with $f|_{U_i} \pitchfork W$
is open and dense, since $\pitchfork(M, N; W)$ is the countable intersection
of these sets.

\medskip\noindent
\textbf{Step 2: Parametric families.}
Working in local coordinates, one constructs a smooth family
$F \colon M \times S \to N$ (where $S$ is an open subset of some $\R^p$)
such that $F_s = F(\cdot, s)$ gives a ``perturbation'' of $f$ for each
$s \in S$, and such that $F$ itself is a submersion near $F^{-1}(W)$
(and hence $F \pitchfork W$).

\medskip\noindent
\textbf{Step 3: Application of Sard's theorem.}
By Theorem~\ref{thm:transversality}, $F^{-1}(W)$ is a submanifold.
Consider the projection $\pi \colon F^{-1}(W) \to S$. By Sard's theorem
(\ref{thm:sard-recall} below), the set of regular values of $\pi$ has
full measure in $S$, hence is dense. One verifies that $s$ is a regular
value of $\pi$ if and only if $F_s \pitchfork W$. The density of
regular values gives the density of transverse maps.

\medskip\noindent
\textbf{Step 4: Openness.}
When $W$ is closed, the openness follows from the stability of
transversality (Theorem~\ref{thm:stability-transversality}), applied
locally and patched using compactness of the relevant subsets.
\end{proof}

\begin{theorem}[Sard's Theorem --- recalled]\label{thm:sard-recall}
\index{Sard's theorem}
Let $f \colon M \to N$ be a smooth map. The set of critical values of $f$
has measure zero in $N$. In particular, the set of regular values is dense.
\end{theorem}

\begin{remark}\label{rem:thom-meaning}
The philosophical content of the Thom transversality theorem is profound:
\emph{any smooth map can be perturbed by an arbitrarily small amount to
become transverse to any given submanifold}. This means that transversality
is the ``expected'' (generic) situation, and non-transverse intersections
are ``exceptional.''
\end{remark}

\begin{corollary}\label{cor:generic-transverse-submanifolds}
Let $A \subseteq M$ be a compact submanifold and $B \subseteq M$ a closed
submanifold. Then $A$ can be made transverse to $B$ by an arbitrarily
small perturbation: there exist embeddings $\iota' \colon A \to M$
arbitrarily $C^\infty$-close to the inclusion $\iota \colon A
\hookrightarrow M$ such that $\iota'(A) \pitchfork B$.
\end{corollary}

% --- TikZ figure: perturbing to transversality ---
\begin{figure}[ht]
\centering
\begin{tikzpicture}[scale=1.2]
  % Before perturbation
  \begin{scope}[xshift=-3.5cm]
    \draw[thick, blue!70!black] (-1.5,0.02)
      .. controls (-0.5,0.5) and (0.5,0.5) .. (1.5,0.02);
    \draw[thick, red!70!black] (-1.5,-0.02)
      .. controls (-0.5,-0.5) and (0.5,-0.5) .. (1.5,-0.02);
    \node[blue!70!black] at (1.7,0.3) {\small $A$};
    \node[red!70!black] at (1.7,-0.3) {\small $B$};
    \node[below] at (0,-1.0) {Non-transverse};
  \end{scope}
  % Arrow
  \draw[->,thick,decorate,decoration={snake,amplitude=1.5pt,segment length=6pt}]
    (-1.2,0) -- (0.3,0)
    node[midway,above=4pt]{\small perturb};
  % After perturbation
  \begin{scope}[xshift=3.5cm]
    \draw[thick, blue!70!black, dashed, opacity=0.4] (-1.5,0.02)
      .. controls (-0.5,0.5) and (0.5,0.5) .. (1.5,0.02);
    \draw[thick, blue!70!black] (-1.5,0.3)
      .. controls (-0.5,0.1) and (0.5,-0.3) .. (1.5,0.2);
    \draw[thick, red!70!black] (-1.5,-0.02)
      .. controls (-0.5,-0.5) and (0.5,-0.5) .. (1.5,-0.02);
    \node[blue!70!black] at (1.7,0.5) {\small $A'$};
    \node[red!70!black] at (1.7,-0.3) {\small $B$};
    \filldraw[black] (0.25,-0.17) circle (1.5pt);
    \node[below] at (0,-1.0) {Transverse: $A' \pitchfork B$};
  \end{scope}
\end{tikzpicture}
\caption{A small perturbation of $A$ achieves transversality with $B$.}
\label{fig:perturbation-transversality}
\end{figure}

% ============================================================
\section{Genericity via Sard's theorem}\label{sec:genericity-sard}
% ============================================================

The paradigm of ``genericity via Sard'' pervades differential topology.
We present some consequences that illustrate this principle.

\begin{proposition}\label{prop:generic-immersion-low-dim}
\index{generic immersion}
Let $M^m$ be a compact smooth manifold. For $n \geq 2m$, a generic smooth
map $f \colon M \to \R^n$ is an immersion. For $n \geq 2m + 1$, a generic
smooth map is an embedding.
\end{proposition}

\begin{proof}[Proof sketch]
One applies the Thom transversality theorem in jet space. Consider the
$1$-jet extension $j^1 f \colon M \to J^1(M, \R^n)$. The set
$\Sigma \subset J^1(M, \R^n)$ of $1$-jets of rank less than $m$ is a
submanifold of codimension $(n - m + 1)$ (by the determinantal variety
structure). If $n \geq 2m$, this codimension exceeds $m = \dim M$,
so $j^1 f \pitchfork \Sigma$ implies $(j^1 f)^{-1}(\Sigma) = \varnothing$
by dimension count, meaning $f$ is an immersion. The embedding statement
for $n \geq 2m + 1$ uses additionally that self-intersections can be
eliminated.
\end{proof}

\begin{proposition}[Generic maps have only transverse self-intersections]
\label{prop:generic-transverse-selfintersections}
For $n \geq 2m$, a generic immersion $f \colon M^m \to \R^n$ has only
transverse double points and no higher-order multiple points.
\end{proposition}

\begin{corollary}\label{cor:generic-curves-R2}
A generic immersion of a closed curve in $\R^2$ has only transverse
double points (simple crossings). This is the generic picture of
plane curves.
\end{corollary}

% ============================================================
\section{Intersection numbers}\label{sec:intersection-numbers}
% ============================================================

Transversality allows us to define intersection numbers, which are
fundamental invariants in differential topology. Throughout this section,
all manifolds are assumed to be oriented unless stated otherwise.

\begin{definition}[Orientation at a transverse intersection point]
\label{def:intersection-orientation}
\index{intersection number!orientation}
Let $M^n$ be an oriented manifold and let $A^a, B^b \subseteq M$ be
compact oriented submanifolds with $A \pitchfork B$ and $a + b = n$.
At each point $x \in A \cap B$, we have
$T_x A \oplus T_x B \cong T_x M$ (as a consequence of transversality and
dimension count). Define the \textbf{sign} $\varepsilon(x) = +1$ if the
isomorphism $T_x A \oplus T_x B \to T_x M$ is orientation-preserving, and
$\varepsilon(x) = -1$ otherwise.
\end{definition}

\begin{definition}[Intersection number]\label{def:intersection-number}
\index{intersection number}
With the setup above, the \textbf{intersection number} of $A$ and $B$
in $M$ is
\[
  I(A, B) \;=\; \sum_{x \in A \cap B} \varepsilon(x) \;\in\; \Z.
\]
Since $A$ is compact and $A \pitchfork B$, the intersection $A \cap B$ is a
compact $0$-dimensional manifold, hence a finite set, so this sum is finite.
\end{definition}

\begin{proposition}\label{prop:intersection-number-properties}
The intersection number $I(A, B)$ satisfies:
\begin{enumerate}[label=(\roman*)]
  \item \textbf{Homotopy invariance}: $I(A, B)$ depends only on the
    homology classes $[A] \in H_a(M; \Z)$ and $[B] \in H_b(M; \Z)$.
  \item \textbf{Symmetry}: $I(A, B) = (-1)^{ab}\, I(B, A)$, where
    $a = \dim A$ and $b = \dim B$.
  \item \textbf{Linearity}: $I(A_1 \sqcup A_2, B) = I(A_1, B) + I(A_2, B)$.
\end{enumerate}
\end{proposition}

\begin{proof}
(i) If $A_0$ and $A_1$ are homologous (there exists a compact oriented
$(a+1)$-manifold $W$ with $\partial W = A_1 - A_0$), then after perturbing
to achieve transversality with $B$, the signed count of intersection points
of $W$ with $B$ gives a cobordism between $A_0 \cap B$ and $A_1 \cap B$,
preserving the total count.

(ii) Swapping the order in the direct sum $T_x A \oplus T_x B$ introduces
a sign $(-1)^{ab}$, from the permutation of basis vectors.

(iii) This follows directly from the definition.
\end{proof}

\begin{example}\label{ex:intersection-torus}
\index{intersection number!on the torus}
On the torus $T^2 = S^1 \times S^1$, let $A = S^1 \times \set{1}$ and
$B = \set{1} \times S^1$. These are transverse and meet in a single point
$(1,1)$. With the standard orientations, $I(A, B) = +1$. This reflects
the fact that the ``meridian'' and ``longitude'' of the torus link once.
\end{example}

% --- TikZ figure: intersection on the torus (schematic) ---
\begin{figure}[ht]
\centering
\begin{tikzpicture}[scale=1.0]
  % Flat torus representation
  \draw[thick] (0,0) rectangle (3,3);
  \draw[thick, blue!70!black, ->-/.style={decoration={markings,
    mark=at position 0.55 with {\arrow{>}}}, postaction={decorate}}]
    (0,1.5) -- (3,1.5) node[midway, above]{\small $A$};
  \draw[thick, red!70!black, ->-/.style={decoration={markings,
    mark=at position 0.55 with {\arrow{>}}}, postaction={decorate}}]
    (1.5,0) -- (1.5,3) node[midway, right]{\small $B$};
  \filldraw[black] (1.5,1.5) circle (2pt) node[above right=2pt]{\small $+1$};
  % Identification arrows
  \draw[thick, ->] (0,0) -- (0,3);
  \draw[thick, ->] (0,0) -- (3,0);
  \draw[thick, ->] (3,0) -- (3,3);
  \draw[thick, ->] (0,3) -- (3,3);
  \node[below] at (1.5,-0.3) {$T^2$ (fundamental domain)};
\end{tikzpicture}
\caption{Intersection of the two standard circles on the torus: $I(A,B)=+1$.}
\label{fig:intersection-torus}
\end{figure}

% ============================================================
\section{Self-intersection number}\label{sec:self-intersection}
% ============================================================

\begin{definition}[Self-intersection number]\label{def:self-intersection}
\index{self-intersection number}
Let $M^n$ be a compact oriented manifold and $A^a \subseteq M$ a compact
oriented submanifold with $2a = n$. The \textbf{self-intersection number}
of $A$ is defined as $I(A, A)$, computed by perturbing one copy of $A$ to
$A'$ so that $A \pitchfork A'$, and setting
\[
  I(A, A) \;=\; I(A, A').
\]
By homotopy invariance, this is independent of the perturbation.
\end{definition}

\begin{proposition}\label{prop:self-intersection-euler}
\index{Euler class!and self-intersection}
The self-intersection number $I(A, A)$ equals the Euler number of the
normal bundle $\nu(A \hookrightarrow M)$:
\[
  I(A, A) = e(\nu_A),
\]
where $e(\nu_A) \in \Z$ is the evaluation of the Euler class on the
fundamental class of $A$.
\end{proposition}

\begin{example}\label{ex:self-intersection-sphere-sphere}
Consider $S^n$ as the equatorial sphere in $S^{2n}$ (via the standard
inclusion $\R^{n+1} \hookrightarrow \R^{2n+1}$). The normal bundle of
$S^n$ in $S^{2n}$ is the tangent bundle $TS^n$ (this can be seen via the
diagonal embedding $\Delta \colon S^n \to S^n \times S^n$). Therefore
$I(S^n, S^n) = \chi(S^n) = 1 + (-1)^n$.
\end{example}

\begin{example}\label{ex:self-intersection-zero-section}
The zero section $M \hookrightarrow TM$ (viewing $TM$ as a $2n$-dimensional
manifold) has self-intersection number equal to $\chi(M)$, the Euler
characteristic of $M$. This provides a differential-topological proof of
the Poincar\'e--Hopf theorem.
\end{example}

% ============================================================
\section{The Lefschetz fixed-point theorem}\label{sec:lefschetz}
% ============================================================

We briefly state the connection between intersection theory and fixed-point
theory.

\begin{definition}[Lefschetz number]\label{def:lefschetz-number}
\index{Lefschetz number}
Let $f \colon M \to M$ be a smooth map on a compact oriented manifold $M$.
The \textbf{Lefschetz number} of $f$ is
\[
  L(f) \;=\; \sum_{k=0}^{\dim M} (-1)^k \operatorname{tr}\bigl(
    f_* \colon H_k(M; \Q) \to H_k(M; \Q)\bigr).
\]
\end{definition}

\begin{theorem}[Lefschetz Fixed-Point Theorem]\label{thm:lefschetz}
\index{Lefschetz fixed-point theorem}
Let $f \colon M \to M$ be a smooth map on a compact oriented manifold.
If $L(f) \neq 0$, then $f$ has a fixed point. More precisely, if the graph
$\Gamma_f = \set{(x, f(x)) : x \in M}$ is transverse to the diagonal
$\Delta = \set{(x,x) : x \in M}$ in $M \times M$, then
\[
  L(f) = I(\Gamma_f, \Delta) = \sum_{\substack{x \in M \\ f(x)=x}}
    \operatorname{sign}\det(\Id - \dd f_x).
\]
\end{theorem}

\begin{remark}\label{rem:lefschetz-euler}
When $f = \Id_M$, we have $L(\Id_M) = \chi(M)$. After perturbing the
identity to a map with only non-degenerate fixed points (a Lefschetz
map), we recover the Poincar\'e--Hopf theorem: the sum of indices of
a generic vector field equals $\chi(M)$. This connects the self-intersection
theory of Section~\ref{sec:self-intersection} to the Lefschetz theory.
\end{remark}

% ============================================================
\section{Exercises}\label{sec:exercises-transversality}
% ============================================================

\begin{exercise}\label{exo:transverse-basic}
\index{transversality!exercises}
Show that two lines in $\R^3$ are transverse if and only if they are not
coplanar (i.e., they are skew). Conclude that two lines in $\R^3$ that
intersect in a point are \emph{never} transverse.
\end{exercise}

\begin{exercise}\label{exo:transverse-codimension}
Let $A^a$ and $B^b$ be submanifolds of $M^n$ with $A \pitchfork B$ and
$a + b < n$. Show that $A \cap B = \varnothing$.
\end{exercise}

\begin{exercise}\label{exo:submersion-transverse}
Let $f \colon M \to N$ be a smooth map and $W \subseteq N$ a submanifold.
\begin{enumerate}[label=(\alph*)]
  \item Show that if $f$ is a submersion, then $f \pitchfork W$.
  \item Give an example showing that the converse is false.
\end{enumerate}
\end{exercise}

\begin{exercise}\label{exo:transverse-product}
Let $f \colon M \to N$ and $g \colon M' \to N'$ be smooth maps, and let
$W \subseteq N$, $W' \subseteq N'$ be submanifolds. Show that if
$f \pitchfork W$ and $g \pitchfork W'$, then
$(f \times g) \pitchfork (W \times W')$.
\end{exercise}

\begin{exercise}\label{exo:intersection-number-mod2}
Let $M$ be a compact manifold (not necessarily oriented) and
$A, B \subseteq M$ compact submanifolds with complementary dimensions
and $A \pitchfork B$. Define the \textbf{mod 2 intersection number}\index{intersection number!mod 2}
$I_2(A, B) = \#(A \cap B) \mod 2$. Show that $I_2(A, B)$ is
a homotopy invariant.
\end{exercise}

\begin{exercise}\label{exo:linking-number}
\index{linking number}
Let $\gamma_1, \gamma_2 \colon S^1 \to \R^3$ be disjoint smooth embeddings.
The \textbf{linking number} $\operatorname{lk}(\gamma_1, \gamma_2)$ can be
defined as follows: let $\Sigma$ be a compact oriented surface with
$\partial \Sigma = \gamma_1(S^1)$, chosen transverse to $\gamma_2$.
Set $\operatorname{lk}(\gamma_1, \gamma_2) = I(\Sigma, \gamma_2)$.
\begin{enumerate}[label=(\alph*)]
  \item Show that this is independent of the choice of $\Sigma$.
  \item Compute the linking number of the Hopf link.
  \item Show that $\operatorname{lk}(\gamma_1, \gamma_2) =
    \operatorname{lk}(\gamma_2, \gamma_1)$.
\end{enumerate}
\end{exercise}

\begin{exercise}\label{exo:self-intersection-RP}
Consider $\RP^1 \subset \RP^2$ as the submanifold defined by the last
homogeneous coordinate being zero. Compute the self-intersection number
$I(\RP^1, \RP^1)$ in $\RP^2$ (working mod $2$, since $\RP^2$ is
non-orientable).
\end{exercise}

\begin{exercise}\label{exo:transversality-jet}
Let $M^m$ and $N^n$ be smooth manifolds with $n \geq 2m + 1$. Using the
Thom transversality theorem applied to $1$-jet extensions, show that a
generic smooth map $f \colon M \to N$ is an injective immersion. If $M$ is
compact, show it is an embedding.
\end{exercise}

\begin{exercise}\label{exo:lefschetz-application}
Let $f \colon S^{2n} \to S^{2n}$ be a smooth map of degree $d$.
\begin{enumerate}[label=(\alph*)]
  \item Compute the Lefschetz number $L(f)$.
  \item Show that if $d \neq (-1)^{2n+1} = -1$, then $f$ has a fixed point.
  \item Deduce that the antipodal map on $S^{2n}$ has a fixed point.
    Why does this not contradict the fact that the antipodal map has no
    fixed point? Resolve the apparent paradox.
\end{enumerate}
\end{exercise}

\begin{exercise}\label{exo:euler-selfintersection-diagonal}
Let $M$ be a compact oriented manifold of dimension $n$. Show that the
self-intersection number of the diagonal $\Delta \subset M \times M$ equals
the Euler characteristic:
\[
  I(\Delta, \Delta) = \chi(M).
\]
\textit{Hint}: Relate this to the Lefschetz number of the identity.
\end{exercise}


%%%%%%%%%%%%%%%%%%%%%%%%%%%%%%%%%%%%%%%%%%%%%%%%%%%%%%%%%%%%%
%%%%%%%%%%%%%%%%%%%%%%%%%%%%%%%%%%%%%%%%%%%%%%%%%%%%%%%%%%%%%
%%
%%  CHAPTER 6 --- DIFFERENTIAL FORMS AND INTEGRATION ON MANIFOLDS
%%
%%%%%%%%%%%%%%%%%%%%%%%%%%%%%%%%%%%%%%%%%%%%%%%%%%%%%%%%%%%%%
%%%%%%%%%%%%%%%%%%%%%%%%%%%%%%%%%%%%%%%%%%%%%%%%%%%%%%%%%%%%%

\chapter{Differential Forms and Integration on Manifolds}
\label{ch:differential-forms}
\minitoc\vspace{1cm}

Differential forms provide the natural language for integration on manifolds.
They unify and generalize many constructions from vector calculus---gradient,
curl, divergence, line integrals, surface integrals, and the fundamental
theorems of Green, Gauss, and Stokes---into a single elegant framework.
In this chapter, we develop the theory of exterior algebra, differential forms,
the exterior derivative, and integration, culminating in the machinery needed
for de Rham cohomology (treated in subsequent chapters).

% ============================================================
\section{Exterior algebra}\label{sec:exterior-algebra}
% ============================================================

We begin with the algebraic foundations: the exterior algebra of a
finite-dimensional real vector space.

\begin{definition}[Alternating multilinear form]\label{def:alternating-form}
\index{alternating form}
Let $V$ be a real vector space of dimension $n$. A \textbf{$k$-linear
alternating form} (or \textbf{exterior $k$-form}) on $V$ is a multilinear
map
\[
  \omega \colon \underbrace{V \times \cdots \times V}_{k} \longrightarrow \R
\]
that is alternating: $\omega(v_1, \ldots, v_k) = 0$ whenever two of the
$v_i$ are equal. Equivalently, for any permutation $\sigma \in \mathfrak{S}_k$,
\[
  \omega(v_{\sigma(1)}, \ldots, v_{\sigma(k)})
    = \operatorname{sgn}(\sigma)\, \omega(v_1, \ldots, v_k).
\]
\end{definition}

\begin{definition}[Exterior power]\label{def:exterior-power}
\index{exterior algebra}\index{exterior power}
The space of all $k$-linear alternating forms on $V$ is denoted
$\Lambda^k(V^*)$ or $\Lambda^k V^*$. By convention, $\Lambda^0 V^* = \R$
and $\Lambda^1 V^* = V^*$. We have $\Lambda^k V^* = 0$ for $k > \dim V$.
The dimension is
\[
  \dim \Lambda^k V^* = \binom{n}{k}.
\]
\end{definition}

\begin{definition}[Wedge product]\label{def:wedge-product}
\index{wedge product}
The \textbf{wedge product} (or \textbf{exterior product})
\[
  \wedge \colon \Lambda^k V^* \times \Lambda^\ell V^*
    \longrightarrow \Lambda^{k+\ell} V^*
\]
is defined by
\[
  (\alpha \wedge \beta)(v_1, \ldots, v_{k+\ell})
    = \frac{1}{k!\,\ell!} \sum_{\sigma \in \mathfrak{S}_{k+\ell}}
      \operatorname{sgn}(\sigma)\,
      \alpha(v_{\sigma(1)}, \ldots, v_{\sigma(k)})\,
      \beta(v_{\sigma(k+1)}, \ldots, v_{\sigma(k+\ell)}).
\]
Equivalently, summing only over $(k,\ell)$-shuffles (permutations
$\sigma$ with $\sigma(1) < \cdots < \sigma(k)$ and
$\sigma(k+1) < \cdots < \sigma(k+\ell)$):
\[
  (\alpha \wedge \beta)(v_1, \ldots, v_{k+\ell})
    = \sum_{\text{$(k,\ell)$-shuffles } \sigma}
      \operatorname{sgn}(\sigma)\,
      \alpha(v_{\sigma(1)}, \ldots, v_{\sigma(k)})\,
      \beta(v_{\sigma(k+1)}, \ldots, v_{\sigma(k+\ell)}).
\]
\end{definition}

\begin{proposition}[Properties of the wedge product]\label{prop:wedge-properties}
\index{wedge product!properties}
For $\alpha \in \Lambda^k V^*$, $\beta \in \Lambda^\ell V^*$, and
$\gamma \in \Lambda^p V^*$:
\begin{enumerate}[label=(\roman*)]
  \item \textbf{Associativity}: $(\alpha \wedge \beta) \wedge \gamma
    = \alpha \wedge (\beta \wedge \gamma)$.
  \item \textbf{Graded commutativity}: $\alpha \wedge \beta
    = (-1)^{k\ell}\, \beta \wedge \alpha$.
  \item \textbf{Bilinearity}: $\wedge$ is bilinear in each argument.
  \item If $e^1, \ldots, e^n$ is a basis of $V^*$, then
    $\set{e^{i_1} \wedge \cdots \wedge e^{i_k} : 1 \leq i_1 < \cdots < i_k \leq n}$
    is a basis of $\Lambda^k V^*$.
\end{enumerate}
\end{proposition}

\begin{proof}
(i) Both sides equal the $(k+\ell+p)$-linear alternating form given by
summing over all $(k,\ell,p)$-shuffles.

(ii) Transposing the first $k$ arguments past the last $\ell$ arguments
requires $k\ell$ transpositions, each contributing a factor of $-1$.

(iii) Clear from the definition.

(iv) One checks that these $\binom{n}{k}$ elements are linearly independent
and span $\Lambda^k V^*$ by expressing any alternating form in terms of them.
\end{proof}

\begin{definition}[Exterior algebra]\label{def:exterior-algebra}
\index{exterior algebra}
The \textbf{exterior algebra} of $V^*$ is the graded algebra
\[
  \Lambda^*(V^*) = \bigoplus_{k=0}^{n} \Lambda^k(V^*)
\]
equipped with the wedge product. It has total dimension $2^n$.
\end{definition}

\begin{example}\label{ex:exterior-R3}
For $V = \R^3$ with dual basis $e^1, e^2, e^3$:
\begin{itemize}
  \item $\Lambda^0 = \R$, dimension $1$.
  \item $\Lambda^1 = \operatorname{span}\set{e^1, e^2, e^3}$, dimension $3$.
  \item $\Lambda^2 = \operatorname{span}\set{e^1 \wedge e^2,\;
    e^1 \wedge e^3,\; e^2 \wedge e^3}$, dimension $3$.
  \item $\Lambda^3 = \operatorname{span}\set{e^1 \wedge e^2 \wedge e^3}$,
    dimension $1$.
\end{itemize}
The isomorphisms $\Lambda^1 \cong \Lambda^2$ and $\Lambda^0 \cong \Lambda^3$
underlie the vector calculus identities relating gradient, curl, and divergence.
\end{example}

% ============================================================
\section{Differential forms on manifolds}\label{sec:differential-forms}
% ============================================================

\begin{definition}[Differential $k$-form]\label{def:differential-form}
\index{differential form}
Let $M$ be a smooth manifold. A \textbf{differential $k$-form}\index{$k$-form}
on $M$ is a smooth section of the $k$-th exterior power of the cotangent
bundle:
\[
  \omega \in \Gamma(\Lambda^k T^*M).
\]
That is, $\omega$ assigns to each $p \in M$ an alternating $k$-linear form
$\omega_p \in \Lambda^k(T_p^*M)$, and this assignment is smooth. The space
of all smooth $k$-forms on $M$ is denoted
\[
  \Omega^k(M) = \Gamma(\Lambda^k T^*M).
\]
By convention, $\Omega^0(M) = C^\infty(M)$.
\end{definition}

\begin{remark}\label{rem:local-expression-forms}
In local coordinates $(x^1, \ldots, x^n)$ on an open set $U \subseteq M$,
any $\omega \in \Omega^k(M)$ can be written as
\[
  \omega\big|_U = \sum_{1 \leq i_1 < \cdots < i_k \leq n}
    \omega_{i_1 \cdots i_k}\, \dd x^{i_1} \wedge \cdots \wedge \dd x^{i_k},
\]
where $\omega_{i_1 \cdots i_k} \in C^\infty(U)$ are smooth coefficient
functions and $\dd x^1, \ldots, \dd x^n$ are the coordinate $1$-forms
(a local frame for $T^*M$).
\end{remark}

\begin{definition}[Algebra of differential forms]\label{def:form-algebra}
The \textbf{de Rham algebra}\index{de Rham algebra} of $M$ is the graded
$\R$-algebra
\[
  \Omega^*(M) = \bigoplus_{k=0}^{n} \Omega^k(M)
\]
with the wedge product defined pointwise:
$(\alpha \wedge \beta)_p = \alpha_p \wedge \beta_p$.
It is a graded-commutative algebra: for $\alpha \in \Omega^k(M)$ and
$\beta \in \Omega^\ell(M)$,
\[
  \alpha \wedge \beta = (-1)^{k\ell}\, \beta \wedge \alpha.
\]
\end{definition}

\begin{example}\label{ex:1-forms-R3}
On $\R^3$, a $1$-form is $\omega = P\,\dd x + Q\,\dd y + R\,\dd z$
where $P, Q, R \in C^\infty(\R^3)$. A $2$-form is
$\eta = A\,\dd y \wedge \dd z + B\,\dd z \wedge \dd x + C\,\dd x \wedge \dd y$.
A $3$-form is $\mu = f\,\dd x \wedge \dd y \wedge \dd z$.
\end{example}

% ============================================================
\section{The exterior derivative}\label{sec:exterior-derivative}
% ============================================================

\begin{theorem}[Existence and uniqueness of the exterior derivative]
\label{thm:exterior-derivative}
\index{exterior derivative}
There exists a unique family of $\R$-linear maps
\[
  \dd \colon \Omega^k(M) \to \Omega^{k+1}(M), \quad k = 0, 1, \ldots, n,
\]
satisfying the following properties:
\begin{enumerate}[label=(\roman*)]
  \item \textbf{Degree $+1$}: $\dd$ maps $k$-forms to $(k+1)$-forms.
  \item \textbf{$\dd^2 = 0$}: For all $\omega \in \Omega^k(M)$,
    $\dd(\dd\omega) = 0$.
  \item \textbf{Graded Leibniz rule}: For $\alpha \in \Omega^k(M)$ and
    $\beta \in \Omega^\ell(M)$,
    \[
      \dd(\alpha \wedge \beta) = (\dd\alpha) \wedge \beta
        + (-1)^k\, \alpha \wedge (\dd\beta).
    \]
  \item \textbf{Agreement on functions}: For $f \in \Omega^0(M)
    = C^\infty(M)$, $\dd f$ is the usual differential of $f$, i.e.,
    $(\dd f)_p(v) = v(f)$ for $v \in T_p M$.
\end{enumerate}
\end{theorem}

\begin{proof}
\textbf{Uniqueness.} Properties (iii) and (iv) force the action of $\dd$ on
local coordinate expressions. In coordinates $(x^1, \ldots, x^n)$, for
$\omega = \sum_I \omega_I\, \dd x^I$ (where $I = (i_1 < \cdots < i_k)$
and $\dd x^I = \dd x^{i_1} \wedge \cdots \wedge \dd x^{i_k}$), we must have
\[
  \dd\omega = \sum_I \dd\omega_I \wedge \dd x^I
    = \sum_I \sum_{j=1}^{n} \frac{\partial \omega_I}{\partial x^j}\,
      \dd x^j \wedge \dd x^I,
\]
since $\dd(\dd x^j) = 0$ by property (ii) applied to $\dd x^j = \dd(x^j)$,
and the Leibniz rule gives
$\dd(\omega_I\, \dd x^I) = \dd\omega_I \wedge \dd x^I
  + \omega_I\, \dd(\dd x^I) = \dd\omega_I \wedge \dd x^I$.
Here $\dd(\dd x^I) = 0$ follows by induction using (ii) and (iii).

\medskip\noindent
\textbf{Existence.} Define $\dd$ by the formula above in each coordinate chart.
To show this is well-defined (independent of coordinates), one can verify
directly that the formula transforms correctly under coordinate changes, or
one can give a coordinate-free formula: for $\omega \in \Omega^k(M)$ and
smooth vector fields $X_0, \ldots, X_k$,
\begin{align*}
  (\dd\omega)(X_0, \ldots, X_k) &= \sum_{i=0}^{k} (-1)^i\,
    X_i\bigl(\omega(X_0, \ldots, \widehat{X_i}, \ldots, X_k)\bigr) \\
  &\quad + \sum_{0 \leq i < j \leq k} (-1)^{i+j}\,
    \omega\bigl([X_i, X_j], X_0, \ldots, \widehat{X_i}, \ldots,
      \widehat{X_j}, \ldots, X_k\bigr),
\end{align*}
where $\widehat{X_i}$ denotes omission. This is manifestly coordinate-free.
One verifies that this formula defines a $(k+1)$-form and satisfies
all four properties.

The property $\dd^2 = 0$ follows in coordinates from the symmetry of
second partial derivatives:
\[
  \dd^2 \omega = \sum_I \sum_{j,\ell}
    \frac{\partial^2 \omega_I}{\partial x^\ell \partial x^j}\,
    \dd x^\ell \wedge \dd x^j \wedge \dd x^I = 0,
\]
since $\frac{\partial^2 \omega_I}{\partial x^\ell \partial x^j}$ is
symmetric in $\ell, j$ while $\dd x^\ell \wedge \dd x^j$ is antisymmetric.
\end{proof}

\begin{example}\label{ex:exterior-derivative-R3}
On $\R^3$, the exterior derivative recovers the classical vector calculus
operators:
\begin{itemize}
  \item For $f \in \Omega^0(\R^3)$: $\dd f = \frac{\partial f}{\partial x}\,
    \dd x + \frac{\partial f}{\partial y}\,\dd y
    + \frac{\partial f}{\partial z}\,\dd z$ (gradient).
  \item For $\omega = P\,\dd x + Q\,\dd y + R\,\dd z \in \Omega^1(\R^3)$:
    \[
      \dd\omega = \Bigl(\frac{\partial R}{\partial y}
        - \frac{\partial Q}{\partial z}\Bigr)\dd y \wedge \dd z
      + \Bigl(\frac{\partial P}{\partial z}
        - \frac{\partial R}{\partial x}\Bigr)\dd z \wedge \dd x
      + \Bigl(\frac{\partial Q}{\partial x}
        - \frac{\partial P}{\partial y}\Bigr)\dd x \wedge \dd y
    \]
    (curl).
  \item For $\eta = A\,\dd y \wedge \dd z + B\,\dd z \wedge \dd x
    + C\,\dd x \wedge \dd y \in \Omega^2(\R^3)$:
    $\dd\eta = \bigl(\frac{\partial A}{\partial x}
      + \frac{\partial B}{\partial y}
      + \frac{\partial C}{\partial z}\bigr)\dd x \wedge \dd y \wedge \dd z$
    (divergence).
\end{itemize}
The identity $\dd^2 = 0$ encodes both $\operatorname{curl}(\operatorname{grad} f) = 0$
and $\operatorname{div}(\operatorname{curl} F) = 0$.
\end{example}

% --- TikZ diagram: de Rham complex ---
\begin{figure}[ht]
\centering
\begin{tikzcd}[column sep=large]
  \Omega^0(M) \arrow[r, "\dd"] &
  \Omega^1(M) \arrow[r, "\dd"] &
  \Omega^2(M) \arrow[r, "\dd"] &
  \cdots \arrow[r, "\dd"] &
  \Omega^n(M) \arrow[r, "\dd"] &
  0
\end{tikzcd}
\caption{The de Rham complex $(\Omega^*(M), \dd)$.
  The condition $\dd^2 = 0$ makes this a cochain complex.}
\label{fig:derham-complex}
\end{figure}

% ============================================================
\section{Pullback of differential forms}\label{sec:pullback-forms}
% ============================================================

\begin{definition}[Pullback]\label{def:pullback-form}
\index{pullback!of differential forms}
Let $f \colon M \to N$ be a smooth map. The \textbf{pullback} is the
$\R$-linear map $f^* \colon \Omega^k(N) \to \Omega^k(M)$ defined by
\[
  (f^*\omega)_p(v_1, \ldots, v_k)
    = \omega_{f(p)}(\dd f_p(v_1), \ldots, \dd f_p(v_k))
\]
for $\omega \in \Omega^k(N)$, $p \in M$, and $v_1, \ldots, v_k \in T_p M$.
\end{definition}

\begin{proposition}[Properties of pullback]\label{prop:pullback-properties}
\index{pullback!properties}
Let $f \colon M \to N$ and $g \colon N \to P$ be smooth maps.
\begin{enumerate}[label=(\roman*)]
  \item \textbf{Functoriality}: $(g \circ f)^* = f^* \circ g^*$.
  \item \textbf{Compatibility with wedge product}:
    $f^*(\alpha \wedge \beta) = (f^*\alpha) \wedge (f^*\beta)$.
  \item \textbf{Identity}: $(\Id_M)^* = \Id_{\Omega^*(M)}$.
  \item \textbf{On functions}: For $h \in \Omega^0(N) = C^\infty(N)$,
    $f^* h = h \circ f$.
\end{enumerate}
\end{proposition}

\begin{theorem}[Naturality of the exterior derivative]\label{thm:naturality-d}
\index{exterior derivative!naturality}\index{naturality of $\dd$}
Let $f \colon M \to N$ be a smooth map. Then the pullback commutes with the
exterior derivative:
\[
  f^* \circ \dd = \dd \circ f^*.
\]
That is, for every $\omega \in \Omega^k(N)$,
$f^*(\dd\omega) = \dd(f^*\omega)$.
\end{theorem}

\begin{proof}
On $0$-forms (functions): for $h \in C^\infty(N)$ and $v \in T_p M$,
\[
  (f^*\dd h)_p(v) = (\dd h)_{f(p)}(\dd f_p(v))
  = (\dd f_p(v))(h)
  = v(h \circ f)
  = (\dd(h \circ f))_p(v)
  = (\dd(f^*h))_p(v).
\]
In local coordinates $(y^1, \ldots, y^n)$ on $N$, any $k$-form can be
written as $\omega = \sum_I \omega_I\, \dd y^I$. Then
\[
  f^*(\dd\omega) = f^*\Bigl(\sum_I \dd\omega_I \wedge \dd y^I\Bigr)
  = \sum_I f^*(\dd\omega_I) \wedge f^*(\dd y^I)
  = \sum_I \dd(f^*\omega_I) \wedge \dd(f^*y^I),
\]
using the $0$-form case and compatibility with wedge products. On the other
hand,
\[
  \dd(f^*\omega) = \dd\Bigl(\sum_I (f^*\omega_I)\, f^*(\dd y^I)\Bigr)
  = \sum_I \dd(f^*\omega_I) \wedge f^*(\dd y^I)
    + \sum_I (f^*\omega_I)\, \dd(f^*(\dd y^I)).
\]
Now $f^*(\dd y^{i_j}) = \dd(f^* y^{i_j}) = \dd(y^{i_j} \circ f)$, so
$\dd(f^*(\dd y^{i_j})) = \dd^2(y^{i_j} \circ f) = 0$. By the Leibniz rule
for the wedge product, $\dd(f^*(\dd y^I)) = 0$. Hence both sides agree.
\end{proof}

% --- Commutative diagram: naturality ---
\begin{figure}[ht]
\centering
\begin{tikzcd}
  \Omega^k(N) \arrow[r, "\dd"] \arrow[d, "f^*"'] &
  \Omega^{k+1}(N) \arrow[d, "f^*"] \\
  \Omega^k(M) \arrow[r, "\dd"'] &
  \Omega^{k+1}(M)
\end{tikzcd}
\caption{Naturality of the exterior derivative: $f^* \circ \dd = \dd \circ f^*$.}
\label{fig:naturality-d}
\end{figure}

% ============================================================
\section{Interior product and Lie derivative}\label{sec:interior-lie}
% ============================================================

\begin{definition}[Interior product]\label{def:interior-product}
\index{interior product}
Let $X$ be a smooth vector field on $M$. The \textbf{interior product}
(or \textbf{contraction}) with $X$ is the linear map
$\iota_X \colon \Omega^k(M) \to \Omega^{k-1}(M)$ defined by
\[
  (\iota_X \omega)(v_1, \ldots, v_{k-1})
    = \omega(X, v_1, \ldots, v_{k-1}).
\]
By convention, $\iota_X$ acts as zero on $\Omega^0(M)$.
\end{definition}

\begin{proposition}\label{prop:interior-product-properties}
\index{interior product!properties}
The interior product satisfies:
\begin{enumerate}[label=(\roman*)]
  \item $\iota_X$ has degree $-1$: it maps $\Omega^k$ to $\Omega^{k-1}$.
  \item $\iota_X^2 = 0$.
  \item \textbf{Graded Leibniz rule}: For $\alpha \in \Omega^k(M)$ and
    $\beta \in \Omega^\ell(M)$,
    \[
      \iota_X(\alpha \wedge \beta) = (\iota_X \alpha) \wedge \beta
        + (-1)^k\, \alpha \wedge (\iota_X \beta).
    \]
    That is, $\iota_X$ is an antiderivation of degree $-1$.
  \item $\iota_{fX} = f\, \iota_X$ for $f \in C^\infty(M)$.
\end{enumerate}
\end{proposition}

\begin{proof}
(ii) For $\omega \in \Omega^k(M)$,
$(\iota_X^2 \omega)(v_1, \ldots, v_{k-2})
= \omega(X, X, v_1, \ldots, v_{k-2}) = 0$ by the alternating property.

(iii) Follows by a direct computation using the definition of the wedge
product, tracking the position of $X$ among the arguments.
\end{proof}

\begin{definition}[Lie derivative of a form]\label{def:lie-derivative-form}
\index{Lie derivative!of differential forms}
Let $X$ be a smooth vector field on $M$ with flow $\phi_t$. The
\textbf{Lie derivative} of $\omega \in \Omega^k(M)$ along $X$ is
\[
  \mathcal{L}_X \omega = \frac{\dd}{\dd t}\bigg|_{t=0} \phi_t^* \omega
  = \lim_{t \to 0} \frac{\phi_t^* \omega - \omega}{t}.
\]
\end{definition}

\begin{theorem}[Cartan's magic formula]\label{thm:cartan-formula}
\index{Cartan formula}\index{magic formula}
For any smooth vector field $X$ on $M$ and any $\omega \in \Omega^k(M)$,
\[
  \mathcal{L}_X \omega = \iota_X(\dd\omega) + \dd(\iota_X \omega)
  = (\iota_X \circ \dd + \dd \circ \iota_X)\omega.
\]
In shorthand: $\mathcal{L}_X = \iota_X \dd + \dd\, \iota_X$.
\end{theorem}

\begin{proof}
We verify the formula by checking that both sides agree on functions and
exact $1$-forms, then extend by the Leibniz rule.

\medskip\noindent
\textbf{On functions} ($k = 0$): For $f \in C^\infty(M)$,
$\mathcal{L}_X f = X(f)$ (the directional derivative). On the right side,
$\iota_X(\dd f) + \dd(\iota_X f) = (\dd f)(X) + 0 = X(f)$.

\medskip\noindent
\textbf{On exact $1$-forms} ($k = 1$, $\omega = \dd f$):
$\mathcal{L}_X(\dd f) = \dd(\mathcal{L}_X f) = \dd(Xf)$ (since
$\mathcal{L}_X$ commutes with $\dd$). On the right side,
$\iota_X(\dd^2 f) + \dd(\iota_X \dd f) = 0 + \dd(Xf)$.

\medskip\noindent
\textbf{Leibniz rule}: Both $\mathcal{L}_X$ and $\iota_X \dd + \dd\, \iota_X$
are derivations of degree $0$ on $\Omega^*(M)$ (i.e., they satisfy the
ungraded Leibniz rule $D(\alpha \wedge \beta) = (D\alpha) \wedge \beta
+ \alpha \wedge (D\beta)$). For $\mathcal{L}_X$, this follows from
$\phi_t^*(\alpha \wedge \beta) = (\phi_t^*\alpha) \wedge (\phi_t^*\beta)$.
For $\iota_X \dd + \dd\, \iota_X$, this follows by expanding and using
the graded Leibniz rules for $\dd$ and $\iota_X$.

Since differential forms are locally generated (over $C^\infty$) by functions
and exact $1$-forms $\dd x^i$, agreement on these generators plus the
Leibniz rule implies agreement everywhere.
\end{proof}

\begin{corollary}\label{cor:lie-derivative-commutes-d}
The Lie derivative commutes with the exterior derivative:
$\mathcal{L}_X \circ \dd = \dd \circ \mathcal{L}_X$.
\end{corollary}

\begin{proof}
$\mathcal{L}_X \dd\omega = \iota_X \dd^2\omega + \dd\iota_X \dd\omega
= \dd\iota_X \dd\omega$. Also
$\dd\mathcal{L}_X\omega = \dd\iota_X\dd\omega + \dd^2\iota_X\omega
= \dd\iota_X\dd\omega$.
\end{proof}

\begin{proposition}\label{prop:lie-derivative-properties}
\index{Lie derivative!properties}
For vector fields $X, Y$ and $\omega, \eta \in \Omega^*(M)$:
\begin{enumerate}[label=(\roman*)]
  \item $\mathcal{L}_X(f\omega) = (Xf)\omega + f\,\mathcal{L}_X\omega$
    for $f \in C^\infty(M)$.
  \item $[\mathcal{L}_X, \iota_Y] = \iota_{[X,Y]}$, i.e.,
    $\mathcal{L}_X \iota_Y - \iota_Y \mathcal{L}_X = \iota_{[X,Y]}$.
  \item $[\mathcal{L}_X, \mathcal{L}_Y] = \mathcal{L}_{[X,Y]}$.
\end{enumerate}
\end{proposition}

\begin{remark}\label{rem:graded-algebra-operators}
The three operators $\dd$ (degree $+1$), $\iota_X$ (degree $-1$), and
$\mathcal{L}_X$ (degree $0$) satisfy the relations of a graded Lie algebra,
with $\dd^2 = 0$, $\iota_X^2 = 0$, and
$\mathcal{L}_X = [\dd, \iota_X]$ (where $[\cdot,\cdot]$ denotes the graded
commutator). This algebraic structure is central to the Weil model of
equivariant cohomology.
\end{remark}

% --- TikZ: summary diagram of operators on forms ---
\begin{figure}[ht]
\centering
\begin{tikzpicture}[>=Stealth, node distance=3.5cm]
  \node (Ok-1) {$\Omega^{k-1}(M)$};
  \node (Ok) [right of=Ok-1] {$\Omega^k(M)$};
  \node (Ok+1) [right of=Ok] {$\Omega^{k+1}(M)$};
  \draw[->] (Ok-1) -- (Ok) node[midway,above]{\small $\dd$};
  \draw[->] (Ok) -- (Ok+1) node[midway,above]{\small $\dd$};
  \draw[->] (Ok) to[bend right=30] node[midway,below]{\small $\iota_X$} (Ok-1);
  \draw[->] (Ok+1) to[bend right=30] node[midway,below]{\small $\iota_X$} (Ok);
  \draw[->,dashed] (Ok) to[loop above, looseness=8, out=110, in=70,
    min distance=15mm] node[above]{\small $\mathcal{L}_X = \iota_X\dd + \dd\iota_X$} (Ok);
\end{tikzpicture}
\caption{The operators $\dd$, $\iota_X$, and $\mathcal{L}_X$ acting on $\Omega^*(M)$.}
\label{fig:operators-forms}
\end{figure}

% ============================================================
\section{Orientation and volume forms}\label{sec:orientation}
% ============================================================

\begin{definition}[Orientation of a manifold]\label{def:orientation-manifold}
\index{orientation!of a manifold}
An \textbf{orientation} of a smooth manifold $M$ of dimension $n$ is a
choice of a maximal atlas $\set{(U_\alpha, \varphi_\alpha)}$ such that all
transition functions $\varphi_\beta \circ \varphi_\alpha^{-1}$ have positive
Jacobian determinant. Equivalently, an orientation is a choice of
connected component of the space of nowhere-vanishing $n$-forms on $M$.
A manifold is \textbf{orientable}\index{orientable manifold} if such a choice
exists.
\end{definition}

\begin{proposition}\label{prop:orientation-nform}
\index{orientation!via $n$-form}
A connected smooth $n$-manifold $M$ is orientable if and only if there
exists a nowhere-vanishing $n$-form $\mu \in \Omega^n(M)$. Such a form
is called a \textbf{volume form}\index{volume form}. Two volume forms
$\mu$ and $\mu'$ determine the same orientation if and only if
$\mu' = f\mu$ for some $f \in C^\infty(M)$ with $f > 0$ everywhere.
\end{proposition}

\begin{example}\label{ex:orientation-Rn}
The standard orientation of $\R^n$ is given by the volume form
$\dd x^1 \wedge \cdots \wedge \dd x^n$.
\end{example}

\begin{example}\label{ex:orientation-sphere}
\index{orientation!of $S^n$}
The unit sphere $S^n \subset \R^{n+1}$ is orientable. The volume form induced
by the outward unit normal $\nu = x$ is
\[
  \mu = \iota_\nu(\dd x^1 \wedge \cdots \wedge \dd x^{n+1})\big|_{S^n}
  = \sum_{i=1}^{n+1} (-1)^{i-1} x^i\, \dd x^1 \wedge \cdots
    \wedge \widehat{\dd x^i} \wedge \cdots \wedge \dd x^{n+1}\big|_{S^n}.
\]
\end{example}

\begin{example}\label{ex:nonorientable-mobius}
\index{M\"obius band}
The M\"obius band and the Klein bottle are non-orientable: they admit no
nowhere-vanishing top-degree form. Equivalently, they have transition
functions with negative Jacobian determinant that cannot be eliminated by
choosing a different atlas.
\end{example}

% ============================================================
\section{Integration of differential forms}\label{sec:integration-forms}
% ============================================================

\subsection{Integration on $\R^n$}

\begin{definition}\label{def:integration-Rn}
\index{integration!on $\R^n$}
Let $\omega = f\,\dd x^1 \wedge \cdots \wedge \dd x^n$ be a compactly
supported $n$-form on an open subset $U \subseteq \R^n$. We define
\[
  \int_U \omega = \int_U f\,\dd x^1 \cdots \dd x^n,
\]
where the right side is the Lebesgue integral.
\end{definition}

\begin{lemma}[Change of variables for forms]\label{lem:change-variables-forms}
\index{change of variables!for differential forms}
Let $\varphi \colon U \to V$ be an orientation-preserving diffeomorphism
between open subsets of $\R^n$. For any compactly supported $n$-form
$\omega$ on $V$,
\[
  \int_U \varphi^*\omega = \int_V \omega.
\]
If $\varphi$ is orientation-reversing, the integral picks up a sign:
$\int_U \varphi^*\omega = -\int_V \omega$.
\end{lemma}

\begin{proof}
Write $\omega = f\,\dd y^1 \wedge \cdots \wedge \dd y^n$. Then
\[
  \varphi^*\omega = (f \circ \varphi)\,
    \varphi^*(\dd y^1 \wedge \cdots \wedge \dd y^n)
  = (f \circ \varphi)\, \det(D\varphi)\, \dd x^1 \wedge \cdots \wedge \dd x^n.
\]
Therefore
\[
  \int_U \varphi^*\omega = \int_U (f \circ \varphi)\, \det(D\varphi)\,
    \dd x^1 \cdots \dd x^n = \int_V f\, \dd y^1 \cdots \dd y^n = \int_V \omega,
\]
where the second equality is the change of variables formula for Lebesgue
integrals (the Jacobian $\abs{\det D\varphi}$ appears, and
$\det D\varphi > 0$ since $\varphi$ is orientation-preserving).
\end{proof}

\subsection{Integration on manifolds}

\begin{definition}[Integration of $n$-forms on oriented manifolds]
\label{def:integration-manifold}
\index{integration!on manifolds}
Let $M$ be an oriented smooth $n$-manifold, and let $\omega \in \Omega^n(M)$
be a compactly supported $n$-form. Choose a locally finite oriented atlas
$\set{(U_\alpha, \varphi_\alpha)}$ and a subordinate partition of
unity\index{partition of unity} $\set{\rho_\alpha}$. Define
\[
  \int_M \omega = \sum_\alpha \int_{\varphi_\alpha(U_\alpha)}
    (\varphi_\alpha^{-1})^*(\rho_\alpha\, \omega).
\]
\end{definition}

\begin{theorem}\label{thm:integration-well-defined}
The integral $\int_M \omega$ is well-defined: it is independent of the
choice of oriented atlas and partition of unity.
\end{theorem}

\begin{proof}
Let $\set{(V_\beta, \psi_\beta)}$ be another oriented atlas with subordinate
partition of unity $\set{\sigma_\beta}$. Then
\begin{align*}
  \sum_\alpha \int_{\varphi_\alpha(U_\alpha)} (\varphi_\alpha^{-1})^*(\rho_\alpha\, \omega)
  &= \sum_\alpha \sum_\beta \int_{\varphi_\alpha(U_\alpha)} (\varphi_\alpha^{-1})^*(\rho_\alpha\, \sigma_\beta\, \omega) \\
  &= \sum_\alpha \sum_\beta \int_{\psi_\beta(V_\beta)} (\psi_\beta^{-1})^*(\rho_\alpha\, \sigma_\beta\, \omega)
\end{align*}
where the second equality uses Lemma~\ref{lem:change-variables-forms} applied
to the transition map $\psi_\beta \circ \varphi_\alpha^{-1}$
(which is orientation-preserving). Summing over $\alpha$ gives
$\sum_\beta \int (\psi_\beta^{-1})^*(\sigma_\beta\, \omega)$.
\end{proof}

\begin{proposition}[Properties of integration]\label{prop:integration-properties}
\index{integration!properties}
Let $M$ be an oriented compact $n$-manifold.
\begin{enumerate}[label=(\roman*)]
  \item \textbf{Linearity}: $\int_M (a\omega + b\eta) = a\int_M \omega + b\int_M \eta$
    for $a, b \in \R$.
  \item \textbf{Orientation}: If $\overline{M}$ denotes $M$ with reversed
    orientation, then $\int_{\overline{M}} \omega = -\int_M \omega$.
  \item \textbf{Diffeomorphism invariance}: If $f \colon M \to N$ is an
    orientation-preserving diffeomorphism, then
    $\int_M f^*\omega = \int_N \omega$.
\end{enumerate}
\end{proposition}

\begin{example}\label{ex:integration-circle}
\index{integration!on $S^1$}
Parametrize $S^1$ by $\theta \mapsto (\cos\theta, \sin\theta)$
for $\theta \in [0, 2\pi)$. The $1$-form
$\omega = -y\,\dd x + x\,\dd y$ restricted to $S^1$ pulls back to
$\dd\theta$, so $\int_{S^1} \omega = \int_0^{2\pi} \dd\theta = 2\pi$.
\end{example}

\begin{example}\label{ex:integration-S2}
\index{integration!on $S^2$}
On $S^2 \subset \R^3$, the standard area form is
$\mu = x\,\dd y \wedge \dd z + y\,\dd z \wedge \dd x + z\,\dd x \wedge \dd y$.
Its integral $\int_{S^2} \mu = 4\pi$, which is the surface area of the
unit sphere.
\end{example}

\subsection{Compact support}\label{subsec:compact-support}

\begin{definition}[Compactly supported forms]\label{def:compact-support-forms}
\index{compact support}\index{differential form!compact support}
Let $M$ be a smooth manifold. A differential form $\omega \in \Omega^k(M)$
has \textbf{compact support} if the closure of
$\set{p \in M : \omega_p \neq 0}$ is compact. The space of compactly
supported $k$-forms is denoted $\Omega_c^k(M)$.
\end{definition}

\begin{remark}\label{rem:compact-support-integration}
Integration $\int_M \colon \Omega_c^n(M) \to \R$ is well-defined on any
oriented manifold (not necessarily compact): the sum in
Definition~\ref{def:integration-manifold} is finite since the support of
$\omega$ is compact and the atlas is locally finite.
\end{remark}

\begin{remark}\label{rem:compact-support-complex}
The exterior derivative preserves compact support:
$\dd \colon \Omega_c^k(M) \to \Omega_c^{k+1}(M)$. Hence
$(\Omega_c^*(M), \dd)$ is a subcomplex of $(\Omega^*(M), \dd)$. Its
cohomology, called \textbf{compactly supported de Rham cohomology}%
\index{de Rham cohomology!compactly supported} and denoted
$H_c^*(M)$, is important in Poincar\'e duality.
\end{remark}

% --- TikZ figure: compact support ---
\begin{figure}[ht]
\centering
\begin{tikzpicture}[scale=1.0]
  % Manifold
  \draw[thick, fill=blue!5, rounded corners=20pt]
    (-3,-1.2) -- (3,-1.2) -- (3,1.2) -- (-3,1.2) -- cycle;
  \node at (2.4,0.8) {$M$};
  % Support region
  \draw[thick, dashed, fill=orange!15, rounded corners=8pt]
    (-1.2,-0.6) -- (1.2,-0.6) -- (1.2,0.6) -- (-1.2,0.6) -- cycle;
  \node at (0,0) {\small $\operatorname{supp}(\omega)$};
  % Bump function profile below
  \begin{scope}[yshift=-2.5cm]
    \draw[->] (-3,0) -- (3,0) node[right]{\small $M$};
    \draw[->] (0,0) -- (0,1.4);
    \draw[thick, blue!70!black] (-3,0) -- (-1.3,0)
      .. controls (-0.8,0) and (-0.6,1) .. (0,1)
      .. controls (0.6,1) and (0.8,0) .. (1.3,0) -- (3,0);
    \node[blue!70!black, right] at (1.5,0.5) {\small $\omega$};
  \end{scope}
\end{tikzpicture}
\caption{A compactly supported form on $M$: the form vanishes outside a
  compact region.}
\label{fig:compact-support}
\end{figure}

% ============================================================
\section{Summary: dictionary between vector calculus and forms}
\label{sec:dictionary}
% ============================================================

\begin{figure}[ht]
\centering
\renewcommand{\arraystretch}{1.4}
\begin{tabular}{@{}ccc@{}}
\toprule
\textbf{Vector Calculus ($\R^3$)} & & \textbf{Differential Forms} \\
\midrule
scalar field $f$ & $\longleftrightarrow$ & $0$-form $f \in \Omega^0$ \\
vector field $\mathbf{F}$ & $\longleftrightarrow$ & $1$-form or $2$-form \\
$\nabla f$ (gradient) & $\longleftrightarrow$ & $\dd f$ \\
$\nabla \times \mathbf{F}$ (curl) & $\longleftrightarrow$ & $\dd\omega^1$ (as $2$-form) \\
$\nabla \cdot \mathbf{F}$ (divergence) & $\longleftrightarrow$ & $\dd\omega^2$ (as $3$-form) \\
$\operatorname{curl}(\operatorname{grad}) = 0$ & $\longleftrightarrow$ & $\dd^2 = 0$ \\
$\operatorname{div}(\operatorname{curl}) = 0$ & $\longleftrightarrow$ & $\dd^2 = 0$ \\
\bottomrule
\end{tabular}
\caption{Correspondence between vector calculus on $\R^3$ and differential forms.}
\label{fig:dictionary-forms}
\end{figure}

% ============================================================
\section{Exercises}\label{sec:exercises-forms}
% ============================================================

\begin{exercise}\label{exo:wedge-computation}
\index{wedge product!exercises}
Let $\alpha = 3\,\dd x - 2\,\dd y + \dd z$ and
$\beta = \dd x + 4\,\dd y - 2\,\dd z$ be $1$-forms on $\R^3$.
Compute $\alpha \wedge \beta$ and verify that
$\alpha \wedge \beta = -\beta \wedge \alpha$.
\end{exercise}

\begin{exercise}\label{exo:exterior-derivative-computation}
\index{exterior derivative!exercises}
Compute the exterior derivative of the following forms on $\R^3$:
\begin{enumerate}[label=(\alph*)]
  \item $\omega = xy^2\,\dd x + e^z\,\dd y + xz\,\dd z$.
  \item $\eta = x^2\,\dd y \wedge \dd z + y^2\,\dd z \wedge \dd x
    + z^2\,\dd x \wedge \dd y$.
  \item $f = \sin(xyz)$.
\end{enumerate}
Verify that $\dd^2 = 0$ for each.
\end{exercise}

\begin{exercise}\label{exo:pullback-computation}
\index{pullback!exercises}
Let $f \colon \R^2 \to \R^3$ be defined by
$f(u,v) = (u^2, uv, v^2)$. Compute $f^*\omega$ where
$\omega = x\,\dd y - y\,\dd z$.
\end{exercise}

\begin{exercise}\label{exo:cartan-formula-verification}
\index{Cartan formula!exercise}
On $\R^3$, let $X = x\,\frac{\partial}{\partial y}
- y\,\frac{\partial}{\partial x}$ (the rotational vector field) and
$\omega = x\,\dd x + y\,\dd y + z\,\dd z$. Verify Cartan's formula
$\mathcal{L}_X \omega = \iota_X \dd\omega + \dd(\iota_X \omega)$
by computing both sides explicitly.
\end{exercise}

\begin{exercise}\label{exo:closed-not-exact}
\index{closed form}\index{exact form}
Consider the $1$-form on $\R^2 \setminus \set{0}$:
\[
  \omega = \frac{-y\,\dd x + x\,\dd y}{x^2 + y^2}.
\]
\begin{enumerate}[label=(\alph*)]
  \item Show that $\dd\omega = 0$ (i.e., $\omega$ is closed).
  \item Show that $\int_{S^1} \omega = 2\pi$, where $S^1$ is the unit circle.
  \item Conclude that $\omega$ is not exact on $\R^2 \setminus \set{0}$.
\end{enumerate}
\end{exercise}

\begin{exercise}\label{exo:volume-form-Sn}
\index{volume form!on $S^n$}
Show that the $n$-form on $\R^{n+1}$,
\[
  \mu = \sum_{i=1}^{n+1} (-1)^{i-1} x^i\, \dd x^1 \wedge \cdots
    \wedge \widehat{\dd x^i} \wedge \cdots \wedge \dd x^{n+1},
\]
restricts to a volume form on $S^n$, and compute $\int_{S^2} \mu$.
\end{exercise}

\begin{exercise}\label{exo:lie-derivative-computation}
\index{Lie derivative!exercises}
Let $X = x\,\frac{\partial}{\partial x} + y\,\frac{\partial}{\partial y}$
on $\R^2$ and $\omega = \dd x \wedge \dd y$.
\begin{enumerate}[label=(\alph*)]
  \item Compute $\mathcal{L}_X \omega$ using Cartan's formula.
  \item Compute $\mathcal{L}_X \omega$ directly from the flow $\phi_t$ of $X$.
  \item Verify that the results agree.
\end{enumerate}
\end{exercise}

\begin{exercise}\label{exo:integration-form}
\index{integration!exercises}
Let $\omega = (x^2 + y^2)\,\dd x \wedge \dd y$ on $\R^2$ and let
$D = \set{(x,y) \in \R^2 : x^2 + y^2 \leq 1}$ be the closed unit disk.
Compute $\int_D \omega$.
\end{exercise}

\begin{exercise}\label{exo:orientation-nform}
\index{orientation!exercises}
Show that the M\"obius band is non-orientable by showing that no
nowhere-vanishing $2$-form exists. \textit{Hint}: consider a loop
traversing the core circle once and track the orientation of a local
frame.
\end{exercise}

\begin{exercise}\label{exo:interior-product-computation}
\index{interior product!exercises}
Let $\omega = \dd x \wedge \dd y \wedge \dd z$ on $\R^3$ and
$X = y\,\frac{\partial}{\partial x} + z\,\frac{\partial}{\partial y}
+ x\,\frac{\partial}{\partial z}$.
\begin{enumerate}[label=(\alph*)]
  \item Compute $\iota_X \omega$.
  \item Compute $\iota_X(\iota_X \omega)$ and verify it is zero.
\end{enumerate}
\end{exercise}

% === END OF CHUNK ch5-6 ===
% === START OF CHUNK ch7-8 ===

%%%%%%%%%%%%%%%%%%%%%%%%%%%%%%%%%%%%%%%%%%%%%%%%%%%%%%%%%%%%%%%%%%%%
%%  CHAPTER 7 — STOKES' THEOREM
%%%%%%%%%%%%%%%%%%%%%%%%%%%%%%%%%%%%%%%%%%%%%%%%%%%%%%%%%%%%%%%%%%%%
\chapter{Stokes' Theorem}\label{ch:stokes}
\minitoc\vspace{1em}

The classical integral theorems of vector calculus---Green's theorem,
the divergence theorem, and the classical Stokes theorem in~$\R^3$---are
all special cases of a single, elegant statement on oriented manifolds
with boundary.  In this chapter we develop the theory of integration on
chains, prove the general Stokes theorem, and draw fundamental
consequences for de~Rham cohomology.\index{Stokes' theorem}

%--------------------------------------------------------------------
\section{Smooth singular chains}\label{sec:chains}
%--------------------------------------------------------------------

\begin{definition}[Standard $k$-simplex]\label{def:standard-simplex}
\index{simplex!standard}
The \emph{standard $k$-simplex} is the subset
\[
  \Delta^k = \set{(t_0,t_1,\dots,t_k)\in\R^{k+1} \;\big|\;
    t_i\ge 0,\; \sum_{i=0}^{k} t_i = 1}.
\]
Equivalently, using coordinates $(t_1,\dots,t_k)$ with $t_0 = 1 - t_1 - \cdots - t_k$, we may
identify $\Delta^k$ with the convex hull of the standard basis vectors
$e_0,e_1,\dots,e_k$ in~$\R^{k+1}$.
\end{definition}

\begin{definition}[Smooth singular simplex]\label{def:sing-simplex}
\index{singular simplex}
Let $M$ be a smooth manifold.  A \emph{smooth singular $k$-simplex} in~$M$
is a smooth map $\sigma\colon \Delta^k \to M$ (smooth meaning it
extends to a smooth map on an open neighbourhood of $\Delta^k$
in the affine hyperplane containing~$\Delta^k$).
\end{definition}

\begin{definition}[Chain group]\label{def:chain-group}
\index{chain!singular}
The \emph{smooth singular chain group} $C_k(M)$ is the free abelian
group generated by all smooth singular $k$-simplices in~$M$.  An
element $c = \sum_i a_i \sigma_i$ with $a_i\in\Z$ is called a
\emph{smooth singular $k$-chain}.
\end{definition}

\begin{definition}[Boundary operator]\label{def:boundary-op}
\index{boundary operator}
For a smooth singular $k$-simplex $\sigma\colon\Delta^k\to M$, define
the \emph{boundary} $\partial\sigma\in C_{k-1}(M)$ by
\[
  \partial\sigma = \sum_{i=0}^{k} (-1)^i\, \sigma\circ F_i^k,
\]
where $F_i^k\colon \Delta^{k-1}\to \Delta^k$ is the $i$-th face map,
\[
  F_i^k(t_0,\dots,t_{k-1}) = (t_0,\dots,t_{i-1},0,t_i,\dots,t_{k-1}).
\]
Extend $\partial$ linearly to all of $C_k(M)$.
\end{definition}

\begin{proposition}\label{prop:dd-zero}
$\partial\circ\partial = 0$.
\end{proposition}

\begin{proof}
For any smooth singular $k$-simplex $\sigma$ one computes
\[
  \partial(\partial\sigma)
  = \sum_{i=0}^k\sum_{j=0}^{k-1} (-1)^{i+j}\, \sigma\circ F_i^k\circ F_j^{k-1}.
\]
The standard identity $F_i^k\circ F_j^{k-1} = F_j^k\circ F_{i-1}^{k-1}$
for $j < i$ shows that each term appears twice with opposite signs, hence the
sum vanishes.
\end{proof}

%--------------------------------------------------------------------
\section{Integration of forms on chains}\label{sec:integration-chains}
%--------------------------------------------------------------------

\begin{definition}[Integral over a simplex]\label{def:integral-simplex}
\index{integration!on chains}
Let $\omega$ be a smooth $k$-form on~$M$ and $\sigma\colon\Delta^k\to M$
a smooth singular $k$-simplex.  The \emph{integral of~$\omega$ over~$\sigma$}
is
\[
  \int_\sigma \omega = \int_{\Delta^k} \sigma^*\omega,
\]
where the right-hand side is the ordinary Lebesgue integral of the
pulled-back form on~$\Delta^k$.
\end{definition}

\begin{definition}[Integral over a chain]\label{def:integral-chain}
For a $k$-chain $c = \sum_i a_i\sigma_i$, define
\[
  \int_c \omega = \sum_i a_i \int_{\sigma_i} \omega.
\]
\end{definition}

\begin{remark}\label{rem:chain-integration-linear}
The map $\omega\mapsto \int_c\omega$ is $\R$-linear in~$\omega$,
and the map $c\mapsto \int_c\omega$ is $\Z$-linear (and extends
to $\R$-linear on real chains $C_k(M;\R)$).
\end{remark}

\begin{proposition}[Stokes for a single simplex]\label{prop:stokes-simplex}
\index{Stokes' theorem!for simplices}
Let $\omega\in\Omega^{k-1}(M)$ and $\sigma\colon\Delta^k\to M$ a smooth
singular $k$-simplex.  Then
\[
  \int_\sigma \dd\omega = \int_{\partial\sigma} \omega.
\]
\end{proposition}

\begin{proof}
By the naturality of pullback and the fact that $\dd$ commutes with
pullback, $\sigma^*(\dd\omega) = \dd(\sigma^*\omega)$.  Write
$\sigma^*\omega = f(t)\,\dd t_1\wedge\cdots\wedge \dd t_{k-1}$ on
$\Delta^k$ (using an appropriate coordinate representation).  The
classical Stokes theorem for $\Delta^k\subset\R^k$ (a domain with
piecewise smooth boundary) then gives
\[
  \int_{\Delta^k} \dd(\sigma^*\omega)
  = \int_{\partial\Delta^k} \sigma^*\omega
  = \sum_{i=0}^k (-1)^i \int_{\Delta^{k-1}} (F_i^k)^*(\sigma^*\omega)
  = \int_{\partial\sigma}\omega. \qedhere
\]
\end{proof}

\begin{corollary}[Stokes for chains]\label{cor:stokes-chain}
For any smooth $k$-chain $c$ and $\omega\in\Omega^{k-1}(M)$,
\[
  \int_c \dd\omega = \int_{\partial c} \omega.
\]
\end{corollary}

%--------------------------------------------------------------------
\section{Oriented manifolds and integration}\label{sec:oriented-integration}
%--------------------------------------------------------------------

We briefly recall the integration theory on oriented manifolds to state
Stokes' theorem in its manifold form.

\begin{definition}[Integration on an oriented manifold]\label{def:integration-manifold}
\index{integration!on manifolds}
Let $M$ be an oriented smooth $n$-manifold (possibly with boundary),
and let $\omega\in\Omega^n_c(M)$ be a compactly supported $n$-form.
Choose a locally finite atlas $\set{(U_\alpha,\varphi_\alpha)}$ of
positively oriented charts and a subordinate partition of
unity~$\set{\rho_\alpha}$.  Define
\[
  \int_M \omega = \sum_\alpha \int_{\varphi_\alpha(U_\alpha)}
    (\varphi_\alpha^{-1})^*(\rho_\alpha\,\omega).
\]
This is independent of all choices.
\end{definition}

\begin{definition}[Induced boundary orientation]\label{def:boundary-orientation}
\index{boundary!orientation}
Let $M$ be an oriented $n$-manifold with boundary.  The
\emph{induced orientation} on $\partial M$ is defined by declaring
that an ordered basis $(v_1,\dots,v_{n-1})$ of $T_p(\partial M)$ is
positively oriented if and only if $(-\nu, v_1,\dots,v_{n-1})$ is a
positively oriented basis of $T_pM$, where $\nu$ is the outward-pointing
normal vector at~$p$.
\end{definition}

%--------------------------------------------------------------------
\section{Stokes' theorem: statement and proof}\label{sec:stokes-proof}
%--------------------------------------------------------------------

\begin{theorem}[Stokes' theorem]\label{thm:stokes}
\index{Stokes' theorem!general}
Let $M$ be an oriented smooth $n$-manifold with boundary (possibly
$\partial M = \emptyset$), and let $\omega\in\Omega^{n-1}_c(M)$ be a
compactly supported $(n{-}1)$-form.  Then
\[
  \boxed{\int_M \dd\omega = \int_{\partial M} \omega,}
\]
where $\partial M$ carries the induced boundary orientation.
\end{theorem}

\begin{proof}
The proof proceeds in three steps.

\medskip\noindent\textbf{Step 1: Local statement in~$\R^n$ and in the
half-space $\mathbb{H}^n$.}\index{half-space}

Write $\mathbb{H}^n = \set{x\in\R^n \mid x_n\ge 0}$.  We first prove the
theorem for compactly supported forms on $\R^n$ and on $\mathbb{H}^n$.

\emph{Case~1: $\omega$ supported in $\R^n$.}\;  Write
\[
  \omega = \sum_{i=1}^n (-1)^{i-1} f_i(x)\, \dd x_1\wedge\cdots
  \wedge\widehat{\dd x_i}\wedge\cdots\wedge \dd x_n.
\]
Then
\[
  \dd\omega = \left(\sum_{i=1}^n \frac{\partial f_i}{\partial x_i}\right)
  \dd x_1\wedge\cdots\wedge \dd x_n.
\]
By Fubini's theorem,
\[
  \int_{\R^n}\dd\omega
  = \sum_{i=1}^n \int_{\R^{n-1}}\left(
    \int_{-\infty}^{+\infty} \frac{\partial f_i}{\partial x_i}\,\dd x_i
  \right)\dd x_1\cdots\widehat{\dd x_i}\cdots\dd x_n.
\]
Each inner integral equals $f_i\big|_{x_i=-\infty}^{x_i=+\infty} = 0$
since $f_i$ is compactly supported.  Thus $\int_{\R^n}\dd\omega = 0$,
and the right-hand side $\int_{\partial\R^n}\omega = 0$ trivially
(since $\partial\R^n = \emptyset$).

\emph{Case~2: $\omega$ supported in $\mathbb{H}^n$.}\;  The same
argument applies to all terms with $i < n$: the integrals vanish because
$f_i$ is compactly supported in the $x_i$-direction.  For the last term
($i = n$),
\[
  \int_{\mathbb{H}^n} \frac{\partial f_n}{\partial x_n}
  \,\dd x_1\cdots \dd x_n
  = \int_{\R^{n-1}}\left(
    \int_0^{+\infty}\frac{\partial f_n}{\partial x_n}\,\dd x_n
  \right) \dd x_1\cdots \dd x_{n-1}.
\]
Since $f_n$ is compactly supported, $f_n(x',+\infty) = 0$, and hence
\[
  \int_0^{+\infty}\frac{\partial f_n}{\partial x_n}\,\dd x_n
  = -f_n(x_1,\dots,x_{n-1},0).
\]
Therefore
\[
  \int_{\mathbb{H}^n}\dd\omega
  = (-1)^n\cdot(-1)\int_{\R^{n-1}} f_n(x',0)\,\dd x_1\cdots \dd x_{n-1}.
\]
On the other hand, $\partial\mathbb{H}^n = \R^{n-1}\times\set{0}$
with the induced orientation (the outward normal is $-e_n$), so
\[
  \int_{\partial\mathbb{H}^n}\omega
  = (-1)^{n-1}\int_{\R^{n-1}} f_n(x',0)\,\dd x_1\cdots \dd x_{n-1},
\]
which matches.

\medskip\noindent\textbf{Step 2: Partition of unity.}\index{partition of unity}

Choose a locally finite atlas $\set{(U_\alpha,\varphi_\alpha)}_\alpha$
of positively oriented charts for~$M$ such that each $\varphi_\alpha(U_\alpha)$
is either an open subset of~$\R^n$ (for interior charts) or an open
subset of~$\mathbb{H}^n$ (for boundary charts).  Let $\set{\rho_\alpha}$
be a subordinate smooth partition of unity.  Then
$\omega = \sum_\alpha \rho_\alpha\omega$ (locally finite sum), and each
$\rho_\alpha\omega$ is compactly supported in~$U_\alpha$.

\medskip\noindent\textbf{Step 3: Reassembly.}

Since $\dd(\rho_\alpha\omega)$ is also compactly supported in $U_\alpha$,
Step~1 gives
\[
  \int_M \dd(\rho_\alpha\omega) = \int_{\partial M} \rho_\alpha\omega
\]
for each~$\alpha$.  Summing over~$\alpha$ (the sum is locally finite,
so we may interchange summation and integration):
\[
  \int_M \dd\omega
  = \int_M \dd\!\left(\sum_\alpha \rho_\alpha\omega\right)
  = \sum_\alpha \int_M \dd(\rho_\alpha\omega)
  = \sum_\alpha \int_{\partial M}\rho_\alpha\omega
  = \int_{\partial M}\omega. \qedhere
\]
\end{proof}

\begin{remark}\label{rem:stokes-closed-manifolds}
When $M$ is a closed manifold (compact, without boundary), Stokes'
theorem gives $\int_M\dd\omega = 0$ for every
$\omega\in\Omega^{n-1}(M)$.  This is the starting point for de~Rham
cohomology.
\end{remark}

%--------------------------------------------------------------------
\section{Classical special cases}\label{sec:stokes-classical}
%--------------------------------------------------------------------

\begin{corollary}[Green's theorem]\label{cor:green}
\index{Green's theorem}
Let $D\subset\R^2$ be a compact region with smooth boundary
$\partial D$ oriented counterclockwise, and let $P,Q\colon D\to\R$
be smooth.  Then
\[
  \oint_{\partial D}(P\,\dd x + Q\,\dd y)
  = \iint_D \left(\frac{\partial Q}{\partial x}
    - \frac{\partial P}{\partial y}\right)\dd x\,\dd y.
\]
\end{corollary}

\begin{proof}
Set $\omega = P\,\dd x + Q\,\dd y\in\Omega^1(D)$.  Then
$\dd\omega = \bigl(\frac{\partial Q}{\partial x} - \frac{\partial P}{\partial y}\bigr)\dd x\wedge\dd y$.
The result is 
ef{thm:stokes} with $M = D$, $n = 2$.
\end{proof}

\begin{corollary}[Divergence theorem]\label{cor:divergence}
\index{divergence theorem}
Let $\Omega\subset\R^3$ be a compact domain with smooth boundary
$\partial\Omega$, and let $\mathbf{F} = (F_1,F_2,F_3)$ be a smooth
vector field on~$\Omega$.  Then
\[
  \iiint_\Omega \operatorname{div}\mathbf{F}\;\dd V
  = \oiint_{\partial\Omega} \mathbf{F}\cdot\mathbf{n}\;\dd S,
\]
where $\mathbf{n}$ is the outward unit normal.
\end{corollary}

\begin{proof}
Define the $2$-form
$\omega = F_1\,\dd y\wedge\dd z + F_2\,\dd z\wedge\dd x + F_3\,\dd x\wedge\dd y$.
Then $\dd\omega = \operatorname{div}\mathbf{F}\,\dd x\wedge\dd y\wedge\dd z$.
Apply 
ef{thm:stokes}.
\end{proof}

\begin{corollary}[Classical Stokes theorem in $\R^3$]\label{cor:classical-stokes}
\index{Stokes' theorem!classical}
Let $S\subset\R^3$ be a compact oriented surface with smooth
boundary~$\partial S$, and let $\mathbf{F}$ be a smooth vector field.
Then
\[
  \iint_S (\nabla\times\mathbf{F})\cdot\mathbf{n}\;\dd S
  = \oint_{\partial S} \mathbf{F}\cdot\dd\mathbf{r}.
\]
\end{corollary}

\begin{proof}
Define $\omega = F_1\,\dd x + F_2\,\dd y + F_3\,\dd z$.  Then $\dd\omega$
corresponds to the curl $\nabla\times\mathbf{F}$ under the standard
identification of $2$-forms with vector fields via the Hodge star.
Apply 
ef{thm:stokes} with $M = S$.
\end{proof}

\begin{example}\label{ex:stokes-torus}
Let $T^2$ be the $2$-torus with the standard orientation, and let
$\omega\in\Omega^1(T^2)$.  Since $\partial T^2 = \emptyset$,
Stokes' theorem gives
$\int_{T^2}\dd\omega = 0$.
In particular, if $\omega = f\,\dd\theta_1 + g\,\dd\theta_2$ in
the standard angular coordinates, then
\[
  \int_{T^2}\left(\frac{\partial g}{\partial\theta_1}
  - \frac{\partial f}{\partial\theta_2}\right)\dd\theta_1\wedge\dd\theta_2 = 0.
\]
\end{example}

\begin{example}[Area via Green's theorem]\label{ex:area-green}
\index{Green's theorem!area computation}
Let $D\subset\R^2$ be a compact region with smooth boundary $\partial D$
oriented counterclockwise.  Choosing $P = -y/2$ and $Q = x/2$ in

ef{cor:green}, we obtain the area formula
\[
  \operatorname{Area}(D) = \frac{1}{2}\oint_{\partial D}(x\,\dd y - y\,\dd x).
\]
For the ellipse $\frac{x^2}{a^2}+\frac{y^2}{b^2}=1$, parametrise
$\gamma(t) = (a\cos t, b\sin t)$ for $t\in[0,2\pi]$.  Then
\[
  \operatorname{Area}(D) = \frac{1}{2}\int_0^{2\pi}
  (a\cos t\cdot b\cos t + b\sin t\cdot a\sin t)\,\dd t
  = \frac{ab}{2}\int_0^{2\pi}\dd t = \pi ab.
\]
\end{example}

\begin{example}[Divergence theorem: flux computation]\label{ex:divergence-flux}
Let $\mathbf{F}(x,y,z) = (x^3, y^3, z^3)$ and let $\Omega$ be the unit
ball in~$\R^3$.  Then $\operatorname{div}\mathbf{F} = 3(x^2+y^2+z^2) = 3r^2$.
By the divergence theorem,
\[
  \oiint_{S^2}\mathbf{F}\cdot\mathbf{n}\,\dd S
  = \iiint_\Omega 3r^2\,\dd V
  = 3\int_0^1 r^2\cdot 4\pi r^2\,\dd r
  = 12\pi\int_0^1 r^4\,\dd r = \frac{12\pi}{5}.
\]
\end{example}

%--------------------------------------------------------------------
\section{De Rham cohomology}\label{sec:derham}
%--------------------------------------------------------------------

\begin{definition}[Closed and exact forms]\label{def:closed-exact}
\index{closed form}\index{exact form}
A differential form $\omega\in\Omega^k(M)$ is \emph{closed} if
$\dd\omega = 0$, and \emph{exact} if $\omega = \dd\eta$ for some
$\eta\in\Omega^{k-1}(M)$.  Since $\dd^2 = 0$, every exact form is
closed.  We write
\[
  Z^k(M) = \ker(\dd\colon\Omega^k(M)\to\Omega^{k+1}(M)),\qquad
  B^k(M) = \im(\dd\colon\Omega^{k-1}(M)\to\Omega^k(M)).
\]
\end{definition}

\begin{definition}[De Rham cohomology]\label{def:derham-cohomology}
\index{de Rham cohomology}
The \emph{$k$-th de Rham cohomology group} of $M$ is the quotient
vector space
\[
  H^k_{\mathrm{dR}}(M) = \frac{Z^k(M)}{B^k(M)}
  = \frac{\ker\dd}{\im\dd}.
\]
The equivalence class $[\omega]$ of a closed form $\omega$ is called
its \emph{cohomology class}.
\end{definition}

\begin{remark}\label{rem:derham-functor}
If $f\colon M\to N$ is smooth, then pullback $f^*\colon\Omega^k(N)\to\Omega^k(M)$
satisfies $f^*\circ\dd = \dd\circ f^*$, hence $f^*$ descends to a
linear map $f^*\colon H^k_{\mathrm{dR}}(N)\to H^k_{\mathrm{dR}}(M)$.
Moreover, $(\Id_M)^* = \Id$ and $(g\circ f)^* = f^*\circ g^*$, so
$H^k_{\mathrm{dR}}$ is a contravariant functor from smooth manifolds
to real vector spaces.
\end{remark}

\begin{example}\label{ex:derham-point}
For $M = \set{\mathrm{pt}}$, the only nonzero differential forms are
$0$-forms (i.e., constants), and $\dd f = 0$ for any constant~$f$.
Hence $H^0_{\mathrm{dR}}(\mathrm{pt}) \cong \R$ and
$H^k_{\mathrm{dR}}(\mathrm{pt}) = 0$ for $k\ge 1$.
\end{example}

\begin{proposition}\label{prop:H0-components}
\index{de Rham cohomology!degree zero}
For any smooth manifold~$M$,
$H^0_{\mathrm{dR}}(M) \cong \R^{\pi_0(M)}$,
where $\pi_0(M)$ denotes the set of connected components of~$M$.
\end{proposition}

\begin{proof}
A $0$-form $f\in C^\infty(M)$ is closed if and only if $\dd f = 0$, i.e.,
$f$ is locally constant.  On a connected manifold a locally constant
function is constant, so $H^0_{\mathrm{dR}}(M) = Z^0(M)\cong\R$.
For a disconnected manifold the result follows by taking one constant
per component.
\end{proof}

\begin{remark}[Mayer--Vietoris sequence]\label{rem:mayer-vietoris}
\index{Mayer--Vietoris sequence}
If $M = U\cup V$ with $U,V$ open, there is a long exact sequence
\[
  \cdots \to H^{k-1}_{\mathrm{dR}}(U\cap V)
  \xrightarrow{\;\delta\;}
  H^k_{\mathrm{dR}}(M)
  \xrightarrow{\;(i_U^*,i_V^*)\;}
  H^k_{\mathrm{dR}}(U)\oplus H^k_{\mathrm{dR}}(V)
  \xrightarrow{\;j_U^*-j_V^*\;}
  H^k_{\mathrm{dR}}(U\cap V) \to \cdots
\]
where $i_U\colon U\hookrightarrow M$, $i_V\colon V\hookrightarrow M$,
$j_U\colon U\cap V\hookrightarrow U$, $j_V\colon U\cap V\hookrightarrow V$
are the inclusions.  This is an extremely powerful computational tool:
combined with homotopy invariance, it allows one to compute the
de~Rham cohomology of many manifolds by induction.
\end{remark}

\begin{example}[Cohomology of $S^1$]\label{ex:cohomology-S1}
\index{de Rham cohomology!circle}
Cover $S^1$ by two open arcs $U$ and $V$, each diffeomorphic to~$\R$,
with $U\cap V$ the disjoint union of two open intervals (each
diffeomorphic to~$\R$).  The Mayer--Vietoris sequence in degree~$0$ reads
\[
  0 \to H^0_{\mathrm{dR}}(S^1) \to \R\oplus\R
  \xrightarrow{\;\alpha\;} \R\oplus\R
  \xrightarrow{\;\delta\;} H^1_{\mathrm{dR}}(S^1) \to 0.
\]
The map $\alpha(a,b) = (a-b, a-b)$, so $\ker\alpha\cong\R$ and
$\im\alpha\cong\R$.  Hence $H^0_{\mathrm{dR}}(S^1)\cong\R$ and
$H^1_{\mathrm{dR}}(S^1)\cong\R^2/\R\cong\R$.  Since $\dim S^1 = 1$,
$H^k_{\mathrm{dR}}(S^1) = 0$ for $k\ge 2$.
\end{example}

%--------------------------------------------------------------------
\section{The Poincar\'e lemma}\label{sec:poincare-lemma}
%--------------------------------------------------------------------

\begin{definition}[Star-shaped domain]\label{def:star-shaped}
\index{star-shaped}
A subset $U\subset\R^n$ is \emph{star-shaped with respect to a
point~$p\in U$} if for every $x\in U$, the line segment from $p$ to~$x$
lies entirely in~$U$.
\end{definition}

\begin{theorem}[Poincar\'e lemma]\label{thm:poincare-lemma}
\index{Poincar\'e lemma}
Let $U\subset\R^n$ be a star-shaped open set.  Then every closed $k$-form
on~$U$ with $k\ge 1$ is exact.  Equivalently,
\[
  H^k_{\mathrm{dR}}(U) = 0 \qquad\text{for all } k\ge 1.
\]
\end{theorem}

\begin{proof}
Without loss of generality assume $U$ is star-shaped with respect to the
origin.  We construct a \emph{chain homotopy operator}\index{chain homotopy}
$K\colon\Omega^k(U)\to\Omega^{k-1}(U)$ satisfying
\begin{equation}\label{eq:homotopy-formula}
  \dd K + K\dd = \Id \quad\text{on } \Omega^k(U),\quad k\ge 1.
\end{equation}
For a $k$-form $\omega = \sum_{I} a_I(x)\,\dd x_{i_1}\wedge\cdots
\wedge\dd x_{i_k}$ (the sum over increasing multi-indices~$I$), define
\[
  (K\omega)(x) = \sum_I \sum_{j=1}^k (-1)^{j-1}
  \left(\int_0^1 t^{k-1}\, a_I(tx)\,\dd t\right)
  x_{i_j}\,\dd x_{i_1}\wedge\cdots\wedge\widehat{\dd x_{i_j}}\wedge
  \cdots\wedge\dd x_{i_k}.
\]
A direct (if lengthy) computation using the Leibniz rule and
$\frac{\dd}{\dd t}[a_I(tx)] = \sum_\ell x_\ell \frac{\partial a_I}{\partial x_\ell}(tx)$
verifies~\eqref{eq:homotopy-formula}.

Now if $\omega$ is closed ($\dd\omega = 0$) and $k\ge 1$,
then $\omega = (\dd K + K\dd)\omega = \dd(K\omega)$, so $\omega$
is exact with primitive $\eta = K\omega$.
\end{proof}

\begin{corollary}\label{cor:poincare-Rn}
$H^k_{\mathrm{dR}}(\R^n) = 0$ for all $k\ge 1$, and
$H^0_{\mathrm{dR}}(\R^n)\cong\R$.
\end{corollary}

\begin{example}[Explicit primitive via homotopy operator]\label{ex:explicit-primitive}
Consider the closed $2$-form $\omega = 2xy\,\dd x\wedge\dd y$ on
$\R^2$ (star-shaped with respect to the origin).  The homotopy operator
from the proof of 
ef{thm:poincare-lemma} gives
\[
  (K\omega)(x,y) = \left(\int_0^1 t\cdot 2(tx)(ty)\,\dd t\right)
  \bigl(x\,\dd y - y\,\dd x\bigr)
  = \left(\int_0^1 2t^3 xy\,\dd t\right)(x\,\dd y - y\,\dd x).
\]
Evaluating the integral: $\int_0^1 2t^3\,\dd t = \frac{1}{2}$, so
\[
  K\omega = \frac{xy}{2}(x\,\dd y - y\,\dd x)
  = \frac{x^2y}{2}\,\dd y - \frac{xy^2}{2}\,\dd x.
\]
One verifies: $\dd(K\omega) = \dd\!\left(\frac{x^2y}{2}\,\dd y
- \frac{xy^2}{2}\,\dd x\right)
= \left(\frac{2xy}{2}+\frac{y^2}{2}+\frac{x^2}{2}-\frac{2xy}{2}\right)
\dd x\wedge\dd y$\,... which does not simplify to $\omega$ because
one must use the correct formula more carefully.  In fact, let us
apply the formula directly for $\omega = f(x,y)\,\dd x\wedge\dd y$
with $f = 2xy$:
\begin{align*}
  K\omega &= \left(\int_0^1 t\cdot f(tx,ty)\,\dd t\right)x\,\dd y
  - \left(\int_0^1 t\cdot f(tx,ty)\,\dd t\right)y\,\dd x \\
  &= x\left(\int_0^1 2t^3xy\,\dd t\right)\dd y
  - y\left(\int_0^1 2t^3xy\,\dd t\right)\dd x \\
  &= \frac{x^2y}{2}\,\dd y - \frac{xy^2}{2}\,\dd x.
\end{align*}
Computing $\dd(K\omega)$:
\[
  \dd\!\left(-\frac{xy^2}{2}\,\dd x + \frac{x^2y}{2}\,\dd y\right)
  = \left(\frac{2xy}{2}+\frac{2xy}{2}\right)\dd x\wedge\dd y
  = 2xy\,\dd x\wedge\dd y = \omega. \checkmark
\]
\end{example}

\begin{example}\label{ex:angle-form}
\index{angle form}
Let $\omega = \frac{-y\,\dd x + x\,\dd y}{x^2+y^2}$ on
$\R^2\setminus\set{0}$.  One checks $\dd\omega = 0$, yet $\omega$ is
not exact: $\int_{S^1}\omega = 2\pi \ne 0$ (and Stokes' theorem would give
$0$ if $\omega = \dd f$).  This shows $H^1_{\mathrm{dR}}(\R^2\setminus\set{0})\ne 0$.
In fact, $H^1_{\mathrm{dR}}(\R^2\setminus\set{0})\cong\R$, generated
by $[\omega/2\pi]$.
\end{example}

%--------------------------------------------------------------------
\section{Homotopy invariance}\label{sec:homotopy-invariance}
%--------------------------------------------------------------------

\begin{theorem}[Homotopy invariance of de Rham cohomology]
  \label{thm:homotopy-invariance}
\index{homotopy invariance!de Rham cohomology}
If $f,g\colon M\to N$ are smoothly homotopic (i.e., there exists a smooth
map $H\colon M\times[0,1]\to N$ with $H(\cdot,0) = f$ and $H(\cdot,1)=g$),
then $f^* = g^*\colon H^k_{\mathrm{dR}}(N)\to H^k_{\mathrm{dR}}(M)$ for
all~$k$.
\end{theorem}

\begin{proof}
We construct a \emph{cochain homotopy} $h\colon\Omega^k(N)\to\Omega^{k-1}(M)$
satisfying
\[
  g^* - f^* = \dd\circ h + h\circ\dd.
\]
Let $\iota_t\colon M\hookrightarrow M\times[0,1]$, $\iota_t(p) = (p,t)$.
For $\omega\in\Omega^k(N)$, set
\[
  h(\omega) = \int_0^1 \iota_t^*\bigl(\iota_{\partial/\partial t}
  H^*\omega\bigr)\,\dd t,
\]
where $\iota_{\partial/\partial t}$ denotes interior multiplication by
$\partial/\partial t$.  Cartan's magic formula gives
\[
  \mathcal{L}_{\partial/\partial t}(H^*\omega)
  = \dd(\iota_{\partial/\partial t} H^*\omega)
  + \iota_{\partial/\partial t}\dd(H^*\omega),
\]
and integrating from $t=0$ to $t=1$ yields the desired identity.
Hence if $\dd\omega = 0$, then $g^*\omega - f^*\omega = \dd(h\omega)$,
so $[g^*\omega] = [f^*\omega]$.
\end{proof}

\begin{example}\label{ex:homotopy-invariance-Rn-minus-0}
\index{de Rham cohomology!punctured space}
The inclusion $\iota\colon S^{n-1}\hookrightarrow\R^n\setminus\set{0}$
and the retraction $r\colon\R^n\setminus\set{0}\to S^{n-1}$,
$r(x) = x/\abs{x}$, satisfy $r\circ\iota = \Id_{S^{n-1}}$ and
$\iota\circ r\simeq\Id_{\R^n\setminus\set{0}}$ (via the straight-line
homotopy $H(x,t) = (1-t)x + t\,x/\abs{x}$, which stays in
$\R^n\setminus\set{0}$).  Hence by homotopy invariance,
\[
  H^k_{\mathrm{dR}}(\R^n\setminus\set{0})
  \cong H^k_{\mathrm{dR}}(S^{n-1})
  \cong \begin{cases} \R & \text{if } k = 0 \text{ or } k = n-1, \\
  0 & \text{otherwise.}\end{cases}
\]
For $n = 2$, this recovers $H^1_{\mathrm{dR}}(\R^2\setminus\set{0})\cong\R$,
generated by the angle form of 
ef{ex:angle-form}.
\end{example}

\begin{corollary}[Homotopy equivalence]\label{cor:htpy-equiv}
\index{homotopy equivalence}
If $f\colon M\to N$ is a smooth homotopy equivalence, then
$f^*\colon H^k_{\mathrm{dR}}(N)\xrightarrow{\;\sim\;} H^k_{\mathrm{dR}}(M)$
is an isomorphism for all~$k$.
\end{corollary}

\begin{corollary}\label{cor:contractible}
If $M$ is contractible, then $H^k_{\mathrm{dR}}(M) = 0$ for $k\ge 1$
and $H^0_{\mathrm{dR}}(M)\cong\R$.  In particular, the Poincar\'e
lemma is a special case of homotopy invariance.
\end{corollary}

%--------------------------------------------------------------------
\section{Degree via integration}\label{sec:degree-integration}
%--------------------------------------------------------------------

\begin{definition}[Degree of a map]\label{def:degree-integration}
\index{degree!via integration}
Let $M,N$ be compact, connected, oriented smooth $n$-manifolds without
boundary, and let $f\colon M\to N$ be smooth.  Choose any
$\omega\in\Omega^n(N)$ with $\int_N\omega \ne 0$.  The \emph{degree}
of~$f$ is
\[
  \deg(f) = \frac{\int_M f^*\omega}{\int_N\omega}.
\]
\end{definition}

\begin{theorem}\label{thm:degree-well-defined}
The degree $\deg(f)$ is a well-defined integer, independent of the choice
of~$\omega$.  Moreover, $\deg(f)$ is a homotopy invariant: if
$f\simeq g$, then $\deg f = \deg g$.
\end{theorem}

\begin{proof}
The map $f^*\colon H^n_{\mathrm{dR}}(N)\to H^n_{\mathrm{dR}}(M)$ is
a linear map between one-dimensional vector spaces (since $M$~and~$N$
are compact, connected, and oriented, $H^n_{\mathrm{dR}}\cong\R$).
Hence $f^*[\omega] = \lambda[\mu]$ for some $\lambda\in\R$,
where $[\mu]$ generates $H^n_{\mathrm{dR}}(M)$.  The ratio
$\int_M f^*\omega / \int_N\omega$ equals this constant~$\lambda$,
independent of~$\omega$.  That $\lambda\in\Z$ follows from the local
analysis: at a regular value $q\in N$, $f^{-1}(q)$ is finite and
$\deg f = \sum_{p\in f^{-1}(q)}\operatorname{sign}\det\dd f_p$,
which is manifestly an integer.  Homotopy invariance follows from

ef{thm:homotopy-invariance}.
\end{proof}

\begin{example}\label{ex:degree-circle}
The map $f\colon S^1\to S^1$, $f(z) = z^k$, has $\deg f = k$.
Indeed, if $\omega = \dd\theta/(2\pi)$ is the normalised volume form,
then $f^*\omega = k\,\dd\theta/(2\pi)$ and
$\int_{S^1}f^*\omega = k$.
\end{example}

\begin{example}\label{ex:degree-antipodal}
The antipodal map $A\colon S^n\to S^n$, $A(x) = -x$, has
$\deg A = (-1)^{n+1}$.  In particular, $A$ is homotopic to the identity
if and only if $n$ is odd.
\end{example}

\begin{proposition}[Degree and regular values]\label{prop:degree-regular}
\index{degree!via regular values}
Let $f\colon M^n\to N^n$ be a smooth map between compact connected
oriented manifolds.  If $q\in N$ is a regular value of~$f$, then
$f^{-1}(q)$ is a finite set and
\[
  \deg f = \sum_{p\in f^{-1}(q)} \operatorname{sign}\det(\dd f_p),
\]
where $\operatorname{sign}\det(\dd f_p) = +1$ if $\dd f_p$ preserves
orientation and $-1$ if it reverses orientation.
\end{proposition}

\begin{proof}
Since $q$ is a regular value, $f^{-1}(q)$ is a $0$-dimensional
submanifold of the compact manifold~$M$, hence a finite set
$\set{p_1,\dots,p_r}$.  By the inverse function theorem, there exist
disjoint open neighbourhoods $U_i$ of $p_i$ mapped diffeomorphically
by~$f$ onto a common neighbourhood $V$ of~$q$.  Choose an $n$-form
$\eta$ supported in~$V$ with $\int_V\eta = 1$.  Then
\[
  \int_M f^*\eta = \sum_{i=1}^r \int_{U_i} f^*\eta
  = \sum_{i=1}^r \operatorname{sign}\det(\dd f_{p_i})
  \int_V\eta
  = \sum_{i=1}^r\operatorname{sign}\det(\dd f_{p_i}),
\]
since $f|_{U_i}$ is a diffeomorphism onto~$V$.  Dividing by
$\int_N\eta = 1$ gives the result.
\end{proof}

\begin{example}\label{ex:degree-hopf}
\index{Hopf map}
The Hopf map $h\colon S^3\to S^2$ is not a map between manifolds
of the same dimension, so the degree formula above does not directly
apply.  However, one can define the \emph{Hopf invariant} using
linking numbers and de~Rham cohomology.  We mention this as motivation
for the rich interplay between topology and differential forms.
\end{example}

%--------------------------------------------------------------------
\section{Poincar\'e duality}\label{sec:poincare-duality}
%--------------------------------------------------------------------

\begin{remark}[Poincar\'e duality -- informal statement]
  \label{rem:poincare-duality}
\index{Poincar\'e duality}
For a closed oriented $n$-manifold~$M$, the wedge product followed by
integration defines a nondegenerate bilinear pairing
\[
  H^k_{\mathrm{dR}}(M) \times H^{n-k}_{\mathrm{dR}}(M) \longrightarrow \R,
  \qquad ([\alpha],[\beta])\longmapsto \int_M \alpha\wedge\beta.
\]
This is the \emph{Poincar\'e duality} isomorphism
$H^k_{\mathrm{dR}}(M)\cong \bigl(H^{n-k}_{\mathrm{dR}}(M)\bigr)^*$.
When each $H^k_{\mathrm{dR}}(M)$ is finite-dimensional (as is the case
for compact manifolds), this gives
$\dim H^k_{\mathrm{dR}}(M) = \dim H^{n-k}_{\mathrm{dR}}(M)$.
A full proof requires compactly supported cohomology and the
Mayer--Vietoris sequence; we refer the reader to Bott--Tu
\emph{Differential Forms in Algebraic Topology}.
\end{remark}

\begin{example}\label{ex:poincare-duality-sphere}
For $S^n$ ($n\ge 1$),
$H^k_{\mathrm{dR}}(S^n)\cong\begin{cases}\R & k=0,n,\\ 0 &
\text{otherwise.}\end{cases}$
Poincar\'e duality is reflected in the symmetry $H^0\cong H^n\cong\R$.
\end{example}

%--------------------------------------------------------------------
\section{Exercises}\label{sec:stokes-exercises}
%--------------------------------------------------------------------

\begin{exercise}\label{exo:stokes-annulus}
Let $A = \set{(x,y)\in\R^2 \mid 1\le x^2+y^2\le 4}$ with the
standard orientation.  Compute $\int_A \dd\omega$ directly and via
Stokes' theorem, where
$\omega = \frac{x\,\dd y - y\,\dd x}{x^2+y^2}$.
\end{exercise}

\begin{exercise}\label{exo:cohomology-cylinder}
Compute $H^k_{\mathrm{dR}}(S^1\times\R)$ for all~$k$ using homotopy
invariance.
\end{exercise}

\begin{exercise}\label{exo:poincare-explicit}
Let $\omega = e^{xyz}(yz\,\dd x + xz\,\dd y + xy\,\dd z)$ on~$\R^3$.
Verify that $\dd\omega = 0$ and find an explicit primitive $\eta$
with $\dd\eta = \omega$ using the homotopy operator from the Poincar\'e
lemma.
\end{exercise}

\begin{exercise}\label{exo:degree-computation}
Let $f\colon S^2\to S^2$ be defined by $f(x,y,z) = (2xz, 2yz, 2z^2-1)$
(this is the map induced by $z\mapsto z^2$ under stereographic
projection).  Compute $\deg f$ by integrating the pullback of the
area form.
\end{exercise}

\begin{exercise}\label{exo:closed-not-exact-torus}
\index{torus!cohomology}
Let $T^2 = S^1\times S^1$ with angular coordinates $(\theta_1,\theta_2)$.
Show that $\omega_1 = \dd\theta_1$ and $\omega_2 = \dd\theta_2$ are
closed but not exact, and that $\set{[\omega_1],[\omega_2]}$ is a basis
for $H^1_{\mathrm{dR}}(T^2)$.
\end{exercise}

\begin{exercise}\label{exo:mayer-vietoris-hint}
Using the fact that $S^n = U_+\cup U_-$ where $U_\pm = S^n\setminus\set{\mp e_{n+1}}$
are each diffeomorphic to~$\R^n$, and the Mayer--Vietoris sequence
\[
  \cdots\to H^{k-1}_{\mathrm{dR}}(U_+\cap U_-)\to H^k_{\mathrm{dR}}(S^n)
  \to H^k_{\mathrm{dR}}(U_+)\oplus H^k_{\mathrm{dR}}(U_-)
  \to H^k_{\mathrm{dR}}(U_+\cap U_-)\to\cdots
\]
compute $H^k_{\mathrm{dR}}(S^n)$ by induction on~$n$.
\end{exercise}

\begin{exercise}\label{exo:hairy-ball}
\index{hairy ball theorem}
Use the degree of the antipodal map and the notion of degree to prove
that $S^{2n}$ admits no nowhere-vanishing smooth vector field (the
hairy ball theorem).  \emph{Hint:} A nowhere-vanishing vector field
gives a homotopy from $\Id$ to the antipodal map.
\end{exercise}

\begin{exercise}\label{exo:flux-integral}
Let $\Sigma\subset\R^3$ be a closed oriented surface enclosing a
bounded domain~$\Omega$, and let $\mathbf{F}$ be a smooth vector field
on~$\R^3$ with $\operatorname{div}\mathbf{F} = 1$ everywhere.  Show
that $\oiint_\Sigma \mathbf{F}\cdot\mathbf{n}\,\dd S = \operatorname{Vol}(\Omega)$.
\end{exercise}

\begin{exercise}\label{exo:stokes-moebius}
Explain why Stokes' theorem (in its manifold form) does not directly
apply to the M\"obius band.  What additional structure would be
needed to define integration of $2$-forms on a M\"obius band?
\end{exercise}

\begin{exercise}\label{exo:cohomology-RPn}
\index{projective space!cohomology}
Using the fact that $S^n\to\RP^n$ is a $2$-sheeted covering space,
show that
\[
  H^k_{\mathrm{dR}}(\RP^n)\cong\begin{cases}
  \R & \text{if } k = 0, \\
  \R & \text{if } k = n \text{ and } n \text{ is odd}, \\
  0  & \text{otherwise}.
  \end{cases}
\]
\emph{Hint:} The pullback of forms via the covering map identifies
$\Omega^k(\RP^n)$ with the $\Z/2$-invariant forms on~$S^n$.
\end{exercise}

\begin{exercise}\label{exo:stokes-sphere-exact}
Let $\omega = x\,\dd y\wedge\dd z + y\,\dd z\wedge\dd x + z\,\dd x\wedge\dd y$
be the standard area form on~$S^2$ (viewed as the restriction of a
$2$-form on~$\R^3$).  Compute $\int_{S^2}\omega$ using the parametrisation
by spherical coordinates and verify that $[\omega]\ne 0$ in
$H^2_{\mathrm{dR}}(S^2)$.
\end{exercise}

\begin{exercise}\label{exo:degree-fund-thm-algebra}
\index{fundamental theorem of algebra}
Use the notion of degree to give a topological proof that every
nonconstant polynomial $p\colon\C\to\C$ has a root.
\emph{Hint:} Consider $f_t(z) = p(tz)/\abs{p(tz)}$ for $z\in S^1$
and use the homotopy invariance of degree.
\end{exercise}


%%%%%%%%%%%%%%%%%%%%%%%%%%%%%%%%%%%%%%%%%%%%%%%%%%%%%%%%%%%%%%%%%%%%
%%  CHAPTER 8 — SARD'S THEOREM
%%%%%%%%%%%%%%%%%%%%%%%%%%%%%%%%%%%%%%%%%%%%%%%%%%%%%%%%%%%%%%%%%%%%
\chapter{Sard's Theorem}\label{ch:sard}
\minitoc\vspace{1em}

Sard's theorem\index{Sard's theorem} is one of the most fundamental
results in differential topology.  It asserts that the set of critical
values of any smooth map has measure zero, so that ``almost every''
value is a regular value.  This seemingly analytic statement has
profound topological consequences: the existence of regular values
leads to the non-existence of smooth retractions $D^n\to S^{n-1}$
and hence to the Brouwer fixed-point theorem, as well as to
transversality results that underpin much of modern differential
topology.

%--------------------------------------------------------------------
\section{Critical and regular values}\label{sec:critical-regular}
%--------------------------------------------------------------------

\begin{definition}[Critical and regular points]\label{def:critical-pt}
\index{critical point}\index{regular point}
Let $f\colon M^m\to N^n$ be a smooth map between smooth manifolds.
A point $p\in M$ is a \emph{critical point}\index{critical point} of~$f$
if the differential $\dd f_p\colon T_pM\to T_{f(p)}N$ is not surjective
(i.e., $\operatorname{rank}\dd f_p < n$).  If $\dd f_p$ is surjective,
$p$ is a \emph{regular point}.
\end{definition}

\begin{definition}[Critical and regular values]\label{def:critical-val}
\index{critical value}\index{regular value}
A point $q\in N$ is a \emph{critical value} of~$f$ if
$q = f(p)$ for some critical point $p$ of~$f$.  A point $q\in N$
is a \emph{regular value} if it is not a critical value, i.e., for
every $p\in f^{-1}(q)$ the differential $\dd f_p$ is surjective.
\end{definition}

\begin{remark}\label{rem:regular-value-not-in-image}
A point $q\in N\setminus f(M)$ is automatically a regular value
(vacuously: there are no points in $f^{-1}(q)$ to fail surjectivity).
\end{remark}

\begin{remark}\label{rem:submersion}
\index{submersion}
If every point of $M$ is a regular point, then $f$ is called a
\emph{submersion}.  By the implicit function theorem, a submersion
$f\colon M^m\to N^n$ (with $m\ge n$) has the property that $f^{-1}(q)$
is a smooth submanifold of~$M$ of dimension $m-n$ for every $q\in f(M)$.
\end{remark}

\begin{example}\label{ex:height-function}
Consider the height function $f\colon S^2\to\R$, $f(x,y,z) = z$.
The critical points are the north and south poles, $p_\pm = (0,0,\pm 1)$,
where $\dd f_{p_\pm} = 0$.  The critical values are $\set{-1,+1}$;
every $c\in(-1,1)$ is a regular value and $f^{-1}(c)$ is a circle.
\end{example}

\begin{example}\label{ex:critical-projection}
Consider the smooth map $f\colon\R^3\to\R^2$ defined by
$f(x,y,z) = (x^2-y^2, 2xy)$.  The Jacobian is
\[
  \dd f_{(x,y,z)} = \begin{pmatrix} 2x & -2y & 0 \\ 2y & 2x & 0 \end{pmatrix}.
\]
This has rank $<2$ precisely when $x = y = 0$, i.e., the critical set
is the $z$-axis $\set{(0,0,z)\mid z\in\R}$.  The set of critical
values is $f(\set{(0,0,z)}) = \set{(0,0)}$, a single point---which
indeed has measure zero in~$\R^2$.
\end{example}

\begin{example}\label{ex:whitney-umbrella}
\index{Whitney umbrella}
The \emph{Whitney umbrella} is the image of $f\colon\R^2\to\R^3$,
$f(u,v) = (uv, u, v^2)$.  The Jacobian is
$\dd f_{(u,v)} = \bigl(\begin{smallmatrix} v & u \\ 1 & 0 \\
0 & 2v\end{smallmatrix}\bigr)$,
which has rank~$2$ unless $v = 0$ (giving $\operatorname{rank}= 1$).
Since $m = 2 < n = 3$, every point is critical (the differential is
never surjective), so $f(\R^2)$ has measure zero in~$\R^3$ by

ef{lem:measure-zero-props}(iii).
\end{example}

%--------------------------------------------------------------------
\section{Measure zero in $\R^n$}\label{sec:measure-zero}
%--------------------------------------------------------------------

\begin{definition}[Measure zero]\label{def:measure-zero}
\index{measure zero}
A subset $A\subset\R^n$ has \emph{(Lebesgue) measure zero} if for
every $\varepsilon > 0$ there exists a countable collection of
open cubes $\set{Q_j}_{j=1}^\infty$ in~$\R^n$ such that
$A\subset\bigcup_j Q_j$ and $\sum_j \operatorname{vol}(Q_j) < \varepsilon$.
\end{definition}

\begin{remark}\label{rem:measure-zero-equiv}
One may equivalently use open balls, rectangles, or cubes.  The notion
is independent of the particular covering shapes used.
\end{remark}

\begin{lemma}\label{lem:measure-zero-props}
\index{measure zero!properties}
The following hold:
\begin{enumerate}[label=(\roman*)]
  \item A countable union of sets of measure zero has measure zero.
  \item If $A\subset\R^n$ has measure zero and $m > n$, then
    $A\times\set{0}\subset\R^m$ has measure zero (as a subset of~$\R^m$).
  \item If $A\subset\R^n$ has measure zero and $g\colon U\to\R^n$
    is a smooth map defined on an open set $U\supset A$, then
    $g(A)$ has measure zero.
  \item $\R^k\times\set{0}\subset\R^n$ has measure zero for $k < n$.
\end{enumerate}
\end{lemma}

\begin{proof}
(i) is a standard $\varepsilon/2^j$ argument.  (ii) follows because
a cube of side $\ell$ in $\R^n$ times $\set{0}$ can be covered by
a thin slab of volume~$\varepsilon$.  For (iii), on any compact subset
$K\subset U$, $g$ is Lipschitz with constant~$C$, and $g$ maps a cube
of side~$\ell$ into a cube of side $C\sqrt{n}\,\ell$, hence of volume
$(C\sqrt{n}\,\ell)^n = C^n n^{n/2}\ell^n$.  The covering of $A\cap K$
by cubes of total volume $<\varepsilon$ gives a covering of
$g(A\cap K)$ by cubes of total volume $<C^n n^{n/2}\varepsilon$.
Since $A = \bigcup_{j=1}^\infty (A\cap K_j)$ for an exhaustion by
compact sets, (i)~completes the proof.
(iv) follows from~(i) and~(ii).
\end{proof}

\begin{definition}[Measure zero in manifolds]\label{def:measure-zero-manifold}
\index{measure zero!in manifolds}
A subset $A$ of a smooth $n$-manifold~$N$ has \emph{measure zero} if
for every chart $(U,\varphi)$ of~$N$, $\varphi(A\cap U)$ has measure
zero in~$\R^n$.  By 
ef{lem:measure-zero-props}(iii), this is
well-defined (independent of the atlas).
\end{definition}

\begin{lemma}[Fubini-type lemma for measure zero]\label{lem:fubini-measure-zero}
\index{Fubini!measure zero}
Let $A\subset\R^{n+1}$ be a subset such that for every $t\in\R$, the
slice $A_t = \set{x\in\R^n \mid (x,t)\in A}$ has measure zero in~$\R^n$.
If furthermore the set $T = \set{t\in\R \mid A_t\ne\emptyset}$ has measure
zero in~$\R$, then $A$ has measure zero in~$\R^{n+1}$.
\end{lemma}

\begin{proof}
For each $t\in T$, cover $A_t$ by cubes in $\R^n$ of total
$n$-dimensional volume less than $\varepsilon_t$.  Taking products
with thin intervals around $t$ and using a careful covering of~$T$
by intervals of total length $<\varepsilon$, one obtains a covering
of $A$ by $(n+1)$-cubes of total volume $<C\varepsilon$.  Details
are left as an exercise.
\end{proof}

More generally, we have the following useful strengthening.

\begin{lemma}[Fubini-type lemma, general form]\label{lem:fubini-general}
Let $A\subset\R^{n+k}$.  If for almost every $y\in\R^k$, the slice
$A_y = \set{x\in\R^n \mid (x,y)\in A}$ has measure zero in~$\R^n$,
then $A$ has measure zero in~$\R^{n+k}$.
\end{lemma}

%--------------------------------------------------------------------
\section{Sard's theorem: statement and proof}\label{sec:sard-proof}
%--------------------------------------------------------------------

\begin{theorem}[Sard's theorem, 1942]\label{thm:sard}
\index{Sard's theorem}
Let $f\colon M^m\to N^n$ be a smooth map between smooth manifolds.
Then the set of critical values of~$f$ has measure zero in~$N$.
\end{theorem}

The proof is by induction on~$m$.  We first treat the Euclidean case.

\begin{theorem}[Sard's theorem, Euclidean version]\label{thm:sard-euclidean}
Let $U\subset\R^m$ be open and $f\colon U\to\R^n$ smooth.  Let
$C = \set{x\in U\mid \operatorname{rank}\dd f_x < n}$ be the set of
critical points.  Then $f(C)$ has measure zero in~$\R^n$.
\end{theorem}

\begin{proof}
We proceed by induction on~$m$.

\medskip\noindent\textbf{Base case $m = 0$:} $f$ is a constant map
and $f(U)$ is a single point (or empty).  If $n\ge 1$, a single point
has measure zero.  If $n = 0$, there are no critical points.

\medskip\noindent\textbf{Inductive step:} Assume the theorem holds for
all smooth maps from open subsets of $\R^{m-1}$.  We stratify the
critical set~$C$ as follows.  For $j\ge 1$, define
\[
  C_j = \set{x\in U \;\Big|\; \frac{\partial^\alpha f}{\partial x^\alpha}(x) = 0
  \text{ for all multi-indices } \abs{\alpha}\le j}.
\]
(Here we adopt the convention that $C_0 = C$.)
We have a descending filtration
\[
  C = C_0 \supset C_1 \supset C_2 \supset \cdots \supset C_{m-n+1} \supset\cdots
\]
where each $C_j$ is closed in~$U$.

The proof is completed in three claims.

\medskip\noindent\textbf{Claim 1:} $f(C\setminus C_1)$ has measure zero.

Let $x_0\in C\setminus C_1$.  Then some first partial derivative of a
component of~$f$ is nonzero at~$x_0$ but $\operatorname{rank}\dd f_{x_0}<n$.
Without loss of generality (reordering coordinates), assume
$\frac{\partial f_1}{\partial x_1}(x_0)\ne 0$.  Define
$g\colon U\to\R^m$ by
\[
  g(x) = (f_1(x), x_2, \dots, x_m).
\]
By the inverse function theorem, $g$ is a diffeomorphism on a
neighbourhood $V$ of~$x_0$.  Then $h = f\circ g^{-1}\colon g(V)\to\R^n$
has the form $h(y_1,y_2,\dots,y_m) = (y_1, h_2(y),\dots,h_n(y))$.
A point $y\in g(V)$ is a critical point of~$h$ if and only if the map
$(y_2,\dots,y_m)\mapsto (h_2(y_1,y_2,\dots,y_m),\dots,h_n(y_1,y_2,\dots,y_m))$
has non-surjective differential (since the first component is just
$y_1$).

For fixed $y_1 = c$, define $h^c\colon\R^{m-1}\to\R^{n-1}$ by
$h^c(y_2,\dots,y_m) = (h_2(c,y_2,\dots,y_m),\dots,h_n(c,y_2,\dots,y_m))$.
Then $y$ is a critical point of~$h$ if and only if $(y_2,\dots,y_m)$ is
a critical point of~$h^c$.  By the induction hypothesis (applied in
dimension $m-1$), $h^c(\text{crit.\ pts.\ of } h^c)$ has measure zero
in~$\R^{n-1}$ for each~$c$.  Thus, by 
ef{lem:fubini-measure-zero},
$h(\text{crit.\ pts.\ of }h\text{ in }g(V))$ has measure zero in~$\R^n$.
Since $f(C\cap V) = h(g(C\cap V))$, and $C\setminus C_1$ is covered
by countably many such neighbourhoods~$V$, we conclude
$f(C\setminus C_1)$ has measure zero.

\medskip\noindent\textbf{Claim 2:} $f(C_j\setminus C_{j+1})$ has
measure zero for each $j\ge 1$.

Let $x_0\in C_j\setminus C_{j+1}$.  Then all partial derivatives of
all components of~$f$ up to order~$j$ vanish at $x_0$, but some
$(j{+}1)$-st partial derivative is nonzero.  Without loss of generality,
assume $\frac{\partial^{j+1} f_1}{\partial x_1\partial x^\beta}(x_0)\ne 0$
for some multi-index $\beta$ with $\abs{\beta} = j$.  Set
$w(x) = \frac{\partial^j f_1}{\partial x^\beta}(x)$; then $w(x_0)=0$
and $\frac{\partial w}{\partial x_1}(x_0)\ne 0$.

Define $g(x) = (w(x), x_2,\dots,x_m)$.  By the inverse function
theorem, $g$ is a local diffeomorphism near~$x_0$.  On the set
$g(C_j\cap V)$, we have $y_1 = 0$ (since $w$ vanishes on~$C_j$ locally).
Thus $g(C_j\cap V)\subset \set{0}\times\R^{m-1}$, and we can apply
the induction hypothesis to the restricted map
$f\circ g^{-1}\big|_{\set{0}\times\R^{m-1}}$
(an $(m-1)$-dimensional problem).  By induction, $f(C_j\cap V)$
has measure zero, and the result follows by a countable cover.

\medskip\noindent\textbf{Claim 3:} $f(C_k)$ has measure zero for
$k\ge m-n+1$ (in fact for $k$ sufficiently large).

We use a Taylor expansion argument.  Let $x_0\in C_k$ and work
in a small cube $Q$ of side $\delta$ centered at~$x_0$.  By the
vanishing of all derivatives up to order~$k$, Taylor's theorem gives
\[
  \abs{f(x) - f(x_0)} \le C\abs{x-x_0}^{k+1}
  \le C'\delta^{k+1}
\]
for $x\in Q$.  Hence $f(C_k\cap Q)$ is contained in a cube of
side~$O(\delta^{k+1})$ in~$\R^n$, with $n$-dimensional volume
$O(\delta^{n(k+1)})$.

Now subdivide $Q$ into $N^m$ subcubes of side $\delta/N$.  The image
of $C_k$ intersected with each subcube lies in a cube of volume
$O((\delta/N)^{n(k+1)})$.  The total volume is bounded by
\[
  N^m \cdot O\!\left(\left(\frac{\delta}{N}\right)^{n(k+1)}\right)
  = O(N^{m - n(k+1)}).
\]
If $k+1 > m/n$, i.e., $k \ge \lceil m/n\rceil$, then $m - n(k+1) < 0$
and this tends to $0$ as $N\to\infty$.  In particular, the condition
$k \ge m - n + 1$ suffices (since $m-n+1 \ge m/n$ when $n\ge 1$).
Hence $f(C_k)$ has measure zero.

\medskip
Combining Claims~1, 2, and~3:
\[
  f(C) = f(C\setminus C_1) \cup f(C_1\setminus C_2) \cup \cdots
  \cup f(C_{k-1}\setminus C_k) \cup f(C_k)
\]
is a finite union of sets of measure zero, hence has measure zero.
\end{proof}

\begin{proof}[Proof of 
ef{thm:sard} (manifold version)]
Cover $M$ by countably many charts $(U_\alpha,\varphi_\alpha)$ and $N$ by
charts $(V_\beta,\psi_\beta)$.  The critical values of~$f$ are
\[
  f(C_f) = \bigcup_{\alpha,\beta} \psi_\beta^{-1}\!\left(
  (\psi_\beta\circ f\circ\varphi_\alpha^{-1})
  (C_{\psi_\beta\circ f\circ\varphi_\alpha^{-1}})\right).
\]
By 
ef{thm:sard-euclidean}, each set
$(\psi_\beta\circ f\circ\varphi_\alpha^{-1})(C_{\cdots})$ has measure zero.
Since $\psi_\beta^{-1}$ is a diffeomorphism, the preimage has measure zero
in~$N$ (
ef{lem:measure-zero-props}(iii)).  A countable union of measure
zero sets has measure zero.
\end{proof}

%--------------------------------------------------------------------
\section{Regular values are dense}\label{sec:regular-dense}
%--------------------------------------------------------------------

\begin{corollary}[Regular values are dense]\label{cor:regular-dense}
\index{regular value!density}
Let $f\colon M\to N$ be smooth.  Then the set of regular values
of~$f$ is dense in~$N$.  More precisely, it is a residual set
(countable intersection of open dense sets) and in particular dense
by the Baire category theorem.
\end{corollary}

\begin{proof}
A set of measure zero in~$\R^n$ has empty interior (since open sets
have positive measure).  Hence its complement is dense.  The set of
regular values, being the complement of a measure-zero set, is dense
in~$N$.

For the residual statement, write $M = \bigcup_j K_j$ with $K_j$
compact.  The set of critical values contributed by~$K_j$ is compact
(as the image of a compact set) and of measure zero, hence closed and
nowhere dense.  Its complement is open and dense.  The set of regular
values contains $\bigcap_j(N\setminus f(C_f\cap K_j))$, a countable
intersection of open dense sets.
\end{proof}

\begin{remark}\label{rem:sard-smoothness}
\index{Sard's theorem!regularity}
Sard's theorem requires $f$ to be sufficiently smooth.  For $f\in C^r$
with $r \ge \max(m-n+1, 1)$, the conclusion still holds (this is the
sharp regularity threshold due to Whitney).  In the $C^\infty$ category
we work in, this is never an issue.  The necessity of the smoothness
condition is illustrated by the following: there exist $C^1$ maps
$f\colon[0,1]\to[0,1]$ whose set of critical values has positive
measure (see Whitney's 1935 example).
\end{remark}

\begin{remark}[Sard's theorem and the Brown--de~Rham theorem]
  \label{rem:sard-brown}
\index{Sard's theorem!consequences}
Sard's theorem has far-reaching consequences beyond those discussed here.
Combined with the implicit function theorem, it implies:
\begin{enumerate}[label=(\roman*)]
  \item Every smooth manifold admits a Morse function (functions whose
    critical points are all nondegenerate).
  \item The set of embeddings $M\hookrightarrow\R^N$ is dense for
    $N\ge 2\dim M + 1$ (Whitney embedding theorem).
  \item Every smooth map can be approximated by maps transverse to any
    given submanifold.
\end{enumerate}
These results form the foundation of differential topology.
\end{remark}

%--------------------------------------------------------------------
\section{Existence of regular values}\label{sec:existence-regular}
%--------------------------------------------------------------------

\begin{corollary}\label{cor:existence-regular}
Let $f\colon M^m\to N^n$ be smooth with $m\ge n$.  Then there exists
a regular value $q\in N$.  If $q\in f(M)$, then $f^{-1}(q)$ is a
smooth submanifold of~$M$ of dimension $m - n$.
\end{corollary}

\begin{proof}
By 
ef{cor:regular-dense} regular values are dense, hence exist.
The preimage theorem (implicit function theorem on manifolds) gives
the submanifold conclusion.
\end{proof}

\begin{example}\label{ex:regular-value-sphere}
Consider $f\colon\R^{n+1}\to\R$, $f(x) = \abs{x}^2$.  Then
$\dd f_x = 2x^T$, which is surjective for $x\ne 0$.  Hence every
$c > 0$ is a regular value and $f^{-1}(c) = S^n(\sqrt{c})$ is a
smooth $n$-manifold.
\end{example}

\begin{example}\label{ex:regular-value-matrix}
\index{orthogonal group}
Consider $f\colon\GL(n,\R)\to\operatorname{Sym}(n,\R)$ defined by
$f(A) = A^TA$.  One shows that the identity matrix $I$ is a regular
value of~$f$.  Hence $f^{-1}(I) = O(n)$ is a smooth submanifold of
$\GL(n,\R)$ of dimension $n^2 - \frac{n(n+1)}{2} = \frac{n(n-1)}{2}$.
\end{example}

%--------------------------------------------------------------------
\section{No retraction theorem and Brouwer fixed point}\label{sec:brouwer}
%--------------------------------------------------------------------

\begin{lemma}[Classification of compact $1$-manifolds]\label{lem:compact-1-mfd}
\index{classification!1-manifolds}
Every compact connected smooth $1$-manifold with boundary is
diffeomorphic to either $S^1$ (if $\partial M = \emptyset$) or $[0,1]$
(if $\partial M\ne\emptyset$).  In particular, every compact
$1$-manifold with boundary has an even number of boundary points.
\end{lemma}

\begin{proof}[Proof sketch]
Let $M$ be a compact connected $1$-manifold with boundary.  Choose a
nonvanishing vector field~$X$ on $M$ (which exists since $M$ is
$1$-dimensional).  The flow of~$X$ starting at any interior point
traces out an injective curve; by compactness this curve must
eventually return to its starting point (giving $S^1$) or reach a
boundary point (giving $[0,1]$).  Since each connected component
diffeomorphic to $[0,1]$ contributes exactly $2$ boundary points and
each component diffeomorphic to $S^1$ contributes $0$, the total
number of boundary points is even.
\end{proof}

\begin{definition}[Retraction]\label{def:retraction}
\index{retraction}
A \emph{smooth retraction} of $D^n$ onto $S^{n-1}$ is a smooth map
$r\colon D^n\to S^{n-1}$ such that $r(x) = x$ for all $x\in S^{n-1}$.
\end{definition}

\begin{theorem}[No retraction theorem]\label{thm:no-retraction}
\index{no retraction theorem}
There is no smooth retraction of $D^n$ onto $S^{n-1}$ for any $n\ge 1$.
\end{theorem}

\begin{proof}
Suppose for contradiction that $r\colon D^n\to S^{n-1}$ is a smooth
retraction.  By Sard's theorem (
ef{thm:sard}), $r$ has a regular
value $q\in S^{n-1}$.

Consider $r^{-1}(q)$.  Since $q$ is a regular value, $r^{-1}(q)$ is
a smooth $1$-dimensional submanifold of~$D^n$ (since
$\dim D^n - \dim S^{n-1} = n - (n-1) = 1$).  Moreover, $r^{-1}(q)$
is compact (as a closed subset of the compact set~$D^n$), so it
consists of finitely many circles and arcs.

Since $r|_{S^{n-1}} = \Id_{S^{n-1}}$, the point $q\in S^{n-1}$ has
exactly one preimage on the boundary, namely $q$ itself.  The arcs have
their endpoints on $\partial D^n = S^{n-1}$, and the total number of
boundary endpoints must be even (each arc contributes $0$ or $2$
boundary points, except for arcs with one interior and one boundary
endpoint).

More precisely, $r^{-1}(q)$ is a compact $1$-manifold with boundary
$r^{-1}(q)\cap S^{n-1}$.  By the classification of compact
$1$-manifolds with boundary, each component is either a circle or a
closed interval.  The boundary points of $r^{-1}(q)$ (as a manifold
with boundary) are exactly the points of $r^{-1}(q)\cap S^{n-1}$,
which is just $\set{q}$ (since $r|_{S^{n-1}} = \Id$).  But a compact
$1$-manifold with boundary has an \emph{even} number of boundary
points (each interval component contributes exactly~$2$).
This contradicts the fact that $r^{-1}(q)\cap S^{n-1} = \set{q}$
has exactly one point.
\end{proof}

\begin{theorem}[Brouwer fixed-point theorem, smooth version]
  \label{thm:brouwer-smooth}
\index{Brouwer fixed-point theorem}
Every smooth map $f\colon D^n\to D^n$ has a fixed point.
\end{theorem}

\begin{proof}
Suppose $f$ has no fixed point, i.e., $f(x)\ne x$ for all $x\in D^n$.
Define $r\colon D^n\to S^{n-1}$ by letting $r(x)$ be the point where the
ray from $f(x)$ through~$x$ intersects~$S^{n-1}$.  Explicitly,
$r(x) = x + t(x)(x - f(x))$ where $t(x)\ge 0$ is determined by
$\abs{r(x)} = 1$.

One checks that $t$ depends smoothly on~$x$ (by the implicit function
theorem applied to $\abs{x+t(x-f(x))}^2 = 1$), so $r$ is smooth.
For $x\in S^{n-1}$: since $\abs{x} = 1$ and $\abs{f(x)}\le 1$ with
$f(x)\ne x$, the ray from $f(x)$ through $x$ exits the disk at~$x$
itself, so $r(x) = x$.

Thus $r$ is a smooth retraction $D^n\to S^{n-1}$, contradicting

ef{thm:no-retraction}.
\end{proof}

\begin{corollary}[Brouwer fixed-point theorem, continuous version]
  \label{cor:brouwer-continuous}
Every continuous map $f\colon D^n\to D^n$ has a fixed point.
\end{corollary}

\begin{proof}
By the Weierstrass approximation theorem, we can approximate any
continuous $f\colon D^n\to D^n$ uniformly by smooth maps
$f_k\colon D^n\to\R^n$ with $f_k\to f$ uniformly.  For $k$
sufficiently large, $g_k(x) = f_k(x)/\max(1,\abs{f_k(x)})$
maps $D^n$ to~$D^n$ and is smooth (or at worst Lipschitz).  One can
refine this argument via a smooth approximation that stays in~$D^n$,
and then apply 
ef{thm:brouwer-smooth} to obtain fixed
points $x_k$ of~$g_k$.  By compactness, $x_k\to x$ (along a
subsequence), and continuity gives $f(x) = x$.
\end{proof}

\begin{remark}\label{rem:brouwer-alternative}
An alternative approach uses Stokes' theorem directly.  If
$r\colon D^n\to S^{n-1}$ were a smooth retraction, then the volume
form $\omega$ of $S^{n-1}$ would satisfy
$\int_{S^{n-1}}\omega = \int_{S^{n-1}} r^*\omega = \int_{D^n}\dd(r^*\omega)
= \int_{D^n} r^*(\dd\omega) = 0$ (since $\dd\omega = 0$ as a form on
$S^{n-1}$, which is a top-degree form).  But $\int_{S^{n-1}}\omega\ne 0$,
giving the contradiction.  This avoids the classification of
$1$-manifolds but uses the integration theory from 
ef{ch:stokes}.
\end{remark}

%--------------------------------------------------------------------
\section{Transversality}\label{sec:transversality}
%--------------------------------------------------------------------

\begin{definition}[Transversality]\label{def:transversality}
\index{transversality}
Let $f\colon M\to N$ be smooth and $W\subset N$ a submanifold.  We say
$f$ is \emph{transverse to~$W$}, written $f\pitchfork W$, if for every
$p\in f^{-1}(W)$,
\[
  \dd f_p(T_pM) + T_{f(p)}W = T_{f(p)}N.
\]
Two submanifolds $X,Y\subset N$ are \emph{transverse}, $X\pitchfork Y$,
if the inclusion $\iota\colon X\hookrightarrow N$ is transverse to~$Y$,
i.e., $T_pX + T_pY = T_pN$ for all $p\in X\cap Y$.
\end{definition}

\begin{remark}\label{rem:transverse-disjoint}
If $f(M)\cap W = \emptyset$, then $f\pitchfork W$ vacuously.
\end{remark}

\begin{theorem}[Preimage under transverse maps]\label{thm:transverse-preimage}
\index{transversality!preimage theorem}
If $f\colon M^m\to N^n$ is smooth and $f\pitchfork W^w\subset N$,
then $f^{-1}(W)$ is a smooth submanifold of~$M$ of dimension
$m - n + w = m - \operatorname{codim}(W)$.
\end{theorem}

\begin{proof}
Let $p\in f^{-1}(W)$ and $q = f(p)\in W$.  Choose a submanifold chart
$(V,\psi)$ for~$W$ near~$q$ such that $\psi(W\cap V) = \psi(V)\cap(\R^w\times\set{0})$.
Let $\pi\colon\R^n\to\R^{n-w}$ be the projection onto the last $n-w$
coordinates.  Then $g = \pi\circ\psi\circ f$ satisfies $g^{-1}(0) = f^{-1}(W)$
near~$p$, and transversality ensures $\dd g_p$ is surjective.  By the
implicit function theorem, $g^{-1}(0)$ is a submanifold of dimension
$m - (n-w)$.
\end{proof}

\begin{example}\label{ex:transverse-curves}
Two curves in $\R^3$ are generically disjoint (since transversality
requires $T_pC_1 + T_pC_2 = T_p\R^3$, which fails when
$\dim T_pC_1 + \dim T_pC_2 = 2 < 3$, so transversality forces
$C_1\cap C_2 = \emptyset$).  Two surfaces in $\R^3$ generically
intersect in a curve.
\end{example}

\begin{example}\label{ex:transverse-diagonal}
Let $f\colon M\to N$ be smooth and consider the graph
$\Gamma_f = \set{(p,f(p))\mid p\in M}\subset M\times N$.
The map $f$ is transverse to $W\subset N$ if and only if
$\Gamma_f\pitchfork (M\times W)$ in $M\times N$.
\end{example}

\begin{theorem}[Transversality theorem, parametric version]
  \label{thm:transversality-parametric}
\index{transversality!theorem}
Let $F\colon M\times S\to N$ be smooth, and let $W\subset N$ be a
submanifold.  If $F\pitchfork W$, then for almost every $s\in S$
(i.e., all $s$ outside a set of measure zero), the map
$f_s = F(\cdot,s)\colon M\to N$ satisfies $f_s\pitchfork W$.
\end{theorem}

\begin{proof}
By 
ef{thm:transverse-preimage}, $Z = F^{-1}(W)$ is a submanifold
of $M\times S$.  Let $\pi\colon Z\to S$ be the projection.  By Sard's
theorem, almost every $s\in S$ is a regular value of~$\pi$.  One checks
that $s$ is a regular value of $\pi$ if and only if $f_s\pitchfork W$.
\end{proof}

\begin{corollary}[Density of transverse maps]\label{cor:transverse-dense}
\index{transversality!density}
Let $W\subset N$ be a closed submanifold.  Then the set of smooth maps
$f\colon M\to N$ that are transverse to~$W$ is dense (in the strong
$C^\infty$ topology) in $C^\infty(M,N)$.
\end{corollary}

\begin{proof}[Proof sketch]
By the tubular neighbourhood theorem, embed a neighbourhood of $W$ in~$N$
as the total space of the normal bundle $\nu W$.  Given $f\colon M\to N$,
define $F\colon M\times B_\varepsilon^k\to N$ (where $k = \operatorname{codim}W$
and $B_\varepsilon^k$ is a small ball) by translating $f$ in the normal
directions.  One shows $F\pitchfork W$; then

ef{thm:transversality-parametric} gives $f_s\pitchfork W$ for
almost every~$s$, and choosing $s$ small makes $f_s$ $C^\infty$-close
to~$f$.
\end{proof}

\begin{remark}[Transversality and stability]\label{rem:transversality-stability}
\index{transversality!stability}
If $W\subset N$ is a closed submanifold and $f\colon M\to N$ satisfies
$f\pitchfork W$, then any map sufficiently $C^1$-close to~$f$ also
satisfies transversality to~$W$.  In other words, the set of maps
transverse to a closed submanifold is not only dense but also open
in $C^\infty(M,N)$.  Combined with density, transversality is
\emph{generic} (it holds on an open dense set).
\end{remark}

\begin{example}[Intersection numbers via transversality]\label{ex:intersection-number}
\index{intersection number}
Let $M^n$ be a compact oriented manifold and $X^p, Y^q\subset M$ be
compact oriented submanifolds with $p + q = n$ and $X\pitchfork Y$.
Then $X\cap Y$ is a finite set of points, and one defines the
\emph{intersection number}
\[
  X\cdot Y = \sum_{p\in X\cap Y} \varepsilon(p),
\]
where $\varepsilon(p) = +1$ if the orientation of $T_pX\oplus T_pY$
agrees with that of $T_pM$, and $\varepsilon(p) = -1$ otherwise.
This integer depends only on the homology classes of $X$ and~$Y$.
\end{example}

%--------------------------------------------------------------------
\section{Exercises}\label{sec:sard-exercises}
%--------------------------------------------------------------------

\begin{exercise}\label{exo:critical-det}
Let $f\colon M_n(\R)\to\R$ be the determinant map $f(A) = \det A$.
Find the critical points of~$f$ and describe $f^{-1}(1) = \SL(n,\R)$
as a smooth submanifold.
\end{exercise}

\begin{exercise}\label{exo:measure-zero-image}
Let $f\colon\R^m\to\R^n$ be smooth with $m < n$.  Prove that $f(\R^m)$
has measure zero in~$\R^n$.  \emph{Hint:} Every point of $\R^m$ is a
critical point since $\dd f_x$ cannot be surjective.
\end{exercise}

\begin{exercise}\label{exo:sard-polynomial}
Let $p\colon\R^n\to\R$ be a polynomial.  Show that the set of critical
values of~$p$ is finite.  \emph{Hint:} The critical set is an algebraic
variety; combine with Sard's theorem and the structure of
semi-algebraic sets.
\end{exercise}

\begin{exercise}\label{exo:retraction-annulus}
Show that there \emph{does} exist a smooth retraction from the annulus
$A = \set{x\in\R^2\mid 1\le\abs{x}\le 2}$ onto the inner circle
$S^1 = \set{x\mid\abs{x}=1}$.  Why does this not contradict

ef{thm:no-retraction}?
\end{exercise}

\begin{exercise}\label{exo:brouwer-hex}
Use the Brouwer fixed-point theorem to prove that the game of Hex
has a winning strategy for the first player.  \emph{Hint:} Show that
every Hex game has a winner (there are no draws), which is equivalent
to a topological connectivity statement related to Brouwer.
\end{exercise}

\begin{exercise}\label{exo:transverse-intersection}
Let $X = \set{(x,y,z)\in\R^3 \mid x^2+y^2 = 1}$ (a cylinder) and
$Y = \set{(x,y,z)\in\R^3\mid x+z = 1}$ (a plane).  Show that
$X\pitchfork Y$ and describe the intersection $X\cap Y$ as a
submanifold.
\end{exercise}

\begin{exercise}\label{exo:normal-bundle-transverse}
\index{normal bundle}
Let $X,Y\subset N$ be transverse submanifolds.  Show that the normal
bundle of $X\cap Y$ in~$N$ satisfies
$\nu_{X\cap Y}^N \cong \nu_{X\cap Y}^X \oplus \nu_{X\cap Y}^Y$.
\end{exercise}

\begin{exercise}\label{exo:morse-regular}
\index{Morse function}
A smooth function $f\colon M\to\R$ on a compact manifold is called a
\emph{Morse function}\index{Morse function} if all its critical points are
nondegenerate (the Hessian matrix in local coordinates is nonsingular).
Use Sard's theorem to prove that Morse functions are dense in
$C^\infty(M,\R)$.  \emph{Hint:} For an embedding
$\iota\colon M\hookrightarrow\R^N$, consider the family of functions
$f_a(x) = \abs{x-a}^2$ parametrised by $a\in\R^N$.
\end{exercise}

\begin{exercise}\label{exo:whitney-embedding}
\index{Whitney embedding theorem}
Let $M^m$ be a compact smooth manifold.  Using Sard's theorem and
transversality, show that $M$ can be immersed in~$\R^{2m}$ and
embedded in~$\R^{2m+1}$.
\emph{Hint:} Start with an embedding $M\hookrightarrow\R^N$
for some large~$N$ (e.g., via Whitney) and project to lower-dimensional
subspaces, using Sard to find projections that preserve injectivity
and immersion properties.
\end{exercise}

\begin{exercise}\label{exo:compact-one-manifolds}
\index{classification!1-manifolds}
Prove the classification of compact connected $1$-manifolds with
boundary: every such manifold is diffeomorphic to either $S^1$ or
$[0,1]$.  Use this to give a self-contained proof that a compact
$1$-manifold with boundary has an even number of boundary points.
\end{exercise}

\begin{exercise}\label{exo:sard-space-filling}
Show that there is no smooth surjection $f\colon\R\to\R^2$.
\emph{Hint:} Since $1 < 2$, every point of $\R$ is a critical point.
Contrast this with the existence of continuous (Peano) space-filling
curves.
\end{exercise}

\begin{exercise}\label{exo:transverse-grassmannian}
\index{Grassmannian}
Let $\operatorname{Gr}_k(\R^n)$ denote the Grassmannian of $k$-planes
in~$\R^n$.  Consider the smooth map $f\colon\R^n\to\R^{n-k}$ given by
projection onto a $(n-k)$-dimensional subspace $V\in\operatorname{Gr}_{n-k}(\R^n)$.
Show that if $M^m\subset\R^n$ is a smooth submanifold, then for almost
every $V$, the restriction $f|_M\colon M\to V\cong\R^{n-k}$ is
transverse to $\set{0}\subset\R^{n-k}$.
\end{exercise}

\begin{exercise}\label{exo:linking-number}
\index{linking number}
Let $\gamma_1\colon S^1\to\R^3$ and $\gamma_2\colon S^1\to\R^3$ be
two disjoint smooth embeddings.  The \emph{linking number} is defined as
\[
  \operatorname{lk}(\gamma_1,\gamma_2)
  = \frac{1}{4\pi}\oint_{\gamma_1}\oint_{\gamma_2}
  \frac{(\mathbf{r}_1-\mathbf{r}_2)\cdot(\dd\mathbf{r}_1\times\dd\mathbf{r}_2)}
  {\abs{\mathbf{r}_1-\mathbf{r}_2}^3}.
\]
Show that this equals the degree of the Gauss map
$G\colon S^1\times S^1\to S^2$, $G(s,t) = \frac{\gamma_1(s)-\gamma_2(t)}
{\abs{\gamma_1(s)-\gamma_2(t)}}$, and that it is an integer.
\end{exercise}

% === END OF CHUNK ch7-8 ===
% === START OF CHUNK ch9-11+end ===

% %%%%%%%%%%%%%%%%%%%%%%%%%%%%%%%%%%%%%%%%%%%%%%%%%%%%%%%%%%%%
%  CHAPTER 9 — MORSE THEORY: INTRODUCTION
% %%%%%%%%%%%%%%%%%%%%%%%%%%%%%%%%%%%%%%%%%%%%%%%%%%%%%%%%%%%%

\chapter{Morse Theory --- Introduction}\label{ch:morse}
\minitoc\vspace{1cm}

Morse theory establishes a profound connection between the analysis
of smooth real-valued functions on a manifold and the topology of the
manifold itself. The key insight, due to Marston Morse in the 1920s
and subsequently refined by Smale, Thom, Milnor, and others, is that
the critical points of a ``generic'' smooth function encode
enough information to reconstruct the homotopy type of the manifold.

The philosophy is simple yet powerful: as we ``sweep'' through the
level sets $f^{-1}(t)$ of a smooth function $f\colon M \to \R$, the
topology of the sublevel set $f^{-1}((-\infty,t])$ changes only when $t$
passes through a critical value, and the nature of the change is
controlled by the \emph{index} of the critical point. This idea leads
to a cell decomposition of $M$, to fundamental inequalities relating
critical points and Betti numbers, and ultimately to the Morse complex,
a precursor to Floer homology and modern symplectic topology.

This chapter develops the foundations: Morse functions, the Morse lemma,
handle attachments, and the weak and strong Morse inequalities.

% ============================================================
\section{Morse Functions and Critical Points}\label{sec:morse-functions}
% ============================================================

\begin{definition}[Critical point, critical value]\label{def:critical-point-morse}
\index{critical point}\index{critical value}
Let $M$ be a smooth manifold and $f\colon M \to \R$ a smooth function.
A point $p \in M$ is a \textbf{critical point} of $f$ if
$\dd f_p = 0$, i.e., in any local coordinates $(x_1,\dots,x_n)$
centered at $p$,
\[
  \frac{\partial f}{\partial x_i}(p) = 0, \qquad i = 1,\dots,n.
\]
The value $f(p) \in \R$ is then called a \textbf{critical value}.
A point of $\R$ that is not a critical value is a \textbf{regular value}.
\end{definition}

\begin{definition}[Non-degenerate critical point, Hessian]\label{def:nondegenerate}
\index{non-degenerate critical point}\index{Hessian}
A critical point $p$ of $f\colon M \to \R$ is \textbf{non-degenerate}
if the \textbf{Hessian matrix}
\[
  H_f(p) = \Bigl(\frac{\partial^2 f}{\partial x_i \partial x_j}(p)\Bigr)_{1 \le i,j \le n}
\]
is non-singular (i.e., $\det H_f(p) \neq 0$) in some (hence any)
local coordinate system centered at $p$.
\end{definition}

\begin{remark}\label{rem:hessian-well-defined}
Although the Hessian matrix $H_f(p)$ depends on the choice of coordinates,
its non-degeneracy does not: a change of coordinates
$y = \varphi(x)$ transforms $H_f(p)$ by congruence
$H \mapsto J^T H J$ where $J$ is the Jacobian of $\varphi$ at $p$,
and congruence preserves non-singularity.
\end{remark}

\begin{definition}[Index of a critical point]\label{def:morse-index}
\index{index!of a critical point}
The \textbf{index} of a non-degenerate critical point $p$ of
$f\colon M \to \R$ is the number of negative eigenvalues of $H_f(p)$
(counted with multiplicity). We denote it $\lambda(p)$ or $\operatorname{ind}(p)$.
\end{definition}

\begin{definition}[Morse function]\label{def:morse-function}
\index{Morse function}
A smooth function $f\colon M \to \R$ is a \textbf{Morse function}
if every critical point of $f$ is non-degenerate.
If $M$ is compact, a Morse function has finitely many critical points.
\end{definition}

\begin{example}[Height function on $S^n$]\label{ex:height-sphere}
\index{height function}
Consider $S^n = \set{x \in \R^{n+1} : \abs{x}^2 = 1}$ and
$f\colon S^n \to \R$, $f(x_0,\dots,x_n) = x_n$. The critical points
are the north pole $N = (0,\dots,0,1)$ and south pole $S = (0,\dots,0,-1)$.

To verify non-degeneracy, use stereographic coordinates at $S$:
$\varphi(x_0,\dots,x_n) = \frac{1}{1+x_n}(x_0,\dots,x_{n-1})$
with inverse giving $x_n = \frac{\abs{u}^2-1}{\abs{u}^2+1}$.
Then $f\circ\varphi^{-1}(u) = \frac{\abs{u}^2-1}{\abs{u}^2+1}$,
whose Hessian at $u=0$ is $\frac{4}{4}I_n = I_n > 0$,
confirming index $0$ at $S$. Similarly, $N$ has index $n$.
Thus $f$ is a Morse function with exactly two critical points.
\end{example}

\begin{example}[Height function on the torus $T^2$]\label{ex:height-torus}
\index{torus!Morse function}
Embed the torus $T^2$ in $\R^3$ in the standard way:
\[
  (u,v) \mapsto \bigl((2+\cos v)\cos u,\;(2+\cos v)\sin u,\;\sin v\bigr).
\]
The height function $f(x,y,z) = z$ restricted to $T^2$ gives
$f(u,v) = \sin v$. This has four critical points:
\begin{itemize}
  \item $v = -\pi/2$ (minimum, index $0$),
  \item two saddle points at $v = 0$ with suitable $u$ values (index $1$),
  \item $v = \pi/2$ (maximum, index $2$).
\end{itemize}

\begin{center}
\begin{tikzpicture}[scale=0.9]
  % Draw torus outline
  \draw[thick] (0,0) ellipse (3 and 1.2);
  \draw[thick] (-1.2,0) arc (180:360:1.2 and 0.5);
  \draw[thick,dashed] (-1.2,0) arc (180:0:1.2 and 0.5);
  % Critical points
  \fill[red] (0,1.2) circle (2.5pt) node[above right]{max, $\lambda=2$};
  \fill[red] (0,-1.2) circle (2.5pt) node[below right]{min, $\lambda=0$};
  \fill[blue] (-1.2,0) circle (2.5pt) node[left]{saddle, $\lambda=1$};
  \fill[blue] (1.2,0) circle (2.5pt) node[right]{saddle, $\lambda=1$};
  % Height arrow
  \draw[->,thick,gray] (4,-1.5) -- (4,1.8) node[above]{$z$};
  % Level sets
  \foreach \y in {-0.8,-0.3,0.3,0.8} {
    \draw[thin,green!50!black,dashed] (-3.5,\y) -- (3.5,\y);
  }
\end{tikzpicture}
\captionof{figure}{The torus $T^2$ with height function. Critical points are shown
with their indices. Dashed lines indicate level sets.}\label{fig:torus-height}
\end{center}
\end{example}

\begin{example}[Height on $\RP^2$]\label{ex:height-rp2}
\index{real projective plane!Morse function}
Realize $\RP^2$ as the quotient of $S^2$ by the antipodal map.
The function $f[x_0:x_1:x_2] = \frac{x_2^2}{x_0^2+x_1^2+x_2^2}$
is a well-defined Morse function on $\RP^2$ with three critical points
of indices $0$, $1$, and $2$. This is consistent with
$\chi(\RP^2) = 1 - 1 + 1 = 1$ and the CW structure
$\RP^2 = e^0 \cup e^1 \cup e^2$.
\end{example}

\begin{example}[A non-Morse function]\label{ex:non-morse}
Consider $f\colon \R^2 \to \R$, $f(x,y) = x^2 y^2$. The origin
is a critical point with Hessian
$H_f(0) = \begin{psmallmatrix} 0 & 0 \\ 0 & 0 \end{psmallmatrix}$,
which is degenerate. The origin is not an isolated critical point
(the entire coordinate axes are critical), confirming that $f$
is not a Morse function. By contrast, $g(x,y) = x^2 + y^2 + x^2y^2$
has a non-degenerate minimum at the origin.
\end{example}

% ============================================================
\section{The Morse Lemma}\label{sec:morse-lemma}
% ============================================================

The Morse lemma states that near a non-degenerate critical point,
a Morse function has a canonical quadratic form up to smooth
change of coordinates.

\begin{theorem}[Morse Lemma]\label{thm:morse-lemma}
\index{Morse lemma}
Let $f\colon M \to \R$ be smooth and let $p$ be a non-degenerate
critical point of $f$ with index $\lambda$. Then there exist local
coordinates $(y_1,\dots,y_n)$ centered at $p$ such that
\[
  f(y) = f(p) - y_1^2 - \cdots - y_\lambda^2 + y_{\lambda+1}^2 + \cdots + y_n^2.
\]
\end{theorem}

\begin{proof}
We may assume $p = 0 \in \R^n$ and $f(0) = 0$. By Taylor's theorem
with integral remainder, since $\dd f_0 = 0$, we can write
\[
  f(x) = \sum_{i,j=1}^n x_i x_j \, h_{ij}(x),
\]
where $h_{ij} \colon U \to \R$ are smooth functions with
$h_{ij}(0) = \frac{1}{2}\frac{\partial^2 f}{\partial x_i \partial x_j}(0)$.
By symmetrizing we may assume $h_{ij} = h_{ji}$.

We proceed by induction. Suppose that after a coordinate change
we have achieved
\[
  f(x) = \pm x_1^2 \pm \cdots \pm x_{k-1}^2 +
  \sum_{i,j \ge k} x_i x_j \, g_{ij}(x)
\]
with $g_{ij} = g_{ji}$ smooth.

\textbf{Case 1:} $g_{kk}(0) \neq 0$.
After possibly changing the sign, assume $g_{kk}(0) > 0$.
Since $g_{kk}$ is continuous, $g_{kk}(x) > 0$ near $0$.
Define
\[
  y_k = \sqrt{g_{kk}(x)}\Bigl(x_k + \sum_{j>k}\frac{g_{kj}(x)}{g_{kk}(x)}x_j\Bigr),
  \qquad y_i = x_i \text{ for } i \neq k.
\]
Then $\frac{\partial y_k}{\partial x_k}(0) = \sqrt{g_{kk}(0)} \neq 0$,
so by the inverse function theorem $(y_1,\dots,y_n)$ is a valid
coordinate system near $0$. A direct computation shows
\[
  f = \pm x_1^2 \pm \cdots \pm x_{k-1}^2 + y_k^2 +
  \sum_{i,j > k} x_i x_j\, g'_{ij}(x)
\]
for new smooth functions $g'_{ij}$, completing the induction step.

\textbf{Case 2:} $g_{kk}(0) = 0$ but $g_{k\ell}(0) \neq 0$ for some $\ell > k$.
The linear substitution $x_k \mapsto x_k + x_\ell$,
$x_\ell \mapsto x_k - x_\ell$ produces a nonzero diagonal coefficient,
reducing to Case~1.

Since $H_f(0)$ is non-degenerate, at each step we can always enter
Case~1 or Case~2. After $n$ steps we obtain the desired form.
The number of negative signs equals the index $\lambda$ by
Sylvester's law of inertia.
\end{proof}

\begin{corollary}\label{cor:morse-isolated}
Non-degenerate critical points are isolated.
\end{corollary}

\begin{proof}
In Morse coordinates $(y_1,\dots,y_n)$, the gradient of
$f(y) = f(p) - y_1^2 - \cdots - y_\lambda^2 + y_{\lambda+1}^2 + \cdots + y_n^2$
vanishes only at $y = 0$.
\end{proof}

\begin{corollary}\label{cor:morse-compact-finite}
If $M$ is compact, every Morse function $f\colon M \to \R$ has
finitely many critical points.
\end{corollary}

\begin{proof}
The set of critical points is a closed subset of the compact
manifold $M$, hence compact. But each critical point is isolated
by \autoref{cor:morse-isolated}. A compact discrete set is finite.
\end{proof}

\begin{example}[Morse lemma in dimension 1]\label{ex:morse-lemma-dim1}
Let $f\colon \R \to \R$ have a non-degenerate critical point at $0$
with $f(0)=0$ and $f''(0) \neq 0$. The Morse lemma gives a
coordinate $y$ with $f(y) = \pm y^2$. If $f''(0)>0$ (minimum),
then $f = y^2$ (index $0$); if $f''(0)<0$ (maximum), then $f = -y^2$ (index $1$).
The coordinate change $y = y(x)$ is determined by
$f(x) = x^2 h(x) = (x\sqrt{h(x)})^2$, so $y = x\sqrt{h(x)}$
where $h(0) = f''(0)/2 > 0$.
\end{example}

\begin{example}[Morse lemma in dimension 2]\label{ex:morse-lemma-dim2}
For $f\colon \R^2 \to \R$ with a saddle point at the origin (index $1$),
the Morse lemma gives coordinates $(y_1,y_2)$ such that
$f = -y_1^2 + y_2^2$. The level sets $f^{-1}(c)$ near the origin are:
\begin{itemize}
  \item for $c = 0$: a pair of lines $y_1 = \pm y_2$ (the ``cross''),
  \item for $c > 0$: hyperbolas opening along the $y_2$-axis,
  \item for $c < 0$: hyperbolas opening along the $y_1$-axis.
\end{itemize}
This local picture governs the topology change at a saddle critical value.
\end{example}

\begin{figure}[ht]
\centering
\begin{tikzpicture}[scale=0.75]
  % Level sets near a saddle point
  \begin{scope}[shift={(-5,0)}]
    \node at (0,-2.8) {$c < 0$};
    \draw[thick,blue] (-2,-0.5) .. controls (-0.5,-0.3) and (-0.3,-0.5) .. (-0.3,-2);
    \draw[thick,blue] (-2,0.5) .. controls (-0.5,0.3) and (-0.3,0.5) .. (-0.3,2);
    \draw[thick,blue] (2,-0.5) .. controls (0.5,-0.3) and (0.3,-0.5) .. (0.3,-2);
    \draw[thick,blue] (2,0.5) .. controls (0.5,0.3) and (0.3,0.5) .. (0.3,2);
    \draw[->] (-2.5,0) -- (2.5,0) node[right]{\small $y_1$};
    \draw[->] (0,-2.5) -- (0,2.5) node[above]{\small $y_2$};
  \end{scope}
  % Level set c=0
  \begin{scope}
    \node at (0,-2.8) {$c = 0$};
    \draw[thick,red] (-2,-2) -- (2,2);
    \draw[thick,red] (-2,2) -- (2,-2);
    \fill[red] (0,0) circle (2pt);
    \draw[->] (-2.5,0) -- (2.5,0) node[right]{\small $y_1$};
    \draw[->] (0,-2.5) -- (0,2.5) node[above]{\small $y_2$};
  \end{scope}
  % Level sets c>0
  \begin{scope}[shift={(5,0)}]
    \node at (0,-2.8) {$c > 0$};
    \draw[thick,blue] (-0.5,-2) .. controls (-0.3,-0.5) and (-0.5,-0.3) .. (-2,-0.3);
    \draw[thick,blue] (0.5,-2) .. controls (0.3,-0.5) and (0.5,-0.3) .. (2,-0.3);
    \draw[thick,blue] (-0.5,2) .. controls (-0.3,0.5) and (-0.5,0.3) .. (-2,0.3);
    \draw[thick,blue] (0.5,2) .. controls (0.3,0.5) and (0.5,0.3) .. (2,0.3);
    \draw[->] (-2.5,0) -- (2.5,0) node[right]{\small $y_1$};
    \draw[->] (0,-2.5) -- (0,2.5) node[above]{\small $y_2$};
  \end{scope}
\end{tikzpicture}
\caption{Level sets of $f = -y_1^2 + y_2^2$ near a saddle point (index $1$).
At $c=0$ the level set degenerates into two crossing lines.}\label{fig:saddle-level-sets}
\end{figure}

% ============================================================
\section{Density of Morse Functions}\label{sec:morse-dense}
% ============================================================

\begin{theorem}[Morse functions are generic]\label{thm:morse-generic}
\index{Morse function!density}
Let $M$ be a smooth manifold (without boundary). Then the set of Morse
functions is dense (and in fact residual) in $C^\infty(M,\R)$ with
the strong (Whitney) $C^2$ topology.

More precisely, if $f\colon M \to \R$ is any smooth function and
$\varepsilon > 0$, there exists a Morse function $g$ with
$\abs{f(x)-g(x)} < \varepsilon$ for all $x \in M$.
\end{theorem}

\begin{proof}[Proof sketch]
If $M \subset \R^N$ is embedded, consider for $a \in \R^N$ the function
$f_a(x) = f(x) + \langle a, x\rangle$. The map
$F\colon M \times \R^N \to T^*M$ given by $F(x,a) = \dd(f_a)_x$ is
a submersion. By the parametric transversality theorem
(\autoref{thm:parametric-transversality}),
for almost every $a \in \R^N$ (in the sense of Lebesgue measure),
$f_a$ is a Morse function. Choosing $\abs{a}$ small gives
the desired approximation.
\end{proof}

\begin{remark}\label{rem:morse-proper}
On a compact manifold $M$, every Morse function is automatically proper.
On non-compact manifolds, one typically requires Morse functions to be
proper and to have distinct critical values, which can also be arranged
by a small perturbation.
\end{remark}

% ============================================================
\section{Passing through Critical Levels}\label{sec:critical-levels}
% ============================================================

\begin{definition}[Sublevel set]\label{def:sublevel}
\index{sublevel set}
For $f\colon M \to \R$ and $a \in \R$, define the \textbf{sublevel set}
$M^a = f^{-1}((-\infty,a]) = \set{x \in M : f(x) \le a}$.
\end{definition}

\begin{theorem}[Regular interval theorem]\label{thm:regular-interval}
\index{regular interval theorem}
Let $f\colon M \to \R$ be a smooth proper function on a manifold
without boundary. If $f^{-1}([a,b])$ is compact and contains no
critical points of $f$, then $M^a$ is diffeomorphic to $M^b$.
Moreover, $M^a$ is a deformation retract of $M^b$.
\end{theorem}

\begin{proof}
Since $[a,b]$ contains no critical values, $\nabla f \neq 0$
on $f^{-1}([a,b])$ (using any Riemannian metric). Define the
vector field
\[
  X = \frac{\nabla f}{\abs{\nabla f}^2}
\]
on a neighborhood of $f^{-1}([a,b])$, multiplied by a suitable
bump function to have compact support. The flow of $X$ moves
the level set $f^{-1}(a)$ to $f^{-1}(b)$ in time $b-a$,
since $X(f) = 1$. The time-$(b-a)$ map of the flow restricts to
a diffeomorphism $M^a \xrightarrow{\;\sim\;} M^b$.
The retraction is given by flowing backwards along $X$.
\end{proof}

% ============================================================
\section{Handle Attachments}\label{sec:handles}
% ============================================================

\begin{definition}[$\lambda$-handle]\label{def:handle}
\index{handle attachment}
A \textbf{$\lambda$-handle} of dimension $n$ is a copy of
$D^\lambda \times D^{n-\lambda}$, where $D^k$ denotes the closed
unit disk in $\R^k$. We say that $M^b$ is obtained from $M^a$ by
\textbf{attaching a $\lambda$-handle} if there is an embedding
$\varphi\colon S^{\lambda-1}\times D^{n-\lambda} \hookrightarrow \partial M^a$
and a diffeomorphism
\[
  M^b \cong M^a \cup_\varphi (D^\lambda \times D^{n-\lambda}).
\]
\end{definition}

\begin{theorem}[Handle attachment at a critical point]\label{thm:handle-attachment}
\index{handle attachment!Morse theory}
Let $f\colon M \to \R$ be a Morse function and $p$ a critical point
of index $\lambda$ with $f(p) = c$. Suppose $f^{-1}([c-\varepsilon,c+\varepsilon])$
is compact and contains no critical point other than $p$.
Then for sufficiently small $\varepsilon > 0$,
$M^{c+\varepsilon}$ is diffeomorphic to $M^{c-\varepsilon}$ with a
$\lambda$-handle attached. In particular, $M^{c+\varepsilon}$ has the
homotopy type of $M^{c-\varepsilon}$ with a $\lambda$-cell $e^\lambda$ attached.
\end{theorem}

\begin{proof}[Proof sketch]
By the Morse lemma, in a neighborhood of $p$ we have coordinates
$(y_1,\dots,y_n)$ with
\[
  f = c - y_1^2 - \cdots - y_\lambda^2 + y_{\lambda+1}^2 + \cdots + y_n^2.
\]
Write $u = (y_1,\dots,y_\lambda)$ and $v = (y_{\lambda+1},\dots,y_n)$,
so $f = c - \abs{u}^2 + \abs{v}^2$. The region
$f^{-1}([c-\varepsilon,c+\varepsilon])$ in Morse coordinates becomes
$\set{(u,v) : -\varepsilon \le -\abs{u}^2+\abs{v}^2 \le \varepsilon}$.
One constructs a modified function $\tilde{f}$ that agrees with $f$
outside a small neighborhood, such that $\tilde{f}^{-1}((-\infty,c-\varepsilon])$
is diffeomorphic to $M^{c-\varepsilon} \cup (D^\lambda \times D^{n-\lambda})$.
See Milnor \cite{milnor-morse} for the complete construction.
\end{proof}

\begin{figure}[ht]
\centering
\begin{tikzpicture}[scale=1.1]
  % 0-handle (disk)
  \begin{scope}[shift={(-4,0)}]
    \fill[blue!15] (0,0) circle (1);
    \draw[thick,blue!70!black] (0,0) circle (1);
    \node at (0,-1.6) {$0$-handle};
    \fill (0,0) circle (1.5pt);
  \end{scope}
  % 1-handle
  \begin{scope}[shift={(0,0)}]
    \fill[blue!15] (-1,-0.4) rectangle (1,0.4);
    \draw[thick,blue!70!black] (-1,-0.4) rectangle (1,0.4);
    \draw[thick,red] (-1,-0.4) -- (-1,0.4);
    \draw[thick,red] (1,-0.4) -- (1,0.4);
    \node at (0,-1.6) {$1$-handle};
    \node[red,font=\small] at (-1.5,0) {$S^0\times D^1$};
  \end{scope}
  % 2-handle
  \begin{scope}[shift={(4,0)}]
    \fill[blue!15] (0,0) circle (1);
    \draw[thick,red] (0,0) circle (1);
    \draw[thick,blue!70!black] (0,0) circle (0.02);
    \node at (0,-1.6) {$2$-handle};
    \node[red,font=\small] at (1.6,0) {$S^1\times D^0$};
  \end{scope}
\end{tikzpicture}
\caption{Handles in dimension $n=2$: a $0$-handle (disk), a $1$-handle
(strip attached along two arcs), and a $2$-handle (disk glued along its boundary circle).}\label{fig:handles-2d}
\end{figure}

% ============================================================
\section{CW Structure from a Morse Function}\label{sec:cw-morse}
% ============================================================

\begin{theorem}[CW decomposition]\label{thm:cw-morse}
\index{CW structure!from Morse function}
Let $M$ be a compact smooth manifold and $f\colon M \to \R$ a Morse
function with critical points $p_1,\dots,p_k$ of indices
$\lambda_1,\dots,\lambda_k$. Then $M$ has the homotopy type of a
CW complex with one cell of dimension $\lambda_i$ for each critical point $p_i$.
\end{theorem}

\begin{proof}
Arrange the critical values $c_1 < c_2 < \cdots < c_k$ (by a small
perturbation we may assume all critical values are distinct).
Choose $a_0 < c_1 < a_1 < c_2 < \cdots < a_{k-1} < c_k < a_k$.
Then:
\begin{enumerate}[label=(\roman*)]
  \item $M^{a_0} = \emptyset$,
  \item by \autoref{thm:handle-attachment}, $M^{a_i}$ has the homotopy type of
    $M^{a_{i-1}}$ with a cell $e^{\lambda_i}$ attached,
  \item $M^{a_k} = M$.
\end{enumerate}
By induction, $M$ has the homotopy type of a CW complex with cells
$e^{\lambda_1},\dots,e^{\lambda_k}$.
\end{proof}

\begin{example}\label{ex:cw-sphere}
The sphere $S^n$ with the height function (\autoref{ex:height-sphere})
has two critical points of indices $0$ and $n$. Hence $S^n$ has the
homotopy type of a CW complex $e^0 \cup e^n$, recovering the standard
CW structure $S^n = D^n / \partial D^n$.
\end{example}

\begin{example}\label{ex:cw-torus}
The torus $T^2$ with the height function (\autoref{ex:height-torus})
has critical points of indices $0$, $1$, $1$, $2$.
Hence $T^2 \simeq e^0 \cup e^1 \cup e^1 \cup e^2$, consistent with
$\chi(T^2) = 1 - 2 + 1 = 0$.
\end{example}

\begin{example}[CW structure on $\CP^n$]\label{ex:cw-cpn}
\index{complex projective space!CW structure}
Consider the function $f\colon \CP^n \to \R$ defined by
\[
  f[z_0:\cdots:z_n] = \frac{\sum_{k=0}^n k\,\abs{z_k}^2}{\sum_{k=0}^n \abs{z_k}^2}.
\]
This is a Morse function with $n+1$ critical points $p_k = [0:\cdots:0:1:0:\cdots:0]$
(the $1$ in position $k$) of index $2k$. Hence $\CP^n$ has the homotopy
type of a CW complex with one cell in each even dimension $0,2,4,\dots,2n$:
\[
  \CP^n \simeq e^0 \cup e^2 \cup e^4 \cup \cdots \cup e^{2n}.
\]
This explains why $H_{2k}(\CP^n;\Z) \cong \Z$ and
$H_{\mathrm{odd}}(\CP^n;\Z) = 0$.
\end{example}

\begin{remark}[Handle decomposition vs.\ CW structure]\label{rem:handle-vs-cw}
The handle decomposition of a compact manifold from a Morse function is
stronger than just a CW structure: it gives a smooth structure on each
``thickened cell.'' In particular, the handles can be rearranged and
cancelled in pairs (if a $\lambda$-handle and a $(\lambda+1)$-handle
meet along a single $\lambda$-sphere), leading to simplified
decompositions. This is the starting point for the $h$-cobordism theorem
and Smale's proof of the generalized Poincar\'{e} conjecture in
dimensions $\ge 5$.
\end{remark}

% ============================================================
\section{Morse Inequalities}\label{sec:morse-inequalities}
% ============================================================

Let $f\colon M \to \R$ be a Morse function on a compact manifold $M$.
Denote by $c_\lambda$ the number of critical points of index $\lambda$,
and by $b_\lambda = \dim H_\lambda(M;\R)$ the $\lambda$-th Betti number.

\begin{theorem}[Weak Morse inequalities]\label{thm:weak-morse}
\index{Morse inequalities!weak}
For each $\lambda \ge 0$,
\[
  c_\lambda \ge b_\lambda.
\]
In particular, $\sum_\lambda c_\lambda \ge \sum_\lambda b_\lambda = \chi(M) + 2\sum_{\text{odd } \lambda} b_\lambda$.
\end{theorem}

\begin{proof}
By \autoref{thm:cw-morse}, $M$ has the homotopy type of a CW complex
with $c_\lambda$ cells of dimension $\lambda$. The cellular chain group
$C_\lambda$ is a free abelian group of rank $c_\lambda$, so
\[
  b_\lambda = \dim H_\lambda(M;\R)
  \le \dim(C_\lambda \otimes \R) = c_\lambda. \qedhere
\]
\end{proof}

\begin{theorem}[Strong Morse inequalities]\label{thm:strong-morse}
\index{Morse inequalities!strong}
For each $k \ge 0$,
\[
  \sum_{j=0}^{k} (-1)^{k-j}\, c_j \;\ge\; \sum_{j=0}^{k} (-1)^{k-j}\, b_j.
\]
Equality holds when $k = n = \dim M$, giving the Euler characteristic formula:
\[
  \chi(M) = \sum_{j=0}^{n} (-1)^j\, c_j = \sum_{j=0}^{n} (-1)^j\, b_j.
\]
\end{theorem}

\begin{proof}
Let $C_\bullet$ be the cellular chain complex from the CW structure,
with $\partial_j\colon C_j \to C_{j-1}$. Set $d_j = \operatorname{rank}(\partial_j)$.
Then $c_j = b_j + d_j + d_{j+1}$ (the rank-nullity relation in each degree).
For the alternating sum:
\[
  \sum_{j=0}^{k}(-1)^{k-j} c_j
  = \sum_{j=0}^{k}(-1)^{k-j}(b_j+d_j+d_{j+1})
  = \sum_{j=0}^{k}(-1)^{k-j}b_j + d_{k+1},
\]
where the telescoping of the $d_j$ terms leaves only $d_{k+1} \ge 0$.
This gives the inequality. When $k = n$, $d_{n+1} = 0$, yielding equality.
\end{proof}

\begin{example}[Morse inequalities for $T^2$]\label{ex:morse-ineq-torus}
For the torus $T^2$, the Betti numbers are $b_0=1$, $b_1=2$, $b_2=1$
and the height function yields $c_0=1$, $c_1=2$, $c_2=1$.
The weak inequalities $c_\lambda \ge b_\lambda$ are all equalities.
Such a Morse function is called \textbf{perfect}\index{perfect Morse function}.
\end{example}

\begin{corollary}\label{cor:morse-lower-bound}
Every Morse function on a compact manifold $M$ has at least
$\sum_\lambda b_\lambda(M)$ critical points.
\end{corollary}

% ============================================================
\section{The Morse Complex}\label{sec:morse-complex}
% ============================================================

\begin{remark}[Connection to the Morse complex]\label{rem:morse-complex}
\index{Morse complex}
Given a Morse function $f$ and a Riemannian metric $g$ on $M$,
one can refine the CW structure into the \textbf{Morse complex}
$(C_\bullet^{\mathrm{Morse}},\partial)$ where:
\begin{itemize}
  \item $C_k^{\mathrm{Morse}}$ is the free abelian group generated by
    critical points of index $k$,
  \item the boundary operator counts (with signs) gradient flow lines
    of $-\nabla_g f$ connecting critical points of consecutive index.
\end{itemize}
For a \textbf{Morse--Smale pair}\index{Morse--Smale pair} $(f,g)$
(where stable and unstable manifolds intersect transversally),
the homology of this complex is isomorphic to the singular homology of $M$.
This construction is the finite-dimensional precursor to Floer homology.
\end{remark}

\begin{figure}[ht]
\centering
\begin{tikzpicture}[scale=1.2]
  % Schematic of gradient flow on torus
  \node[circle,fill=red,inner sep=2pt,label=above:{max}] (max) at (0,3) {};
  \node[circle,fill=blue,inner sep=2pt,label=left:{saddle}] (s1) at (-1.5,1.5) {};
  \node[circle,fill=blue,inner sep=2pt,label=right:{saddle}] (s2) at (1.5,1.5) {};
  \node[circle,fill=green!60!black,inner sep=2pt,label=below:{min}] (min) at (0,0) {};
  % Flow lines
  \draw[->,thick,red!60!black] (max) to[out=200,in=80] (s1);
  \draw[->,thick,red!60!black] (max) to[out=340,in=100] (s2);
  \draw[->,thick,blue!60!black] (s1) to[out=280,in=160] (min);
  \draw[->,thick,blue!60!black] (s2) to[out=260,in=20] (min);
  \draw[->,thick,red!60!black] (max) to[out=210,in=90] (-2,1.5) to[out=270,in=170] (s1);
  \draw[->,thick,blue!60!black] (s1) to[out=310,in=150] (min);
  \draw[->,thick,red!60!black] (max) to[out=330,in=90] (2,1.5) to[out=270,in=10] (s2);
  \draw[->,thick,blue!60!black] (s2) to[out=230,in=30] (min);
  % Labels
  \node at (4,1.5) {
    $\begin{aligned}
      &\partial_2(\text{max}) = \text{flow to saddles}\\
      &\partial_1(\text{saddle}) = \text{flow to min}
    \end{aligned}$
  };
\end{tikzpicture}
\caption{Schematic gradient flow lines for the height function on $T^2$,
illustrating the Morse complex boundary operator.}\label{fig:morse-complex-flow}
\end{figure}

% ============================================================
\section{Exercises}\label{sec:exo-morse}
% ============================================================

\begin{exercise}\label{exo:morse-1}
Show that $f\colon \R^n \to \R$ given by $f(x) = \abs{x}^4$ has a
single critical point at the origin which is degenerate. Conclude
that $f$ is not a Morse function.
\end{exercise}

\begin{exercise}\label{exo:morse-2}
Let $A$ be a real symmetric $n\times n$ matrix and define
$f\colon S^{n-1}\to\R$ by $f(x) = x^T A x$. Show that $f$ is a
Morse function if and only if $A$ has distinct eigenvalues. Determine
the critical points and their indices.
\end{exercise}

\begin{exercise}\label{exo:morse-3}
Let $M = S^1 \times S^1$ and $f(e^{i\theta},e^{i\varphi}) = \cos\theta + \cos\varphi$.
Verify that $f$ is a Morse function, find all critical points, compute
their indices, and verify the Euler characteristic formula.
\end{exercise}

\begin{exercise}\label{exo:morse-4}
Prove that the Morse lemma implies non-degenerate critical points are
isolated, directly from the normal form.
\end{exercise}

\begin{exercise}\label{exo:morse-5}
Show that $\RP^n$ admits a Morse function with exactly $n+1$ critical
points, one of each index $0,1,\dots,n$.
\emph{Hint:} use $f[x_0:\cdots:x_n] = \frac{\sum_i a_i x_i^2}{\sum_i x_i^2}$
with distinct $a_i$.
\end{exercise}

\begin{exercise}\label{exo:morse-6}
Let $\Sigma_g$ be the closed orientable surface of genus $g$.
Determine the minimum number of critical points for any Morse function
on $\Sigma_g$. Is a perfect Morse function always achievable?
\end{exercise}

\begin{exercise}\label{exo:morse-7}
\index{Reeb theorem}
(Reeb's theorem) Let $M$ be a compact manifold admitting a Morse function
with exactly two critical points. Prove that $M$ is homeomorphic to $S^n$.
\emph{Hint:} the two critical points must have indices $0$ and $n$; use the
regular interval theorem to show $M \cong D^n \cup_{S^{n-1}} D^n$.
\end{exercise}

\begin{exercise}\label{exo:morse-8}
Verify the strong Morse inequalities for $\CP^n$ equipped with the Morse
function arising from the norm-square of a generic linear functional.
\emph{Hint:} $\CP^n$ has Betti numbers $b_{2k}=1$ for $0 \le k \le n$
and $b_{\text{odd}}=0$.
\end{exercise}

\begin{exercise}\label{exo:morse-9}
Let $f\colon M \to \R$ be Morse on a compact manifold. Show that if all
critical points have even index, then $M$ has the homotopy type of a CW
complex with no odd-dimensional cells. Deduce that $H_{\text{odd}}(M;\Z)$
is torsion-free.
\end{exercise}

\begin{exercise}\label{exo:morse-10}
Let $f\colon M \to \R$ be a Morse function on a compact manifold and
let $-f$ denote the function $p \mapsto -f(p)$. Show that $-f$ is
also a Morse function, and that $p$ has index $\lambda$ for $f$ if
and only if $p$ has index $n - \lambda$ for $-f$ (where $n = \dim M$).
Deduce the Poincar\'{e} duality relation $c_\lambda(f) = c_{n-\lambda}(-f)$
for the number of critical points.
\end{exercise}

\begin{exercise}\label{exo:morse-11}
(Gradient-like vector fields)\index{gradient-like vector field}
A vector field $X$ on $M$ is \textbf{gradient-like} for a Morse function
$f$ if $X(f) > 0$ outside the critical set and $X$ agrees with $-\nabla_g f$
in a Morse-coordinate neighborhood of each critical point.
Show that gradient-like vector fields exist for any Morse function
(use a partition of unity).
\end{exercise}


% %%%%%%%%%%%%%%%%%%%%%%%%%%%%%%%%%%%%%%%%%%%%%%%%%%%%%%%%%%%%
%  CHAPTER 10 — DEGREE OF A SMOOTH MAP AND APPLICATIONS
% %%%%%%%%%%%%%%%%%%%%%%%%%%%%%%%%%%%%%%%%%%%%%%%%%%%%%%%%%%%%

\chapter{Degree of a Smooth Map and Applications}\label{ch:degree}
\minitoc\vspace{1cm}

The notion of degree assigns an integer to a smooth map between
compact manifolds of the same dimension, measuring ``how many times''
the domain wraps around the target. This is one of the most fundamental
invariants in topology: it is computable, homotopy-invariant, and yet
powerful enough to prove major theorems.

Degree theory provides a unifying framework for many classical results,
including the Brouwer fixed point theorem, the hairy ball theorem,
the Borsuk--Ulam theorem, and the fundamental theorem of algebra.
We develop both the mod~$2$ and the oriented (integer-valued) degree,
prove their key properties, and showcase a wide range of applications.

% ============================================================
\section{Mod 2 Degree}\label{sec:mod2-degree}
% ============================================================

\begin{definition}[Mod 2 degree]\label{def:mod2-degree}
\index{degree!mod 2}
Let $f\colon M \to N$ be a smooth map between compact connected manifolds
of the same dimension $n$ (without boundary). For a regular value
$q \in N$, define the \textbf{mod 2 degree} of $f$ as
\[
  \deg_2(f) = \#f^{-1}(q) \mod 2 \;\in\; \Z/2\Z.
\]
\end{definition}

\begin{theorem}[Well-definedness of mod 2 degree]\label{thm:mod2-welldefined}
The mod 2 degree $\deg_2(f)$ is independent of the choice of regular
value $q$.
\end{theorem}

\begin{proof}
Let $q_0$ and $q_1$ be two regular values of $f$. By Sard's theorem
(\autoref{thm:sard}), the set of regular values is dense and open in $N$.
Since $N$ is connected, there is a smooth path $\gamma\colon [0,1]\to N$
from $q_0$ to $q_1$ consisting entirely of regular values (after a small
perturbation of $\gamma$ using Sard's theorem applied to the map
$f \times \Id_{[0,1]}$).

Consider the pullback $W = \set{(p,t)\in M\times[0,1] : f(p) = \gamma(t)}$.
Since $\gamma(t)$ is a regular value for each $t$, $W$ is a compact
$1$-manifold with boundary
\[
  \partial W = f^{-1}(q_0)\times\set{0} \;\sqcup\; f^{-1}(q_1)\times\set{1}.
\]
A compact $1$-manifold with boundary has an even number of boundary points,
so $\#f^{-1}(q_0) + \#f^{-1}(q_1) \equiv 0 \pmod{2}$, whence
$\#f^{-1}(q_0) \equiv \#f^{-1}(q_1) \pmod{2}$.
\end{proof}

\begin{remark}\label{rem:mod2-compact}
The compactness of $M$ ensures that $f^{-1}(q)$ is a finite set for
any regular value $q$ (it is a discrete subset of a compact space).
Without compactness, the cardinality could be infinite.
\end{remark}

\begin{proposition}[Homotopy invariance of mod 2 degree]\label{prop:mod2-homotopy}
If $f,g\colon M \to N$ are smoothly homotopic, then $\deg_2(f) = \deg_2(g)$.
\end{proposition}

\begin{proof}
Let $F\colon M \times [0,1] \to N$ be a smooth homotopy with $F_0 = f$,
$F_1 = g$. By Sard's theorem, choose a regular value $q$ of $F$ (which
is then a regular value of both $f$ and $g$). Then $F^{-1}(q)$ is a compact
$1$-manifold with boundary $f^{-1}(q)\sqcup g^{-1}(q)$, so
$\#f^{-1}(q) \equiv \#g^{-1}(q) \pmod{2}$.
\end{proof}

% ============================================================
\section{Oriented Degree}\label{sec:oriented-degree}
% ============================================================

\begin{definition}[Oriented degree]\label{def:oriented-degree}
\index{degree!oriented}
Let $f\colon M \to N$ be a smooth map between compact connected
\emph{oriented} manifolds of the same dimension $n$. For a regular value
$q \in N$, define the \textbf{degree} of $f$ as
\[
  \deg(f) = \sum_{p \in f^{-1}(q)} \operatorname{sgn}(\det \dd f_p),
\]
where $\operatorname{sgn}(\det \dd f_p) = +1$ if $\dd f_p\colon T_pM \to T_qN$
preserves orientation, and $-1$ if it reverses orientation.
\end{definition}

\begin{theorem}[Well-definedness and homotopy invariance]\label{thm:degree-welldefined}
\index{degree!homotopy invariance}
The integer $\deg(f)$ is independent of the regular value $q$ and is
a homotopy invariant: if $f \simeq g$ (smoothly homotopic), then
$\deg(f) = \deg(g)$.
\end{theorem}

\begin{proof}
\textbf{Independence of regular value.}
Let $q_0, q_1$ be regular values of $f$. As in the mod~2 case, connect
them by a path of regular values and consider the $1$-manifold
$W = \set{(p,t): f(p) = \gamma(t)}$. Now $W$ inherits an orientation
from the orientations of $M$ and $N$. The signed count of boundary
points in an oriented compact $1$-manifold is zero, giving
$\sum_{p\in f^{-1}(q_0)}\operatorname{sgn} - \sum_{p\in f^{-1}(q_1)}\operatorname{sgn} = 0$.

\textbf{Homotopy invariance.}
Given a smooth homotopy $F\colon M\times[0,1]\to N$, choose $q$ a
regular value of $F$. Then $F^{-1}(q)$ is a compact oriented
$1$-manifold with boundary. The boundary contributions with signs
give $\deg(g)-\deg(f) = 0$.
\end{proof}

\begin{remark}\label{rem:degree-properties}
Key properties of the degree:
\begin{enumerate}[label=(\alph*)]
  \item $\deg(\Id_M) = 1$.
  \item If $\deg(f) \neq 0$, then $f$ is surjective.
  \item Reversing the orientation of $M$ or $N$ changes the sign of $\deg(f)$.
  \item For non-orientable manifolds, only the mod~2 degree is defined.
\end{enumerate}
\end{remark}

% ============================================================
\section{Degree via Differential Forms}\label{sec:degree-forms}
% ============================================================

\begin{theorem}[Degree formula via integration]\label{thm:degree-forms}
\index{degree!via differential forms}
Let $f\colon M \to N$ be a smooth map between compact connected oriented
$n$-manifolds. For any $n$-form $\omega$ on $N$ with $\int_N \omega \neq 0$,
\[
  \deg(f) = \frac{\int_M f^*\omega}{\int_N \omega}.
\]
\end{theorem}

\begin{proof}[Proof sketch]
By de Rham theory, $H^n_{\mathrm{dR}}(N) \cong \R$ for a compact connected
oriented $n$-manifold. The pullback $f^*$ acts on $H^n$ by multiplication
by an integer (the degree). Since $\int_N$ provides an isomorphism
$H^n(N) \xrightarrow{\sim} \R$, we get
$\int_M f^*\omega = \deg(f)\cdot\int_N\omega$.
One checks that this integer agrees with the signed count of preimages.
\end{proof}

% ============================================================
\section{Computations}\label{sec:degree-computations}
% ============================================================

\begin{example}[Degree of $z^n$]\label{ex:degree-zn}
\index{degree!of $z^n$}
The map $f\colon S^1 \to S^1$, $f(z) = z^n$ (for $n \in \Z$) has
$\deg(f) = n$. Indeed, for any regular value $q \in S^1$, the preimage
$f^{-1}(q)$ consists of $\abs{n}$ points, each with sign
$\operatorname{sgn}(n)$.

Alternatively, letting $\omega = \dd\theta/(2\pi)$ be the standard
$1$-form on $S^1$, we have $f^*\omega = n\,\dd\theta/(2\pi)$,
and $\int_{S^1}f^*\omega = n$.
\end{example}

\begin{example}[Antipodal map]\label{ex:degree-antipodal}
\index{antipodal map!degree}
The antipodal map $\alpha\colon S^n \to S^n$, $\alpha(x) = -x$,
has $\deg(\alpha) = (-1)^{n+1}$. This is because $\alpha$ is the
composition of $(n+1)$ reflections, each of degree $-1$.
In particular, $\deg(\alpha) = -1$ when $n$ is even, and
$\deg(\alpha) = +1$ when $n$ is odd.
\end{example}

\begin{example}[Degree of a covering map]\label{ex:degree-covering}
\index{degree!covering map}
Let $\pi\colon \tilde{M}\to M$ be a smooth $k$-sheeted covering map
between compact connected oriented manifolds. Then $\deg(\pi)=k$
(if $\pi$ preserves orientation). Indeed, every point $q\in M$ is a
regular value of $\pi$ (since $\pi$ is a local diffeomorphism),
and $\#\pi^{-1}(q)=k$ with all signs positive.

For instance, the double cover $\pi\colon S^n \to \RP^n$ has
$\deg(\pi) = 2$ when $n$ is odd (so $\RP^n$ is orientable and
$\pi$ preserves orientation). When $n$ is even, $\RP^n$ is
non-orientable and only $\deg_2(\pi) = 0$ is defined.
\end{example}

\begin{example}[Degree and winding number]\label{ex:winding-number}
\index{winding number}
For a smooth map $f\colon S^1 \to S^1$, the degree $\deg(f)$
coincides with the classical \textbf{winding number}:
\[
  \deg(f) = \frac{1}{2\pi}\int_0^{2\pi} f'(\theta)\,\dd\theta
\]
where we identify $f$ with a function $\R\to\R$ satisfying
$f(\theta+2\pi) = f(\theta) + 2\pi\deg(f)$. This is the archetypal
example of the degree-via-forms formula.
\end{example}

\begin{proposition}\label{prop:degree-composition}
For smooth maps $f\colon M \to N$ and $g\colon N \to P$ between
compact connected oriented manifolds of the same dimension,
$\deg(g \circ f) = \deg(g)\cdot\deg(f)$.
\end{proposition}

\begin{proof}
Choose a regular value $q$ of $g \circ f$ that is also a regular value
of $g$ (using Sard). Then
\[
  \deg(g\circ f) = \sum_{p\in(g\circ f)^{-1}(q)}
  \operatorname{sgn}\det\dd(g\circ f)_p
  = \sum_{r\in g^{-1}(q)}\operatorname{sgn}\det\dd g_r
  \sum_{p\in f^{-1}(r)}\operatorname{sgn}\det\dd f_p
  = \deg(g)\cdot\deg(f). \qedhere
\]
\end{proof}

% ============================================================
\section{Applications of Degree Theory}\label{sec:degree-applications}
% ============================================================

\begin{theorem}[Brouwer fixed point theorem]\label{thm:brouwer-degree}
\index{Brouwer fixed point theorem}
Every continuous map $f\colon D^n \to D^n$ has a fixed point.
\end{theorem}

\begin{proof}
Suppose $f$ has no fixed point. Define a retraction $r\colon D^n \to S^{n-1}$
by sending $x$ to the point where the ray from $f(x)$ through $x$ hits
$S^{n-1}$. Then $r$ is continuous and $r|_{S^{n-1}} = \Id_{S^{n-1}}$.
Consider $i\colon S^{n-1}\hookrightarrow D^n$ the inclusion. Then
$r\circ i = \Id$, so $\deg(r\circ i) = 1$. But $r$ extends over $D^n$,
so by homotopy invariance (since $D^n$ is contractible)
$r$ is homotopic to a constant, hence $\deg(r\circ i) = 0$, a contradiction.
\end{proof}

\begin{theorem}[Hairy ball theorem]\label{thm:hairy-ball}
\index{hairy ball theorem}
There is no continuous nowhere-vanishing tangent vector field on $S^{2k}$
(for any $k \ge 1$).
\end{theorem}

\begin{proof}
Suppose $v\colon S^{2k}\to \R^{2k+1}$ is a continuous nowhere-vanishing
tangent vector field. Normalize it to $\hat{v}(x) = v(x)/\abs{v(x)}$,
so $\hat{v}(x)\perp x$ and $\abs{\hat{v}(x)}=1$. Define the homotopy
\[
  H_t(x) = (\cos\pi t)\,x + (\sin\pi t)\,\hat{v}(x), \qquad t\in[0,1].
\]
One checks $\abs{H_t(x)} = 1$ for all $t$ (since $x \perp \hat{v}(x)$ and
both are unit vectors). Thus $H_t\colon S^{2k}\to S^{2k}$, with
$H_0 = \Id$ and $H_1 = -\Id$ (the antipodal map). By homotopy invariance,
$\deg(\Id) = \deg(-\Id)$, i.e., $1 = (-1)^{2k+1} = -1$, a contradiction.
\end{proof}

\begin{theorem}[Borsuk--Ulam theorem]\label{thm:borsuk-ulam}
\index{Borsuk--Ulam theorem}
For every continuous map $f\colon S^n \to \R^n$, there exists a
point $x \in S^n$ with $f(x) = f(-x)$.
\end{theorem}

\begin{proof}
Suppose $f(x) \neq f(-x)$ for all $x$. Define
$g\colon S^n \to S^{n-1}$ by
\[
  g(x) = \frac{f(x)-f(-x)}{\abs{f(x)-f(-x)}}.
\]
Then $g(-x) = -g(x)$, i.e., $g$ is an odd (equivariant) map.

We show by induction on $n$ that an odd map $S^n \to S^{n-1}$ cannot exist.
For $n=1$: an odd map $g\colon S^1 \to S^0 = \set{\pm 1}$ would satisfy
$g(-x)=-g(x)$, but $S^1$ is connected and $S^0$ is not, so $g$ must be
constant, contradicting oddness.

For $n \ge 2$: restrict $g$ to the equatorial $S^{n-1}\subset S^n$.
This gives an odd map $g|_{S^{n-1}}\colon S^{n-1}\to S^{n-1}$.
An odd map $S^{n-1}\to S^{n-1}$ has odd degree (a standard result
using the homology of the quotient $\RP^{n-1}$), so
$\deg(g|_{S^{n-1}})$ is odd, hence nonzero.
But $g|_{S^{n-1}}$ extends to the upper hemisphere $D^n_+$, so
it is null-homotopic, giving $\deg(g|_{S^{n-1}})=0$, a contradiction.
\end{proof}

\begin{corollary}[Ham sandwich theorem]\label{cor:ham-sandwich}
\index{ham sandwich theorem}
Given $n$ bounded measurable sets in $\R^n$, there exists a hyperplane
that bisects each of them simultaneously.
\end{corollary}

% ============================================================
\section{Lefschetz Fixed Point Theory}\label{sec:lefschetz}
% ============================================================

\begin{definition}[Lefschetz number]\label{def:lefschetz-number}
\index{Lefschetz number}
Let $f\colon M \to M$ be a continuous map on a compact manifold.
The \textbf{Lefschetz number} of $f$ is
\[
  L(f) = \sum_{k=0}^{n} (-1)^k \operatorname{tr}\bigl(f_*\colon H_k(M;\Q)\to H_k(M;\Q)\bigr).
\]
\end{definition}

\begin{theorem}[Lefschetz fixed point theorem]\label{thm:lefschetz}
\index{Lefschetz fixed point theorem}
If $L(f) \neq 0$, then $f$ has a fixed point.
\end{theorem}

\begin{proof}[Proof sketch]
One shows that $L(f)$ equals the intersection number of the graph
$\Gamma_f = \set{(x,f(x))}$ with the diagonal $\Delta = \set{(x,x)}$
in $M \times M$. If $f$ has no fixed point, then $\Gamma_f \cap \Delta = \emptyset$,
so the intersection number is zero, hence $L(f)=0$.
For a transversal intersection, each fixed point $p$ contributes
$\pm 1$ to $L(f)$, and the sign is the \textbf{local Lefschetz index},
which equals $\operatorname{sgn}\det(\Id - \dd f_p)$.
\end{proof}

\begin{corollary}\label{cor:lefschetz-brouwer}
Since $L(\Id_M) = \chi(M)$, any map homotopic to the identity on a
manifold with $\chi(M)\neq 0$ has a fixed point. In particular, this
recovers the Brouwer theorem for contractible manifolds.
\end{corollary}

\begin{example}[Lefschetz number computations]\label{ex:lefschetz-computations}
\begin{enumerate}[label=(\alph*)]
  \item For $f = \Id_{S^{2k}}\colon S^{2k}\to S^{2k}$, we have
    $L(f) = 1 + (-1)^{2k}\cdot 1 = 2 = \chi(S^{2k}) \neq 0$,
    so $\Id$ has a fixed point (obviously).
    More interestingly, \emph{every} map $f\colon S^{2k}\to S^{2k}$
    with $\deg(f)\neq(-1)^{2k+1} = -1$ has $L(f) = 1+\deg(f)\neq 0$
    and hence has a fixed point.
  \item For the torus $T^2$, $\chi(T^2)=0$, so $L(\Id)=0$; the Lefschetz
    theorem gives no information. Indeed, translations on $T^2$ have
    no fixed points.
  \item For the antipodal map $\alpha$ on $S^n$,
    $L(\alpha)=1+(-1)^n(-1)^{n+1} = 1-1 = 0$ when $n$ is even (consistent
    with $\alpha$ having no fixed points on $S^{2k}$).
\end{enumerate}
\end{example}

% ============================================================
\section{Intersection Number as Degree}\label{sec:intersection-degree}
% ============================================================

\begin{definition}[Intersection number]\label{def:intersection-number-deg}
\index{intersection number}
Let $M$ be a compact oriented $n$-manifold, and let
$X^k, Y^{n-k}\subset M$ be compact oriented submanifolds
intersecting transversally. The \textbf{intersection number}
$I(X,Y) \in \Z$ is defined as
\[
  I(X,Y) = \sum_{p \in X \cap Y} \varepsilon(p),
\]
where $\varepsilon(p) = +1$ if the orientation of $T_pX \oplus T_pY$
agrees with that of $T_pM$, and $\varepsilon(p) = -1$ otherwise.
\end{definition}

\begin{proposition}\label{prop:intersection-homotopy}
The intersection number $I(X,Y)$ depends only on the homology classes
$[X] \in H_k(M;\Z)$ and $[Y]\in H_{n-k}(M;\Z)$, and agrees with the
homological intersection pairing
\[
  H_k(M;\Z)\otimes H_{n-k}(M;\Z) \to H_0(M;\Z)\cong\Z.
\]
\end{proposition}

\begin{proposition}\label{prop:intersection-degree}
Let $f\colon M^n \to N^n$ be a smooth map between compact oriented
manifolds. For a regular value $q \in N$,
\[
  \deg(f) = I(M, \set{q})
\]
where the intersection is computed in the graph: the intersection number
of the graph $\Gamma_f \subset M \times N$ with $M \times \set{q}$
equals $\deg(f)$.
\end{proposition}

\begin{example}[Self-intersection number]\label{ex:self-intersection}
\index{self-intersection number}
Consider $\CP^1 \subset \CP^2$ (a complex line). Its self-intersection
number $I(\CP^1,\CP^1) = +1$. This can be computed by perturbing
one copy slightly: two generic complex lines in $\CP^2$ intersect
in exactly one point, transversally, with positive sign.
More generally, two algebraic curves of degrees $d$ and $e$ in
$\CP^2$ have intersection number $de$ (B\'{e}zout's theorem).
\end{example}

% ============================================================
\section{Exercises}\label{sec:exo-degree}
% ============================================================

\begin{exercise}\label{exo:degree-1}
Compute $\deg_2(f)$ and $\deg(f)$ for $f\colon S^1\to S^1$,
$f(e^{i\theta}) = e^{3i\theta}$.
\end{exercise}

\begin{exercise}\label{exo:degree-2}
Let $f\colon S^n \to S^n$ be a smooth map with no fixed points.
Show that $\deg(f) = (-1)^{n+1}$. \emph{Hint:} $f$ is homotopic
to the antipodal map via
$H_t(x)=\frac{(1-t)f(x)-tx}{\abs{(1-t)f(x)-tx}}$.
\end{exercise}

\begin{exercise}\label{exo:degree-3}
Show that $\deg\colon [S^n,S^n] \to \Z$ is a group isomorphism
(where $[S^n,S^n]$ denotes homotopy classes of based maps with the
group structure given by the co-$H$-space structure of $S^n$).
Deduce that $\pi_n(S^n)\cong\Z$.
\end{exercise}

\begin{exercise}\label{exo:degree-4}
Let $p(z) = z^n + a_{n-1}z^{n-1}+\cdots+a_0$ be a complex polynomial.
Show that for $R$ sufficiently large, $p/\abs{p}\colon S^1_R \to S^1$
has degree $n$ (where $S^1_R$ is the circle of radius $R$). Deduce the
fundamental theorem of algebra.
\end{exercise}

\begin{exercise}\label{exo:degree-5}
Use degree theory to prove that every map $f\colon S^{2n}\to S^{2n}$
has either a fixed point or sends some point to its antipode.
\end{exercise}

\begin{exercise}\label{exo:degree-6}
Let $f\colon T^n \to T^n$ be a smooth map of the $n$-torus. Express
$\deg(f)$ in terms of the matrix $A \in M_n(\Z)$ induced by $f_*$
on $H_1(T^n;\Z)\cong\Z^n$. Verify for the map $(z_1,z_2)\mapsto(z_1^2 z_2, z_1 z_2^3)$ on $T^2$.
\end{exercise}

\begin{exercise}\label{exo:degree-7}
Show that the Hopf map $h\colon S^3\to S^2$, $h(z_1,z_2)=[z_1:z_2]\in\CP^1\cong S^2$,
has the property that $\#h^{-1}(q)$ is constant for all regular values $q$,
but the degree formula does not apply directly. Why?
\end{exercise}

\begin{exercise}\label{exo:degree-8}
Let $f\colon \Sigma_g \to S^2$ be a smooth map from a surface of genus
$g \ge 1$. Show that $\deg(f) = 0$ if and only if $f$ is null-homotopic.
\emph{Hint:} use the cohomological degree formula and $H^2(\Sigma_g;\Z)\cong\Z$.
\end{exercise}

\begin{exercise}\label{exo:degree-9}
(Gauss map and degree) Let $S \subset \R^3$ be a compact oriented surface
and $\nu\colon S \to S^2$ its Gauss map. Show that
$\deg(\nu) = \frac{1}{2}\chi(S)$.
\emph{Hint:} use the Gauss--Bonnet theorem.
\end{exercise}


% %%%%%%%%%%%%%%%%%%%%%%%%%%%%%%%%%%%%%%%%%%%%%%%%%%%%%%%%%%%%
%  CHAPTER 11 — WHITNEY EMBEDDING THEOREM
% %%%%%%%%%%%%%%%%%%%%%%%%%%%%%%%%%%%%%%%%%%%%%%%%%%%%%%%%%%%%

\chapter{Whitney Embedding Theorem}\label{ch:whitney}
\minitoc\vspace{1cm}

A fundamental question in differential topology asks: can every abstract
smooth manifold be realized as a submanifold of some Euclidean space?
The Whitney embedding theorem provides a definitive affirmative answer,
with sharp dimension bounds.

Historically, manifolds were first studied as subsets of Euclidean space
(Riemann, Poincar\'{e}). The intrinsic definition of a manifold via
charts and atlases, due to Weyl and Whitney, raised the question of
whether the abstract viewpoint is truly more general. Whitney's landmark
1936 paper showed that it is not: every abstract smooth manifold admits
a smooth embedding into some $\R^N$.

This chapter proves the ``easy'' Whitney
theorem (compact $n$-manifold embeds in $\R^{2n+1}$) in full detail,
states the ``hard'' Whitney theorem (embedding in $\R^{2n}$),
develops the Whitney approximation theorems, and discusses applications
including the classification of surfaces and exotic phenomena.

% ============================================================
\section{Review of Embeddings and Immersions}\label{sec:embed-review}
% ============================================================

\begin{definition}[Immersion]\label{def:immersion-review}
\index{immersion}
A smooth map $f\colon M^m \to N^n$ is an \textbf{immersion} if
$\dd f_p\colon T_pM \to T_{f(p)}N$ is injective for every $p \in M$.
Necessarily $m \le n$.
\end{definition}

\begin{definition}[Embedding]\label{def:embedding-review}
\index{embedding}
A smooth map $f\colon M \to N$ is a (smooth) \textbf{embedding} if:
\begin{enumerate}[label=(\roman*)]
  \item $f$ is an immersion,
  \item $f$ is a topological embedding, i.e., $f\colon M \to f(M)$
    is a homeomorphism (where $f(M)$ has the subspace topology).
\end{enumerate}
If $M$ is compact, condition (ii) is equivalent to $f$ being injective.
\end{definition}

\begin{definition}[Proper map]\label{def:proper-review}
\index{proper map}
A continuous map $f\colon M \to N$ is \textbf{proper} if the preimage
of every compact set is compact. A proper injective immersion is an embedding.
\end{definition}

\begin{remark}\label{rem:embed-submanifold}
If $f\colon M \hookrightarrow N$ is an embedding, then $f(M)$ is a
(regular/embedded) submanifold of $N$, and $f$ is a diffeomorphism
onto $f(M)$.
\end{remark}

\begin{example}[Immersions that are not embeddings]\label{ex:immersion-not-embedding}
\begin{enumerate}[label=(\alph*)]
  \item The figure-eight $\gamma\colon (-\pi,\pi)\to\R^2$,
    $\gamma(t) = (\sin 2t, \sin t)$, is an injective immersion but
    not an embedding (the image in $\R^2$ has a self-intersection
    at the origin in the closure).
  \item The map $\iota\colon \R\to T^2$ given by $t\mapsto(e^{it},e^{i\alpha t})$
    with $\alpha$ irrational is an injective immersion whose image is
    dense in $T^2$. This is an immersion that is far from being an embedding.
\end{enumerate}
\end{example}

% ============================================================
\section{The Weak Whitney Embedding Theorem}\label{sec:weak-whitney}
% ============================================================

\begin{theorem}[Weak Whitney embedding theorem]\label{thm:weak-whitney}
\index{Whitney embedding theorem!weak}
Every compact smooth $n$-manifold can be embedded in $\R^{2n+1}$.
\end{theorem}

We break the proof into several steps.

\begin{lemma}[Embedding in Euclidean space of some dimension]\label{lem:embed-some-RN}
Every compact smooth $n$-manifold $M$ can be embedded in $\R^N$
for some $N$.
\end{lemma}

\begin{proof}
Since $M$ is compact, choose a finite atlas $\set{(U_\alpha,\varphi_\alpha)}_{\alpha=1}^k$
and a subordinate partition of unity $\set{\rho_\alpha}$.
Define $F\colon M \to \R^{k(n+1)}$ by
\[
  F(p) = \bigl(\rho_1(p)\varphi_1(p),\;\rho_1(p),\;\dots,\;
  \rho_k(p)\varphi_k(p),\;\rho_k(p)\bigr),
\]
where $\rho_\alpha(p)\varphi_\alpha(p)$ is extended by zero outside $U_\alpha$.
We claim $F$ is an embedding:
\begin{itemize}
  \item \textbf{Injectivity:} If $F(p) = F(q)$, then $\rho_\alpha(p) = \rho_\alpha(q)$
    for all $\alpha$. Pick $\alpha$ with $\rho_\alpha(p) > 0$. Then
    $\rho_\alpha(p)\varphi_\alpha(p) = \rho_\alpha(q)\varphi_\alpha(q)$
    implies $\varphi_\alpha(p) = \varphi_\alpha(q)$, so $p = q$.
  \item \textbf{Immersion:} If $v \in T_pM\setminus\set{0}$, pick $\alpha$
    with $\rho_\alpha(p) > 0$. The component $\rho_\alpha\varphi_\alpha$ has
    differential $\rho_\alpha(p)\,\dd(\varphi_\alpha)_p + \varphi_\alpha(p)\,\dd(\rho_\alpha)_p$,
    and by choosing appropriate $\alpha$ this is nonzero on $v$.
  \item Since $M$ is compact, an injective immersion is an embedding.
    \qedhere
\end{itemize}
\end{proof}

\begin{lemma}[Dimension reduction by projection]\label{lem:dimension-reduction}
\index{Whitney!dimension reduction}
If $M^n \subset \R^N$ is a compact embedded submanifold with $N > 2n+1$,
then there exists a unit vector $v \in S^{N-1}$ such that the orthogonal
projection $\pi_v\colon \R^N \to v^\perp \cong \R^{N-1}$ restricts to
an embedding $\pi_v|_M\colon M \hookrightarrow \R^{N-1}$.
\end{lemma}

\begin{proof}
We must show that for a generic direction $v$, the projection $\pi_v|_M$
is both injective and an immersion.

\textbf{Step 1: Injectivity.}
Define $\alpha\colon M\times M\setminus\Delta \to S^{N-1}$ by
\[
  \alpha(p,q) = \frac{p-q}{\abs{p-q}}.
\]
The projection $\pi_v|_M$ fails to be injective if and only if $v$
lies in the image of $\alpha$. The domain has dimension $2n$,
and the codomain has dimension $N-1$. Since $N > 2n+1$, we have
$2n < N-1$, so by Sard's theorem the image of $\alpha$ has measure
zero in $S^{N-1}$.

\textbf{Step 2: Immersion.}
Define $\beta\colon TM\setminus\set{0\text{-section}} \to S^{N-1}$ by
$\beta(p,w) = w/\abs{w}$ (viewing $T_pM \subset \R^N$ via the embedding).
The projection $\pi_v|_M$ fails to be an immersion at $p$ if and only
if $v \in T_pM$, i.e., $v$ is in the image of $\beta$ restricted to
the unit sphere bundle of $TM$, which has dimension $2n-1$.
Since $2n-1 < N-1$, the image of $\beta$ also has measure zero
by Sard's theorem.

\textbf{Step 3: Conclusion.}
The set of ``bad'' directions $v$ (where $\pi_v|_M$ fails to be
an embedding) is contained in $\operatorname{im}(\alpha)\cup\operatorname{im}(\beta)$,
a set of measure zero in $S^{N-1}$. Hence for almost every $v$,
$\pi_v|_M$ is an embedding into $\R^{N-1}$.
\end{proof}

\begin{proof}[Proof of \autoref{thm:weak-whitney}]
By \autoref{lem:embed-some-RN}, embed $M^n \hookrightarrow \R^N$ for some $N$.
If $N > 2n+1$, apply \autoref{lem:dimension-reduction} repeatedly to reduce
the ambient dimension one step at a time until we reach $\R^{2n+1}$.
\end{proof}

\begin{corollary}[Immersion in $\R^{2n}$]\label{cor:weak-immersion}
Every compact smooth $n$-manifold can be immersed in $\R^{2n}$.
\end{corollary}

\begin{proof}
For immersion we only need Step~2 of \autoref{lem:dimension-reduction}:
the ``bad'' directions for the immersion condition form a set of
dimension $\le 2n-1$ in $S^{N-1}$. We can project to $\R^{N-1}$
preserving the immersion property as long as $2n-1 < N-1$, i.e.,
$N > 2n$. Starting from some $N$ and projecting repeatedly, we reach
$\R^{2n}$.
\end{proof}

\begin{remark}\label{rem:immersion-sharpness}
The immersion dimension $2n$ from \autoref{cor:weak-immersion} is not sharp.
The sharp result is the \textbf{Hirsch--Smale immersion theorem}\index{Hirsch--Smale theorem},
combined with obstruction theory, which shows that every compact
$n$-manifold immerses in $\R^{2n-1}$ (for $n>1$). For specific manifolds,
even lower dimensions are possible: $\RP^2$ immerses in $\R^3$
(Boy's surface), and all orientable surfaces immerse in $\R^3$.
\end{remark}

% ============================================================
\section{The Strong Whitney Embedding Theorem}\label{sec:strong-whitney}
% ============================================================

\begin{theorem}[Strong Whitney embedding theorem]\label{thm:strong-whitney}
\index{Whitney embedding theorem!strong}
Every smooth $n$-manifold (with $n > 0$) admits a proper embedding
into $\R^{2n}$.
\end{theorem}

\begin{remark}\label{rem:strong-whitney}
The dimension $2n$ is sharp: the non-orientable manifold
$\RP^n$ (for $n$ a power of $2$) cannot be embedded in $\R^{2n-1}$
by Stiefel--Whitney class obstructions.

The proof of the strong Whitney theorem is substantially harder than
the weak version. The key new ingredient is the \textbf{Whitney trick}
\index{Whitney trick}: pairs of intersection points of opposite sign can
be cancelled by an isotopy, provided $2n \ge 5$ (i.e., $n \ge 3$).
For $n = 1$ the result is trivial (embed $S^1 \hookrightarrow \R^2$),
and $n = 2$ requires separate (difficult) arguments. We refer to
\cite{whitney-embed} and \cite{adachi} for the full proof.
\end{remark}

% ============================================================
\section{Whitney Approximation Theorems}\label{sec:whitney-approx}
% ============================================================

\begin{theorem}[Whitney approximation --- functions]\label{thm:whitney-approx-func}
\index{Whitney approximation theorem}
Let $M$ be a smooth manifold (possibly non-compact) and
$f\colon M \to \R^N$ a continuous map. Then $f$ can be uniformly
approximated by smooth maps. More precisely, for any continuous
$\varepsilon\colon M \to (0,\infty)$, there exists a smooth map
$g\colon M \to \R^N$ with $\abs{f(p)-g(p)} < \varepsilon(p)$ for all $p$.
If $f$ is already smooth on a closed set $A$, then $g$ can be chosen
to agree with $f$ on $A$.
\end{theorem}

\begin{theorem}[Whitney approximation --- maps between manifolds]\label{thm:whitney-approx-mfld}
Every continuous map between smooth manifolds is homotopic to a
smooth map. If two smooth maps are continuously homotopic, they
are smoothly homotopic.
\end{theorem}

\begin{remark}\label{rem:smooth-homotopy-category}
\autoref{thm:whitney-approx-mfld} has the profound consequence that
the \textbf{smooth and continuous homotopy categories of manifolds are equivalent}.
In particular, any homotopy-theoretic invariant (fundamental group,
homology, homotopy groups) can be computed using either smooth or
continuous maps. This justifies the free passage between smooth and
continuous settings throughout differential topology.
\end{remark}

\begin{definition}[Isotopy]\label{def:isotopy}
\index{isotopy}
An \textbf{isotopy} of embeddings $f_0,f_1\colon M\hookrightarrow N$
is a smooth map $F\colon M\times[0,1]\to N$ such that $F_0=f_0$,
$F_1=f_1$, and $F_t$ is an embedding for each $t \in [0,1]$.
An \textbf{ambient isotopy} is a smooth family of diffeomorphisms
$\Phi_t\colon N \to N$ with $\Phi_0 = \Id$ and $\Phi_1\circ f_0 = f_1$.
\end{definition}

\begin{theorem}[Isotopy extension theorem]\label{thm:isotopy-extension}
\index{isotopy extension theorem}
Let $M$ be compact and $N$ a manifold without boundary.
Every isotopy of embeddings $f_t\colon M \hookrightarrow N$ extends
to an ambient isotopy $\Phi_t\colon N \to N$.
\end{theorem}

% ============================================================
\section{Applications and Examples}\label{sec:whitney-applications}
% ============================================================

\begin{example}[Standard embeddings]\label{ex:standard-embeddings}
\index{embedding!examples}
Some classical embeddings:
\begin{enumerate}[label=(\alph*)]
  \item $S^n \hookrightarrow \R^{n+1}$ as the unit sphere (codimension $1$).
  \item $T^2 = S^1 \times S^1 \hookrightarrow \R^3$:
    $(e^{i\theta},e^{i\varphi}) \mapsto
    \bigl((2+\cos\varphi)\cos\theta,\;(2+\cos\varphi)\sin\theta,\;\sin\varphi\bigr)$.
  \item $T^n = (S^1)^n \hookrightarrow \R^{2n}$ via
    $(e^{i\theta_1},\dots,e^{i\theta_n})\mapsto(\cos\theta_1,\sin\theta_1,\dots,\cos\theta_n,\sin\theta_n)$,
    which is an embedding (``flat torus'').
\end{enumerate}
\end{example}

\begin{example}[$\RP^2$ in $\R^4$]\label{ex:rp2-r4}
\index{real projective plane!embedding}
The real projective plane $\RP^2$ does not embed in $\R^3$
(since it is non-orientable and compact without boundary). However,
it embeds in $\R^4$. An explicit embedding is given by
\[
  [x:y:z] \mapsto (xy,\; xz,\; y^2 - z^2,\; 2yz),
\]
restricted to $S^2 \to \R^4$, which descends to a well-defined
embedding $\RP^2 \hookrightarrow \R^4$.
\end{example}

\begin{example}[Boy's surface --- an immersion of $\RP^2$ in $\R^3$]\label{ex:boy-surface}
\index{Boy's surface}
Although $\RP^2$ cannot be embedded in $\R^3$, it can be \emph{immersed}
in $\R^3$. The famous \textbf{Boy surface} is an immersion
$\RP^2 \looparrowright \R^3$ with a single triple point and no pinch points.
It was discovered by Werner Boy (a student of Hilbert) in 1901,
disproving Hilbert's conjecture that no such immersion exists.
An explicit parameterization can be given using the Apéry formula;
see \cite{guillemin-pollack} for a discussion.
\end{example}

\begin{example}[Non-orientable surfaces cannot embed in $\R^3$]\label{ex:nonorientable-no-r3}
No compact non-orientable surface (without boundary) embeds in $\R^3$.
This follows from the \textbf{Jordan--Brouwer separation theorem}:
\index{Jordan--Brouwer separation theorem}
a compact connected hypersurface $\Sigma^{n-1}\subset\R^n$ separates
$\R^n$ into two connected components (a bounded ``inside'' and an
unbounded ``outside''). The outward normal then gives a continuous
nowhere-vanishing normal field, which orients $\Sigma$.
\end{example}

\begin{proposition}[Embedding dimension and Stiefel--Whitney classes]\label{prop:embed-obstruction}
\index{Stiefel--Whitney class}
If a compact $n$-manifold $M$ embeds in $\R^{n+k}$, then the
$(k+1)$-st through $n$-th Stiefel--Whitney classes vanish:
$\overline{w}_i(M) = 0$ for $i > k$ (where $\overline{w}$ denotes the
dual Stiefel--Whitney classes). This provides obstructions to
low-codimension embeddings.
\end{proposition}

\begin{example}\label{ex:rp-embed-obstruction}
The non-immersion $\RP^{2^k}\not\looparrowright\R^{2^{k+1}-2}$ can
be proved using Stiefel--Whitney classes. For instance,
$\RP^4$ does not immerse in $\R^6$, showing that Whitney's bound
of $\R^{2n}$ is close to optimal for this family.
\end{example}

\begin{figure}[ht]
\centering
\begin{tikzpicture}[scale=0.8]
  % Schematic cross-sections for the torus embedding in R^3
  \draw[thick,blue!70!black] (0,0) ellipse (3 and 1.5);
  \draw[thick,blue!70!black] (0,0) ellipse (1 and 0.5);
  \draw[thick,blue!70!black,dashed] (-3,0) arc (180:360:3 and 0.8);
  \draw[thick,blue!70!black] (-3,0) arc (180:0:3 and 0.8);
  \draw[thick,blue!70!black,dashed] (-1,0) arc (180:360:1 and 0.3);
  \draw[thick,blue!70!black] (-1,0) arc (180:0:1 and 0.3);
  \node at (0,-2.5) {$T^2 \hookrightarrow \R^3$};
\end{tikzpicture}
\caption{Schematic of the standard embedding of the torus $T^2$ in $\R^3$.}\label{fig:torus-r3}
\end{figure}

% ============================================================
\section{Classification of Compact Surfaces}\label{sec:surface-classification}
% ============================================================

\begin{theorem}[Classification of compact surfaces]\label{thm:surface-classification}
\index{surface classification}
Every compact connected surface (2-manifold) without boundary is
diffeomorphic to exactly one of the following:
\begin{enumerate}[label=(\roman*)]
  \item The sphere $S^2$ (genus $0$, orientable),
  \item The connected sum $\Sigma_g = T^2 \# \cdots \# T^2$ ($g$ copies)
    for $g \ge 1$ (genus $g$, orientable),
  \item The connected sum $N_k = \RP^2 \# \cdots \# \RP^2$ ($k$ copies)
    for $k \ge 1$ (non-orientable).
\end{enumerate}
The invariants are the orientability and the Euler characteristic
$\chi(\Sigma_g) = 2-2g$ and $\chi(N_k) = 2-k$.
\end{theorem}

\begin{remark}\label{rem:surface-embed}
The embedding dimensions for compact surfaces are:
\begin{itemize}
  \item Every orientable surface $\Sigma_g$ embeds in $\R^3$.
  \item Every non-orientable surface $N_k$ embeds in $\R^4$
    (but not in $\R^3$). In particular, the Klein bottle $N_2$
    embeds in $\R^4$.
  \item By the weak Whitney theorem, all surfaces embed in $\R^5$;
    by the strong theorem, in $\R^4$.
\end{itemize}
\end{remark}

% ============================================================
\section{Exotic Spheres --- A Preview}\label{sec:exotic-spheres}
% ============================================================

\begin{remark}[Exotic spheres]\label{rem:exotic-spheres}
\index{exotic sphere}
A natural question is whether the smooth structure on $S^n$ is unique.
Remarkably, Milnor \cite{milnor-exotic} showed in 1956 that $S^7$ admits
exotic smooth structures: there exist smooth manifolds homeomorphic but
not diffeomorphic to $S^7$. The group $\Theta_n$ of exotic $n$-spheres
(up to orientation-preserving diffeomorphism, under connected sum) has been
computed for many $n$. Notable facts:
\begin{itemize}
  \item $\Theta_7 \cong \Z/28\Z$ (there are $28$ smooth structures on $S^7$),
  \item $\Theta_n = 0$ for $n \le 6$,
  \item $\Theta_4$ is unknown (the smooth Poincar\'{e} conjecture in
    dimension $4$ is open).
\end{itemize}
The construction of exotic spheres uses Morse theory and handle
decompositions in an essential way: Milnor's original examples arise as
boundaries of certain disk bundles over $S^4$, analyzed via Morse-theoretic
handle attachments.
\end{remark}

% ============================================================
\section{Exercises}\label{sec:exo-whitney}
% ============================================================

\begin{exercise}\label{exo:whitney-1}
Show that the Klein bottle $K = \RP^2 \# \RP^2$ does not embed in $\R^3$.
\emph{Hint:} a compact surface without boundary embedded in $\R^3$
separates $\R^3$ (Jordan--Brouwer) and hence is orientable.
\end{exercise}

\begin{exercise}\label{exo:whitney-2}
Construct an explicit embedding of $S^1 \times S^1$ in $\R^3$ and
verify it is an embedding (injective immersion, homeomorphism onto
its image).
\end{exercise}

\begin{exercise}\label{exo:whitney-3}
\index{Veronese embedding}
The \textbf{Veronese embedding} $\nu\colon \RP^2 \to \R^6$ is defined by
$\nu[x:y:z] = (x^2,y^2,z^2,\sqrt{2}\,xy,\sqrt{2}\,xz,\sqrt{2}\,yz)$
(on $S^2$, modulo the antipodal identification). Verify that $\nu$ is a
well-defined smooth embedding.
\end{exercise}

\begin{exercise}\label{exo:whitney-4}
Prove that any compact $1$-manifold (without boundary) is diffeomorphic
to a finite disjoint union of circles $S^1$. Deduce that it embeds in $\R^2$.
\end{exercise}

\begin{exercise}\label{exo:whitney-5}
Let $M^n$ be a compact parallelizable manifold (i.e., $TM$ is trivial).
Show that $M$ immerses in $\R^{n+1}$.
\emph{Hint:} use the trivialization to define a map to $S^n$
(Gauss map) and the immersion theorem.
\end{exercise}

\begin{exercise}\label{exo:whitney-6}
Show that $\RP^n$ can be embedded in $\R^{2n}$ for all $n\ge 1$.
For $n=1$, give an explicit embedding $\RP^1 \cong S^1 \hookrightarrow\R^2$.
\end{exercise}

\begin{exercise}\label{exo:whitney-7}
Prove that if $M^n$ embeds in $\R^{n+1}$ (codimension $1$), then $M$
is orientable if and only if the normal bundle is trivial.
\end{exercise}

\begin{exercise}\label{exo:whitney-8}
Show that the connected sum $M_1 \# M_2$ of two compact $n$-manifolds
embeds in $\R^{2n+1}$ if both $M_1$ and $M_2$ do. Does the same
hold for $\R^{2n}$?
\end{exercise}

\begin{exercise}\label{exo:whitney-9}
(Whitney--Graustein theorem)\index{Whitney--Graustein theorem}
Two immersions $\gamma_0,\gamma_1\colon S^1\looparrowright\R^2$ are
regularly homotopic if and only if they have the same
\textbf{turning number} $\tau(\gamma) = \deg(\gamma'/\abs{\gamma'}\colon S^1\to S^1)$.
Verify this for the figure-eight curve ($\tau=0$) and the standard
circle ($\tau=1$).
\end{exercise}


% %%%%%%%%%%%%%%%%%%%%%%%%%%%%%%%%%%%%%%%%%%%%%%%%%%%%%%%%%%%%
%  END MATTER
% %%%%%%%%%%%%%%%%%%%%%%%%%%%%%%%%%%%%%%%%%%%%%%%%%%%%%%%%%%%%

\appendix

% %%%%%%%%%%%%%%%%%%%%%%%%%%%%%%%%%%%%%%%%%%%%%%%%%%%%%%%%%%%%
%  APPENDIX A — REVIEW OF ANALYSIS AND TOPOLOGY
% %%%%%%%%%%%%%%%%%%%%%%%%%%%%%%%%%%%%%%%%%%%%%%%%%%%%%%%%%%%%

\chapter{Review of Analysis and Topology}\label{app:analysis}

This appendix collects standard results from analysis and point-set
topology used throughout the text. Proofs are generally omitted;
see \cite{lee-smooth} or \cite{rudin-analysis} for details.

% ============================================================
\section{Inverse and Implicit Function Theorems}\label{sec:app-ift}
% ============================================================

\begin{theorem}[Inverse function theorem]\label{thm:app-ift}
\index{inverse function theorem}
Let $U \subset \R^n$ be open and $F\colon U \to \R^n$ be $C^k$
($k \ge 1$). If $\det(\dd F_p) \neq 0$ at some $p \in U$, then there
exist open neighborhoods $V$ of $p$ and $W$ of $F(p)$ such that
$F|_V\colon V \xrightarrow{\sim} W$ is a $C^k$ diffeomorphism.
\end{theorem}

\begin{theorem}[Implicit function theorem]\label{thm:app-implicit}
\index{implicit function theorem}
Let $U \subset \R^n\times\R^m$ be open and $F\colon U \to \R^m$ be $C^k$.
Suppose $F(a,b)=0$ and the partial derivative $\frac{\partial F}{\partial y}(a,b)$
is invertible. Then there exist neighborhoods $V\ni a$ in $\R^n$ and
$W\ni b$ in $\R^m$, and a unique $C^k$ map $g\colon V\to W$ with
$g(a)=b$ and $F(x,g(x))=0$ for all $x\in V$.
\end{theorem}

% ============================================================
\section{Sard's Theorem}\label{sec:app-sard}
% ============================================================

\begin{theorem}[Sard's theorem]\label{thm:sard}
\index{Sard's theorem}
Let $f\colon M^m \to N^n$ be a smooth map between smooth manifolds.
The set of critical values of $f$ has Lebesgue measure zero in $N$.
Equivalently, the set of regular values is dense in $N$
(in fact, residual in the sense of Baire).
\end{theorem}

% ============================================================
\section{Partitions of Unity}\label{sec:app-partition}
% ============================================================

\begin{theorem}[Existence of partitions of unity]\label{thm:app-partition}
\index{partition of unity}
Let $M$ be a smooth manifold and $\set{U_\alpha}_{\alpha\in A}$ an
open cover of $M$. There exists a smooth partition of unity
$\set{\rho_\alpha}$ subordinate to $\set{U_\alpha}$, i.e.,
$\operatorname{supp}(\rho_\alpha)\subset U_\alpha$,
$\rho_\alpha \ge 0$, and $\sum_\alpha \rho_\alpha = 1$
(locally finite sum).
\end{theorem}

% ============================================================
\section{Transversality}\label{sec:app-transversality}
% ============================================================

\begin{definition}[Transversality]\label{def:app-transversality}
\index{transversality}
A smooth map $f\colon M \to N$ is \textbf{transverse} to a submanifold
$S \subset N$ (written $f \pitchfork S$) if for every $p \in f^{-1}(S)$,
\[
  \dd f_p(T_pM) + T_{f(p)}S = T_{f(p)}N.
\]
Two submanifolds $X,Y \subset N$ are \textbf{transverse} if the inclusion
$X \hookrightarrow N$ is transverse to $Y$.
\end{definition}

\begin{theorem}[Transversality theorem]\label{thm:app-transversality}
If $f\pitchfork S$, then $f^{-1}(S)$ is a submanifold of $M$
with $\operatorname{codim}(f^{-1}(S)) = \operatorname{codim}(S)$.
\end{theorem}

\begin{theorem}[Parametric transversality theorem]\label{thm:parametric-transversality}
\index{transversality!parametric}
Let $F\colon M\times S \to N$ be smooth and $F\pitchfork W$ for a
submanifold $W\subset N$. Then for almost every $s \in S$
(a residual set of full measure), the map $F_s = F(-,s)\colon M\to N$
satisfies $F_s \pitchfork W$.
\end{theorem}

% ============================================================
\section{Compact Manifolds and Exhaustion Functions}\label{sec:app-exhaustion}
% ============================================================

\begin{theorem}\label{thm:app-exhaustion}
Every smooth manifold $M$ admits a proper smooth function
$\rho\colon M \to [0,\infty)$ (an \textbf{exhaustion function}).
If $M$ is compact, $\rho$ can be taken to be any smooth function.
\end{theorem}

\begin{theorem}[Ehresmann's fibration theorem]\label{thm:app-ehresmann}
\index{Ehresmann fibration theorem}
Let $f\colon M\to N$ be a proper smooth submersion between manifolds
without boundary. Then $f$ is a locally trivial fiber bundle.
\end{theorem}

% ============================================================
\section{Smooth Bump Functions and Cutoffs}\label{sec:app-bump}
% ============================================================

\begin{theorem}[Existence of bump functions]\label{thm:app-bump}
\index{bump function}
For any point $p$ in a smooth manifold $M$ and any open neighborhood
$U \ni p$, there exists a smooth function $\rho\colon M \to [0,1]$
with $\rho(p)=1$ and $\operatorname{supp}(\rho)\subset U$.
\end{theorem}

\begin{remark}\label{rem:app-bump-consequences}
The existence of bump functions has several important consequences:
\begin{enumerate}[label=(\roman*)]
  \item Smooth manifolds are normal topological spaces (Urysohn-type
    separation via smooth functions).
  \item Any closed subset of a smooth manifold is the zero set of some
    smooth non-negative function.
  \item Smooth functions separate points: for any $p \neq q$ in $M$,
    there exists $\rho \in C^\infty(M)$ with $\rho(p) \neq \rho(q)$.
\end{enumerate}
\end{remark}

% ============================================================
\section{Tubular Neighborhoods}\label{sec:app-tubular}
% ============================================================

\begin{theorem}[Tubular neighborhood theorem]\label{thm:app-tubular}
\index{tubular neighborhood}
Let $S \subset M$ be a compact embedded submanifold. Then $S$ has a
\textbf{tubular neighborhood}: there exist an open neighborhood $U$
of $S$ in $M$ and a diffeomorphism $\Phi\colon \nu(S) \xrightarrow{\sim} U$
from (a neighborhood of the zero section in) the normal bundle $\nu(S)$
to $U$, with $\Phi|_S = \Id_S$.
\end{theorem}

\begin{remark}\label{rem:tubular-uniqueness}
The tubular neighborhood is unique up to isotopy: any two tubular
neighborhood embeddings $\Phi_0,\Phi_1\colon \nu(S)\to M$ that agree
on $S$ are isotopic through tubular neighborhood embeddings.
This is a consequence of the isotopy extension theorem.
\end{remark}

% ============================================================
%  BIBLIOGRAPHY
% ============================================================

\begin{thebibliography}{99}

\bibitem{milnor-morse}
\index{Milnor, J.}
J.~Milnor,
\textit{Morse Theory},
Annals of Mathematics Studies 51,
Princeton University Press, 1963.

\bibitem{milnor-exotic}
J.~Milnor,
\textit{On manifolds homeomorphic to the 7-sphere},
Annals of Mathematics \textbf{64} (1956), 399--405.

\bibitem{milnor-topology}
J.~Milnor,
\textit{Topology from the Differentiable Viewpoint},
Princeton University Press, 1965.

\bibitem{guillemin-pollack}
\index{Guillemin, V.}\index{Pollack, A.}
V.~Guillemin and A.~Pollack,
\textit{Differential Topology},
Prentice-Hall, 1974; AMS Chelsea Publishing, 2010.

\bibitem{hirsch}
\index{Hirsch, M.}
M.W.~Hirsch,
\textit{Differential Topology},
Graduate Texts in Mathematics 33,
Springer-Verlag, 1976.

\bibitem{lee-smooth}
\index{Lee, J.M.}
J.M.~Lee,
\textit{Introduction to Smooth Manifolds},
Graduate Texts in Mathematics 218,
Springer, 2nd edition, 2013.

\bibitem{broecker-jaenich}
\index{Br\"{o}cker, Th.}\index{J\"{a}nich, K.}
Th.~Br\"{o}cker and K.~J\"{a}nich,
\textit{Introduction to Differential Topology},
Cambridge University Press, 1982.

\bibitem{bott-tu}
\index{Bott, R.}\index{Tu, L.W.}
R.~Bott and L.W.~Tu,
\textit{Differential Forms in Algebraic Topology},
Graduate Texts in Mathematics 82,
Springer-Verlag, 1982.

\bibitem{whitney-embed}
\index{Whitney, H.}
H.~Whitney,
\textit{The self-intersections of a smooth $n$-manifold in $2n$-space},
Annals of Mathematics \textbf{45} (1944), 220--246.

\bibitem{adachi}
\index{Adachi, M.}
M.~Adachi,
\textit{Embeddings and Immersions},
Translations of Mathematical Monographs 124,
American Mathematical Society, 1993.

\bibitem{matsumoto}
\index{Matsumoto, Y.}
Y.~Matsumoto,
\textit{An Introduction to Morse Theory},
Translations of Mathematical Monographs 208,
American Mathematical Society, 2002.

\bibitem{kosinski}
\index{Kosinski, A.}
A.~Kosinski,
\textit{Differential Manifolds},
Academic Press, 1993; Dover, 2007.

\bibitem{rudin-analysis}
\index{Rudin, W.}
W.~Rudin,
\textit{Principles of Mathematical Analysis},
McGraw-Hill, 3rd edition, 1976.

\bibitem{Lee2013}
\index{Lee, J.M.}
J.M.~Lee,
\textit{Introduction to Smooth Manifolds},
Graduate Texts in Mathematics 218,
Springer, 2nd edition, 2013.

\bibitem{GP2010}
\index{Guillemin, V.}\index{Pollack, A.}
V.~Guillemin and A.~Pollack,
\textit{Differential Topology},
Prentice-Hall, 1974; AMS Chelsea Publishing, 2010.

\bibitem{Tu2011}
\index{Tu, L.W.}
L.W.~Tu,
\textit{An Introduction to Manifolds},
Universitext, Springer, 2nd edition, 2011.

\bibitem{Hirsch1976}
\index{Hirsch, M.}
M.W.~Hirsch,
\textit{Differential Topology},
Graduate Texts in Mathematics 33,
Springer-Verlag, 1976.

\end{thebibliography}

% ============================================================
%  INDEX
% ============================================================

\printindex

\end{document}
% === END OF CHUNK ch9-11+end ===
